\section{The normalized primitive gauge of a normal tensor}

A normal tensor need not be trace-preserving or have spectral radius one.
Standard treatments normalize the tensor by scaling the spectral radius of the
transfer map to one and applying a gauge transformation by the square root of a
positive fixed point~\cite{Nachtergaele1996Spectral,PerezGarcia2007Matrix}. This
gauge reduction extends the finite-range spectral gap to arbitrary normal
tensors.

\begin{lemma}[Rescaling invariance of block injectivity and normality]
    \label{thm:normality_rescaling_invariance}
    \lean{MPSTensor.wordSpan_smul_eq, MPSTensor.isNBlkInjective_smul_iff,
        MPSTensor.isNormal_smul_iff}
    \leanok
    \uses{def:word_span, def:l_blk_injective, def:normal}
    Let $\zeta\neq0$ and let $A$ be a tensor.  For every $N$ the length-$N$
    word spans of $A$ and of $\zeta A$ agree, so $\zeta A$ is $N$-block
    injective exactly when $A$ is, and $\zeta A$ is normal exactly when $A$ is.
\end{lemma}

\begin{proof}
    \leanok
    \uses{lem:eval_word_smul}
    Every length-$N$ word of $\zeta A$ is the corresponding word of $A$
    multiplied by $\zeta^{N}$, and $\zeta^{N}$ is invertible, so each span
    contains the other.
\end{proof}

\begin{lemma}[A normal tensor has a nonzero matrix]
    \label{thm:normal_tensor_nonzero_matrix}
    \lean{MPSTensor.exists_apply_ne_zero_of_isNormal}
    \leanok
    \uses{def:normal}
    Let $D>0$ and let $A$ be a normal tensor.  Then $A^i\neq0$ for at least
    one $i$.
\end{lemma}

\begin{proof}
    \leanok
    \uses{def:l_blk_injective}
    If every $A^i$ vanished, then every word of positive length would vanish,
    so the span of the length-$N$ words would be $\{0\}$ for the length $N>0$
    at which $A$ is block injective.  That span is $\MN{D}$, which contains
    $\Id\neq0$ because $D>0$.
\end{proof}

\begin{lemma}[Trace pairing of a conjugated word against a boundary matrix]
    \label{lem:trace_gauge_boundary}
    \lean{MPSTensor.trace_gauge_boundary}
    \leanok
    Let $X$ be invertible and let $E$ and $Y$ be $D\times D$ matrices.  Then
    \begin{align}
        \tr\bigl(XEX^{-1}Y\bigr) & =\tr\bigl(E\,X^{-1}YX\bigr). \notag
    \end{align}
\end{lemma}

\begin{proof}
    \leanok
    Cyclicity of the trace moves the leading factor $X$ to the end of the
    product.
\end{proof}

\begin{lemma}[Gauge invariance of the local ground space]
    \label{thm:local_ground_space_gauge_invariance}
    \lean{MPSTensor.GaugeEquiv.groundSpace_le,
        MPSTensor.GaugeEquiv.groundSpace_eq}
    \leanok
    \uses{def:gauge_equiv, def:ground_space_parent}
    Let $A$ and $B$ be gauge equivalent.  For every $L$ one has
    $G_L(A)=G_L(B)$.
\end{lemma}

\begin{proof}
    \leanok
    \uses{lem:eval_word_gauge, lem:trace_gauge_boundary}
    By Lemma~\ref{lem:eval_word_gauge}, $B^i=XA^iX^{-1}$ implies
    $B^\sigma=XA^\sigma X^{-1}$ for every length-$L$ configuration $\sigma$, so
    Lemma~\ref{lem:trace_gauge_boundary} identifies the boundary
    parametrization of $B$ at $Y$ with that of $A$ at $X^{-1}YX$.  This gives
    $G_L(B)\subseteq G_L(A)$, and the symmetric argument applied to
    $A^i=X^{-1}B^iX$ gives the reverse inclusion.
\end{proof}

\begin{lemma}[Rescaling invariance of the local ground space]
    \label{thm:local_ground_space_rescaling}
    \lean{MPSTensor.groundSpace_smul_eq}
    \leanok
    \uses{def:ground_space_parent}
    Let $\zeta\neq0$.  For every $L$ one has $G_L(\zeta A)=G_L(A)$.
\end{lemma}

\begin{proof}
    \leanok
    \uses{lem:eval_word_smul}
    The boundary parametrization of $\zeta A$ at length $L$ sends $X$ to the
    value of the boundary parametrization of $A$ at $\zeta^{L}X$, and
    multiplication by $\zeta^{L}$ is a bijection of $\MN{D}$.
\end{proof}

\begin{lemma}[Local ground spaces determine the parent Hamiltonian]
    \label{thm:parent_hamiltonian_from_ground_space}
    \lean{MPSTensor.parentInteraction_eq_of_groundSpace_eq,
        MPSTensor.localTerm_eq_of_parentInteraction_eq,
        MPSTensor.parentHamiltonian_eq_of_groundSpace_eq,
        MPSTensor.parentHamiltonianES_eq_of_groundSpace_eq}
    \leanok
    \uses{def:ground_space_parent, def:parent_interaction,
        def:parent_hamiltonian, def:parent_hamiltonian_es}
    Let $A$ and $B$ have the same bond dimension and satisfy $G_L(A)=G_L(B)$.
    Then their canonical parent interactions at length $L$ agree, hence so do
    all their translated local terms and their parent Hamiltonians
    $H_{N,L}(A)=H_{N,L}(B)$ and $H_{N,L}^{\mathrm{ES}}(A)
        =H_{N,L}^{\mathrm{ES}}(B)$ on every periodic chain.
\end{lemma}

\begin{proof}
    \leanok
    \uses{def:ground_space_es, def:parent_interaction, def:local_term_parent,
        def:parent_hamiltonian, def:parent_hamiltonian_es}
    The parent interaction at length $L$ is the orthogonal projector onto
    $G_L(A)^\perp$, so it depends on $G_L(A)$ alone.  Every translated local
    term is built from that single projector, and the parent Hamiltonian is
    their sum.
\end{proof}

\begin{lemma}[Normality alone gives an irreducible tensor]
    \label{thm:normal_irreducible_tensor}
    \lean{MPSTensor.isPrimitivePaper_of_isNormal,
        MPSTensor.isIrreducibleTensor_of_isPrimitivePaper}
    \leanok
    \uses{def:normal, def:paper_primitive, def:has_inv_proj,
        def:irreducible_tensor}
    Let $A$ be a normal tensor, with no normalization imposed on $A$.  Under
    Definition~\ref{def:paper_primitive} and
    Definition~\ref{def:irreducible_tensor}, $A$ is primitive and has no
    nontrivial invariant projection, so $A$ is irreducible.
\end{lemma}

\begin{proof}
    \leanok
    \uses{def:eventually_full_kraus_rank, rem:wolf_6_8_pairwise}
    Normality is eventual full Kraus rank
    (Remark~\ref{rem:wolf_6_8_pairwise}), so $S_N(A)=\MN{D}$ for some $N>0$.
    Applying the length-$N$ words to a nonzero $\varphi$ then spans $\C^D$,
    which is the spreading condition.  A subspace invariant under all matrices
    of $A$ and containing a nonzero vector therefore contains the span of the
    images of that vector, namely all of $\C^D$; so no invariant projection is
    nontrivial.
\end{proof}

\begin{lemma}[Normalized primitive gauge of a normal tensor]
    \label{thm:normalized_primitive_gauge}
    \lean{MPSTensor.exists_isPrimitiveMPS_gauge_of_isNormal}
    \leanok
    \uses{def:normal, def:primitive_mps, def:gauge_equiv, def:mpv,
        def:ground_space_parent, def:parent_interaction,
        def:chain_ground_space}
    Let $D>0$ and let $A$ be a normal tensor.  There exist a scalar
    $\zeta\neq0$, a tensor $B$, and a matrix $\rho>0$ such that $\zeta A$ and
    $B$ are gauge equivalent, the pair $(B,\rho)$ is primitive, and
    \begin{align}
        V^{(N)}(B)_{\sigma}
         & =\zeta^{N}V^{(N)}(A)_{\sigma} \notag
    \end{align}
    for every chain length $N$ and every configuration $\sigma$.  The local
    ground spaces, the parent interactions, and the periodic chain ground
    spaces of $A$ and $B$ also coincide:
    \begin{align}
        G_L(A)              & =G_L(B), \notag             \\
        h_L(A)              & =h_L(B), \notag             \\
        \mathcal G_{N,L}(A) & =\mathcal G_{N,L}(B) \notag
    \end{align}
    for all $L$ and $N$, where $h_L$ is the canonical parent interaction and
    $\mathcal G_{N,L}$ the periodic chain ground space.

    Following the normalization taken without loss of generality in
    \cite[equations~(3.1)--(3.2b), lines~1394--1435]{Nachtergaele1996Spectral}
    and at the start of the proof of
    \cite[Theorem~4, lines~765--770]{PerezGarcia2007Matrix}, the gauge here is
    built from the Perron eigenvector of the adjoint transfer map, giving the
    trace-preserving orientation.  That orientation is what conditions~1 and~2
    of \cite[Theorem~4, lines~752--758]{PerezGarcia2007Matrix},
    $\sum_iA^i(A^i)^\dagger=\Id$ and $\sum_i(A^i)^\dagger\Lambda A^i=\Lambda$
    with $\Lambda>0$, become after the further gauge by $\Lambda^{1/2}$:
    condition~2 turns into $\sum_i(B^i)^\dagger B^i=\Id$ and condition~1 into
    the statement that $\Lambda$ is a positive definite fixed point of the
    transfer map of $B$.
\end{lemma}

\begin{proof}
    \leanok
    \uses{thm:normal_irreducible_tensor,
        thm:normalized_normal_strongly_irreducible,
        thm:primitive_from_strongly_irred,
        thm:normality_rescaling_invariance,
        thm:local_ground_space_rescaling, thm:normal_tensor_nonzero_matrix,
        thm:tp_data_irreducible, thm:gauge_invariance_normal,
        thm:gauge_invariance, thm:local_ground_space_gauge_invariance,
        lem:eval_word_smul,
        thm:parent_hamiltonian_from_ground_space,
        lem:parent_spanning_of_local_space_eq}
    Lemma~\ref{thm:normal_irreducible_tensor} makes $A$ an irreducible
    tensor, with no normalization imposed on $A$, and
    Lemma~\ref{thm:normal_tensor_nonzero_matrix} gives the nonzero matrix
    that the Perron--Frobenius eigenvector requires.
    Lemma~\ref{thm:tp_data_irreducible} therefore returns a
    positive definite $\sigma$ and an eigenvalue $r>0$ of the adjoint transfer
    map.  With $\zeta=r^{-1/2}$ and $B^i=\zeta\,\sigma^{1/2}A^i\sigma^{-1/2}$,
    one has $\sum_i(B^i)^\dagger B^i=\Id$, and $\zeta A$ and $B$ are gauge
    equivalent.  Because rescaling preserves normality by
    Lemma~\ref{thm:normality_rescaling_invariance} and gauging preserves it
    by Lemma~\ref{thm:gauge_invariance_normal}, $B$ is normal.  Since $B$
    is trace preserving,
    Lemma~\ref{thm:normalized_normal_strongly_irreducible} makes $B$
    strongly irreducible, and
    Lemma~\ref{thm:primitive_from_strongly_irred} then returns a positive
    definite $\rho$ with $(B,\rho)$ primitive.  By
    Theorem~\ref{thm:gauge_invariance} and Lemma~\ref{lem:eval_word_smul},
    gauge equivalence leaves the matrix product vectors unchanged and
    multiplying every matrix by $\zeta$ multiplies every length-$N$ word by
    $\zeta^{N}$, hence every length-$N$ coefficient of $B$ is $\zeta^{N}$ times
    that of $A$.  By Lemma~\ref{thm:local_ground_space_rescaling} together
    with gauge invariance of the local ground space, $G_L(A)=G_L(\zeta A)=G_L(B)$;
    Lemma~\ref{thm:parent_hamiltonian_from_ground_space} and
    Lemma~\ref{lem:parent_spanning_of_local_space_eq} then identify the parent
    interaction and the periodic chain ground space.
\end{proof}

\begin{lemma}[Uniform finite-range gap for normal MPS parent Hamiltonians]
    \label{thm:normal_parent_hamiltonian_finite_range_knabe_gap}
    \lean{MPSTensor.exists_parentHamiltonianES_gap_of_isNormal}
    \leanok
    \uses{def:normal, def:l_blk_injective, def:parent_hamiltonian_es}
    Let $D>0$ and let $A$ be a normal tensor, with no normalization imposed on
    $A$.  There exist $l>1$, $\epsilon_l\geq0$, an integer $m\geq l+1$, and
    $\delta>0$ such that $A$ is $l$-block injective,
    $\epsilon_l<1/\sqrt{l+1}$, and the arbitrary finite-range coefficient
    (\ref{eq:normal_finite_range_knabe_delta}) derived in
    \path{docs/paper-gaps/knabe88_finite_range_coefficient.tex},
    \begin{align}
        \delta
         & =\frac{m\left(1-\epsilon_l\sqrt{l+1}\right)^2-l^2}{m-l},
        \label{eq:normal_finite_range_knabe_delta}
    \end{align}
    realizes the uniform gap of matrix product state parent Hamiltonians
    asserted at \cite[lines~2183--2187]{Cirac2021Matrix}: for every
    $N\geq2m$ and every $v\in(\ker H_{N,l+1}^{\mathrm{ES}}(A))^\perp$,
    \begin{align}
        \delta\lVert v\rVert
         & \leq\left\lVert H_{N,l+1}^{\mathrm{ES}}(A)v\right\rVert.
        \label{eq:normal_finite_range_knabe_gap}
    \end{align}
\end{lemma}

\begin{proof}
    \leanok
    \uses{thm:normalized_primitive_gauge,
        thm:primitive_parent_hamiltonian_finite_range_knabe_gap,
        thm:parent_hamiltonian_from_ground_space,
        thm:normality_rescaling_invariance,
        lem:gauge_invariance_block_injective}
    Applying Theorem~\ref{thm:primitive_parent_hamiltonian_finite_range_knabe_gap}
    to the normalized primitive representative $B$ of
    Lemma~\ref{thm:normalized_primitive_gauge} yields the bound for $B$.
    Since $G_L(A)=G_L(B)$ for every $L$,
    Lemma~\ref{thm:parent_hamiltonian_from_ground_space} makes the two
    parent Hamiltonians equal, so both the kernel and the norm bound pass to
    $A$.  By Lemma~\ref{lem:gauge_invariance_block_injective} applied to the
    inverse gauge and Lemma~\ref{thm:normality_rescaling_invariance}, block
    injectivity of $B$ at length $l$ passes first to $\zeta A$ and then to $A$.
\end{proof}

\begin{lemma}[Positive gap transfer under Loewner domination]
    \label{lem:positive_gap_transfer_loewner}
    \lean{LinearMap.IsPositive.re_inner_ge_of_norm_gap,
        LinearMap.IsPositive.norm_gap_of_le_of_ker_eq}
    \leanok
    Let $P$ and $Q$ be finite-dimensional complex endomorphisms, and let
    $\gamma\geq0$.  Suppose that $P$ is positive and
    \begin{align}
        \gamma\lVert v\rVert
         & \leq\lVert Pv\rVert
    \end{align}
    for every $v\in(\ker P)^\perp$.  Then
    \begin{align}
        \gamma\lVert v\rVert^2
         & \leq\operatorname{Re}\langle Pv,v\rangle
    \end{align}
    on $(\ker P)^\perp$.  If also $P\leq Q$ and $\ker P=\ker Q$, then
    \begin{align}
        \gamma\lVert v\rVert
         & \leq\lVert Qv\rVert
    \end{align}
    for every $v\in(\ker Q)^\perp$.
\end{lemma}

% Source: Rozmán--Molnár--Schuch, arXiv:2607.19078v1, Lemma block-obc
% (statement at main.tex lines 607--651, argument at lines 795--811).  The
% source tex is not vendored under Papers/, and references.bib has no key for
% it; cite by footnote until the paper is vendored and a bib entry is added.
\begin{lemma}[Kernels of positively weighted sums of positive operators]
    \label{thm:weighted_sum_kernel}
    \lean{LinearMap.IsPositive.apply_eq_zero_of_re_inner_eq_zero,
        WeightedPositiveKernel.re_inner_weighted_sum_apply,
        WeightedPositiveKernel.apply_weighted_sum_eq_zero_iff,
        WeightedPositiveKernel.apply_sum_eq_zero_iff,
        WeightedPositiveKernel.ker_weighted_sum_eq_iInf,
        WeightedPositiveKernel.ker_sum_eq_iInf,
        WeightedPositiveKernel.ker_weighted_sum_eq_sum}
    \leanok
    Let $(P_i)_{i\in I}$ be a finite family of positive operators on an
    inner product space.  A positive operator $P$ annihilates a vector $v$
    as soon as $\operatorname{Re}\langle Pv,v\rangle=0$, and consequently,
    for real weights $w_i>0$,
    \begin{align}
        \ker\Bigl(\sum_{i\in I}w_iP_i\Bigr)
        =\ker\Bigl(\sum_{i\in I}P_i\Bigr)
        =\bigcap_{i\in I}\ker P_i,
        \label{eq:ph_weighted_sum_kernel}
    \end{align}
    equivalently $(\sum_{i\in I}w_iP_i)v=0$ exactly when $P_iv=0$ for
    every $i\in I$.  For the empty family both sums are zero and all three
    kernels are the full space.  In the blocked-sum comparison of
    Lemma~\texttt{block-obc} of
    Rozmán--Molnár--Schuch\footnote{arXiv:2607.19078v1.}, where the
    blocked parent Hamiltonian is a weighted sum $\sum_i\gamma_ih_i$ of
    positive local terms with $\gamma_i>0$, this identity provides the
    analytic kernel step.
\end{lemma}

\begin{proof}
    \leanok
    If $(\sum_iw_iP_i)v=0$, the real part of the inner product with $v$
    expands as
    \begin{align}
        \operatorname{Re}\Bigl\langle\Bigl(\sum_iw_iP_i\Bigr)v,v\Bigr\rangle
        =\sum_iw_i\operatorname{Re}\langle P_iv,v\rangle
        =0.
    \end{align}
    Each summand is nonnegative, so each vanishes, and $w_i>0$ forces
    $\operatorname{Re}\langle P_iv,v\rangle=0$.  For a positive operator
    $P$, the quadratic form is nonnegative along the real affine line
    $t\mapsto v+t\cdot(Pv)$ and vanishes at $t=0$; its linear coefficient
    is $2\lVert Pv\rVert^2$, so evaluating at small negative $t$ forces
    $\lVert Pv\rVert=0$, that is $Pv=0$.  The converse inclusion holds
    term by term, and the unweighted comparison is the case $w_i=1$.
\end{proof}

\begin{lemma}[Loewner domination of positively weighted sums]
    \label{thm:weighted_sum_loewner}
    \lean{WeightedPositiveKernel.isPositive_weighted_sum,
        WeightedPositiveKernel.weighted_sum_le_sum}
    \leanok
    Let $(P_i)_{i\in I}$ be a finite family of positive operators on an
    inner product space, with real weights $(w_i)_{i\in I}$.  If
    $w_i\geq0$ for every $i$, then $\sum_{i\in I}w_iP_i$ is positive.  If
    $w_i\leq1$ for every $i$, as in the comparison step of the blocked-sum
    argument in the source cited in Lemma~\ref{thm:weighted_sum_kernel}
    where the coefficients satisfy $0<\gamma_i\leq1$, then
    \begin{align}
        \sum_{i\in I}w_iP_i
         & \leq\sum_{i\in I}P_i
    \end{align}
    in Loewner order.
\end{lemma}

\begin{proof}
    \leanok
    Positivity of the weighted sum follows from positivity of each scalar
    multiple $w_iP_i$ and of finite sums of positive operators.  For the
    domination, the difference is again a weighted sum,
    \begin{align}
        \sum_{i\in I}P_i-\sum_{i\in I}w_iP_i
        =\sum_{i\in I}(1-w_i)P_i,
    \end{align}
    now with nonnegative weights $1-w_i\geq0$, hence positive.
\end{proof}

\begin{lemma}[Gap transfer through a positively weighted sum]
    \label{thm:weighted_sum_gap_transfer}
    \lean{WeightedPositiveKernel.norm_gap_sum_of_weighted_sum}
    \leanok
    Let $(P_i)_{i\in I}$ be a finite family of positive operators on a
    finite-dimensional complex inner product space, with real weights
    $0<w_i\leq1$, and let $\gamma\geq0$.  If
    \begin{align}
        \gamma\lVert v\rVert
         & \leq\Bigl\lVert\Bigl(\sum_{i\in I}w_iP_i\Bigr)v\Bigr\rVert
    \end{align}
    for every $v\in\bigl(\ker\sum_{i\in I}w_iP_i\bigr)^\perp$, then
    \begin{align}
        \gamma\lVert v\rVert
         & \leq\Bigl\lVert\Bigl(\sum_{i\in I}P_i\Bigr)v\Bigr\rVert
    \end{align}
    for every $v\in\bigl(\ker\sum_{i\in I}P_i\bigr)^\perp$.
\end{lemma}

\begin{proof}
    \leanok
    \uses{thm:weighted_sum_kernel, thm:weighted_sum_loewner, lem:positive_gap_transfer_loewner}
    The weighted sum is positive and, by
    Lemma~\ref{thm:weighted_sum_loewner}, dominated by the unweighted
    sum; by the kernel equality~(\ref{eq:ph_weighted_sum_kernel}) of
    Lemma~\ref{thm:weighted_sum_kernel}, the two sums have the same
    kernel.  Lemma~\ref{lem:positive_gap_transfer_loewner} transfers the
    norm gap from the smaller weighted sum to the larger unweighted sum.
\end{proof}

\begin{lemma}[Positive gap at a fixed length]
    \label{thm:parent_hamiltonian_fixed_length_gap}
    \lean{MPSTensor.parentHamiltonianES_exists_gap}
    \leanok
    \uses{def:parent_hamiltonian_es}
    For any tensor $A$, interaction range $L$, and chain length $N$, there
    exists $\delta_N>0$ such that
    $\delta_N\lVert v\rVert\leq\lVert H_{N,L}^{\mathrm{ES}}(A)v\rVert$
    for every $v\in(\ker H_{N,L}^{\mathrm{ES}}(A))^\perp$.
\end{lemma}

\begin{proof}
    \leanok
    The restriction of $H_{N,L}^{\mathrm{ES}}(A)$ to its kernel complement
    is injective, and its inverse on its image is bounded because the domain is
    finite dimensional. Thus $\lVert v\rVert\leq C_N\lVert
        H_{N,L}^{\mathrm{ES}}(A)v\rVert$ for some $C_N>0$, so
    $\delta_N=C_N^{-1}$ gives the desired bound; if the kernel complement is
    zero, any positive constant suffices.
\end{proof}

\begin{lemma}[Extension of an eventual uniform gap]
    \label{thm:parent_hamiltonian_eventual_gap}
    \lean{MPSTensor.parentHamiltonianES_gap_of_eventual_gap}
    \leanok
    \uses{def:parent_hamiltonian_es}
    Fix $A$ and $L$. Suppose $\delta>0$ and $M\in\mathbb N$ satisfy
    $\delta\lVert v\rVert\leq\lVert H_{N,L}^{\mathrm{ES}}(A)v\rVert$
    for every $N\geq M$ and every
    $v\in(\ker H_{N,L}^{\mathrm{ES}}(A))^\perp$.
    Then some $\gamma>0$ satisfies the inequality with $\gamma$ in place of
    $\delta$ for every $N$.
\end{lemma}

\begin{proof}
    \leanok
    \uses{thm:parent_hamiltonian_fixed_length_gap}
    Taking $\gamma$ to be the minimum of $\delta$ and the finitely many
    positive constants $\delta_N$ given by
    Lemma~\ref{thm:parent_hamiltonian_fixed_length_gap} for $N<M$ (or
    equivalently inducting on $M$ and replacing $\delta$ by
    $\min\{\delta,\delta_{M-1}\}$ at each step) extends the bound to every
    $N$.
\end{proof}

\begin{lemma}[Uniform gap for normal MPS parent Hamiltonians]
    \label{thm:normal_mps_parent_hamiltonian_gap}
    \lean{MPSTensor.exists_parentHamiltonianES_uniform_gap_of_isNormal}
    \leanok
    \uses{def:normal, def:l_blk_injective, def:parent_hamiltonian_es}
    Let $D>0$ and let $A$ be a normal tensor. For the normal-tensor case of
    the uniform gap assertion in \cite[lines~2183--2187]{Cirac2021Matrix},
    there exist a suitable interaction range $l+1$ with $l>1$ and a constant
    $\gamma>0$ such that $A$ is $l$-block injective and, for every $N$ and
    every $v\in(\ker H_{N,l+1}^{\mathrm{ES}}(A))^\perp$,
    \begin{align}
        \gamma\lVert v\rVert
         & \leq \lVert H_{N,l+1}^{\mathrm{ES}}(A)v\rVert.
        \notag
    \end{align}
\end{lemma}

\begin{proof}
    \leanok
    \uses{thm:normal_parent_hamiltonian_finite_range_knabe_gap,
        thm:parent_hamiltonian_eventual_gap}
    Lemma~\ref{thm:normal_parent_hamiltonian_finite_range_knabe_gap}
    gives $l>1$, $m\geq l+1$, and $\delta>0$ with the required norm
    inequality for $N\geq2m$, and
    Lemma~\ref{thm:parent_hamiltonian_eventual_gap} applied with
    $L=l+1$ and $M=2m$ extends it to every $N$.
\end{proof}

\begin{lemma}[Gap bound from a cyclic anticommutator estimate]
    \label{thm:anticommutator_gap_bound}
    \leanok
    \uses{lem:parent_hamiltonian_cyclic_overlap_anticommutator_gap}
    Let $L>1$. Suppose that for every $N\ge2L$ and every overlapping
    off-diagonal pair of cyclic length-$L$ windows,
    \begin{align}
        h_{i,\mathrm{ES}}h_{j,\mathrm{ES}}
        +h_{j,\mathrm{ES}}h_{i,\mathrm{ES}}
         & \ge
        -\left(1-\frac{1}{4L}\right)\frac{1}{2(L-1)}
        (h_{i,\mathrm{ES}}+h_{j,\mathrm{ES}})
        \notag
    \end{align}
    in quadratic form. Then the parent Hamiltonian has the explicit gap
    constant $\gamma_L:=1/(4L)>0$.
    Concretely, if $N\ge2L$ and
    $v\in(\mathcal G_{N,L}^{\mathrm{ES}}(A))^\perp$, then
    $\gamma_L\|v\| \le \|H_{N,L}^{\mathrm{ES}}(A)v\|$.
\end{lemma}

\begin{proof}
    \leanok
    The finite-overlap martingale reduction applies with the cyclic overlap
    relation: at most $2(L-1)$ off-diagonal neighbours per row; for
    non-overlapping pairs the disjoint-window estimate
    $\re\langle h_{i,\mathrm{ES}}v,h_{j,\mathrm{ES}}v\rangle\ge 0$ holds;
    and the anticommutator inequality is retained as an explicit hypothesis.
\end{proof}

\begin{lemma}[Conditional spectral gap for MPS parent Hamiltonians]
    \label{thm:parent_hamiltonian_gapped}
    \lean{MPSTensor.parentHamiltonian_gapped_of_anticommutator}
    \leanok
    \uses{thm:anticommutator_gap_bound}
    For a tensor $A$ and a range $L>1$, assume the cyclic-window
    anticommutator estimate of Lemma~\ref{thm:anticommutator_gap_bound}.
    Then there exists $\gamma>0$
    such that for every $N\ge2L$ and every vector
    $v\in(\mathcal G_{N,L}^{\mathrm{ES}}(A))^\perp$ one has
    $\gamma\|v\| \le \|H_{N,L}^{\mathrm{ES}}(A)v\|$.
    By Lemma~\ref{lem:quadratic_form_to_norm_bound}, this implies the usual
    uniform lower bound on every nonzero eigenvalue of the parent Hamiltonian.
\end{lemma}

\begin{proof}
    \leanok
    This follows immediately from the corresponding explicit theorem by taking
    $\gamma$ to be the constant produced there.
\end{proof}

\begin{definition}[Grouped prefix and interval projections]
    \label{def:grouped_open_ground_projections}
    \lean{MPSTensor.groupedNestedGroundProjectionsES,MPSTensor.groupedIntervalGroundProjectionES}
    \leanok
    \uses{def:fixed-volume-prefix-ground-projections,
        def:open_suffix_parent_hamiltonian_es}
    Fix $p\geq1$, $M\geq2$, and an open chain of length $N=pM$.
    For $0\leq n\leq M$, let $G_n$ be the ground projection of the range-$2p$ parent
    Hamiltonian on $[0,pn)$, extended by the identity on the remaining
    sites. For $0\leq n<M$, let $Q_n$ be the ground projection of the same interaction
    restricted to $[\max\{0,p(n-1)\},p(n+1))$.
    Thus $Q_0=I$, while for $n\geq1$ the interval Hamiltonian is
    a single parent interaction.
\end{definition}

\begin{lemma}[Locality and the initial grouped differences]
    \label{lem:grouped_open_ground_projection_locality}
    \lean{MPSTensor.grouped_martingaleDifference_commute,MPSTensor.grouped_martingaleDifference_zero,MPSTensor.groupedIntervalGroundProjectionES_comp_martingaleDifference_one}
    \leanok
    With $E_n=G_n-G_{n+1}$, every pair $0\leq m,n<M$ with $m<n-1$ or $m>n$
    satisfies
    \begin{align}
        [E_m,Q_n] & =0.
    \end{align}
    The initial differences satisfy $E_0=0$ and $Q_1E_1=0$.
\end{lemma}
\begin{proof}
    \leanok
    \uses{def:grouped_open_ground_projections,
        thm:physical-outside-window-martingale-cancellation,
        thm:fixed-volume-c3-boundary-indices,
        thm:fixed-volume-prefix-ground-projections-nested}
    In the first case both prefix Hamiltonians have support disjoint from
    the interval defining $Q_n$. In the second case their kernels are
    contained in the local interval kernel. The corresponding orthogonal
    projections therefore commute in either case.
    The prefixes of lengths zero and $p$ contain no interaction, so
    $G_0=G_1=I$. At the next step $Q_1=G_2$, giving
    $Q_1E_1=G_2(I-G_2)=0$.
\end{proof}

\begin{lemma}[The first two grouped differences need no estimate]
    \label{thm:grouped_c3_from_second_index}
    \lean{MPSTensor.grouped_martingaleDifference_norm_le_of_two_le}
    \leanok
    \uses{def:grouped_open_ground_projections}
    Let $p\geq1$, let $N\geq2p$, and let $M\geq0$. Define $G_n$, $Q_n$, and
    $E_n=G_n-G_{n+1}$ on the $N$-site chain by the formulas of
    Definition~\ref{def:grouped_open_ground_projections}. Let $\epsilon\geq0$
    and suppose that $\|Q_nE_n\|\leq\epsilon$ for $2\leq n<M$. Then
    $\|Q_nE_n\|\leq\epsilon$ for $0\leq n<M$.
\end{lemma}
\begin{proof}
    \leanok
    \uses{lem:grouped_open_ground_projection_locality}
    By Lemma~\ref{lem:grouped_open_ground_projection_locality},
    \begin{align}
        Q_0E_0 & =0, & Q_1E_1 & =0,
    \end{align}
    so both norms are $0\leq\epsilon$. For $2\leq n<M$ the bound is the
    hypothesis.
\end{proof}

\begin{lemma}[Grouped open-chain coercivity]
    \label{thm:grouped_open_parent_gap}
    \lean{MPSTensor.openParentHamiltonianES_norm_gap_of_grouped_c3}
    \leanok
    Let $p\geq1$, $M\geq2$, and $N=pM$. Consider the
    range-$2p$ open parent Hamiltonian $H_N$ and the grouped
    projections above. Let $0\leq\epsilon<1/\sqrt2$ and suppose that
    $\|Q_nE_n\|\leq\epsilon$ for $0\leq n<M$. Then, for every
    $v\perp\ker H_N$,
    \begin{align}
        (1-\epsilon\sqrt2)^2\|v\| & \leq\|H_Nv\|.
    \end{align}
\end{lemma}
\begin{proof}
    \leanok
    \uses{lem:grouped_open_ground_projection_locality,
        thm:nachtergaele_c1_c3_full_range_norm_gap,
        thm:grouped_open_parent_c1, thm:grouped_open_suffix_c2,
        lem:open_parent_hamiltonian_es_positive,
        thm:fixed-volume-c3-boundary-indices,
        thm:fixed-volume-c3-projector-conjugacy}
    Applying the full-range martingale estimate to the grouped filtration with
    window width one, where each local window Hamiltonian is the projection
    $I-Q_n$, yields the result: Lemma~\ref{thm:grouped_open_suffix_c2} gives
    condition C2 with constant one, Lemma~\ref{thm:grouped_open_parent_c1}
    gives condition C1 with constant one because the nonzero local window terms
    are distinct summands of $H_N$, and
    Lemma~\ref{lem:open_parent_hamiltonian_es_positive} gives positivity of
    $H_N$. The locality lemma provides the commutation condition, the assumed
    norm bound provides C3, and the terminal projection is the projection onto
    $\ker H_N$.
\end{proof}

\begin{lemma}[Energy counting for grouped open windows]
    \label{thm:grouped_open_parent_c1}
    \lean{MPSTensor.grouped_openParentHamiltonianES_C1}
    \leanok
    \uses{def:open_suffix_parent_hamiltonian_es, def:open_parent_hamiltonian_es}
    Let $p\geq1$, $M\geq0$, and $N=pM$. For every vector $v$ on the
    $N$-site chain,
    \begin{align}
        0 & \leq \sum_{n=0}^{M-1}
        \re\langle H^{\mathrm{suffix}}_{N,2p;2p,p(n+1)}(A)v,v\rangle
        \leq \re\langle H^{\mathrm{open}}_{N,2p}(A)v,v\rangle.
    \end{align}
\end{lemma}
\begin{proof}
    \leanok
    \uses{thm:open_parent_hamiltonian_c1_full_range,
        lem:fixed-volume-hamiltonians-positive}
    The indices $p(n+1)-1$, for $0\leq n<M$, are distinct and belong to
    $\{0,\ldots,pM-1\}$. Positivity bounds the sum over these indices by
    the full sum, which is bounded by the open-chain energy.
\end{proof}

\begin{lemma}[Unit local gap for grouped open windows]
    \label{thm:grouped_open_suffix_c2}
    \lean{MPSTensor.grouped_openSuffixParentHamiltonianES_C2}
    \leanok
    \uses{def:open_suffix_parent_hamiltonian_es}
    Let $p\geq1$, $0\leq n<M$, and $N=pM$. Write $Q_n$ for the
    orthogonal projection onto the kernel of
    $H^{\mathrm{suffix}}_{N,2p;2p,p(n+1)}(A)$. Then
    \begin{align}
        \lVert(I-Q_n)v\rVert^2
         & \leq\re\langle H^{\mathrm{suffix}}_{N,2p;2p,p(n+1)}(A)v,v\rangle
    \end{align}
    for every vector $v$ on the $N$-site chain.
\end{lemma}
\begin{proof}
    \leanok
    \uses{thm:open_suffix_parent_hamiltonian_c2_full_range}
    Applying the full-range local gap bound with $l=2p-1$ and index
    $p(n+1)-1<N$ yields the inequality.
\end{proof}

\begin{lemma}[Prefix ground spaces with free right sites]
    \label{thm:ph_prefix_kernel_free_right_sites}
    \lean{MPSTensor.mem_range_leftBoundaryMapES_iff_rightFiber,
        MPSTensor.ker_openPrefixParentHamiltonianES_eq_range_leftBoundaryMapES}
    \leanok
    Let $R>0$ and suppose that the length-$n$ open parent Hamiltonian has
    kernel $G_n(A)$. On $n+r$ sites, the kernel of the same interactions
    restricted to the first $n$ sites is $G_n(A)\otimes(\C^d)^{\otimes r}$.
    A vector belongs to this space if and only if each restriction obtained
    by fixing the last $r$ physical coordinates belongs to $G_n(A)$.
\end{lemma}
\begin{proof}
    \leanok
    \uses{thm:open-hamiltonian-right-spectator-conjugacy,
        thm:fixed-volume-c3-projector-conjugacy,
        lem:right-spectator-kernel-projection,
        thm:left_boundary_map_range_eq}
    Decomposing a vector $v$ on $n+r$ sites according to the last $r$
    physical coordinates, with $v_s$ denoting the vector on $n$ sites
    obtained by fixing these coordinates to $s\in\{0,\ldots,d-1\}^r$, shows
    that the prefix Hamiltonian acts on each component $v_s$ by the length-$n$
    open Hamiltonian, whose kernel is $G_n(A)$. Hence
    $v\in G_n(A)\otimes(\C^d)^{\otimes r}$ if and only if, for every
    $s\in\{0,\ldots,d-1\}^r$,
    \begin{align}
        v_s & \in G_n(A).
    \end{align}
\end{proof}

\begin{lemma}[Reassociated suffix ground space]
    \label{thm:ph_reassociated_suffix_kernel}
    \lean{MPSTensor.physicalReassocES_map_submodule,
        MPSTensor.range_reassocTailBoundaryMapES_eq_ker_openSuffixParentHamiltonianES}
    \leanok
    Reassociation identifies any subspace of the physical space on
    $K+(L+Q)$ sites with the same subspace on $(K+L)+Q$ sites.
    In particular, when $L+Q>0$, the kernel of the single parent interaction
    on the last $L+Q$ sites is
    $(\C^d)^{\otimes K}\otimes G_{L+Q}(A)$ in these common coordinates.
\end{lemma}
\begin{proof}
    \leanok
    \uses{thm:fixed-volume-c3-projector-conjugacy,
        thm:spectator_boundary_hilbert_factorization}
    Let $\sigma$ denote the reassociation isometry from the physical space
    on $K+(L+Q)$ sites to that on $(K+L)+Q$ sites. Because $\sigma$ preserves
    the ordered physical coordinates, every family $(S_n)_n$ of subspaces of
    the $n$-site physical spaces satisfies
    \begin{align}
        \sigma\bigl(S_{K+(L+Q)}\bigr) & = S_{(K+L)+Q}.
    \end{align}
    A parent interaction is the orthogonal projection onto the complement
    of the local boundary-condition space, acting trivially on the other
    sites. Its kernel on the last $L+Q$ sites is therefore
    $(\C^d)^{\otimes K}\otimes G_{L+Q}(A)$, and the display carries this
    description to the common coordinates.
\end{proof}

\begin{lemma}[Terminal interval kernels]
    \label{thm:ph_terminal_interval_kernel}
    \lean{MPSTensor.ker_openSuffixParentHamiltonianES_eq_of_local_kernel,
        MPSTensor.ker_openSuffixParentHamiltonianES_eq_range_reassocTailBoundaryMapES}
    \leanok
    Let $0<R\leq W\leq n\leq N$, and suppose that the range-$R$ open parent
    Hamiltonian on $W$ sites has kernel $G_W(A)$. On the $N$-site chain,
    the range-$R$ interaction sum on $[n-W,n)$ and the single parent
    interaction on the same interval have the same kernel. In particular,
    for $W=L+Q$ and a chain of length $K+L+Q$, the kernel of the terminal
    range-$R$ interaction sum coincides with the boundary space
    $(\C^d)^{\otimes K}\otimes G_{L+Q}(A)$ of condition (C3$'$) in
    \cite{Nachtergaele1996Spectral}.
\end{lemma}
\begin{proof}
    \leanok
    \uses{thm:cyclic_active_block_local_conjugacy,
        thm:ph_reassociated_suffix_kernel}
    The physical coordinate isometry identifies both interval operators
    with the corresponding $W$-site operators tensored with the identity on
    the remaining sites. Tensoring with the identity preserves equality of
    kernels, and the $W$-site kernels are both $G_W(A)$. The terminal case
    then follows from the reassociated description of the kernel of the
    single parent interaction on the last $L+Q$ sites.
\end{proof}

\begin{lemma}[Active grouped projector error]
    \label{thm:ph_active_grouped_projector_error}
    \lean{MPSTensor.norm_suffixProjection_comp_prefixDifference_eq_projector_defect}
    \leanok
    Let $R=L+Q>0$ and suppose that the range-$R$ open parent Hamiltonians
    on $K+L$ and $K+L+Q$ sites have their boundary-condition spaces as kernels.
    On $K+L+Q$ sites, let $P$ project onto the ground space of the first
    $K+L$ sites, let $T$ project onto the ground space of the last $L+Q$
    sites, and let $F$ project onto the whole-chain ground space. Then the
    grouped projection product matches the three-interval projector error of
    \cite[condition C3$'$]{Nachtergaele1996Spectral}:
    \begin{align}
        \|T(P-F)\|=\|F-PT\|.
        \notag
    \end{align}
\end{lemma}
\begin{proof}
    \leanok
    \uses{thm:ph_prefix_kernel_free_right_sites,
        thm:ph_reassociated_suffix_kernel,
        thm:fixed-volume-c3-projector-conjugacy, thm:range_ground_space_map_es,
        thm:reassociated_tail_boundary_factorization}
    The whole-chain boundary space lies in the suffix boundary space,
    so $TF=F$. Thus $T(P-F)=TP-F$. Taking adjoints, using that orthogonal
    projections are self-adjoint, and reversing the sign gives the identity.
\end{proof}

\begin{lemma}[Periodic gaps from divisible open intervals]
    \label{thm:periodic_gap_from_divisible_open}
    \lean{MPSTensor.exists_parentHamiltonianES_gap_of_eventually_divisible_openParentHamiltonianES_gap}
    \leanok
    \uses{def:open_parent_hamiltonian_es, def:parent_hamiltonian_es}
    Let $A$ be a tensor, let $p,R\geq 1$ be integers, and let $\gamma>0$.
    Suppose that, for every sufficiently large integer $M$ and every
    $u\perp\ker H^{\mathrm{open}}_{pM,R}(A)$,
    $\gamma\lVert u\rVert\leq\lVert H^{\mathrm{open}}_{pM,R}(A)u\rVert$.
    Then there is $\delta>0$ such that, for every sufficiently large $N$
    and every $v\perp\ker H^{\mathrm{ES}}_{N,R}(A)$,
    $\delta\lVert v\rVert\leq\lVert H^{\mathrm{ES}}_{N,R}(A)v\rVert$.
\end{lemma}
\begin{proof}
    \leanok
    \uses{thm:parent_hamiltonian_gap_from_open_knabe_block}
    Choosing $n$ with $(R-1)^2<n\gamma$ and $M\geq n+2R$ for which the
    open-chain bound holds, and setting $m=pM-R+1$, gives $m\geq\max\{n,R\}$
    and active open volume $m+R-1=pM$. The finite-range Knabe comparison then
    yields the positive constant $\delta=(m\gamma-(R-1)^2)/(m-R+1)$ and the
    asserted bound for every $N\geq2m$.
\end{proof}

\begin{lemma}[Spectator extension of a grouped projection difference]
    \label{thm:grouped_spectator_norm}
    \lean{MPSTensor.norm_suffixGroundProjection_comp_prefixDifference_le_active}
    \leanok
    \uses{def:fixed-volume-prefix-ground-projections,
        def:open_suffix_parent_hamiltonian_es}
    Let the interaction range satisfy $R>0$, let $p\leq n$, and add
    $r\geq0$ sites to the right of an active interval of length $n$.
    In each ambient volume $s\in\{n,n+r\}$, write $P^{(s)}_a$ for the
    kernel projection of the prefix Hamiltonian ending at $a$, and write
    $Q^{(s)}$ for the kernel projection of the suffix Hamiltonian of length
    $l$ ending at $n$. Then
    \begin{align}
        \bigl\lVert Q^{(n+r)}(P^{(n+r)}_p-P^{(n+r)}_n)\bigr\rVert
         & \leq
        \bigl\lVert Q^{(n)}(P^{(n)}_p-P^{(n)}_n)\bigr\rVert.
    \end{align}
\end{lemma}
\begin{proof}
    \leanok
    \uses{thm:open-hamiltonian-right-spectator-conjugacy,
        lem:right-spectator-kernel-projection, lem:right-spectator-extension}
    Let $\Phi$ be the isometry $\ell^2([d]^{n+r})\to\ell^2([d]^n\times[d]^r)$
    splitting a configuration into its active and spectator coordinates, and
    write $T^{\oplus[d]^r}$ for the fiberwise extension of an operator $T$
    on $n$ sites. By
    Lemma~\ref{thm:open-hamiltonian-right-spectator-conjugacy}, $\Phi$
    conjugates each of the three Hamiltonians to its fiberwise extension.
    Lemma~\ref{lem:right-spectator-kernel-projection} transfers this to
    their kernel projections: for $a\in\{p,n\}$,
    \begin{align}
        \Phi P^{(n+r)}_a\Phi^{-1} & =\bigl(P^{(n)}_a\bigr)^{\oplus[d]^r}, &
        \Phi Q^{(n+r)}\Phi^{-1}   & =\bigl(Q^{(n)}\bigr)^{\oplus[d]^r}.
    \end{align}
    Since fiberwise extension respects composition and subtraction
    (Lemma~\ref{lem:right-spectator-extension}),
    \begin{align}
        \Phi\,Q^{(n+r)}\bigl(P^{(n+r)}_p-P^{(n+r)}_n\bigr)\Phi^{-1}
         & =\Bigl(Q^{(n)}\bigl(P^{(n)}_p-P^{(n)}_n\bigr)\Bigr)^{\oplus[d]^r}.
    \end{align}
    Conjugation by $\Phi$ preserves the operator norm, and
    $\lVert T^{\oplus[d]^r}\rVert\leq\lVert T\rVert$ by
    Lemma~\ref{lem:right-spectator-extension}.
\end{proof}

\begin{lemma}[Grouped C3 from a three-interval estimate]
    \label{thm:grouped_c3_from_projector_defect}
    \lean{MPSTensor.grouped_martingaleDifference_norm_le_of_projector_defect}
    \leanok
    Let $p\geq1$. Suppose that the range-$2p$ open parent
    Hamiltonian of a tensor has its boundary-condition space as kernel
    at every length at least $2p$. For every $K\geq0$, let
    $P_K$, $T_K$, and $F_K$ be the boundary-space projections
    on the first $K+p$ sites, the last $2p$ sites, and all
    $K+2p$ sites, respectively, with free sites elsewhere.
    Suppose that $\epsilon\geq0$ and that, for every $K\geq0$,
    \begin{align}
        \|F_K-P_KT_K\|\leq\epsilon.
        \notag
    \end{align}
    Then on every chain of length $pM$ with $M\geq2$, the grouped
    martingale projections satisfy
    \begin{align}
        \|Q_nE_n\|\leq\epsilon
        \notag
    \end{align}
    for every $0\leq n<M$.
\end{lemma}
\begin{proof}
    \leanok
    \uses{lem:grouped_open_ground_projection_locality,
        thm:grouped_c3_from_second_index,
        thm:ph_active_grouped_projector_error,thm:grouped_spectator_norm}
    The products at indices zero and one vanish, so by
    Lemma~\ref{thm:grouped_c3_from_second_index} it suffices to treat
    $n\geq2$. For such $n$, setting $K=p(n-1)$ gives an active interval of
    length $p(n+1)$ whose preceding prefix has length $pn$. The kernel
    hypothesis and the active projector identity yield the desired bound on
    this interval, and adjoining the remaining $pM-p(n+1)$ sites does not
    increase the operator norm.
\end{proof}

\begin{lemma}[Open-chain gaps on divisible lengths for primitive block sums]
    \label{thm:primitive_block_sum_open_gap_subsequence}
    \lean{MPSTensor.%
        exists_ge_openParentHamiltonianES_toTensorFromBlocks_gap_of_isPrimitiveMPS}
    \leanok
    Let $A_0,\ldots,A_{r-1}$ be normalized primitive tensors with
    common physical dimension $d\geq1$, positive bond dimensions,
    and positive-definite invariant matrices. Assume that no two distinct
    blocks are related by a virtual similarity and a phase. Let
    $B=\bigoplus_{j=0}^{r-1}\mu_jA_j$, with every $\mu_j\neq0$.
    By the grouped martingale estimate of
    \cite[Theorem~2.1(ii) and Section~6]{Nachtergaele1996Spectral},
    for every $R_*\in\mathbb N$ there are $p\geq\max\{1,R_*\}$ and
    $\gamma>0$ such that
    \begin{align}
        \gamma\|v\|\leq
        \|H^{\mathrm{open}}_{pM,2p}(B)v\|\notag
    \end{align}
    for every $M\geq2$ and every
    $v\perp\ker H^{\mathrm{open}}_{pM,2p}(B)$.
\end{lemma}
\begin{proof}
    \leanok
    \uses{thm:ph_primitive_blocks_open_kernel,
        thm:ph_block_diagonal_projector_decay,
        thm:grouped_c3_from_projector_defect,thm:grouped_open_parent_gap}
    Choosing $p$ above $R_*$ and the open-kernel threshold, and sufficiently
    large that the three-interval projector error is at most $1/2$ for
    every pair of outer lengths at least $p$, ensures that the grouped C3
    estimate holds with $\epsilon=1/2<1/\sqrt2$. The grouped open-chain
    bound then gives $\gamma=(1-\sqrt2/2)^2>0$ at every length $pM$, $M\geq2$.
\end{proof}

\begin{lemma}[Open-chain gaps at every length for primitive block sums]
    \label{thm:primitive_block_sum_open_gap_all_lengths}
    \lean{MPSTensor.%
        exists_ge_openParentHamiltonianES_toTensorFromBlocks_gap_all_of_isPrimitiveMPS}
    \leanok
    Under the hypotheses of
    Lemma~\ref{thm:primitive_block_sum_open_gap_subsequence}, as in the
    open-chain form of \cite[Theorem~1.2]{Nachtergaele1996Spectral}, for every
    $R_*\in\mathbb N$ there are $p\geq\max\{1,R_*\}$ and $\delta>0$ such that
    \begin{align}
        \delta\|v\|\leq
        \|H^{\mathrm{open}}_{N,2p}(B)v\|\notag
    \end{align}
    for every chain length $N$ and every
    $v\perp\ker H^{\mathrm{open}}_{N,2p}(B)$ (for $N<2p$ the Hamiltonian is
    zero and the assertion is vacuous).
\end{lemma}
\begin{proof}
    \leanok
    \uses{thm:primitive_block_sum_open_gap_subsequence,
        thm:ph_primitive_blocks_open_kernel,
        thm:ph_block_diagonal_projector_decay,
        thm:ph_prefix_kernel_free_right_sites,
        thm:ph_terminal_interval_kernel,
        thm:open_interval_gap_transport,
        thm:open_parent_overlapping_interval_gap}
    Fix $p$ as in
    Lemma~\ref{thm:primitive_block_sum_open_gap_subsequence}, above the
    open-kernel threshold and so large that the three-interval projector
    error is at most $1/2$, and write $N=pM+q$ with $0\leq q<p$ and $M\geq3$.
    If $q=0$ the grouped bound applies. Otherwise, the prefix of length
    $pM$ carries the grouped gap $\gamma$, and the terminal interval of
    length $2p+q$ carries a gap $\eta>0$, the minimum of the finitely many
    positive gaps at lengths below $3p$. The two intervals overlap in $2p$
    sites, their kernels are the boundary spaces of the three-interval
    estimate, and combining the interval gaps gives the bound
    $\min\{\gamma,\eta\}/16$ at every $N\geq3p$. The finitely many
    remaining lengths have positive gaps on their kernel complements, so a
    uniform $\delta>0$ exists.
\end{proof}

\begin{lemma}[Uniformly gapped interaction ranges for primitive block sums]
    \label{thm:primitive_block_sum_parent_gap}
    \lean{MPSTensor.exists_ge_parentHamiltonianES_toTensorFromBlocks_gap_of_isPrimitiveMPS,
        MPSTensor.exists_parentHamiltonianES_toTensorFromBlocks_gap_of_isPrimitiveMPS}
    \leanok
    Let $A_0,\ldots,A_{r-1}$ be normalized primitive tensors with
    common physical dimension $d\geq1$, positive bond dimensions,
    and positive-definite invariant matrices. Assume that no two blocks
    are related by a virtual similarity and a phase. Let
    $\mu_j\ne0$ for every $j$, and set
    $B=\bigoplus_{j=0}^{r-1}\mu_jA_j$.
    For every prescribed lower bound $R_*$, there exist
    $R\geq\max\{1,R_*\}$, $\gamma>0$, and $N_0$ such that
    the periodic range-$R$ parent Hamiltonian satisfies
    \begin{align}
        \gamma\|v\|\leq\|H_{R,N}v\|
        \notag
    \end{align}
    for every $N\geq N_0$ and every $v\perp\ker H_{R,N}$.
    In particular, the block sum admits a parent interaction with a
    uniform positive periodic gap.
\end{lemma}
\begin{proof}
    \leanok
    \uses{thm:primitive_block_sum_open_gap_subsequence,
        thm:periodic_gap_from_divisible_open}
    Choosing $p\geq\max\{1,R_*\}$ by
    Lemma~\ref{thm:primitive_block_sum_open_gap_subsequence} and setting
    $R=2p$, the uniform open-chain gap on lengths $pM$, $M\geq2$, yields a
    positive periodic gap at every sufficiently large length by
    Lemma~\ref{thm:periodic_gap_from_divisible_open}.
\end{proof}

\begin{lemma}[Operator order after adjoining free sites]
    \label{thm:ph_spectator_operator_order}
    \lean{ContinuousLinearMap.inner_rightFiberwiseMap,
        ContinuousLinearMap.isPositive_rightFiberwiseMap,
        ContinuousLinearMap.rightFiberwiseMap_mono,
        ContinuousLinearMap.rightFiberwiseMap_smul,
        LinearIsometryEquiv.conj_le_conj_iff}
    \leanok
    Let $T$ be an operator on a finite-dimensional physical space and
    let $(e_s)_s$ be an orthonormal basis for an additional physical space.
    For $x=\sum_s x_s\otimes e_s$ and $y=\sum_s y_s\otimes e_s$,
    \begin{align}
        \langle (T\otimes\Id)x,y\rangle
        =\sum_s\langle Tx_s,y_s\rangle.
        \notag
    \end{align}
    Consequently $T\ge0$ implies $T\otimes\Id\ge0$, and $T\le U$ implies
    $T\otimes\Id\le U\otimes\Id$. Scalar multiplication commutes with
    extension: $(cT)\otimes\Id=c(T\otimes\Id)$.
    If $V$ is a unitary identification of physical spaces, then
    $VTV^{-1}\le VUV^{-1}$ if and only if $T\le U$.
\end{lemma}
\begin{proof}
    \leanok
    Expanding the inner product in the product orthonormal basis and summing
    first over the additional coordinates establishes the inner-product
    identity. Positivity follows term by term, and applying positivity to
    $U-T$ gives the order assertion. The scalar identity holds on each
    coordinate, and unitary conjugation preserves and reflects positivity and
    hence operator order.
\end{proof}

\begin{lemma}[Comparison on a finite interval]
    \label{lem:parent-finite-interval-comparison}
    \lean{MPSTensor.exists_pos_parentInteractionES_openParentHamiltonianES_comparison}
    \leanok
    \uses{def:open_parent_hamiltonian_es}
    Let $A$ be an MPS tensor, and let $R,W$ be interaction and interval lengths.
    Write $K_{R,W}=H^{\mathrm{open}}_{W,R}(A)$ for the open parent Hamiltonian
    of interaction length $R$ on $W$ sites, and suppose that its kernel is
    $\mathcal G_W(A)$. Writing $h_W$ for the orthogonal projection onto
    $\mathcal G_W(A)^\perp$, there are positive constants $\kappa,C$ such that
    \begin{align}
        \kappa h_W \leq K_{R,W} \leq C h_W.
        \notag
    \end{align}
\end{lemma}
\begin{proof}
    \leanok
    \uses{lem:positive-kernel-projection-comparison}
    The operator $K_{R,W}$ is a sum of orthogonal projections, hence
    positive. By hypothesis its kernel is $\mathcal G_W(A)$, so the
    orthogonal projection onto $(\ker K_{R,W})^\perp$ is $h_W$.
    Lemma~\ref{lem:positive-kernel-projection-comparison} applied to
    $K_{R,W}$ gives the two bounds.
\end{proof}

\begin{lemma}[Comparison with the kernel projection]
    \label{lem:positive-kernel-projection-comparison}
    \lean{LinearMap.exists_pos_mul_norm_le_of_mem_orthogonal_ker,
        LinearMap.IsPositive.smul_orthogonal_ker_projection_le_of_norm_gap,
        LinearMap.IsPositive.exists_pos_smul_orthogonal_ker_projection_le,
        LinearMap.IsPositive.exists_pos_le_smul_orthogonal_ker_projection}
    \leanok
    Let $H$ be a positive operator on a finite-dimensional complex Hilbert
    space. If $P$ is the orthogonal projection onto $(\ker H)^\perp$,
    then $aP\leq H\leq bP$ for some $a,b>0$. More precisely, a norm gap
    $\gamma\geq0$ on $(\ker H)^\perp$ implies $\gamma P\leq H$.
\end{lemma}
\begin{proof}\leanok
    The restriction of $H$ to $(\ker H)^\perp$ is injective and thus has
    a positive norm lower bound. Since $H$ is positive, its eigenvalues
    on this subspace have the same lower bound. The operator norm gives
    the upper bound. Orthogonal projection onto $(\ker H)^\perp$
    leaves the quadratic form of $H$ unchanged.
\end{proof}

\begin{lemma}[Kernels and gaps under positive comparison]
    \label{lem:positive-comparison-kernel-gap}
    \lean{LinearMap.IsPositive.ker_le_of_le, LinearMap.IsPositive.ker_eq_of_smul_le_of_le_smul, LinearMap.IsPositive.norm_gap_of_smul_le_of_le_smul}
    \leanok
    If $0\leq P\leq Q$, then $\ker Q\subseteq\ker P$.
    Consequently, if $P,Q$ are positive operators and $a,b,c>0$ satisfy
    \begin{align}
        aP\leq bQ\leq cP,
        \notag
    \end{align}
    their kernels agree. In finite dimension, a norm gap $\gamma>0$
    for $P$ on its kernel complement implies a norm gap
    $a\gamma/b$ for $Q$ on the same subspace.
\end{lemma}
\begin{proof}\leanok
    A positive quadratic form vanishes only on the kernel of its operator.
    Evaluating the order bounds on kernel vectors proves the inclusions.
    On the common kernel complement, the norm gap for $P$ is equivalent
    to $P\geq\gamma I$. Hence $bQ\geq a\gamma I$, which gives the
    stated norm gap by Cauchy--Schwarz.
\end{proof}

\begin{lemma}[Transfer to a prescribed interaction range]
    \label{thm:parent-gap-fixed-range-transfer}
    \lean{MPSTensor.exists_parentHamiltonianES_gap_of_larger_range}
    \leanok
    \uses{def:open_parent_hamiltonian_es, def:ground_space_es,
        def:parent_hamiltonian_es}
    Let $A$ be an MPS tensor, and let $R,W$ be positive integers with
    $W\geq2R$. Suppose that the open range-$R$ Hamiltonian on $W$
    sites has kernel $\mathcal G_W(A)$. If the periodic parent Hamiltonians
    of range $W$ have a uniform positive gap for all sufficiently large
    chain lengths, then the same holds at range $R$.
\end{lemma}
\begin{proof}
    \leanok
    \uses{lem:parent-finite-interval-comparison,
        thm:periodic-parent-range-comparison,
        lem:positive-comparison-kernel-gap}
    Lemma~\ref{lem:parent-finite-interval-comparison} provides
    $\kappa,C>0$ such that
    $\kappa h_W\leq K_{R,W}\leq Ch_W$, and summing all cyclic translates by
    Lemma~\ref{thm:periodic-parent-range-comparison} gives
    \begin{align}
        \kappa H_{W,N}\leq (W-R+1)H_{R,N}\leq C H_{W,N}
        \notag
    \end{align}
    for sufficiently large $N$. By
    Lemma~\ref{lem:positive-comparison-kernel-gap}, the periodic kernels
    therefore agree, and a gap $\gamma$ at range $W$ transfers to the gap
    $\kappa\gamma/(W-R+1)$ at range $R$.
\end{proof}

\begin{lemma}[The interaction on a full interval]
    \label{lem:full-interval-parent-interaction}
    \lean{MPSTensor.localTermES_self_zero_eq_parentInteractionES}
    \leanok
    For an MPS tensor $A$ and an interval of positive length $W$,
    the length-$W$ parent interaction starting at the first site is the
    orthogonal projection $h_W$ onto $\mathcal G_W(A)^\perp$.
\end{lemma}
\begin{proof}\leanok
    Restriction to the full interval is the identity, so the local interaction
    is precisely the defining orthogonal projection.
\end{proof}

\begin{lemma}[Periodic comparison from an open interval]
    \label{thm:periodic-parent-range-comparison}
    \lean{MPSTensor.parentHamiltonianES_comparison_of_local_open_comparison}
    \leanok
    Let $1\leq R\leq m$, $N\geq 2m$, and $W=m+R-1$.
    Write $K_{R,W}$ for the open parent Hamiltonian of interaction length
    $R$ on $W$ sites, and $H_{R,N}$ for the corresponding periodic
    Hamiltonian on $N$ sites. If real constants $\kappa,C$ satisfy
    \begin{align}
        \kappa h_W\leq K_{R,W}\leq C h_W,
        \notag
    \end{align}
    where $h_W$ projects onto $\mathcal G_W(A)^\perp$, then
    \begin{align}
        \kappa H_{W,N}\leq mH_{R,N}\leq C H_{W,N}.
        \notag
    \end{align}
\end{lemma}
\begin{proof}\leanok
    \uses{lem:full-interval-parent-interaction,
        thm:cyclic_active_block_local_conjugacy,
        thm:cyclic_knabe_block_open_parent_hamiltonian}
    Extending each local inequality by the identity on the complementary sites
    and conjugating by the cyclic identification of the chosen interval
    preserves operator order. Summing over all cyclic origins then gives the
    asserted bounds, since each range-$R$ interaction occurs
    exactly $m$ times, once for each of its possible positions in the
    open interval.
\end{proof}

\begin{lemma}[Uniform gap above an interaction threshold]
    \label{thm:primitive_block_sum_parent_gap_threshold}
    \lean{MPSTensor.exists_parentHamiltonianES_toTensorFromBlocks_gap_threshold_of_isPrimitiveMPS}
    \leanok
    Let $B=\bigoplus_j\mu_j A_j$, where the coefficients are nonzero and
    $(A_j)_j$ is a finite family of normalized primitive tensors of positive
    bond dimension, pairwise inequivalent under similarity and phase. Assume that
    each recorded invariant density is positive definite and that the physical
    dimension is positive. There exists $R_0\geq1$ such that, for every
    $R\geq R_0$, there exist $\gamma_R>0$ and $N_R$ for which
    \begin{align}
        \gamma_R\|v\|\leq\|H_{R,N}(B)v\|
        \notag
    \end{align}
    for every $N\geq N_R$ and every $v\perp\ker H_{R,N}(B)$.
\end{lemma}
\begin{proof}\leanok
    \uses{thm:primitive_block_sum_parent_gap,thm:ph_primitive_blocks_open_kernel,thm:parent-gap-fixed-range-transfer}
    Let $R_0$ be large enough that every open parent Hamiltonian of range at
    least $R_0$ has the expected local ground space. For any fixed $R\geq R_0$
    and any uniformly gapped range $W\geq2R$, the range-$R$ Hamiltonian on $W$
    sites and the orthogonal projection onto the complement of the local
    MPS space have the same kernel. Finite dimensionality gives two-sided
    positive order bounds. Summing their cyclic translates gives two-sided
    bounds between the periodic Hamiltonians of ranges $R$ and $W$.
    These bounds identify their kernels and transfer the positive gap to range
    $R$.
\end{proof}

\begin{lemma}[Gap at the common block-injectivity range]
    \label{thm:primitive-block-parent-gap-three-block-bound}
    \lean{MPSTensor.exists_parentHamiltonianES_toTensorFromBlocks_gap_of_isPrimitiveMPS_of_threeBlock_bound}
    \leanok
    Let $B=\bigoplus_{j=0}^{r-1}\mu_j A_j$, where $r\geq2$, all coefficients
    are nonzero, and the normalized primitive tensors $A_j$ are pairwise
    inequivalent under similarity and phase. Assume positive physical and bond
    dimensions, positive definite recorded invariant densities, and a common
    block-injectivity length $L_0\geq1$. Removing the restriction to the
    range constructed by the overlap argument recorded
    in~\cite{gap:cpgsv21_block_parent_interaction_range}, for every
    $R\geq3(r-1)(L_0+1)+1$ there exist $\gamma_R>0$ and $N_R$ such that
    \begin{align}
        \gamma_R\|v\|\leq\|H_{R,N}(B)v\|
        \notag
    \end{align}
    for every $N\geq N_R$ and every $v\perp\ker H_{R,N}(B)$.
\end{lemma}
\begin{proof}\leanok
    \uses{thm:primitive_block_sum_parent_gap,thm:parent-gap-fixed-range-transfer,
        thm:ph_primitive_blocks_sharp_open_kernel}
    The common block-injectivity hypothesis identifies the kernel of the
    open-chain range-$R$ Hamiltonian on every interval of length at least
    $R$ with the local MPS space. For a uniformly gapped range
    $W\geq2R$, the two-sided comparison on the interval of length $W$
    and its cyclic translates transfers the gap from range $W$ to range
    $R$.
\end{proof}

\begin{lemma}[Gap uniform in every chain length]
    \label{thm:primitive-block-parent-uniform-gap-three-block-bound}
    \lean{MPSTensor.exists_parentHamiltonianES_toTensorFromBlocks_uniform_gap_of_threeBlock_bound}
    \leanok
    Let $B=\bigoplus_{j=0}^{r-1}\mu_j A_j$, where $r\geq2$, all coefficients
    are nonzero, and the normalized primitive tensors $A_j$ are pairwise
    inequivalent under similarity and phase. Assume positive physical and bond
    dimensions, positive definite recorded invariant densities, and a common
    block-injectivity length $L_0\geq1$. For every
    $R\geq3(r-1)(L_0+1)+1$, there exists $\gamma_R>0$ such that, for every
    $N\in\mathbb N$ and every $v\perp\ker H_{R,N}(B)$,
    \begin{align}
        \gamma_R\|v\|\leq\|H_{R,N}(B)v\|.
    \end{align}
\end{lemma}
\begin{proof}\leanok
    \uses{thm:primitive-block-parent-gap-three-block-bound,thm:parent_hamiltonian_eventual_gap}
    The eventual uniform gap leaves only finitely many chain lengths untreated.
    At each such length the Hamiltonian, restricted to its kernel complement,
    has a positive lower norm bound. Taking the minimum of these bounds and the
    eventual uniform constant proves the assertion.
\end{proof}

\begin{lemma}[Nonzero eigenvalues of the transfer remainder]
    \label{thm:fnw-transfer-eigenvalue-from-remainder}
    \lean{MPSTensor.IsPrimitiveMPS.fnwTransfer_hasEigenvalue_of_remainder}
    \leanok
    \uses{def:primitive_mps, def:fnw_transfer_map, def:fnw_limit_map}
    Let $A$ be primitive with nonzero bond dimension and fixed point $\rho$.
    Write $E=E_{\mathrm{FNW},A}$ for its observable transfer map and
    $P=E_\infty^\rho$ for its stationary projection. Every nonzero
    eigenvalue of $E-P$ is an eigenvalue of $E$.
\end{lemma}
\begin{proof}\leanok
    \uses{thm:fnw_limit_map_projection,thm:fnw_limit_absorbs_transfer}
    The identities $PE=P$ and $P^2=P$ imply $P(E-P)=0$. If
    $(E-P)X=\nu X$ with $X\neq0$ and $\nu\neq0$, applying $P$ gives
    $PX=0$, hence $EX=\nu X$.
\end{proof}

\begin{lemma}[FNW decay rates from transfer eigenvalues]
    \label{thm:fnw-transfer-eigenvalue-rate}
    \lean{MPSTensor.IsPrimitiveMPS.fnwWeightedRemainderSpectralRadius_lt_of_transfer_eigenvalues}
    \lean{MPSTensor.IsPrimitiveMPS.exists_fnwMixingQuantity_le_geometric_of_transfer_eigenvalues}
    \leanok
    \uses{def:primitive_mps, def:fnw_transfer_map,
        def:fnw_weighted_remainder_data, def:fnw_mixing_quantity}
    Let $A$ be primitive with nonzero bond dimension and positive-definite
    fixed point $\rho$. Suppose $\lambda>0$ and every eigenvalue
    $\nu\neq1$ of its observable transfer map satisfies $|\nu|<\lambda$.
    Then the spectral radius of the transfer remainder in the $\rho$-weighted
    norm is less than $\lambda$. Following the eigenvalue prescription of
    \cite[Section~6, equation~(6.1)]{Nachtergaele1996Spectral}, if
    $\tr\rho=1$, there is a constant $C_\lambda>0$ such that, for every
    integer $n\geq1$,
    \begin{align}
        a(n)\leq C_\lambda\lambda^n.
    \end{align}
    % The dimension-only constant in the source projector estimate remains separate.
\end{lemma}
\begin{proof}\leanok
    \uses{thm:fnw-transfer-eigenvalue-from-remainder,thm:fnw_remainder_eigenvalue_gap,thm:fnw_mixing_quantity_geometric}
    A nonzero remainder eigenvalue is a transfer eigenvalue, and primitivity
    excludes the eigenvalue $1$. The zero eigenvalue also has modulus less than
    $\lambda$. In finite dimension these strict eigenvalue bounds give a strict
    spectral-radius bound, to which the prescribed-rate mixing estimate applies.
\end{proof}

\begin{lemma}[Prescribed FNW prefactors after an overlap threshold]
    \label{thm:fnw-eventual-prescribed-prefactor}
    \lean{MPSTensor.IsPrimitiveMPS.eventually_fnwMixingQuantity_le_of_pos_prefactor}
    \leanok
    \uses{def:primitive_mps, def:fnw_transfer_map, def:fnw_mixing_quantity}
    Let $A$ be primitive with nonzero bond dimension and a positive-definite
    trace-one fixed point $\rho$. Suppose $\lambda>0$ exceeds the modulus
    of every nonunit transfer eigenvalue. For every $c>0$, there is an integer
    $M$ such that, for every $n\geq M$,
    \begin{align}
        a(n)\leq c\lambda^n.
    \end{align}
    The onset $M$ may depend on the state, rate, and prescribed prefactor.
\end{lemma}
\begin{proof}\leanok
    \uses{thm:fnw-transfer-eigenvalue-rate,thm:fnw_mixing_quantity_geometric}
    For an intermediate rate $\mu$ strictly between the remainder spectral
    radius and $\lambda$, the mixing estimate gives $a(n)\leq C_\mu\mu^n$.
    Since $(\mu/\lambda)^n$ tends to zero, eventually
    $C_\mu(\mu/\lambda)^n\leq c$.
\end{proof}

\begin{lemma}[Eventual projector bound with a prescribed prefactor]
    \label{thm:fnw-eventual-projector-prescribed-prefactor}
    \lean{MPSTensor.IsPrimitiveMPS.eventually_projector_defect_le_of_pos_prefactor}
    \leanok
    \uses{def:primitive_mps, def:fnw_transfer_map,
        def:reassociated_tail_boundary_map, def:spectator_boundary_maps_es,
        def:ground_space_es}
    Under the preceding hypotheses, suppose additionally that $\lambda<1$.
    In the version of \cite[Section~6, equation~(6.1)]{Nachtergaele1996Spectral}
    with the additional onset stated
    explicitly~\cite{gap:cpgsv21_martingale_overlap}, for every $c>0$ there
    exists an onset $M$ (depending on the state, rate, and prefactor, and not
    asserted to be the injectivity threshold) such that, for every overlap
    $m\geq M$, every prefix length $r\geq0$, and every suffix length $\ell\geq1$,
    we have $c\lambda^m<1$ and
    \begin{align}
        \left\|P^{\mathrm{tail}}_{r;m,\ell}P^{\mathrm{left}}_{r+m;\ell}
        -G_{r+m+\ell}\right\|
        \leq \frac{c\lambda^m(1+c\lambda^m)}{1-c\lambda^m}.
    \end{align}
\end{lemma}
\begin{proof}\leanok
    \uses{thm:fnw-eventual-prescribed-prefactor,thm:fnw_projector_defect_geometric,
        thm:normal_from_primitive_posdef}
    Once the onset is large enough that the overlap exceeds an injectivity
    length and $c\lambda^m<1$, the preceding mixing estimate satisfies the
    hypotheses of the fixed-overlap FNW projector estimate, uniformly in the
    two exterior lengths.
\end{proof}

\subsection{The dimension-only prefactor at a prescribed rate}

\begin{definition}[Weighted matrix-unit path vectors]
    \label{def:weighted-matrix-unit-path-vectors}
    \lean{MPSTensor.FNWDimensionConstant.matrixUnitTensor,MPSTensor.FNWDimensionConstant.matrixUnitTensor_apply,MPSTensor.FNWDimensionConstant.pathConfig,MPSTensor.FNWDimensionConstant.pathVector}
    \leanok
    \uses{def:ground_space_map}
    Let $t=(t_{ij})$ and $z=(z_{ij})$ be real matrices indexed by a finite
    set of $k$ vertices. Set $A^{ij}=t_{ij}E_{ij}$. For endpoints $a,f$,
    define the four-site vector
    \begin{align}
        v_{a,f}(z)
         & =\sum_{b,c,e}z_{be}t_{ab}t_{bc}t_{ce}t_{ef}
        |ab,bc,ce,ef\rangle. \notag
    \end{align}
    Write $\mathbf 1$ for the mark with all entries one and, for a real
    function $u$ on the vertices, $u\otimes\mathbf 1$ and $\mathbf 1\otimes u$
    for the marks $z_{be}=u_b$ and $z_{be}=u_e$. For $X\in\MN{k}$, the
    four-site boundary vector of $A$ has the coordinates
    \begin{align}
        \Gamma_4(X)(i_0j_0,i_1j_1,i_2j_2,i_3j_3)
         & =\operatorname{tr}\bigl(A^{i_0j_0}A^{i_1j_1}A^{i_2j_2}A^{i_3j_3}X\bigr).
        \notag
    \end{align}
\end{definition}

\begin{lemma}[Path contractions]
    \label{thm:weighted-matrix-unit-path-contractions}
    \lean{MPSTensor.FNWDimensionConstant.pathConfig_eq_iff,MPSTensor.FNWDimensionConstant.inner_pathVector,MPSTensor.FNWDimensionConstant.groundSpaceMap_pathConfig,MPSTensor.FNWDimensionConstant.inner_pathVector_groundSpaceMap,MPSTensor.FNWDimensionConstant.pathVector_first_mark_apply,MPSTensor.FNWDimensionConstant.pathVector_last_mark_apply}
    \leanok
    Writing $T_{abcef}=t_{ab}t_{bc}t_{ce}t_{ef}$, one has
    \begin{align}
        \langle v_{a,f}(z),v_{a,f}(z')\rangle
         & =\sum_{b,c,e}z_{be}z'_{be}T_{abcef}^{2}, \notag \\
        \langle v_{a,f}(z),\Gamma_4(X)\rangle
         & =X_{fa}\sum_{b,c,e}z_{be}T_{abcef}^{2}. \notag
    \end{align}
    For every configuration $\sigma=(i_0j_0,i_1j_1,i_2j_2,i_3j_3)$ and every
    real function $u$ on the vertices,
    \begin{align}
        v_{a,f}(u\otimes\mathbf 1)(\sigma)
         & =u_{j_0}\,v_{a,f}(\mathbf 1)(\sigma), \notag \\
        v_{a,f}(\mathbf 1\otimes u)(\sigma)
         & =u_{i_3}\,v_{a,f}(\mathbf 1)(\sigma). \notag
    \end{align}
\end{lemma}
\begin{proof}\leanok
    \uses{def:weighted-matrix-unit-path-vectors}
    Distinct vertex paths $(b,c,e)$ give distinct physical basis vectors
    $|ab,bc,ce,ef\rangle$, which gives the first identity. For the second,
    multiplying four consecutive letters yields
    \begin{align}
        A^{ab}A^{b'c}A^{c'e}A^{e'f}
         & =\delta_{bb'}\delta_{cc'}\delta_{ee'}\,T_{abcef}\,E_{af}, \notag \\
        \operatorname{tr}(E_{af}X)
         & =X_{fa}. \notag
    \end{align}
    Hence $\Gamma_4(X)(ab,bc,ce,ef)=T_{abcef}X_{fa}$ on every vertex path.
    On the support of $v_{a,f}$, the first letter $ab$ determines $b=j_0$ and
    the last letter $ef$ determines $e=i_3$, which gives the two mark
    identities.
\end{proof}

\begin{lemma}[Spectator multipliers preserve interval spaces]
    \label{thm:spectator-multiplier-interval-space}
    \lean{MPSTensor.marked_groundSpaceMap_mem_tailBoundaryRange,MPSTensor.marked_groundSpaceMap_mem_leftBoundaryRange}
    \leanok
    \uses{def:ground_space_map, def:spectator_boundary_maps_es}
    Let $K,L\in\N$ and $X\in\MN{D}$. Let $\varphi$ be a complex function of
    the prefix configurations $u\colon\Fin K\to\Fin d$ and $\psi$ a complex
    function of the suffix configurations $\tau\colon\Fin L\to\Fin d$. Then
    the vector $(u,\tau)\mapsto\varphi(u)\,\Gamma_{K+L}(X)(u,\tau)$ belongs to
    the range of $\Gamma^{\mathrm{tail}}_{K,L}$, and the vector
    $(u,\tau)\mapsto\psi(\tau)\,\Gamma_{K+L}(X)(u,\tau)$ belongs to the range
    of $\Gamma^{\mathrm{left}}_{K,L}$.
\end{lemma}
\begin{proof}\leanok
    \uses{def:tail_boundary_map, def:left_boundary_map}
    Factoring the matrix product at the boundary of the interval gives
    \begin{align}
        \varphi(u)\,\Gamma_{K+L}(X)(u,\tau)
         & =\tr\bigl(A^\tau\,\varphi(u)XA^u\bigr)
        =\Gamma^{\mathrm{tail}}_{K,L}(Y)(u,\tau), \notag \\
        \psi(\tau)\,\Gamma_{K+L}(X)(u,\tau)
         & =\tr\bigl(A^u\,\psi(\tau)A^\tau X\bigr)
        =\Gamma^{\mathrm{left}}_{K,L}(Z)(u,\tau), \notag
    \end{align}
    where $Y_u=\varphi(u)XA^u$ and $Z_\tau=\psi(\tau)A^\tau X$.
\end{proof}

\begin{lemma}[Matrix-unit boundary vectors]
    \label{thm:matrix-unit-path-ground-space}
    \lean{MPSTensor.FNWDimensionConstant.pathVector_one_eq_groundSpaceMap}
    \leanok
    \uses{def:weighted-matrix-unit-path-vectors}
    For the weighted matrix-unit tensor $A^{ij}=t_{ij}E_{ij}$ and all
    endpoints $a,f$, one has $v_{a,f}(\mathbf 1)=\Gamma_4(E_{fa})$.
\end{lemma}
\begin{proof}
    \leanok
    \uses{def:weighted-matrix-unit-path-vectors}
    For every configuration,
    \begin{align}
        \operatorname{tr}\bigl(A^{i_0j_0}A^{i_1j_1}A^{i_2j_2}A^{i_3j_3}E_{fa}\bigr)
         & =\delta_{i_0a}\delta_{j_0i_1}\delta_{j_1i_2}\delta_{j_2i_3}\delta_{j_3f}\,
        t_{i_0j_0}t_{i_1j_1}t_{i_2j_2}t_{i_3j_3}. \notag
    \end{align}
    The nonzero coordinates of $\Gamma_4(E_{fa})$ are therefore those of the
    paths $(ab,bc,ce,ef)$, with value $t_{ab}t_{bc}t_{ce}t_{ef}$. These are
    the coordinates of $v_{a,f}(\mathbf 1)$.
\end{proof}

\begin{lemma}[Marked paths in interval spaces]
    \label{thm:fnw-marked-path-interval-spaces}
    \lean{MPSTensor.FNWDimensionConstant.pathVector_first_mark_mem_tailBoundaryRange,MPSTensor.FNWDimensionConstant.pathVector_last_mark_mem_leftBoundaryRange}
    \leanok
    \uses{def:weighted-matrix-unit-path-vectors, def:spectator_boundary_maps_es}
    For every real function $u$ on the vertices and all endpoints $a,f$,
    \begin{align}
        v_{a,f}(u\otimes\mathbf 1)
         & \in\operatorname{range}\Gamma^{\mathrm{tail}}_{1,3}, \notag \\
        v_{a,f}(\mathbf 1\otimes u)
         & \in\operatorname{range}\Gamma^{\mathrm{left}}_{3,1}. \notag
    \end{align}
\end{lemma}
\begin{proof}\leanok
    \uses{thm:weighted-matrix-unit-path-contractions,thm:matrix-unit-path-ground-space,thm:spectator-multiplier-interval-space}
    By Lemmas~\ref{thm:weighted-matrix-unit-path-contractions}
    and~\ref{thm:matrix-unit-path-ground-space}, for
    $\sigma=(i_0j_0,i_1j_1,i_2j_2,i_3j_3)$,
    \begin{align}
        v_{a,f}(u\otimes\mathbf 1)(\sigma)
         & =u_{j_0}\,\Gamma_4(E_{fa})(\sigma), \notag \\
        v_{a,f}(\mathbf 1\otimes u)(\sigma)
         & =u_{i_3}\,\Gamma_4(E_{fa})(\sigma). \notag
    \end{align}
    The first multiplier depends only on the first site and the second only
    on the last site. Lemma~\ref{thm:spectator-multiplier-interval-space}
    with $(K,L)=(1,3)$ and $(K,L)=(3,1)$ gives the two memberships.
\end{proof}

\begin{lemma}[Normalization and orthogonality of marked paths]
    \label{thm:fnw-path-vector-norm-orthogonal}
    \lean{MPSTensor.FNWDimensionConstant.norm_pathVector_eq_one}
    \lean{MPSTensor.FNWDimensionConstant.pathVector_mem_groundSpaceES_orthogonal}
    \leanok
    For a real weighted matrix-unit tensor, fix external vertices $a,f$ and
    write $p_{bce}=t_{ab}t_{bc}t_{ce}t_{ef}$. The path vector with mark $z_{be}$
    has norm one if $\sum_{b,c,e}z_{be}^2p_{bce}^2=1$. It is orthogonal to
    the full four-site ground space if $\sum_{b,c,e}z_{be}p_{bce}^2=0$.
\end{lemma}
\begin{proof}\leanok
    \uses{thm:weighted-matrix-unit-path-contractions}
    Expanding the self-inner product of the path vector and its inner product
    with an arbitrary boundary vector shows that the latter equals the
    indicated scalar sum multiplied by the boundary matrix entry at $(f,a)$.
\end{proof}

\begin{lemma}[A matrix coefficient bound for a projection defect]
    \label{thm:fnw-projection-witness-bound}
    \lean{Submodule.norm_inner_le_norm_projection_defect}
    \lean{Submodule.norm_inner_le_norm_projection_defect_reverse}
    \leanok
    Let $U,V,W$ be subspaces of a complex Hilbert space with orthogonal
    projections $P_U,P_V,P_W$. If $x\in V\cap W^\perp$ and $y\in U$ have
    norm one, then
    \begin{align}
        |\langle y,x\rangle|\leq\|P_UP_V-P_W\|
        =\|P_VP_U-P_W\|.
    \end{align}
\end{lemma}
\begin{proof}\leanok
    Since $P_Vx=x$, $P_Wx=0$, and $P_Uy=y$, the inner product on the left
    equals $\langle y,(P_UP_V-P_W)x\rangle$, so Cauchy--Schwarz and the
    operator norm bound give the inequality; taking adjoints shows that the
    two operator norms agree.
\end{proof}

\begin{lemma}[The dimension-four coefficient at overlap two]
    \label{thm:fnw-dimension-four-coefficient}
    \lean{MPSTensor.FNWDimensionConstant.dimensionFour_coefficient_lt}
    \leanok
    For $c=16$, $\lambda=1/1000$, and $m=2$, one has $c\lambda^m<1$ and
    \begin{align}
        \frac{c\lambda^m(1+c\lambda^m)}{1-c\lambda^m}<\frac1{16}.
    \end{align}
\end{lemma}
\begin{proof}\leanok
    Here $c\lambda^m=1/62500$, and the displayed rational coefficient equals
    $62501/(62500\cdot62499)$.
\end{proof}

\begin{lemma}[A positive Markov matrix with a nilpotent remainder]
    \label{thm:fnw_dimension_markov_matrix}
    \lean{MPSTensor.FNWDimensionConstant.hadamardU, MPSTensor.FNWDimensionConstant.hadamardV, MPSTensor.FNWDimensionConstant.hadamardW, MPSTensor.FNWDimensionConstant.uniformMatrix, MPSTensor.FNWDimensionConstant.markovMatrix, MPSTensor.FNWDimensionConstant.squareMatrix, MPSTensor.FNWDimensionConstant.remainderMatrix, MPSTensor.FNWDimensionConstant.markovMatrix_pos, MPSTensor.FNWDimensionConstant.markovMatrix_sum_identities, MPSTensor.FNWDimensionConstant.markovMatrix_pow_two, MPSTensor.FNWDimensionConstant.markovMatrix_pow_three, MPSTensor.FNWDimensionConstant.remainderMatrix_nilpotent}
    \leanok
    To examine the prescribed-rate coefficient in
    \cite[Section~6, equation~(6.1)]{Nachtergaele1996Spectral}, let
    $e=(1,1,1,1)^{\mathsf T}$ and let $u=(1,1,-1,-1)^{\mathsf T}$,
    $v=(1,-1,1,-1)^{\mathsf T}$, $w=(1,-1,-1,1)^{\mathsf T}$. With
    $Q=ee^{\mathsf T}/4$ and $P=Q+(uv^{\mathsf T}+vw^{\mathsf T})/16$,
    the matrix $P$ is strictly positive and doubly stochastic, and
    \begin{align}
        P^2     & =Q+uw^{\mathsf T}/64, & P^3            & =Q,   \\
        Pu      & =0,                   & w^{\mathsf T}P & =0,   \\
        (P-Q)^3 & =0,                   & (P-Q)^2        & \ne0.
    \end{align}
\end{lemma}
\begin{proof}
    \leanok
    The four displayed sign vectors are pairwise orthogonal and have squared norm
    four. Multiplication gives the identities. Each entry of $P$ belongs to
    $\{1/8,1/4,3/8\}$; its row and column sums are one.
\end{proof}

\begin{lemma}[Scalar path sums]
    \label{thm:fnw_dimension_markov_path_sums}
    \lean{MPSTensor.FNWDimensionConstant.realMarkovMatrix, MPSTensor.FNWDimensionConstant.realHadamardU, MPSTensor.FNWDimensionConstant.realHadamardW, MPSTensor.FNWDimensionConstant.markovAmplitude, MPSTensor.FNWDimensionConstant.realMarkovMatrix_pos, MPSTensor.FNWDimensionConstant.realMarkovMatrix_sum_row, MPSTensor.FNWDimensionConstant.realMarkovMatrix_sum_col, MPSTensor.FNWDimensionConstant.realHadamardU_sq, MPSTensor.FNWDimensionConstant.realHadamardW_sq, MPSTensor.FNWDimensionConstant.realMarkovMatrix_mul_hadamardU, MPSTensor.FNWDimensionConstant.realHadamardW_mul_markovMatrix, MPSTensor.FNWDimensionConstant.realMarkovMatrix_pow_two, MPSTensor.FNWDimensionConstant.realMarkovMatrix_pow_three, MPSTensor.FNWDimensionConstant.realMarkovMatrix_path_sum, MPSTensor.FNWDimensionConstant.realMarkovMatrix_path_hadamardU_sum, MPSTensor.FNWDimensionConstant.realMarkovMatrix_path_hadamardUW_sum, MPSTensor.FNWDimensionConstant.markovAmplitude_sq, MPSTensor.FNWDimensionConstant.markovAmplitude_pos}
    \leanok
    \uses{thm:fnw_dimension_markov_matrix}
    Regard $P$, $u$ and $w$ above as real. Then $P_{ij}>0$,
    $\sum_jP_{ij}=\sum_jP_{ji}=1$ and $u_i^2=w_i^2=1$ for every $i$, and the
    identities $Pu=0$, $w^{\mathsf T}P=0$, $P^2=Q+uw^{\mathsf T}/64$ and
    $P^3=Q$ of Lemma~\ref{thm:fnw_dimension_markov_matrix} hold over $\R$.
    For every pair of endpoints $a,f$, set
    $p_{abcef}=P_{ab}P_{bc}P_{ce}P_{ef}$. Then
    \begin{align}
        \sum_{b,c,e}p_{abcef}        & =\frac14,    &
        \sum_{b,c,e}u_b p_{abcef}    & =0,          &
        \sum_{b,c,e}u_bw_e p_{abcef} & =\frac1{64}.
    \end{align}
    The amplitudes $t_{ij}=\sqrt{P_{ij}}$ are strictly positive and satisfy
    $t_{ij}^2=P_{ij}$.
\end{lemma}
\begin{proof}
    \leanok
    \uses{thm:fnw_dimension_markov_matrix}
    The real identities are those of
    Lemma~\ref{thm:fnw_dimension_markov_matrix}, read in $\R$.
    Applying the identities for $P^2$ and $P^3$ together with $u_b^2=w_e^2=1$
    yields the path sums, and positivity gives the assertions concerning
    square roots.
\end{proof}

\begin{lemma}[A dimension-four physical projection defect]
    \label{thm:fnw-dimension-four-physical-defect}
    \lean{MPSTensor.FNWDimensionConstant.dimensionFour_projection_defect_lower}
    \leanok
    For the transition matrix $P$ and weighted matrix-unit tensor
    $A^{ij}=\sqrt{P_{ij}}E_{ij}$ defined above, let $G_I$ denote the
    orthogonal projection onto the ground space on the interval $I$,
    extended by the identity on the remaining sites. Then
    \begin{align}
        \left\|G_{[1,3]}G_{[0,2]}-G_{[0,3]}\right\|
        \geq \frac1{16}. \notag
    \end{align}
\end{lemma}
\begin{proof}\leanok
    \uses{thm:weighted-matrix-unit-path-contractions,thm:matrix-unit-path-ground-space,thm:fnw-marked-path-interval-spaces,thm:fnw-path-vector-norm-orthogonal,thm:fnw-projection-witness-bound,thm:fnw_dimension_markov_path_sums}
    Fix external vertices $a,f$ and write
    $T_{abcef}=\sqrt{P_{ab}P_{bc}P_{ce}P_{ef}}$. Consider
    \begin{align}
        x & =2\sum_{b,c,e}u_bT_{abcef}|ab,bc,ce,ef\rangle,\notag \\
        y & =2\sum_{b,c,e}w_eT_{abcef}|ab,bc,ce,ef\rangle.\notag
    \end{align}
    With $t_{ij}=\sqrt{P_{ij}}$ these are $x=2v_{a,f}(u\otimes\mathbf 1)$ and
    $y=2v_{a,f}(\mathbf 1\otimes w)$. By
    Lemma~\ref{thm:fnw-marked-path-interval-spaces}, $x$ belongs to the
    range of $G_{[1,3]}$ and $y$ belongs to the range of $G_{[0,2]}$. The path contraction identities give
    $\|x\|=\|y\|=1$, $x\perp\operatorname{range}G_{[0,3]}$, and
    $\langle y,x\rangle=1/16$. Thus the matrix coefficient of
    $G_{[0,2]}G_{[1,3]}-G_{[0,3]}$ between $y$ and $x$ is $1/16$.
    Taking the adjoint gives the asserted order of the two projections.
\end{proof}

\begin{lemma}[Weighted matrix-unit normalization]
    \label{thm:weighted-matrix-unit-normalization}
    \lean{MPSTensor.FNWDimensionConstant.matrixUnitTensor_isInjective,MPSTensor.FNWDimensionConstant.matrixUnitTensor_isTP}
    \leanok
    Let $A_{ij}=t_{ij}E_{ij}$, where $t$ is a real matrix.
    If every entry of $t$ is nonzero, then the letters of $A$ span the full
    matrix algebra. If the squared entries in every column sum to one, then
    \begin{align}
        \sum_{i,j} A_{ij}^*A_{ij}=I.
    \end{align}
\end{lemma}
\begin{proof}
    \leanok
    When $t_{ij}\ne0$, the matrix unit $E_{ij}$ is a scalar multiple of
    $A_{ij}$. The normalization follows from
    $A_{ij}^*A_{ij}=t_{ij}^2E_{jj}$ and the column-sum assumption.
\end{proof}

\begin{lemma}[Transfer map of a weighted matrix-unit tensor]
    \label{thm:fnw_dimension_matrix_unit_transfer}
    \lean{MPSTensor.FNWDimensionConstant.transferMap_matrixUnitTensor}
    \leanok
    Let $t=(t_{ij})$ be a real matrix and define physical letters
    $A_{ij}=t_{ij}E_{ij}$. The associated transfer map satisfies
    \begin{align}
        \mathcal E_A(X)
         & =\operatorname{diag}\left(\sum_j t_{ij}^2X_{jj}\right)_i.
    \end{align}
\end{lemma}
\begin{proof}
    \leanok
    Multiplying the matrix units in $\sum_{i,j}A_{ij}XA_{ij}^{\dagger}$ shows
    that only the diagonal entries of $X$ contribute.
\end{proof}

\begin{lemma}[The positive Markov matrix-unit channel]
    \label{thm:fnw_dimension_markov_channel}
    \lean{MPSTensor.FNWDimensionConstant.markovTensor, MPSTensor.FNWDimensionConstant.markovTensor_isInjective, MPSTensor.FNWDimensionConstant.markovTensor_isTP, MPSTensor.FNWDimensionConstant.markovTensor_transferMap_one, MPSTensor.FNWDimensionConstant.markovTensor_transferMap_cube}
    \leanok
    For the four-dimensional matrix $P$ used to examine
    \cite[Section~6, equation~(6.1)]{Nachtergaele1996Spectral}, let
    $A_{ij}=\sqrt{P_{ij}}E_{ij}$. Then $A$ is injective at length one, its transfer
    map $\mathcal E_A$ is unital and trace preserving, and every matrix $X$ satisfies
    \begin{align}
        \mathcal E_A^3(X) & =\frac{\operatorname{tr}(X)}4 I.
    \end{align}
\end{lemma}
\begin{proof}
    \leanok
    \uses{thm:fnw_dimension_markov_matrix, thm:fnw_dimension_matrix_unit_transfer, thm:weighted-matrix-unit-normalization}
    The positive coefficients ensure that the physical letters span all matrix
    units. Column stochasticity gives trace preservation, and row stochasticity
    gives unitality. On diagonal matrices the transfer map acts by $P$, while it
    annihilates off-diagonal matrix units. Therefore $P^3=ee^{\mathsf T}/4$ gives
    the asserted formula.
\end{proof}

\begin{lemma}[A faithful primitive fixed point]
    \label{thm:fnw-dimension-markov-primitive}
    \lean{MPSTensor.FNWDimensionConstant.markovDensity,MPSTensor.FNWDimensionConstant.markovDensity_posDef,MPSTensor.FNWDimensionConstant.markovDensity_trace,MPSTensor.FNWDimensionConstant.markovDensity_fixed,MPSTensor.FNWDimensionConstant.markovTensor_isPrimitiveMPS,MPSTensor.FNWDimensionConstant.markovTensor_fnwTransfer_eigenvalue_eq_zero}
    \leanok
    For the weighted matrix-unit tensor of the Markov example, the matrix
    $\rho=I/4$ is positive definite, has trace one, and is stationary. The tensor
    has a complementary transfer spectral gap at $\rho$. Every nonunit
    eigenvalue of its FNW observable transfer map is zero.
\end{lemma}
\begin{proof}
    \leanok
    \uses{thm:weighted-matrix-unit-normalization,thm:fnw_dimension_markov_channel}
    The row sums of the Markov matrix give stationarity of $I/4$.
    For a nonunit eigenvalue $\nu$ of the trace-preserving transfer map,
    trace preservation forces an associated eigenvector $X$ to have trace zero.
    The depolarizing third iterate then gives $\nu^3X=0$, and hence $\nu=0$.
    The bilinear trace adjoint has the same eigenvalues, giving the assertion
    for the FNW observable transfer map. One-site injectivity gives
    irreducibility. Together with the unique peripheral eigenvalue and the
    faithful stationary density, this implies the complementary spectral gap.
\end{proof}

\begin{lemma}[Failure of the dimension-only coefficient at a prescribed rate]
    \label{thm:fnw-dimension-only-counterexample}
    \lean{MPSTensor.FNWDimensionConstant.dimensionFour_projection_defect_exceeds_coefficient,MPSTensor.FNWDimensionConstant.exists_dimensionFour_counterexample}
    \leanok
    There exists an injective matrix-product tensor of physical dimension
    $16$ and auxiliary dimension $4$ with a faithful primitive stationary
    density matrix of trace one such that every eigenvalue other than $1$
    of its observable transfer map has modulus smaller than
    $\lambda=1/1000$. Nevertheless, at overlap length $m=2$,
    \begin{align}
        16\lambda^2 & <1,\notag                                                \\
        \left\|G_{[1,3]}G_{[0,2]}-G_{[0,3]}\right\|
                    & >\frac{16\lambda^2(1+16\lambda^2)}{1-16\lambda^2}.\notag
    \end{align}
    As recorded in \cite{gap:cpgsv21_martingale_overlap}, the dimension-only
    coefficient and the prescribed-rate assertion following Nachtergaele's
    bound therefore cannot hold simultaneously at this fixed overlap length.
\end{lemma}
\begin{proof}\leanok
    \uses{thm:fnw-dimension-four-physical-defect,thm:fnw-dimension-four-coefficient,thm:fnw-dimension-markov-primitive}
    For the weighted matrix-unit tensor associated with the transition
    matrix above and the stationary density $I/4$, the transfer map becomes
    completely depolarizing after three steps, so every nonunit eigenvalue
    of the observable transfer map is zero. The rate $1/1000$ is therefore
    admissible. The physical projection defect is at least $1/16$, whereas
    the proposed coefficient is strictly smaller than $1/16$.
\end{proof}

\begin{lemma}[Block injectivity from a simultaneous word span]
    \label{thm:block_injective_of_word_tuple_span}
    \lean{MPSTensor.isNBlkInjective_of_wordTupleSpanTop}
    \leanok
    \uses{def:blockwise_product_word_span, def:l_blk_injective}
    Let $A^0,\ldots,A^{r-1}$ be tensors with bond dimensions
    $D_0,\ldots,D_{r-1}$, and let $L\in\mathbb N$. Suppose that their
    simultaneous word evaluations at length $L$ span
    $\prod_{k=0}^{r-1}\MN{D_k}$. Then every $A^j$ is $L$-block injective.
\end{lemma}
\begin{proof}
    \leanok
    Let $M\in\MN{D_j}$ and let $X$ be the tuple with $X_j=M$ and $X_k=0$ for
    $k\neq j$. By hypothesis,
    \begin{align}
        X=\sum_{|\omega|=L}c_\omega(A^{0,\omega},\ldots,A^{r-1,\omega})\notag
    \end{align}
    for some coefficients $c_\omega$. Its $j$-th component is
    $M=\sum_{|\omega|=L}c_\omega A^{j,\omega}$, so the length-$L$ words of
    $A^j$ span $\MN{D_j}$.
\end{proof}

\begin{lemma}[Separation from simultaneous injectivity]
    \label{thm:block-word-span-separation}
    \lean{MPSTensor.not_gaugePhaseEquiv_of_wordTupleSpanTop}
    \leanok
    Let $A^0,\ldots,A^{r-1}$ be tensors with positive bond dimensions. As a
    consequence of the simultaneous block-injectivity condition in
    \cite[Section IV.C]{Cirac2021Matrix}, if their simultaneous word
    evaluations at length $S$ span the product of their matrix algebras, then
    no two distinct blocks are gauge-phase equivalent.
\end{lemma}
\begin{proof}
    \leanok
    \uses{lem:product_span_block_image_direct_sum, thm:ground_space_map_injective_blk,
        thm:gauge_phase_mpv_scaling, lem:mpv_mem_ground_space,
        lem:mpv_eq_ground_space_map_one, thm:block_injective_of_word_tuple_span}
    The local ground spaces of the blocks at length $S$ are independent. Each
    block is $S$-block injective by
    Lemma~\ref{thm:block_injective_of_word_tuple_span}, so each boundary
    map at that length is injective, and its value at the identity matrix,
    the periodic matrix product vector $V^{(S)}(A^i)=\Gamma^{A^i}_S(\Id)$, is
    nonzero. Suppose that $i\neq j$ and $A^{i,a}=\zeta XA^{j,a}X^{-1}$ for every
    physical index $a$. Then
    \begin{align}
        V^{(S)}(A^i)=\zeta^SV^{(S)}(A^j)\in G_S(A^i)\cap G_S(A^j)=\{0\},\notag
    \end{align}
    so $\Gamma^{A^i}_S(\Id)=0$, contradicting injectivity of the boundary map.
\end{proof}

\begin{lemma}[Propagation of a positive-length simultaneous span]
    \label{thm:block_word_span_propagation}
    \lean{MPSTensor.wordTupleSpanTop_of_ge}
    \leanok
    \uses{def:blockwise_product_word_span}
    Let $A^0,\ldots,A^{r-1}$ have positive bond dimensions and a common
    physical dimension. As in the propagation of block injectivity in
    \cite[Section~IV.C]{Cirac2021Matrix}, if their simultaneous length-$S$
    word tuples span the product matrix algebra with $S>0$, then their word
    tuples span the same algebra at every length $n\geq S$.
\end{lemma}
\begin{proof}
    \leanok
    \uses{thm:block_injective_of_word_tuple_span,
        thm:normalized_primitive_gauge,thm:word_tuple_span_gauge_invariance,
        thm:block_word_span_independent_rescaling,thm:ph_tp_word_span_propagation}
    Each block is $S$-block injective and hence normal. Rescaling and
    conjugating each block into normalized primitive form preserves
    the simultaneous span at $S$, trace preservation propagates that span
    to $n\geq S$, and inverse conjugation and rescaling return the full span
    to the original family at length $n$.
\end{proof}

\begin{lemma}[Open ground spaces at a simultaneous injectivity length]
    \label{thm:block_open_kernel_simultaneous_injectivity}
    \lean{MPSTensor.%
        ker_openParentHamiltonianES_toTensorFromBlocks_eq_groundSpaceES_of_wordTupleSpanTop}
    \leanok
    Let $A^0,\ldots,A^{r-1}$ have positive physical and bond dimensions.
    Suppose that their simultaneous length-$S$ word tuples span the product
    matrix algebra, where $S>0$. Set $B=\bigoplus_j\mu_jA^j$, with
    every $\mu_j\neq0$. By the inversion and regrowth argument in
    \cite[Section~IV.C, lines~2114--2129]{Cirac2021Matrix}, every
    $S+1\leq R\leq N$ satisfies
    $\ker H^{\mathrm{open}}_{N,R}(B)=G_N^{\mathrm{ES}}(B)$.
\end{lemma}
\begin{proof}
    \leanok
    \uses{thm:block_word_span_primitive_normalization,
        thm:ph_tp_word_span_propagation,thm:ph_tp_block_open_ground_space,
        thm:ph_block_replacement_parent_equality}
    Normalizing the blocks preserves their simultaneous span at $S$
    and every component local space. Trace preservation gives the span
    at all lengths at least $S$, so the open-chain intersection theorem
    gives the kernel identity for the normalized block sum. The original
    and normalized sums have equal local spaces at $R$ and $N$; hence
    their open Hamiltonians and their spaces $G_N^{\mathrm{ES}}$ agree.
\end{proof}

\begin{lemma}[Open-chain gaps at a simultaneous injectivity length]
    \label{thm:block_open_gap_all_lengths_simultaneous_injectivity}
    \lean{MPSTensor.%
        exists_openParentHamiltonianES_toTensorFromBlocks_gap_of_wordTupleSpanTop}
    \leanok
    Let $A^0,\ldots,A^{r-1}$ have positive bond dimensions and common
    physical dimension $d\geq1$. Suppose that their simultaneous
    length-$S$ word tuples span the product matrix algebra, where $S>0$.
    Let $B=\bigoplus_j\mu_jA^j$, with every $\mu_j\ne0$.
    With no normalization or inequivalence assumption required on the
    original blocks, the finite-interval conclusion of
    \cite[Theorem~1.2 and Section~6]{Nachtergaele1996Spectral} holds at
    every interaction range $R\geq S+1$: there exists $\delta>0$
    such that
    \begin{align}
        \delta\|v\|\leq\|H^{\mathrm{open}}_{N,R}(B)v\|
        \notag
    \end{align}
    for every $N\geq R$ and every $v\perp G_N(B)$.
\end{lemma}
\begin{proof}\leanok
    \uses{thm:block_word_span_primitive_normalization,
        thm:block-word-span-separation,
        thm:block_open_kernel_simultaneous_injectivity,
        thm:primitive_block_sum_open_gap_all_lengths,
        lem:parent-finite-interval-comparison,
        lem:open_parent_range_comparison}
    Independent normalization of the blocks into primitive form preserves
    their simultaneous span at $S$ and all component local ground spaces, and
    simultaneous injectivity separates these normalized blocks. The
    all-length primitive theorem gives a uniform gap at some range $2p$
    with $p\geq R$. At range $R$, the open kernel is $G_N$ for every
    $N\geq R$. On the interval of length $2p$, finite-dimensional
    comparison gives $\kappa h_{2p}\leq H^{\mathrm{open}}_{2p,R}$ for
    some $\kappa>0$. The interval comparison transfers the long-range
    gap to range $R$ at every $N\geq2p$, with loss $2p-R+1$.
    A finite minimum includes $R\leq N<2p$. Equality of the local spaces
    returns the same Hamiltonians and ground spaces for the original
    weighted block tensor.
\end{proof}

\begin{lemma}[Dimension-bound open gaps for normal blocks]
    \label{thm:normal_block_open_gap_dimension_bound}
    \lean{MPSTensor.%
        exists_openParentHamiltonianES_toTensorFromBlocks_gap_of_isNormalTensor}
    \leanok
    Let $A^j$ be normal tensors with common physical dimension $d\geq1$,
    pairwise inequivalent under virtual similarity and phase. Set
    $D_* = \max\{\sum_j D_j,1\}$ and $B=\bigoplus_j\mu_jA^j$,
    with every $\mu_j\ne0$. For every $R\geq3D_*^5$, there exists
    $\delta>0$ such that $\delta\|v\|\leq\|H^{\mathrm{open}}_{N,R}(B)v\|$
    for every $N\geq R$ and every $v\perp G_N(B)$.
\end{lemma}
\begin{proof}\leanok
    \uses{thm:block_open_gap_all_lengths_simultaneous_injectivity,
        thm:cpsv_bnt_sharp_span_ambient_cap}
    The dimension-bound blocking theorem gives a positive simultaneous
    injectivity length $S$ with $S+1\leq3D_*^5$, so the open-chain gap
    theorem applies at every $R\geq S+1$.
\end{proof}

\begin{lemma}[Open parent gaps for canonical tensors]
    \label{thm:canonical_open_gap_all_lengths}
    \lean{MPSTensor.CPSVCanonicalFormData.openParentHamiltonianES_eq_representatives,
        MPSTensor.CPSVCanonicalFormData.%
        exists_openParentHamiltonianES_gap_of_wordTupleSpanTop,
        MPSTensor.IsCPSVCanonicalForm.exists_bnt_open_gap_of_wordTupleSpanTop,
        MPSTensor.IsCPSVCanonicalForm.%
        exists_openParentHamiltonianES_gap_of_three_bondDim_pow_five_le,
        MPSTensor.CPSVCanonicalFormData.%
        ker_openParentHamiltonianES_eq_groundSpaceES_of_wordTupleSpanTop}
    \leanok
    Let $A$ be a canonical tensor with physical dimension $d\geq1$, and
    choose its distinct normal representatives $B^j$. At every positive
    interaction range, its open parent Hamiltonian agrees with that of
    $\bigoplus_j B^j$, including all multiplicities and the ambient
    reconstruction of $A$. If the representatives are simultaneously
    injective at $S>0$ and $R\geq S+1$, then
    $\ker H^{\mathrm{open}}_{N,R}(A)=G_N(A)$ for every $N\geq R$,
    and there is $\delta>0$ such that
    \begin{align}
        \delta\|v\|\leq\|H^{\mathrm{open}}_{N,R}(A)v\|
        \notag
    \end{align}
    for every $N\geq R$ and every $v\perp G_N(A)$.
    If the bond dimension $D$ is positive, the conclusion holds at every
    $R\geq3D^5$ without a given injectivity length.
\end{lemma}
\begin{proof}\leanok
    \uses{thm:block_open_gap_all_lengths_simultaneous_injectivity,
        thm:block_open_kernel_simultaneous_injectivity,
        thm:canonical_bounded_range_ground_space}
    The canonical reconstruction preserves every positive-length local
    ground space, identifying the open parent Hamiltonians and their
    boundary-condition spaces with those of the distinct representatives, to
    which the block open-chain gap and kernel theorems apply; for the dimension
    bound, the blocking estimate gives $S>0$ with $S+1\leq3D^5$.
\end{proof}

\begin{definition}[Open sums of a local interaction]
    \label{def:open_local_interaction_sum}
    \lean{MPSTensor.openInteractionHamiltonianES}
    \leanok
    For a local operator $h$ on $R$ sites, let
    $H^{\mathrm{open}}_N(h)$ be the sum of its translates over starts
    $i$ satisfying $0\leq i$ and $i+R\leq N$. Each translate acts as
    the identity on the remaining physical sites.
\end{definition}

\begin{lemma}[Properties of open interaction sums]
    \label{lem:open_local_interaction_sum_properties}
    \lean{MPSTensor.openInteractionHamiltonianES_mono,
        MPSTensor.openInteractionHamiltonianES_smul,
        MPSTensor.openInteractionHamiltonianES_zero,
        MPSTensor.openInteractionHamiltonianES_isPositive,
        MPSTensor.openInteractionHamiltonianES_parentInteractionES}
    \leanok
    Open interaction sums preserve order and positivity, commute with
    scalar multiplication, and vanish for the zero interaction.
    When $R>0$, the sum of the canonical parent projection is the
    canonical open parent Hamiltonian.
\end{lemma}
\begin{proof}\leanok
    \uses{def:open_local_interaction_sum}
    Extension by the identity and physical coordinate isometries preserve
    order, positivity, and scalar multiplication. Summation preserves
    these properties. Extending the local parent projection gives each
    canonical parent term.
\end{proof}

\begin{lemma}[Open gaps for fixed positive parent interactions]
    \label{thm:open_positive_parent_gap}
    \lean{MPSTensor.IsParentInteraction.exists_pos_open_comparison,
        MPSTensor.IsParentInteraction.exists_open_uniform_gap_of_canonical_gap,
        MPSTensor.CPSVCanonicalFormData.%
        exists_openInteraction_uniform_gap_of_wordTupleSpanTop}
    \leanok
    Let $R>0$, and let $h\geq0$ have kernel $G_R(A)$.
    There are $\kappa,C>0$, independent of $N$, such that
    \begin{align}
        \kappa H^{\mathrm{open}}_{N,R}(A)
        \leq H^{\mathrm{open}}_N(h)
        \leq C H^{\mathrm{open}}_{N,R}(A).
        \notag
    \end{align}
    Hence a uniform canonical open-chain gap $\gamma>0$ for $N\geq R$
    gives the gap $\kappa\gamma$ for $H^{\mathrm{open}}_N(h)$ on its
    kernel complement. In particular, this applies to every fixed positive
    parent interaction of a canonical tensor when its normal representatives
    are simultaneously injective at $S>0$ and $R\geq S+1$.
\end{lemma}
\begin{proof}\leanok
    \uses{lem:open_local_interaction_sum_properties,
        thm:canonical_open_gap_all_lengths}
    Finite-dimensional positivity and the prescribed local kernel give
    $\kappa h_R(A)\leq h\leq C h_R(A)$, and summing the translated
    comparisons yields the two-sided operator bounds. Because the two positive
    operators bound one another by positive constants, their kernels coincide,
    and the lower bound transfers the canonical gap to the common kernel
    complement.
\end{proof}

\begin{lemma}[Compact comparison for open parent interactions]
    \label{thm:compact_open_positive_parent_gap}
    \lean{MPSTensor.%
        exists_uniform_openInteractionHamiltonianES_gap_of_compact_canonical_gap,
        MPSTensor.CPSVCanonicalFormData.exists_uniform_openInteraction_gap_of_compact,
        MPSTensor.IsCPSVCanonicalForm.%
        exists_uniform_openInteraction_gap_of_compact_of_bounded_range}
    \leanok
    Let $K$ be a compact parameter set. Suppose the local ground projections
    of $A_x$ are continuous, and let $h_x$ be a continuous family of positive
    interactions with kernel $G_R(A_x)$, where $R>0$. If the canonical open
    parent Hamiltonians have one positive gap for every $x\in K$ and every
    $N\geq N_0$, then the interaction sums $H^{\mathrm{open}}_N(h_x)$
    also have one positive gap for those parameters and lengths.

    For a fixed canonical tensor $A$, this conclusion requires no given
    canonical gap: simultaneous injectivity at $S>0$ and $R\geq S+1$
    give one gap uniform over $x\in K$ and every $N\geq R$.
    When its bond dimension $D$ is positive, $R\geq3D^5$ suffices without
    a given injectivity length.
\end{lemma}
\begin{proof}\leanok
    \uses{thm:open_positive_parent_gap,thm:canonical_open_gap_all_lengths}
    Compactness and continuity give a single $\kappa>0$ with
    $\kappa h_R(A_x)\leq h_x$ for every $x\in K$. Each interaction also
    has a finite positive upper comparison constant. The preceding proof
    transfers the given gap with the common lower constant; the upper
    constants identify the kernels and do not affect the resulting gap.
    For a fixed canonical tensor, the all-length open gap theorem provides the
    required canonical gap.
\end{proof}

\begin{lemma}[Compact gaps on bounded open intervals]
    \label{thm:compact_bounded_open_interval_gap}
    \lean{MPSTensor.%
        exists_uniform_openParentHamiltonianES_toTensorFromBlocks_gap_fixed_volume,
        MPSTensor.%
        exists_uniform_openParentHamiltonianES_toTensorFromBlocks_gap_bounded_volumes}
    \leanok
    Let $K$ be a compact parameter set, and let $A_x^j$ be continuous
    block tensors with positive physical and bond dimensions. Suppose that
    their length-$S$ word tuples span the product matrix algebra for every
    parameter, where $S>0$. Set $B_x=\bigoplus_j\mu_j(x)A_x^j$, with
    all weights nonzero; the weights need not be continuous.
    Fix $R\geq S+1$ and a finite upper length $M$.
    There is $\delta>0$ such that
    \begin{align}
        \delta\|v\|\leq\|H^{\mathrm{open}}_{N,R}(B_x)v\|
        \notag
    \end{align}
    for every $x\in K$, every $R\leq N\leq M$, and every
    $v\perp\ker H^{\mathrm{open}}_{N,R}(B_x)$.
    The constant may depend on $M$.
\end{lemma}
\begin{proof}\leanok
    \uses{thm:block_open_kernel_simultaneous_injectivity,
        thm:block_word_span_propagation}
    At each fixed $N\geq R$, simultaneous injectivity gives continuous
    local ground projections at lengths $R$ and $N$. The open Hamiltonian
    depends continuously on the first projection, and its kernel is $G_N$.
    The compact finite-dimensional norm-gap theorem gives one positive
    bound at this $N$. Taking the minimum over $R\leq N\leq M$ gives the
    claim. Nonzero weights preserve every local ground space.
\end{proof}

\begin{definition}[Zero-padded open interactions]
    \label{def:open_knabe_padding}
    \lean{MPSTensor.zmodOpenLocalTermES}
    \leanok
    \uses{def:open_parent_hamiltonian_es}
    Let $R,P>0$. On the Hilbert space of an $N$-site chain, index the
    nonwrapping canonical interactions by $s\in\Z/P\Z$. Writing $\bar s$
    for the representative in $\{0,\ldots,P-1\}$, define
    \begin{align}
        \widetilde h_s & =
        \begin{cases}
            h_{\bar s}^{(N,R)}(A), & \bar s+R\leq N, \\
            0,                     & \bar s+R>N.
        \end{cases}
        \label{eq:open_knabe_padding}
    \end{align}
    The period indexes the interactions; the physical chain still has $N$ sites.
\end{definition}

\begin{lemma}[Projection, sum, and separation of padded interactions]
    \label{thm:open_knabe_padding_properties}
    \lean{MPSTensor.zmodOpenLocalTermES_isSymmetricProjection,
        MPSTensor.zmodOpenLocalTermES_commute_of_oriented_separation,
        MPSTensor.sum_zmodOpenLocalTermES_eq_openParentHamiltonianES}
    \leanok
    \uses{def:open_knabe_padding}
    Every $\widetilde h_s$ is an orthogonal projection. If $R\leq e$ and
    $e+R\leq P$, then $\widetilde h_s\widetilde h_{s+e}
        =\widetilde h_{s+e}\widetilde h_s$ for every $s$. If $N\leq P$,
    $\sum_s\widetilde h_s=H^{\mathrm{open}}_{N,R}(A)$.
\end{lemma}
\begin{proof}\leanok
    \uses{def:open_knabe_padding,lem:transported_es_local_projections,
        lem:transported_es_nonoverlap_positive,
        thm:physical-outside-window-martingale-cancellation}
    Each term is a local projection or zero. At the stated separation,
    either $\bar s+R\leq\overline{s+e}$ or
    $\overline{s+e}+R\leq\bar s$. When both terms are active, their
    nonwrapping intervals are disjoint, so the local terms commute.
    For $P\geq N$, every nonwrapping start occurs exactly once; all
    remaining summands vanish.
\end{proof}

\begin{lemma}[Clipped windows of padded interactions]
    \label{thm:open_knabe_clipped_windows}
    \lean{MPSTensor.cyclicWindowSum_zmodOpenLocalTermES_eq_zero_or_interval,
        MPSTensor.cyclicWindowSum_zmodOpenLocalTermES_quadraticForm_gap}
    \leanok
    \uses{def:open_knabe_padding}
    Let $m>0$, $N\geq R>0$, and $P=N+2m$. Put
    $A_s=\sum_{q=0}^{m-1}\widetilde h_{s+q}$. Each $A_s$ is zero or the
    range-$R$ open Hamiltonian on an interval $[n-W,n)$, extended by the
    identity on the remaining sites, with
    $R\leq W\leq m+R-1$ and $W\leq n\leq N$.
    If every open Hamiltonian on $R\leq W\leq m+R-1$ sites has norm
    gap at least $\gamma>0$ on its kernel complement, then
    $A_s^2\geq\gamma A_s$ for every $s$.
\end{lemma}
\begin{proof}\leanok
    \uses{def:open_knabe_padding,thm:open_interval_gap_transport}
    The active starts are $0,\ldots,N-R$. Padding by $2m$ zeros prevents
    a cyclic window of $m$ starts from meeting two separated parts of this
    interval. Its active starts therefore form one interval $[a,b)$, or
    none. In the first case, with $W=b-a+R-1$ and $n=b+R-1$, the local terms
    sum to the interval Hamiltonian on $[a,n)$.
    Its norm gap is preserved under extension by free sites. Positivity
    converts this norm gap into $A_s^2\geq\gamma A_s$; the zero case is immediate.
\end{proof}

\begin{lemma}[An open-chain finite-size gap criterion]
    \label{thm:open_finite_range_knabe_gap}
    \lean{MPSTensor.openParentHamiltonianES_gap_of_finite_open_gaps}
    \leanok
    \uses{def:open_parent_hamiltonian_es}
    Let $1\leq R\leq m$, and suppose that every range-$R$ open parent
    Hamiltonian on $R\leq W\leq m+R-1$ sites has norm gap at least
    $\gamma$ on its kernel complement. By the cyclic finite-range inequality
    derived in \cite{gap:knabe88_finite_range_coefficient} (rather than the
    printed periodic nearest-neighbor criterion), if $m\gamma>(R-1)^2$, then
    \begin{align}
        \delta & =\frac{m\gamma-(R-1)^2}{m-R+1}>0,
        \label{eq:open_finite_range_knabe_delta}
    \end{align}
    and $\delta\|v\|\leq\|H^{\mathrm{open}}_{N,R}(A)v\|$ for every
    $N\geq R$ and $v\perp\ker H^{\mathrm{open}}_{N,R}(A)$.
\end{lemma}
\begin{proof}\leanok
    \uses{thm:open_knabe_padding_properties,thm:open_knabe_clipped_windows,
        thm:finite_range_cyclic_knabe,lem:quadratic_form_to_norm_bound}
    For each physical volume $N$, with period $P=N+2m$,
    the padded projections have sum $H^{\mathrm{open}}_{N,R}(A)$ and
    satisfy the separated commutation hypothesis. Every cyclic window
    satisfies $A_s^2\geq\gamma A_s$. The cyclic inequality therefore gives
    $(H^{\mathrm{open}}_{N,R}(A))^2\geq\delta H^{\mathrm{open}}_{N,R}(A)$
    with $\delta$ as in~(\ref{eq:open_finite_range_knabe_delta}).
    The positive-operator quadratic-form estimate gives the asserted norm gap.
\end{proof}

\begin{lemma}[Persistence of gaps on finitely many open intervals]
    \label{thm:open_gap_bounded_volume_continuity}
    \lean{MPSTensor.eventually_openParentHamiltonianES_toTensorFromBlocks_gap_fixed_volume,
        MPSTensor.eventually_openParentHamiltonianES_toTensorFromBlocks_gap_bounded_volumes}
    \leanok
    Let $A_x^j$ be continuous block tensors with positive physical and bond
    dimensions. Suppose their length-$S$ word tuples span the product matrix
    algebra at every parameter, where $S>0$, and put
    $B_x=\bigoplus_j\mu_j(x)A_x^j$, with every weight nonzero.
    Fix $R\geq S+1$ and a finite upper volume $M$.
    As in the finite-volume continuity step of
    \cite[Appendix~A, lines~2575--2578]{Schuch2011Phases}, with no continuity
    of the weights required, if at $x_0$ the open Hamiltonians on every
    $R\leq N\leq M$ have norm gap at least $\gamma>0$ on $G_N(B_{x_0})^\perp$,
    then every $\delta<\gamma$ is a common lower bound on $G_N(B_x)^\perp$ for
    those volumes and all $x$ in some neighborhood of $x_0$ (and the same
    assertion holds for one fixed $N\geq R$).
\end{lemma}
\begin{proof}\leanok
    \uses{thm:block_ground_space_continuity,thm:block_word_span_propagation,
        thm:kernel_gap_perturbation}
    The ground projections at $R$ and $N$ are continuous. The open
    Hamiltonian is a finite sum of translates of $I-P_{G_R}$, hence
    continuous. Its ground space contains $G_N$.
    The operator and projection perturbation estimate preserves each strict
    lower bound at fixed $N$, and intersecting the finitely many resulting
    neighborhoods completes the proof.
\end{proof}

\begin{lemma}[Local persistence of an all-length open gap]
    \label{thm:open_gap_all_length_neighborhood}
    \lean{MPSTensor.eventually_openParentHamiltonianES_toTensorFromBlocks_gap_of_uniform_gap}
    \leanok
    Under the block-family hypotheses of
    Lemma~\ref{thm:open_gap_bounded_volume_continuity}, suppose that at
    $x_0$ one $\gamma>0$ bounds the open norm gaps on $G_N(B_{x_0})^\perp$
    for every $N\geq R$. Then some $\delta>0$ bounds the open norm gaps
    on their kernel complements for every $N\geq R$ and every $x$ in
    one neighborhood of $x_0$.
\end{lemma}
\begin{proof}\leanok
    \uses{thm:open_gap_bounded_volume_continuity,thm:open_finite_range_knabe_gap,
        thm:block_open_kernel_simultaneous_injectivity}
    For $m\geq R$ chosen with $m\gamma/2>(R-1)^2$, continuity preserves
    the bound $\gamma/2$ on every $R\leq W\leq m+R-1$ in one
    neighborhood. Simultaneous injectivity identifies each kernel with
    $G_W(B_x)$. The finite-size criterion gives the same positive
    $\delta=(m\gamma/2-(R-1)^2)/(m-R+1)$ for all $N\geq R$ there.
\end{proof}

\begin{lemma}[Compact uniform open gaps for simultaneously injective blocks]
    \label{thm:compact_block_open_gap}
    \lean{MPSTensor.exists_uniform_openParentHamiltonianES_toTensorFromBlocks_gap_of_compact_all_lengths}
    \leanok
    Let $K$ be compact, and let $A_x^j$ be continuous block tensors with
    positive physical and bond dimensions. Suppose their length-$S$ word
    tuples span the product matrix algebra for every $x$, with $S>0$.
    Put $B_x=\bigoplus_j\mu_j(x)A_x^j$, with every weight nonzero.
    Following the compactness argument of \cite[Appendix~A]{Schuch2011Phases}
    at the stated sufficient range, with no continuity of the weights required,
    for each fixed $R\geq S+1$ there is $\delta>0$ such that
    $\delta\|v\|\leq\|H^{\mathrm{open}}_{N,R}(B_x)v\|$ for every
    $x\in K$, every $N\geq R$, and every
    $v\perp\ker H^{\mathrm{open}}_{N,R}(B_x)$.
\end{lemma}
\begin{proof}\leanok
    \uses{thm:block_open_gap_all_lengths_simultaneous_injectivity,
        thm:open_gap_all_length_neighborhood}
    The pointwise all-length open gap gives a positive bound at each
    parameter. The preceding theorem preserves some positive bound on
    a neighborhood, uniformly in $N$. A finite subcover of $K$ and the
    minimum of its finitely many bounds give $\delta>0$.
\end{proof}

\begin{lemma}[Compact dimension-bound open gaps for normal blocks]
    \label{thm:compact_normal_block_open_gap}
    \lean{MPSTensor.exists_uniform_block_openParentHamiltonianES_gap_of_compact_isNormalTensor}
    \leanok
    Let $A_x^j$ be continuous normalized normal blocks, pairwise inequivalent
    under virtual similarity and phase at every parameter, with physical
    dimension $d\geq1$ and fixed bond dimensions $D_j$.
    Put $D_* =\max\{\sum_jD_j,1\}$ and
    $B_x=\bigoplus_j\mu_j(x)A_x^j$, with every weight nonzero.
    For every compact $K$ and every fixed $R\geq3D_*^5$, one positive
    open norm gap holds uniformly for $x\in K$ and all $N\geq R$.
\end{lemma}
\begin{proof}\leanok
    \uses{thm:cpsv_bnt_sharp_span_ambient_cap,thm:block_word_span_propagation,
        thm:compact_block_open_gap}
    The dimension-bound span theorem gives, at each parameter, a positive
    simultaneous injectivity length $L_x$ with $L_x+1\leq3D_*^5$.
    Propagation gives the common length $S=3D_*^5-1>0$, so the compact
    open-gap theorem applies at every $R\geq S+1$.
\end{proof}

\begin{lemma}[Compact uniform open gaps for positive parent interactions]
    \label{thm:compact_block_open_positive_gap}
    \lean{MPSTensor.exists_uniform_block_openParentInteraction_gap_of_compact,
        MPSTensor.exists_uniform_block_openParentInteraction_gap_of_compact_isNormalTensor}
    \leanok
    Under the hypotheses of Lemma~\ref{thm:compact_block_open_gap},
    let $h_x$ be a continuous family of positive range-$R$ interactions
    with $\ker h_x=G_R(B_x)$. Then there is $\delta>0$ such that
    $\delta\|v\|\leq\|H^{\mathrm{open}}_N(h_x)v\|$ for every $x\in K$,
    every $N\geq R$, and every $v\perp\ker H^{\mathrm{open}}_N(h_x)$.
    The same conclusion holds under the normalized normal-block hypotheses
    of Lemma~\ref{thm:compact_normal_block_open_gap} at $R\geq3D_*^5$.
    Nonzero weights need not be continuous.
\end{lemma}
\begin{proof}\leanok
    \uses{thm:compact_block_open_gap,thm:compact_normal_block_open_gap,
        thm:block_ground_space_continuity,thm:block_word_span_propagation,
        thm:cpsv_bnt_sharp_span_ambient_cap,thm:compact_open_positive_parent_gap}
    The local ground projections are continuous at range $R$.
    The canonical open gaps are uniform on $K$ and all $N\geq R$ by
    the two preceding compact results. Compact positive comparison gives
    one $\kappa>0$ with $\kappa h_R(B_x)\leq h_x$ on $K$.
    Summing translates transfers the canonical gap with factor $\kappa$.
    A pointwise upper comparison identifies the kernels.
\end{proof}

\begin{lemma}[Open-chain gaps at a prescribed simultaneous injectivity range]
    \label{thm:block_open_gap_subsequence_simultaneous_injectivity}
    \lean{MPSTensor.%
        exists_openParentHamiltonianES_toTensorFromBlocks_gap_subsequence_of_wordTupleSpanTop}
    \leanok
    Let $A^0,\ldots,A^{r-1}$ have positive physical and bond dimensions.
    Suppose that their simultaneous length-$S$ word tuples span the product
    matrix algebra, where $S>0$. Set $B=\bigoplus_j\mu_jA^j$, with
    every $\mu_j\neq0$. Combining the block-injective local spaces of
    \cite[Section~IV.C]{Cirac2021Matrix} with the martingale estimate of
    \cite[Theorem~2.1(ii) and Section~6]{Nachtergaele1996Spectral}, for any
    fixed $R\geq S+1$ there are $p>0$, with $R\leq2p$, and $\delta>0$ such that
    \begin{align}
        \delta\|v\|\leq\|H^{\mathrm{open}}_{pM,R}(B)v\|\notag
    \end{align}
    for every $M\geq2$ and every $v\perp G_{pM}^{\mathrm{ES}}(B)$.
\end{lemma}
\begin{proof}
    \leanok
    \uses{thm:block_word_span_primitive_normalization,
        thm:block-word-span-separation,
        thm:primitive_block_sum_open_gap_subsequence,
        thm:block_open_kernel_simultaneous_injectivity,
        lem:parent-finite-interval-comparison,lem:open_parent_range_comparison,
        thm:ph_block_replacement_parent_equality}
    Normalizing the blocks preserves all local spaces and the span at $S$,
    and the simultaneous span separates the blocks. Choosing $p\geq R$ and a
    uniform open-chain gap $\gamma>0$ at range $2p$ on every length $pM$,
    $M\geq2$, gives range-$R$ and range-$2p$ open kernels that both equal
    $G_{pM}^{\mathrm{ES}}$. On the interval of length $2p$, finite-dimensional
    comparison gives $\kappa h_{2p}\leq H^{\mathrm{open}}_{2p,R}$ with
    $\kappa>0$. Open-range comparison transfers the gap to
    $\delta=\kappa\gamma/(2p-R+1)>0$ at range $R$. Equality of all local
    spaces carries this bound to the original block sum.
\end{proof}

\begin{lemma}[Strict Knabe intervals at a simultaneous injectivity range]
    \label{thm:block_strict_knabe_window_simultaneous_injectivity}
    \lean{MPSTensor.%
        exists_strict_openParentHamiltonianES_toTensorFromBlocks_window_of_wordTupleSpanTop}
    \leanok
    Let $A^0,\ldots,A^{r-1}$ have positive physical and bond dimensions.
    Suppose that their simultaneous length-$S$ word tuples span the product
    matrix algebra, where $S>0$. Set $B=\bigoplus_j\mu_jA^j$, with
    every $\mu_j\neq0$. For any fixed $R\geq S+1$, there are $m\geq R$ and
    $\gamma>0$ satisfying the strict finite-range threshold $(R-1)^2<m\gamma$
    of Lemma~\ref{thm:parent_hamiltonian_gap_from_open_knabe_block} (with
    the coefficient discussed in
    \path{docs/paper-gaps/knabe88_finite_range_coefficient.tex}) such that
    \begin{align}
        \gamma\|v\|\leq\|H^{\mathrm{open}}_{m+R-1,R}(B)v\|\notag
    \end{align}
    for every $v\perp G_{m+R-1}^{\mathrm{ES}}(B)$.
\end{lemma}
\begin{proof}
    \leanok
    \uses{thm:block_open_gap_subsequence_simultaneous_injectivity}
    Let $p>0$ and $\delta>0$ give the open-chain gap on all lengths
    $pM$, $M\geq2$, and choose $n\in\mathbb N$ with $(R-1)^2<n\delta$.
    With $M=n+2R+2$ and $m=pM+1-R$, one has $m\geq R$, $n\leq m$,
    and $m+R-1=pM$, so $\gamma=\delta$ has both the required gap and
    $(R-1)^2<m\gamma$.
\end{proof}

\begin{lemma}[A uniform gap at the simultaneous injectivity range]
    \label{thm:block-parent-uniform-gap-simultaneous-injectivity}
    \lean{MPSTensor.exists_parentHamiltonianES_toTensorFromBlocks_uniform_gap_of_wordTupleSpanTop}
    \leanok
    Let $A^0,\ldots,A^{r-1}$ have positive physical and bond dimensions.
    Suppose their word tuples of a given positive length $S$ span the full
    product matrix algebra. Let $B=\bigoplus_{j=0}^{r-1}\mu_jA^j$, where every
    $\mu_j$ is nonzero. Realizing the block-injective gap assertion of
    \cite[Section~IV.C]{Cirac2021Matrix} at the simultaneous injectivity range,
    for each interaction range $R\geq S+1$ there is a constant $\gamma_R>0$
    such that
    \begin{align}
        \gamma_R\lVert v\rVert\leq\lVert H_{R,N}v\rVert
    \end{align}
    for every chain length $N\in\mathbb N$ and every
    $v\perp\ker H_{R,N}$, where $H_{R,N}$ is the canonical periodic parent
    Hamiltonian of $B$.
\end{lemma}
\begin{proof}
    \leanok
    \uses{thm:block_word_span_primitive_normalization, thm:block-word-span-separation,
        thm:primitive_block_sum_parent_gap,
        thm:parent-gap-fixed-range-transfer, thm:parent_hamiltonian_eventual_gap,
        thm:ph_tp_word_span_propagation, thm:ph_tp_block_open_ground_space,
        thm:ph_block_replacement_parent_equality}
    Normalizing the blocks retains simultaneous injectivity at $S$ and
    all local ground spaces, and simultaneous injectivity forces distinct
    blocks to be inequivalent. The grouped martingale estimate and finite-range
    periodic comparison provide a gapped interaction range $W\geq2R$.
    Let $B'$ be the weighted sum of the normalized blocks. Trace preservation
    propagates simultaneous injectivity from $S$ to every larger length, so,
    since $R\geq S+1$, the open-chain ground-space theorem gives
    \begin{align}
        \ker H^{\mathrm{open}}_{W,R}(B')=G_W^{\mathrm{ES}}(B').\notag
    \end{align}
    Two-sided comparison of the
    periodic range-$R$ and range-$W$ Hamiltonians therefore transfers the gap
    to $R$. Taking a minimum with the finitely many remaining finite-volume
    gaps makes the bound uniform over every $N$. Finally, normalization left
    the original parent Hamiltonians unchanged.
\end{proof}

\begin{lemma}[Canonical parent Hamiltonians and normal representatives]
    \label{thm:canonical-parent-hamiltonian-representatives}
    \lean{MPSTensor.CPSVCanonicalFormData.parentHamiltonianES_eq_representatives}
    \leanok
    Let $A$ be a tensor in canonical form, with distinct normal
    representatives $B_0,\ldots,B_{g-1}$, and set $B=\bigoplus_j B_j$.
    For every positive interaction range $R$ and every chain length $N$,
    the canonical periodic parent Hamiltonians satisfy
    \begin{align}
        H_{R,N}(A)=H_{R,N}(B).\notag
    \end{align}
    This equality includes the multiplicity copies and the ambient
    reconstruction of the original tensor.
\end{lemma}
\begin{proof}\leanok
    \uses{thm:ph_canonical_representative_local_spaces,
        thm:block_sum_local_ground_space}
    The local ground spaces satisfy
    \begin{align}
        G_R(A)=\bigvee_jG_R(B_j)=G_R(B).\notag
    \end{align}
    Their orthogonal complementary projections therefore agree. Translating
    and summing those projections gives the same Hamiltonian.
\end{proof}

\begin{lemma}[Canonical gap for every canonical decomposition]
    \label{thm:canonical-decomposition-gap-simultaneous-injectivity}
    \lean{MPSTensor.CPSVCanonicalFormData.exists_parentHamiltonianES_uniform_gap_of_wordTupleSpanTop}
    \leanok
    Let $A$ be a tensor in canonical form with positive physical dimension.
    Fix any canonical decomposition of
    $A$, and let $B_0,\ldots,B_{g-1}$ be its distinct normal representatives.
    Let $S$ be a positive integer such that the simultaneous length-$S$ word
    tuples of $B_0,\ldots,B_{g-1}$ span the product of the block matrix
    algebras. Then, for every $R\geq S+1$, there exists $\gamma>0$ such that,
    for every chain length $N$ and every vector $v\perp\ker H_{R,N}(A)$,
    \begin{align}
        \gamma\|v\|\leq\|H_{R,N}(A)v\|.\notag
    \end{align}
\end{lemma}
\begin{proof}\leanok
    \uses{thm:canonical-parent-hamiltonian-representatives,thm:block-parent-uniform-gap-simultaneous-injectivity}
    For $B=\bigoplus_jB_j$, its blocks satisfy the hypotheses of the uniform
    gap at the simultaneous injectivity range, so $H_{R,N}(B)$ has a positive
    uniform gap. Since $H_{R,N}(A)=H_{R,N}(B)$ for every $N$, the same
    constant works for $A$, without changing the range or the chain length.
\end{proof}

\begin{lemma}[Canonical gap at a given simultaneous injectivity length]
    \label{thm:canonical-gap-simultaneous-injectivity}
    \lean{MPSTensor.IsCPSVCanonicalForm.exists_bnt_uniform_gap_of_wordTupleSpanTop}
    \leanok
    Let $A$ be a tensor in canonical form with positive physical dimension.
    Including canonical forms with repeated blocks as in
    \cite[Section~IV.C]{Cirac2021Matrix}, one may choose a basis of normal
    tensors $B_0,\ldots,B_{g-1}$ once, before $S$ and $R$, such that for every
    positive integer $S$ whose simultaneous length-$S$ word tuples span the
    product of the block matrix algebras and every $R\geq S+1$, there exists
    $\gamma>0$ such that, for every chain length $N$ and every vector
    $v\perp\ker H_{R,N}(A)$,
    \begin{align}
        \gamma\|v\|\leq\|H_{R,N}(A)v\|.\notag
    \end{align}
\end{lemma}
\begin{proof}\leanok
    \uses{thm:canonical-decomposition-gap-simultaneous-injectivity}
    Fixing one canonical decomposition of $A$ and taking its distinct normal
    representatives yields a basis of normal tensors for $A$ to which the
    preceding theorem applies for every $S$ and every $R\geq S+1$.
\end{proof}

\begin{lemma}[A dimension bound for canonical parent gaps]
    \label{thm:canonical_parent_gap_bounded_range}
    \lean{MPSTensor.IsCPSVCanonicalForm.%
        exists_parentHamiltonianES_uniform_gap_of_three_bondDim_pow_five_le}
    \leanok
    Let $A$ be a tensor in canonical form with positive physical dimension
    and positive bond dimension $D$. For every $R\geq3D^5$, there is
    $\gamma>0$ (depending on $A$ and $R$, but not on $N$) such that
    \begin{align}
        \gamma\|v\|\leq\|H_{R,N}(A)v\|\notag
    \end{align}
    for every $N\in\mathbb N$ and every $v\perp\ker H_{R,N}(A)$.
    Combining the blocking estimate of
    \cite[lines~317--345]{Cirac2016MPDO_arXiv} with the gap assertion of
    \cite[Section~IV.C]{Cirac2021Matrix}, this sufficient range need not be
    optimal: the sharper sufficient condition $R\geq S+1$ is recorded in
    \cite{gap:cpgsv21_block_parent_interaction_range}, while
    \cite{gap:cpgsv21_short_range_parent_gap} shows that the unrestricted
    assertion fails at some shorter ranges.
\end{lemma}
\begin{proof}
    \leanok
    \uses{thm:canonical_bounded_range_ground_space,
        thm:canonical-decomposition-gap-simultaneous-injectivity}
    For any canonical decomposition of $A$, the dimension estimate gives a
    simultaneous injectivity length $S>0$ for its distinct normal
    representatives with $S+1\leq3D^5\leq R$, and the canonical gap at that
    simultaneous injectivity length gives the assertion.
\end{proof}

\begin{lemma}[Comparison of positive parent interactions]
    \label{thm:parent-interaction-comparison}
    \lean{MPSTensor.IsParentInteraction.exists_pos_comparison}
    \leanok
    Let $h$ be a positive interaction on $R$ sites whose kernel is the local
    MPS space $G_R(A)$, in the sense of
    \cite[Section IV.C]{Cirac2021Matrix}. Write $h_R=I-P_{G_R(A)}$ for the
    canonical parent projection. There are constants $\kappa,C>0$ such that
    \begin{align}
        \kappa h_R\leq h\leq Ch_R.
    \end{align}
\end{lemma}
\begin{proof}\leanok
    \uses{lem:positive-kernel-projection-comparison}
    Applying the finite-dimensional comparison with the projection onto the
    orthogonal complement of the kernel of $h$, which is $h_R$, gives the two
    bounds.
\end{proof}

\begin{definition}[Periodic sum of a fixed interaction]
    \label{def:periodic-fixed-interaction}
    \lean{MPSTensor.periodicLocalInteractionES,MPSTensor.periodicInteractionHamiltonianES}
    \leanok
    Let $h$ act on the physical space of $R$ consecutive sites. For
    $R\leq N$, let $U_i$ be the unitary which places the cyclic interval
    of length $R$ starting at $i$ first and the remaining sites second.
    Define
    \begin{align}
        h_i    & =U_i^*(h\otimes I)U_i,      &
        H_N(h) & =\sum_{i=0}^{N-1}h_i.\notag
    \end{align}
    When $R>N$, set each translated interaction, and hence their sum, to
    zero. This agrees with the convention for the canonical parent Hamiltonian.
\end{definition}

\begin{lemma}[Order and positivity of periodic interaction sums]
    \label{thm:periodic-fixed-interaction-order}
    \lean{MPSTensor.periodicLocalInteractionES_mono,
        MPSTensor.periodicLocalInteractionES_smul,
        MPSTensor.periodicInteractionHamiltonianES_zero,
        MPSTensor.periodicInteractionHamiltonianES_mono,
        MPSTensor.periodicInteractionHamiltonianES_smul,
        MPSTensor.periodicInteractionHamiltonianES_isPositive}
    \leanok
    Cyclic extension and its periodic sum preserve operator inequalities
    and commute with scalar multiplication. In particular, for every
    volume $N$,
    \begin{align}
        h\leq k & \quad\Longrightarrow\quad H_N(h)\leq H_N(k),\notag                      \\
        H_N(ch) & =cH_N(h),                                          & H_N(0) & =0,\notag \\
        h\geq0  & \quad\Longrightarrow\quad H_N(h)\geq0.\notag
    \end{align}
    The scalar identity holds for every complex number $c$.
\end{lemma}
\begin{proof}\leanok
    \uses{def:periodic-fixed-interaction,thm:ph_spectator_operator_order}
    Extension by the identity preserves order and scalar multiplication,
    as does unitary conjugation. Summing over the cyclic positions preserves
    these relations. Positivity is the comparison with zero.
\end{proof}

\begin{lemma}[Canonical parent projections give the canonical Hamiltonian]
    \label{thm:periodic-fixed-interaction-canonical}
    \lean{MPSTensor.periodicLocalInteractionES_parentInteractionES,
        MPSTensor.periodicInteractionHamiltonianES_parentInteractionES}
    \leanok
    For a tensor $A$ and a positive interaction range $R$, the periodic
    sum of the canonical parent projection $h_R=I-P_{G_R(A)}$ is exactly the
    canonical parent Hamiltonian, at every volume:
    \begin{align}
        H_N(h_R)=H_{R,N}(A).\notag
    \end{align}
\end{lemma}
\begin{proof}\leanok
    \uses{def:periodic-fixed-interaction,thm:cyclic_active_block_local_conjugacy,lem:full-interval-parent-interaction}
    In cyclic active-interval coordinates, each translated canonical term
    is the full-interval parent projection extended by the identity on the
    spectator sites. Conjugating back identifies the terms individually.
    When $R>N$, both definitions give zero.
\end{proof}

\begin{lemma}[Uniform comparison of periodic parent Hamiltonians]
    \label{thm:positive-parent-interaction-periodic-comparison}
    \lean{MPSTensor.IsParentInteraction.exists_pos_periodic_comparison}
    \leanok
    Let $R>0$ and let $h$ be a fixed positive operator with kernel
    $G_R(A)$, and let $H_{R,N}(A)$ be the canonical periodic parent
    Hamiltonian. There are $\kappa,C>0$, independent of $N$, such that
    \begin{align}
        \kappa H_{R,N}(A)\leq H_N(h)\leq C H_{R,N}(A).
    \end{align}
\end{lemma}
\begin{proof}
    \leanok
    \uses{thm:parent-interaction-comparison,thm:periodic-fixed-interaction-order,thm:periodic-fixed-interaction-canonical}
    Lemma~\ref{thm:parent-interaction-comparison} gives $\kappa,C>0$ with
    $\kappa h_R\leq h\leq Ch_R$. Each translated term $h_i$ is $h$ extended by
    the identity and conjugated by the unitary $U_i$, so these inequalities
    hold term by term, $\kappa(h_R)_i\leq h_i\leq C(h_R)_i$. Summing over the
    cyclic translates and using the monotonicity and homogeneity of the
    periodic sum from Lemma~\ref{thm:periodic-fixed-interaction-order} gives
    \begin{align}
        \kappa H_N(h_R)\leq H_N(h)\leq CH_N(h_R).\notag
    \end{align}
    Lemma~\ref{thm:periodic-fixed-interaction-canonical} identifies
    $H_N(h_R)$ with $H_{R,N}(A)$.
\end{proof}

\begin{lemma}[Uniform gap for a fixed positive parent interaction]
    \label{thm:positive-parent-interaction-uniform-gap}
    \lean{MPSTensor.IsParentInteraction.exists_uniform_gap_of_canonical_gap}
    \leanok
    Let $A$ be a matrix product tensor and let $R>0$. Suppose that the periodic
    Hamiltonians formed from the orthogonal projection onto $G_R(A)^\perp$
    have a positive gap uniformly in the chain length. Let $h$ be a fixed positive
    operator on $R$ sites with kernel $G_R(A)$, and let $H_N(h)$ be the sum
    of its cyclic translates. Then there is $\delta>0$ such that, for every chain
    length $N$ and every $v\perp\ker H_N(h)$,
    \begin{align}
        \delta\lVert v\rVert\leq\lVert H_N(h)v\rVert.
    \end{align}
    For $N<R$ both periodic Hamiltonians are defined to be zero.
\end{lemma}
\begin{proof}
    \leanok
    \uses{thm:positive-parent-interaction-periodic-comparison,lem:positive-comparison-kernel-gap}
    Following the projector reduction of \cite[Section~IV.C]{Cirac2021Matrix},
    finite-dimensional positivity gives constants $\kappa,C>0$ such that
    $\kappa h_R\leq h\leq C h_R$, where $h_R=I-P_{G_R(A)}$ is the canonical
    parent projection.
    Cyclic extension and summation give the same inequalities for the periodic
    Hamiltonians, with constants independent of $N$. The two inequalities imply
    equality of the kernels, and restriction to their common orthogonal
    complement transfers a canonical gap $\gamma$ to the gap $\kappa\gamma$.
\end{proof}

\begin{lemma}[Positive parent interactions for every canonical decomposition]
    \label{thm:canonical-decomposition-positive-parent-gap-simultaneous-injectivity}
    \lean{MPSTensor.CPSVCanonicalFormData.exists_parentInteraction_uniform_gap_of_wordTupleSpanTop}
    \leanok
    Let $A$ be a tensor in canonical form with positive physical dimension.
    Fix any canonical decomposition of $A$, and let $B_0,\ldots,B_{g-1}$ be
    its distinct normal representatives.
    Suppose that $S>0$ and the tuples of length-$S$ words span the product of
    the block matrix algebras. For every $R\geq S+1$ and
    every fixed positive operator $h$ on $R$ sites with kernel $G_R(A)$,
    there is $\gamma>0$ such that, for every $N$ and every $v\perp\ker H_N(h)$,
    \begin{align}
        \gamma\lVert v\rVert\leq\lVert H_N(h)v\rVert.
    \end{align}
\end{lemma}
\begin{proof}
    \leanok
    \uses{thm:canonical-decomposition-gap-simultaneous-injectivity,thm:positive-parent-interaction-uniform-gap}
    Applying the canonical projector gap of
    Lemma~\ref{thm:canonical-decomposition-gap-simultaneous-injectivity} at
    range $R$ and then the comparison theorem for the fixed interaction $h$
    yields the bound, with constants depending on $h$ but not on $N$.
\end{proof}

\begin{lemma}[Positive parent interactions at the simultaneous injectivity range]
    \label{thm:canonical-positive-parent-gap-simultaneous-injectivity}
    \lean{MPSTensor.IsCPSVCanonicalForm.exists_bnt_parentInteraction_uniform_gap_of_wordTupleSpanTop}
    \leanok
    Let $A$ be a tensor in canonical form with positive physical dimension.
    One may choose a basis of normal tensors $B_0,\ldots,B_{g-1}$ such that
    the following holds. Suppose that $S>0$ and the tuples of length-$S$ words
    span the product of the block matrix algebras. For
    every $R\geq S+1$ and every fixed positive operator $h$ on $R$ sites with
    kernel $G_R(A)$, there is $\gamma>0$ such that, for every $N$ and every
    $v\perp\ker H_N(h)$,
    \begin{align}
        \gamma\lVert v\rVert\leq\lVert H_N(h)v\rVert.
    \end{align}
    The representatives are chosen once, independently of $S$, $R$, and $h$.
\end{lemma}
\begin{proof}
    \leanok
    \uses{thm:canonical-gap-simultaneous-injectivity,thm:positive-parent-interaction-uniform-gap}
    Applying the canonical projector gap at range $R$ and then the comparison
    theorem for the fixed interaction $h$ yields the bound, with constants
    depending on $h$ but not on $N$.
\end{proof}

\begin{lemma}[A dimension bound for fixed positive parent interactions]
    \label{thm:canonical_positive_parent_gap_bounded_range}
    \lean{MPSTensor.IsCPSVCanonicalForm.%
        exists_parentInteraction_uniform_gap_of_three_bondDim_pow_five_le}
    \leanok
    Let $A$ be a tensor in canonical form with positive physical dimension
    and positive bond dimension $D$. Let $R\geq3D^5$ and let $h$ be any
    fixed positive operator on $R$ sites with $\ker h=G_R(A)$. Combining the
    positive-interaction comparison of \cite[Section~IV.C]{Cirac2021Matrix}
    with the sufficient dimension bound of
    \cite[lines~317--345]{Cirac2016MPDO_arXiv}, there is $\gamma>0$
    (depending on $A$, $R$, and $h$, but not on $N$) such that
    \begin{align}
        \gamma\|v\|\leq\|H_N(h)v\|\notag
    \end{align}
    for every $N\in\mathbb N$ and every $v\perp\ker H_N(h)$.
\end{lemma}
\begin{proof}
    \leanok
    \uses{thm:canonical_parent_gap_bounded_range,
        thm:positive-parent-interaction-uniform-gap}
    The canonical projector interaction has a uniform positive periodic
    gap at range $R$. Comparison with the fixed interaction $h$ preserves
    the kernel and gives a positive gap independent of $N$.
\end{proof}

\section{Prescribed gaps at sufficiently long interaction ranges}

Write $H_{R,N}$ for the periodic parent Hamiltonian on $N$ sites with
interaction range $R$. The local interactions in this section are orthogonal
projectors with coefficient one. The range increases with the requested gap; no rescaling
of the Hamiltonian is made. The bounds follow from the projector-decay and
martingale estimates of~\cite[Theorem~2.1 and Section~6]{Nachtergaele1996Spectral}
and the finite-range comparison. They do not improve the gap at a prescribed
shorter interaction range.

\begin{lemma}[A prescribed periodic gap from an open-chain estimate]
    \label{thm:parent-prescribed-gap-window}
    \lean{MPSTensor.parentHamiltonianES_gap_ge_one_sub_of_open_gap}
    \leanok
    \uses{def:parent_hamiltonian_es,def:open_parent_hamiltonian_es}
    Let $0<\eta<1$, $1\le R\le m$, and
    $2(R-1)^2\le m\eta$. Suppose the open parent Hamiltonian on
    $m+R-1$ sites has gap at least $1-\eta/2$. Then, for every
    $N\ge2m$ and every $v\perp\ker H_{R,N}$,
    \begin{align}
        (1-\eta)\|v\| & \le\|H_{R,N}v\|.
    \end{align}
\end{lemma}
\begin{proof}
    \leanok
    \uses{thm:parent_hamiltonian_gap_from_open_knabe_block}
    Put $l=R-1$ and $\gamma=1-\eta/2$. The finite-range comparison
    gives $\delta=(m\gamma-l^2)/(m-l)$. Its denominator is positive,
    and
    \begin{align}
        (m-l)(\delta-(1-\eta))
         & =m\eta/2-l^2+(1-\eta)l \\
         & \ge0.
    \end{align}
    Thus $\delta\ge1-\eta>0$, which also verifies the positive-numerator
    hypothesis of the comparison.
\end{proof}

\begin{lemma}[Preservation of finite-interval gaps]
    \label{thm:open_interval_gap_transport}
    \lean{MPSTensor.openPrefixParentHamiltonianES_norm_gap_of_local_gap,
        MPSTensor.openSuffixParentHamiltonianES_norm_gap_of_local_gap,
        MPSTensor.openSuffixParentHamiltonianES_conj_cyclicActiveBlock}
    \leanok
    Let $R>0$. If the range-$R$ parent Hamiltonian on $n$ sites has norm
    gap $\kappa>0$ on its kernel complement, its extension to the prefix
    of an $(n+q)$-site chain has the same lower bound.
    Likewise, if $R\leq W\leq n\leq N$ and the range-$R$ Hamiltonian
    on $W$ sites has gap at least $\kappa$, its interaction sum on
    $[n-W,n)$, acting on the $N$-site chain, has gap at least $\kappa$.
    The physical coordinate isometry identifies this interval interaction
    sum with the $W$-site Hamiltonian tensored with the identity on the
    remaining sites. These are the finite-interval gaps used in
    condition C2 of~\cite{Nachtergaele1996Spectral}.
\end{lemma}
\begin{proof}\leanok
    \uses{lem:positive-kernel-projection-comparison,
        lem:right-spectator-kernel-projection,
        thm:ph_spectator_operator_order,
        thm:open-hamiltonian-right-spectator-conjugacy,
        thm:cyclic_active_block_local_conjugacy}
    For the finite-interval operator $T$ and its ground projection $P$,
    the given gap gives $\kappa(I-P)\leq T$. Extension by free sites
    preserves this inequality. The kernel projection extends to
    $P\otimes I$, and the physical coordinate isometry carries both
    the Hamiltonian and this projection to the chosen interval.
    Thus the extended operator dominates $\kappa$ times its excitation
    projection. On the kernel complement, the quadratic-form bound and
    Cauchy--Schwarz give the asserted norm bound.
\end{proof}

\begin{lemma}[Sharp gap comparison from two ground spaces]
    \label{thm:sharp_two_ground_projection_gap}
    \lean{FrustrationFree.norm_gap_of_two_ground_projection_bounds}
    \leanok
    Let $H$ be a linear operator on a finite-dimensional complex Hilbert
    space, let $W=\ker H$, and suppose $W\subseteq U\cap V$. Write
    $P_W,P_U,P_V$ for the orthogonal projections onto these subspaces.
    Suppose $\kappa\geq0$, $c<1$,
    $\|P_W-P_UP_V\|\leq c$, and
    $\kappa((I-P_U)+(I-P_V))\leq H$ in the order of self-adjoint
    operators. Then $\kappa(1-c)\|v\|\leq\|Hv\|$ for every $v\perp W$.
    This is a two-subspace consequence of the ground-projection estimate
    in \cite[Section~6, the three-interval commutation lemma,
        part~(ii)]{Nachtergaele1996Spectral}.
\end{lemma}
\begin{proof}\leanok
    \uses{thm:friedrichs_ground_space_excitation_anticommutator,
        lem:quadratic_form_to_norm_bound,lem:positive_gap_transfer_loewner}
    If $z\in U\cap V$ and $z\perp W$, the projection defect gives
    $\|z\|\leq c\|z\|$, hence $z=0$. Subtracting $P_Wv$ from an
    arbitrary $v\in U\cap V$ proves $W=U\cap V$.
    Put $L=(I-P_U)+(I-P_V)$. The projection defect bounds the
    Friedrichs overlap by $c$. The complementary-projection
    anticommutator inequality gives $L^2\geq(1-c)L$, and
    $\ker L=U\cap V=W$. The finite-dimensional spectral theorem gives
    $(1-c)\|v\|^2\leq\langle Lv,v\rangle$ for $v\perp W$.
    Applying $\kappa L\leq H$ and Cauchy--Schwarz yields
    $\kappa(1-c)\|v\|^2\leq\langle Hv,v\rangle
        \leq\|Hv\|\|v\|$. Cancel $\|v\|$ when $v\ne0$.
\end{proof}

\begin{lemma}[Combining gaps on two open intervals]
    \label{thm:open_parent_overlapping_interval_gap}
    \lean{MPSTensor.openPrefixParentHamiltonianES_le_openParentHamiltonianES,
        MPSTensor.openSuffixParentHamiltonianES_le_openParentHamiltonianES,
        MPSTensor.openParentHamiltonianES_norm_gap_of_overlapping_interval_gaps}
    \leanok
    Let $H$ be the range-$R$ parent Hamiltonian on an open chain of length
    $N$, and let $H_L$ and $H_T$ be its interaction sums on a prefix of
    length $n$ and on a terminal interval of length $W$, respectively.
    Write $P_L$, $P_T$, and $P$ for their orthogonal ground projections.
    Suppose $\kappa\geq0$, and suppose both interval Hamiltonians have a
    norm gap at least $\kappa$ on their respective kernel complements.
    If $\|P-P_LP_T\|\leq1/2$, then
    $(\kappa/16)\|v\|\leq\|Hv\|$ for every $v\perp\ker H$.
    The interval interaction sums always satisfy $H_L\leq H$ and
    $H_T\leq H$.
\end{lemma}
\begin{proof}\leanok
    \uses{lem:positive-kernel-projection-comparison}
    Each interval interaction sum is a subfamily of the positive terms of
    $H$, so $H_L+H_T\leq2H$. The interval gaps give
    $\kappa(P_L^\perp+P_T^\perp)\leq H_L+H_T$.
    For $v\perp\ker H$, the projection defect gives
    $\|P_LP_Tv\|\leq\|v\|/2$. Decomposing
    $v=P_LP_Tv+P_LP_T^\perp v+P_L^\perp v$ yields
    $\|v\|\leq2(\|P_L^\perp v\|+\|P_T^\perp v\|)$.
    Squaring gives
    $\|v\|^2\leq8(\|P_L^\perp v\|^2+\|P_T^\perp v\|^2)$.
    Hence $\kappa\|v\|^2\leq16\langle Hv,v\rangle$.
    Cauchy--Schwarz and cancellation of $\|v\|$ prove the claim.
\end{proof}

\begin{lemma}[Prescribed open-chain projector defects]
    \label{thm:primitive-prescribed-projector-defect}
    \lean{MPSTensor.IsPrimitiveMPS.eventually_openChain_groundProjection_defect_mul_sqrt_lt}
    \leanok
    \uses{def:primitive_mps,def:l_blk_injective,
        def:open_chain_ground_subspaces_es,def:ground_space_es}
    Let $A$ be a normalized primitive tensor with positive bond dimension
    and a positive-definite invariant matrix. For every $a>0$, every
    sufficiently large integer $l$ satisfies $l>1$, $l$-block injectivity,
    and a number $\epsilon_l\ge0$ with $\epsilon_l\sqrt{l+1}<a$
    such that, for every prefix length $K\ge0$,
    \begin{align}
        \left\|G_{\lbrack K,K+l+1)}G_{\lbrack 0,K+l)}-G_{\lbrack 0,K+l+1)}\right\|
         & \le\epsilon_l.
    \end{align}
    Here $G_I$ denotes the projection onto the matrix-product boundary
    space on the indicated interval, extended by the identity on its
    complement in the chain of length $K+l+1$.
\end{lemma}
\begin{proof}
    \leanok
    \uses{thm:primitive_fnw_openchain_defect_geometric,thm:wordspan_blk_inj_succ}
    The geometric estimate gives
    $\epsilon_l=c\lambda^l(1+c\lambda^l)/(1-c\lambda^l)$ beyond the
    injectivity threshold, with $c>0$ and $0<\lambda<1$.
    Since $\sqrt{l+1}\lambda^l\to0$ and $c\lambda^l\to0$,
    \begin{align}
        \epsilon_l\sqrt{l+1}
         & =\bigl(\sqrt{l+1}\lambda^l\bigr)
        \frac{c(1+c\lambda^l)}{1-c\lambda^l}\longrightarrow0.
    \end{align}
    Choose $l$ beyond the injectivity threshold, above one, and sufficiently
    large that $c\lambda^l<1$ and this product is below $a$.
\end{proof}

\begin{lemma}[Prescribed gaps for primitive tensors]
    \label{thm:primitive-parent-prescribed-gap}
    \lean{MPSTensor.IsPrimitiveMPS.eventually_parentHamiltonianES_gap_ge_one_sub}
    \leanok
    \uses{def:primitive_mps,def:parent_hamiltonian_es}
    Let $A$ be a normalized primitive tensor with positive bond dimension
    and a positive-definite invariant matrix. For every $0<\eta<1$,
    every sufficiently large integer $l$ has the following property.
    If $m\ge l+1$, $2l^2\le m\eta$, and $N\ge2m$, then
    \begin{align}
        (1-\eta)\|v\| & \le\|H_{l+1,N}v\|
    \end{align}
    for every $v\perp\ker H_{l+1,N}$. One may take
    $m=\max\{l+1,\lceil2l^2/\eta\rceil\}$.
    The threshold for $l$ depends on the tensor; no dimension-only
    bound on this threshold is asserted.
\end{lemma}
\begin{proof}
    \leanok
    \uses{thm:primitive-prescribed-projector-defect,
        thm:literal_fixed_volume_martingale_c3,
        thm:open_parent_hamiltonian_fixed_c3_gap,
        thm:parent-prescribed-gap-window}
    The geometric projector estimate gives
    $\epsilon_l=c\lambda^l(1+c\lambda^l)/(1-c\lambda^l)$ and
    $\epsilon_l\sqrt{l+1}\to0$. Choose $l$ above the injectivity
    threshold and sufficiently large that
    $0\le s_l=\epsilon_l\sqrt{l+1}<\eta/4$. The open-chain
    martingale coefficient satisfies
    \begin{align}
        (1-s_l)^2=1-2s_l+s_l^2\ge1-\eta/2.
    \end{align}
    Apply Lemma~\ref{thm:parent-prescribed-gap-window} with $R=l+1$.
\end{proof}

\begin{lemma}[Prescribed gaps for primitive block sums]
    \label{thm:primitive-block-parent-prescribed-gap}
    \lean{MPSTensor.eventually_toTensorFromBlocks_parentHamiltonianES_gap_ge_one_sub}
    \leanok
    \uses{def:primitive_mps,def:parent_hamiltonian_es,def:gauge_phase_equiv}
    Let $B=\bigoplus_j\mu_jA_j$, where the $A_j$ are normalized
    primitive tensors with common positive physical dimension, positive
    bond dimensions, and positive-definite invariant matrices. Suppose
    $\mu_j\ne0$ and no two blocks are related by similarity and
    multiplication by any nonzero complex scalar.
    For every $0<\eta<1$ and every sufficiently large integer $p$,
    the periodic range-$2p$ parent Hamiltonian has gap at least
    $1-\eta$ at every sufficiently large chain length.
\end{lemma}
\begin{proof}
    \leanok
    \uses{thm:ph_primitive_blocks_open_kernel,
        thm:ph_block_diagonal_projector_decay,
        thm:grouped_c3_from_projector_defect,
        thm:grouped_open_parent_gap,thm:parent-prescribed-gap-window}
    Choose $p$ above the open-kernel threshold and large enough that
    the three-interval projector error is at most
    $\epsilon=\eta/(4\sqrt2)$. The grouped open-chain bound on
    $pM$ sites has coefficient
    $(1-\epsilon\sqrt2)^2=(1-\eta/4)^2\ge1-\eta/2$.
    Choose $M$ large enough that $m=pM-2p+1$ satisfies
    $2p\le m$ and $2(2p-1)^2\le m\eta$.
    The active comparison window has $m+2p-1=pM$ sites, so
    Lemma~\ref{thm:parent-prescribed-gap-window} applies for every
    $N\ge2m$.
\end{proof}

\begin{lemma}[Prescribed gaps without normalization]
    \label{thm:normal-parent-prescribed-gap}
    \lean{MPSTensor.eventually_parentHamiltonianES_gap_ge_one_sub_of_isNormal,
        MPSTensor.eventually_toTensorFromBlocks_parentHamiltonianES_gap_ge_one_sub_of_isNormal}
    \leanok
    \uses{def:normal,def:parent_hamiltonian_es}
    The conclusion of Lemma~\ref{thm:primitive-parent-prescribed-gap}
    holds for every normal tensor of positive bond dimension, without
    a normalization hypothesis. Likewise, the conclusion of
    Lemma~\ref{thm:primitive-block-parent-prescribed-gap} holds for
    normal blocks with nonzero weights and positive physical and bond
    dimensions, provided no two blocks are related by similarity and
    multiplication by any nonzero complex scalar.
\end{lemma}
\begin{proof}
    \leanok
    \uses{thm:normalized_primitive_gauge,
        thm:ph_primitive_representatives_parent_equality,
        thm:primitive-parent-prescribed-gap,
        thm:primitive-block-parent-prescribed-gap}
    Replace the tensor, or each block independently, by its normalized
    primitive representative. The canonical parent Hamiltonians agree
    exactly at every interaction range and every chain length, so the
    prescribed lower bounds transfer unchanged.
\end{proof}

\begin{lemma}[Three normal blocks with the singlet parent interaction]
    \label{lem:three_normal_blocks_singlet_parent}
    \lean{MPSTensor.shortRangeHeisenbergBlock,
        MPSTensor.shortRangeHeisenbergTensor,
        MPSTensor.wordTupleSpanTop_shortRangeHeisenbergBlock_two,
        MPSTensor.not_wordTupleSpanTop_shortRangeHeisenbergBlock_one,
        MPSTensor.shortRangeHeisenbergBlock_normalization,
        MPSTensor.shortRangeHeisenbergBlock_isNormalTensor,
        MPSTensor.shortRangeHeisenbergBlock_isNormal,
        MPSTensor.shortRangeHeisenbergBlock_blocksNotGaugePhaseEquiv,
        MPSTensor.isCPSVBasisOfNormalTensors_shortRangeHeisenbergTensor,
        MPSTensor.mem_groundSpaceES_shortRangeHeisenbergTensor_two_iff,
        MPSTensor.groundSpaceES_shortRangeHeisenbergTensor_two_starProjection,
        MPSTensor.parentInteractionES_shortRangeHeisenbergTensor_two_apply}
    \leanok
    The direct sum of the three bond-one qubit tensors
    $(1,0)$, $(0,1)$ and $(1/\sqrt2,1/\sqrt2)$, with weights one,
    has a basis of three normalized, normal, pairwise inequivalent blocks.
    Its simultaneous injectivity length is exactly two. Its two-site
    local ground space is $\operatorname{Sym}^2(\mathbb C^2)$, and its
    canonical parent interaction is the singlet projection
    $h_2=(I-F)/2$, where $F$ exchanges the two physical sites.
\end{lemma}

\begin{proof}
    \leanok
    The word tuples for $00$, $11$ and $01$ are
    $(1,0,1/2)$, $(0,1,1/2)$ and $(0,0,1/2)$, which span
    $\mathbb C^3$; two one-site words cannot span this space.
    Each scalar block has transfer normalization one. Simultaneous
    spanning separates the blocks and gives eventual independence of
    their periodic vectors. The two-site boundary vectors span
    $|00\rangle$, $|11\rangle$ and $|++\rangle$, exactly the symmetric
    subspace. Its orthogonal complement is the singlet line.
\end{proof}

\begin{lemma}[Failure of the unrestricted parent-Hamiltonian gap assertion]
    \label{thm:short_range_parent_gap_counterexample}
    \lean{SpinChain.oneMagnon_edge_sum,
        ZMod.cycleLaplacian_stdAddChar,
        ZMod.cycleFourierGap_pos,
        ZMod.tendsto_cycleFourierGap_zero,
        MPSTensor.shortRangeHeisenbergMagnon,
        MPSTensor.shortRangeHeisenbergMagnon_ne_zero,
        MPSTensor.parentHamiltonianES_shortRangeHeisenbergMagnon,
        MPSTensor.not_exists_eventual_uniform_gap_shortRangeHeisenbergTensor_two}
    \leanok
    \uses{lem:three_normal_blocks_singlet_parent, def:parent_hamiltonian_es}
    For the preceding tensor, the periodic range-two parent Hamiltonian
    has a nonzero eigenvector of eigenvalue
    $\varepsilon_N=1-\cos(2\pi/N)>0$ for every $N\geq2$.
    These eigenvalues tend to zero. Consequently there are no
    $\delta>0$ and $N_0$ such that
    \[
        \delta\lVert v\rVert\leq\lVert H_{N,2}v\rVert
        \qquad(N\geq N_0,\ v\perp\ker H_{N,2}).
    \]
    Thus the unrestricted assertion in
    \cite[Section~IV.C, lines~2183--2187]{Cirac2021Matrix} is false.
    The sufficient-range theorems remain valid, since this example
    has $R=S=2$. See \cite{gap:cpgsv21_short_range_parent_gap}.
\end{lemma}

\begin{proof}
    \leanok
    \uses{lem:three_normal_blocks_singlet_parent}
    On configurations with one down spin the sum of singlet projections
    acts as $f(j)-(f(j+1)+f(j-1))/2$. The first Fourier character of
    $\mathbb Z/N\mathbb Z$ is therefore an eigenvector with the stated
    eigenvalue. Its coefficient at the zero site is one, so it is nonzero.
    Self-adjointness makes it orthogonal to the kernel. A uniform positive
    norm bound would give $\delta\leq\varepsilon_N$ for every sufficiently
    large $N$, contradicting $\varepsilon_N\to0$.
\end{proof}

\begin{lemma}[Compression through an injective boundary map]
    \label{thm:injective_boundary_compression_norm}
    \lean{ContinuousLinearMap.norm_starProjection_comp_le_inverseGram,
        ContinuousLinearMap.norm_starProjection_comp_sub_smul_le_inverseGram,
        ContinuousLinearMap.tendsto_starProjection_comp_sub_smul_zero_of_inverseGram_bound,
        ContinuousLinearMap.norm_starProjection_apply_sub_smul_le_compression,
        ContinuousLinearMap.norm_inner_sub_scalar_le_compression}
    \leanok
    Let $T:E\to F$ be an injective bounded linear map between real or
    complex Hilbert spaces, with $E$ finite-dimensional. Write
    $K=(T^*T)^{-1}$ and let $P$ be the orthogonal projection onto
    $\operatorname{ran}T$. For every bounded operator $X$ on $F$ and
    every scalar $c$,
    \begin{align}
        \|PXP\|\leq\|K\|\,\|T^*XT\|,
        \qquad
        \|PXP-cP\|\leq\|K\|\,\|T^*XT-cT^*T\|.
        \notag
    \end{align}
    For a family $T_i,X_i,c_i$ indexed by an arbitrary filter, suppose
    $T_i$ is eventually injective and its inverse Gram operators are
    eventually bounded by one constant. Then
    $\|T_i^*X_iT_i-c_iT_i^*T_i\|\to0$ implies
    $\|P_iX_iP_i-c_iP_i\|\to0$. The Hilbert spaces may vary with $i$.
    For any finite-dimensional ground subspace $G$ and unit vector $v\in G$,
    with orthogonal projection $P_G$, one also has
    \begin{align}
        |\langle v,Xv\rangle-c|
        \leq\|P_GXv-cv\|
        \leq\|P_GXP_G-cP_G\|.
        \notag
    \end{align}
\end{lemma}
\begin{proof}\leanok
    Since $P=TKT^*$, one has $PXP=(TK)(T^*XT)(KT^*)$.
    Both outside factors have norm $\sqrt{\|K\|}$. Their product
    therefore gives the first bound. Apply it to $X-cI$ for the second.
    The limiting assertion follows by bounding the compression defect
    by a fixed constant times the pullback defect.
    The unit-vector bounds follow from $P_Gv=v$ and the Cauchy--Schwarz
    inequality.
\end{proof}

\begin{lemma}[Algebra of interior observables]
    \label{thm:bulk_observable_algebra}
    \lean{MPSTensor.bulkObservable_eq_chainWindowOperator,
        MPSTensor.bulkObservable_mul,
        MPSTensor.bulkObservable_conjTranspose,
        MPSTensor.norm_toEuclideanCLM_bulkObservable_le}
    \leanok
    Place a $k$-site observable $X$ between free intervals of lengths
    $\ell,r$, obtaining $X^{(\ell,r)}=I_\ell\otimes X\otimes I_r$.
    Placement preserves products and adjoints. When $k\geq1$, it is
    the usual nonwrapping window inclusion and
    $\|X^{(\ell,r)}\|\leq\|X\|$.
\end{lemma}
\begin{proof}\leanok
    Tensor products multiply factorwise and commute with adjoints.
    The nonwrapping window inclusion is a unital homomorphism preserving
    adjoints between matrix algebras, so it is contractive.
\end{proof}

\begin{lemma}[Ground-space subtraction in the commutator gap estimate]
    \label{thm:ground_projection_commutator_gap}
    \lean{FrustrationFree.energy_gap_sub_ground_projection,
        FrustrationFree.commutator_ground_expectation,
        FrustrationFree.commutator_ground_gap,
        FrustrationFree.commutator_ground_gap_of_ker_eq_span}
    \leanok
    Let $H\geq0$ act on a finite-dimensional complex Hilbert space,
    let $K=\ker H$, and suppose $\gamma\geq0$ and
    $\gamma\|v\|\leq\|Hv\|$ for every $v\perp K$.
    Then $\langle Hv,v\rangle\geq
        \gamma(\|v\|^2-\|P_Kv\|^2)$ for every $v$.
    For every $\psi\in K$ and linear operator $X$,
    \begin{align}
        \langle\psi,X^*[H,X]\psi\rangle
         & =\langle HX\psi,X\psi\rangle
        \geq\gamma\bigl(\|X\psi\|^2-\|P_KX\psi\|^2\bigr).
        \label{eq:ph_ground_commutator}
    \end{align}
    If $K=\C\psi$ and $\|\psi\|=1$, the subtracted term is
    $|\langle\psi,X\psi\rangle|^2$.
    This is the finite-volume estimate for the limiting argument in
    \cite[Section~6]{Nachtergaele1996Spectral}; the full ground-space
    projection is retained as described in
    \cite{gap:nachtergaele96_infinite_volume_ground_projection}.
\end{lemma}
\begin{proof}\leanok
    \uses{lem:positive-kernel-projection-comparison}
    The gap gives $H\geq\gamma P_{K^\perp}$. Pythagoras yields
    $\|P_{K^\perp}v\|^2=\|v\|^2-\|P_Kv\|^2$.
    Since $H\psi=0$, the adjoint identity gives
    $\langle\psi,X^*[H,X]\psi\rangle=\langle HX\psi,X\psi\rangle$.
    Apply the first estimate to $v=X\psi$.
    For a one-dimensional ground space, its projection is
    $P_Kv=\langle\psi,v\rangle\psi$.
\end{proof}

\begin{lemma}[Limiting gap from decay of the ground projection]
    \label{thm:commutator_gap_projection_limit}
    \lean{FrustrationFree.commutator_gap_of_tendsto_ground_projection_zero}
    \leanok
    Let $H_n,X_n$ act on finite-dimensional complex Hilbert spaces which
    may depend on $n$. Suppose that, for all sufficiently large $n$,
    $H_n\geq0$, $\psi_n\in\ker H_n$, and $H_n$ has norm gap at least
    the same constant $\gamma\geq0$ on $(\ker H_n)^\perp$.
    Let $a,e\in\R$. If $\|X_n\psi_n\|^2\to a$,
    $\|P_{\ker H_n}X_n\psi_n\|^2\to0$, and
    $\langle\psi_n,X_n^*[H_n,X_n]\psi_n\rangle\to e$, then
    $\gamma a\leq e$.
\end{lemma}
\begin{proof}\leanok
    \uses{thm:ground_projection_commutator_gap}
    For all sufficiently large $n$, (\ref{eq:ph_ground_commutator})
    gives the required inequality before taking limits.
    Its left side tends to $\gamma(a-0)=\gamma a$ and its right side
    tends to $e$. The order relation on $\R$ is closed.
\end{proof}

\begin{corollary}[Limiting gap for one-dimensional finite-volume ground spaces]
    \label{cor:commutator_gap_rank_one_limit}
    \lean{FrustrationFree.commutator_gap_of_tendsto_mean_zero_of_ker_eq_span}
    \leanok
    Let $H_n,X_n$ act on finite-dimensional complex Hilbert spaces which
    may depend on $n$. Suppose that, for all sufficiently large $n$,
    $H_n\geq0$, $\|\psi_n\|=1$, $\ker H_n=\C\psi_n$, and $H_n$ has
    norm gap at least the same constant $\gamma\geq0$ on
    $(\ker H_n)^\perp$. If $\|X_n\psi_n\|^2\to a$,
    $\langle\psi_n,X_n\psi_n\rangle\to0$, and
    $\langle\psi_n,X_n^*[H_n,X_n]\psi_n\rangle\to e$, then
    $\gamma a\leq e$.
    Finite-volume uniqueness is a hypothesis of this corollary.
\end{corollary}
\begin{proof}\leanok
    \uses{thm:ground_projection_commutator_gap}
    The rank-one projection formula makes the subtracted term the squared
    absolute value of the mean. This tends to zero, so the finite-volume
    estimate passes to the limit.
\end{proof}

\begin{lemma}[Sparse grouped open-chain energy bound]
    \label{thm:sparse_grouped_open_energy_gap}
    \lean{MPSTensor.sparseGroupedOpenParentHamiltonianES,
        MPSTensor.grouped_martingaleDifference_norm_le_of_projector_defect_of_le,
        MPSTensor.sparseGroupedOpenParentHamiltonianES_energy_bound_of_grouped_c3,
        MPSTensor.sparseGroupedOpenParentHamiltonianES_isPositive,
        MPSTensor.sparseGroupedOpenParentHamiltonianES_projection_bound_of_grouped_c3,
        MPSTensor.sparseGroupedOpenParentHamiltonianES_add_terminal_le}
    \leanok
    Let $p\geq1$, $M\geq2$, and $pM\leq N$. Sample the canonical
    range-$2p$ interactions at starts $0,p,\ldots,p(M-2)$, and write
    $H_{\mathrm{sp}}$ for their sum on $N$ sites. Let $U$ be the kernel
    of the full range-$2p$ interaction on the prefix of length $pM$.
    Suppose the grouped martingale projections satisfy
    $\|Q_nE_n\|\leq\epsilon<1/\sqrt2$, with $\epsilon\geq0$.
    Then
    \[
        \langle H_{\mathrm{sp}}v,v\rangle
        \geq(1-\sqrt2\epsilon)^2\|v\|^2
        \qquad(v\perp U),
        \qquad
        H_{\mathrm{sp}}\geq(1-\sqrt2\epsilon)^2(I-P_U).
    \]
    The grouped projector estimate remains valid after adjoining free
    sites on the right. If $N=pM+q$ with $q>0$, adjoining the terminal
    range-$2p$ parent projection $T$ gives
    $H_{\mathrm{sp}}+T\leq H^{\mathrm{open}}_{N,2p}$.
    This is a consequence of
    \cite[Theorem~2.1(ii)]{Nachtergaele1996Spectral}.
\end{lemma}
\begin{proof}\leanok
    \uses{thm:grouped_c3_from_projector_defect,thm:grouped_spectator_norm,
        thm:nachtergaele_c1_c3_full_range_energy_estimate}
    Apply the grouped martingale energy inequality directly to the sampled
    interaction sum. Every sampled term annihilates $U$, so decomposition
    into $U$ and $U^\perp$ gives the operator inequality. Free-site extension
    does not increase the projector norm. When $q>0$, the terminal start
    $N-2p=p(M-2)+q$ is strictly larger than all sampled starts. Thus the
    sampled terms and the terminal term form a subfamily of the full
    positive interaction sum, each counted once.
\end{proof}

\begin{lemma}[Sharp open-chain gap with a residual terminal interval]
    \label{thm:sparse_grouped_terminal_open_gap}
    \lean{MPSTensor.openParentHamiltonianES_gap_of_sparse_grouped_and_terminal}
    \leanok
    Let $p\geq1$, $M\geq2$, $0<q<p$, and $N=pM+q$.
    Assume that every range-$2p$ open Hamiltonian at length at least $2p$
    has its boundary-condition space as kernel. Suppose that the
    three-interval ground-projection defect is at most
    $0\leq\epsilon<1/\sqrt2$ whenever the overlap has at least $p$ sites,
    uniformly in both outer intervals, with the right interval nonempty.
    Then
    \[
        (1-\sqrt2\epsilon)^2(1-\epsilon)\|v\|
        \leq\|H^{\mathrm{open}}_{N,2p}v\|
        \qquad(v\perp\ker H^{\mathrm{open}}_{N,2p}).
    \]
\end{lemma}
\begin{proof}\leanok
    \uses{thm:sparse_grouped_open_energy_gap,
        thm:sharp_two_ground_projection_gap,thm:ph_terminal_interval_kernel,
        thm:ph_prefix_kernel_free_right_sites}
    Let $U$ be the prefix ground space at length $pM$, and let $V$ be
    the terminal ground space on $2p$ sites. Their overlap contains
    $2p-q\geq p$ sites. Put $\kappa=(1-\sqrt2\epsilon)^2\leq1$.
    The sparse energy bound and the terminal projection give
    $H^{\mathrm{open}}_{N,2p}\geq\kappa((I-P_U)+(I-P_V))$.
    The full ground space is contained in both $U$ and $V$, and the
    three-interval estimate bounds its projection defect by $\epsilon$.
    The sharp two-ground-space comparison gives the conclusion.
\end{proof}

\begin{lemma}[Prescribed open-chain gaps for primitive block sums]
    \label{thm:primitive_block_open_prescribed_gap}
    \lean{MPSTensor.eventually_toTensorFromBlocks_openParentHamiltonianES_gap_ge_one_sub}
    \leanok
    Let $A_0,\ldots,A_{r-1}$ be normalized primitive tensors with common
    physical dimension $d\geq1$, positive bond dimensions, and
    positive-definite invariant matrices. Suppose no two distinct blocks
    are related by a virtual similarity and a phase. Put
    $B=\bigoplus_j\mu_jA_j$, with $\mu_j\ne0$ for every $j$.
    For every $0<\eta<1$, there is $p_0$ such that, for every
    $p\geq p_0$ and every $N\geq2p$,
    \[
        (1-\eta)\|v\|\leq\|H^{\mathrm{open}}_{N,2p}(B)v\|
        \qquad(v\perp\ker H^{\mathrm{open}}_{N,2p}(B)).
    \]
    This is a quantitative consequence of
    \cite[Theorem~2.1(ii) and Section~6]{Nachtergaele1996Spectral}.
\end{lemma}
\begin{proof}\leanok
    \uses{thm:sparse_grouped_terminal_open_gap,
        thm:grouped_open_parent_gap,thm:ph_primitive_blocks_open_kernel,
        thm:ph_block_diagonal_projector_decay}
    Choose $\epsilon=\eta/(4\sqrt2)$ and take $p$ above the kernel
    threshold and the projector-decay threshold for every overlap length
    at least $p$. Write $N=pM+q$, with $M\geq2$ and $0\leq q<p$.
    For $q=0$, the grouped gap is $(1-\sqrt2\epsilon)^2$.
    For $q>0$, the residual-interval theorem gives
    $(1-\sqrt2\epsilon)^2(1-\epsilon)$.
    Both coefficients are at least $1-\eta$.
\end{proof}

\begin{lemma}[Prescribed open-chain gaps for normal block sums]
    \label{thm:normal_block_open_prescribed_gap}
    \lean{MPSTensor.eventually_toTensorFromBlocks_openParentHamiltonianES_gap_ge_one_sub_of_isNormal}
    \leanok
    Let $A_0,\ldots,A_{r-1}$ be normal tensors with common physical
    dimension $d\geq1$ and positive bond dimensions. Suppose no two
    distinct blocks are related by similarity and multiplication by a
    nonzero complex scalar, and let $B=\bigoplus_j\mu_jA_j$, with
    $\mu_j\ne0$ for every $j$. For every $0<\eta<1$, there is $p_0$
    such that, for every $p\geq p_0$ and every $N\geq2p$,
    \[
        (1-\eta)\|v\|\leq\|H^{\mathrm{open}}_{N,2p}(B)v\|
        \qquad(v\perp\ker H^{\mathrm{open}}_{N,2p}(B)).
    \]
    Thus the canonical open-chain gap approaches one as the even
    interaction range grows, uniformly over all lengths at least that
    range. This is a quantitative consequence of the projector decay and
    martingale estimates of
    \cite[Theorem~2.1(ii) and Section~6]{Nachtergaele1996Spectral}.
\end{lemma}
\begin{proof}\leanok
    \uses{thm:primitive_block_open_prescribed_gap,
        thm:normal_blocks_primitive_representatives}
    Independent primitive normalization preserves all block ground spaces
    and therefore every combined canonical open Hamiltonian. It also
    preserves inequivalence. Apply the primitive block theorem to the
    normalized representatives and transfer its bound through the exact
    equality of open Hamiltonians.
\end{proof}

\begin{definition}[Insertion expectation]
    \label{def:primitive_insertion_expectation}
    \lean{MPSTensor.observableInsertionExpectation}
    \leanok
    For a tensor $A$, a matrix $\rho$ with nonzero trace, and an observable
    $X$ on $k$ sites, write
    $E_X(Y)=\sum_{\sigma,\tau}X_{\sigma\tau}A^\tau Y(A^\sigma)^*$ and
    $\eta_A(X)=\operatorname{tr}(E_X(\rho))/\operatorname{tr}\rho$.
    This is the local expectation appearing in
    \cite[Section~6]{Nachtergaele1996Spectral}.
\end{definition}

\begin{definition}[Observable in an interior interval]
    \label{def:primitive_bulk_observable}
    \lean{MPSTensor.bulkObservable}
    \leanok
    For an observable $X$ on $k$ sites and nonnegative integers $\ell,r$,
    let $X_{\ell,r}=I_\ell\otimes X\otimes I_r$ on $N=\ell+k+r$ sites,
    with the tensor factors ordered from left to right.
\end{definition}

\begin{lemma}[Boundary compression of an insertion]
    \label{thm:primitive_insertion_boundary_identity}
    \lean{MPSTensor.adjoint_groundSpaceMapES_observable_groundSpaceMapES,
        MPSTensor.physicalObservableTransfer_bulkObservable,
        MPSTensor.adjoint_groundSpaceMapES_bulkObservable_groundSpaceMapES}
    \leanok
    Let $\Gamma_n$ be the Hilbert-space boundary map and let
    $\mathcal R$ denote Gram reshuffling. Then
    \begin{align}
        \Gamma_k^*X\Gamma_k          & =\mathcal R(E_X),           &
        E_{X_{\ell,r}}               & =E^\ell E_XE^r,               \\
        \Gamma_N^*X_{\ell,r}\Gamma_N & =\mathcal R(E^\ell E_XE^r).
    \end{align}
\end{lemma}
\begin{proof}\leanok
    \uses{def:ground_space_map_es, def:gram_reshuffle, def:primitive_bulk_observable}
    Pair both sides of the first identity with matrix-unit boundary vectors.
    The resulting finite sums agree term by term.
    Insertion transfers multiply under consecutive tensor products,
    while the identity insertion on $n$ sites is $E^n$.
    These facts give the remaining identities.
\end{proof}

\begin{lemma}[Primitive insertion limit]
    \label{thm:primitive_insertion_transfer_limit}
    \lean{MPSTensor.fixedPointProj_physicalObservableTransfer_fixedPointProj,
        MPSTensor.IsPrimitiveMPS.transferMap_pow_tendsto_fixedPointProj,
        MPSTensor.IsPrimitiveMPS.physicalObservableTransfer_sandwich_tendsto}
    \leanok
    Let $A$ be a normalized primitive tensor with invariant matrix $\rho$,
    and let $P_\rho(Y)=\operatorname{tr}(Y)\rho/\operatorname{tr}\rho$.
    For every local observable $X$,
    $P_\rho E_XP_\rho=\eta_A(X)P_\rho$ and $E^n\to P_\rho$ in operator norm.
    Along any filter on which $\ell,r\to\infty$,
    $E^\ell E_XE^r\to\eta_A(X)P_\rho$.
\end{lemma}
\begin{proof}\leanok
    \uses{def:primitive_insertion_expectation, thm:primitive_fnw_remainder_power}
    Substitution into the rank-one projection gives the first identity.
    For $n\geq1$, the decomposition
    $E^n=P_\rho+(E-P_\rho)^n$ and decay of the complementary powers
    give the second. Continuity of multiplication gives the insertion limit.
\end{proof}

\begin{lemma}[Limiting boundary compression]
    \label{thm:primitive_insertion_gram_limit}
    \lean{MPSTensor.IsPrimitiveMPS.groundSpaceMapES_bulkObservable_compression_tendsto}
    \leanok
    Under the preceding hypotheses, with $K_\infty=\mathcal R(P_\rho)$,
    $\Gamma_N^*X_{\ell,r}\Gamma_N\to\eta_A(X)K_\infty$ as both
    free intervals tend to infinity.
\end{lemma}
\begin{proof}\leanok
    \uses{thm:primitive_insertion_boundary_identity, thm:primitive_insertion_transfer_limit}
    Apply the continuous linear reshuffling map to the insertion limit.
\end{proof}

\begin{lemma}[Scalar compression on the entire primitive ground space]
    \label{thm:primitive_bulk_ground_compression}
    \lean{MPSTensor.IsPrimitiveMPS.bulkObservable_groundSpace_compression_tendsto_scalar,
        MPSTensor.IsPrimitiveMPS.bulkObservable_groundSpace_compression_tendsto_zero,
        MPSTensor.IsPrimitiveMPS.bulkObservable_groundSpace_projection_eventually_near_expectation}
    \leanok
    Suppose in addition that $\rho$ is positive definite. Let $G_N$ be the
    open-boundary MPS space and $P_N$ its orthogonal projection. Then
    \begin{align}
        \|P_NX_{\ell,r}P_N-\eta_A(X)P_N\| & \longrightarrow0.
    \end{align}
    In particular, centered observables satisfy
    $\|P_NX_{\ell,r}P_N\|\to0$. For every $\varepsilon>0$, eventually
    $\|P_NX_{\ell,r}\psi-\eta_A(X)\psi\|<\varepsilon$ simultaneously
    for all $\psi\in G_N$ with $\|\psi\|\leq1$.
\end{lemma}
\begin{proof}\leanok
    \uses{thm:primitive_insertion_gram_limit, thm:primitive_ground_space_gram_tendsto_fixed_point,
        thm:primitive_ground_space_map_eventual_inverse_gram_bound,
        thm:injective_boundary_compression_norm}
    Subtract $\eta_A(X)\Gamma_N^*\Gamma_N$ from the limiting boundary
    compression. The difference tends to zero. The boundary maps are
    eventually injective and their inverse Gram operators remain bounded.
    The compression estimate therefore gives the operator-norm limit.
    Apply the resulting operator to vectors of norm at most one.
\end{proof}

\begin{lemma}[Centered projection decay and local expectation convergence]
    \label{thm:primitive_bulk_ground_vector_limits}
    \lean{MPSTensor.IsPrimitiveMPS.bulkObservable_groundSpace_projection_norm_sq_tendsto_zero,
        MPSTensor.IsPrimitiveMPS.bulkObservable_groundState_expectation_tendsto}
    \leanok
    With the preceding hypotheses, let $\psi_N$ be any varying family of
    vectors eventually belonging to $G_N$. If $\|\psi_N\|\leq1$ and
    $\eta_A(X)=0$, then $\|P_NX_{\ell,r}\psi_N\|^2\to0$.
    If instead $\|\psi_N\|=1$, then
    $\langle\psi_N,X_{\ell,r}\psi_N\rangle\to\eta_A(X)$.
    These assertions concern one primitive sector; the general pure-GVBS
    infinite-volume statement is discussed in
    \cite{gap:nachtergaele96_infinite_volume_ground_projection}.
\end{lemma}
\begin{proof}\leanok
    \uses{thm:primitive_bulk_ground_compression, thm:injective_boundary_compression_norm}
    The projection norm is bounded by the centered compression norm.
    Squaring preserves its zero limit. For unit vectors, the scalar
    expectation deviation is bounded by the scalar compression norm.
\end{proof}

\begin{lemma}[Positive observable insertions]
    \label{thm:primitive_insertion_positive}
    \lean{MPSTensor.physicalObservableTransfer_posSemidef}
    \leanok
    For any tensor $A$, positive observable $X$ on $k$ sites, and positive
    virtual matrix $\rho$, the insertion $E_X(\rho)$ is positive.
    This is the positivity of the local GVBS contraction in
    \cite[Section~3, equations (3.1)--(3.2b)]{Nachtergaele1996Spectral}.
\end{lemma}
\begin{proof}\leanok
    \uses{def:primitive_insertion_expectation}
    Let $B_{a,(\sigma,b)}=(A^\sigma)_{ab}$, where $\sigma$ runs over
    words of length $k$. The insertion is
    $E_X(\rho)=B(X^{\mathsf T}\otimes\rho)B^*$.
    Transposition, tensor products, and congruence preserve positivity.
\end{proof}

\begin{lemma}[Local state properties and consistency]
    \label{thm:primitive_insertion_local_state}
    \lean{MPSTensor.observableInsertionExpectation_one,
        MPSTensor.observableInsertionExpectation_nonneg,
        MPSTensor.observableInsertionExpectation_add,
        MPSTensor.observableInsertionExpectation_smul,
        MPSTensor.observableInsertionExpectation_appendObservable_one,
        MPSTensor.observableInsertionExpectation_one_appendObservable,
        MPSTensor.IsPrimitiveMPS.observableInsertionExpectation_one,
        MPSTensor.IsPrimitiveMPS.observableInsertionExpectation_nonneg}
    \leanok
    Fix $A$ and a virtual matrix $\rho$ with nonzero trace. For each
    interval length $k$, the functional $\eta_A$ is complex linear.
    If $\rho\geq0$, then $\eta_A(X)\geq0$ for $X\geq0$.
    If $E(\rho)=\rho$, then $\eta_A(I_k)=1$ and adjoining identity
    factors on the right preserves expectations. If
    $\sum_i(A^i)^*A^i=I$, adjoining identity factors on the left
    preserves expectations. Thus a normalized primitive tensor with
    its positive invariant matrix defines compatible states on finite
    intervals, as in
    \cite[Section~3, equations (3.1)--(3.2b)]{Nachtergaele1996Spectral}.
\end{lemma}
\begin{proof}\leanok
    \uses{def:primitive_insertion_expectation, thm:primitive_insertion_positive}
    Linearity follows from linearity of the insertion and the trace.
    The positive numerator and denominator give positivity.
    At a fixed point, $E^m(\rho)=\rho$, so the identity insertion
    $E_{I_m}=E^m$ gives normalization. Consecutive insertions compose,
    giving $E_{X\otimes I_m}=E_XE^m$ and
    $E_{I_m\otimes X}=E^mE_X$. Invariance removes the right factor,
    while trace preservation removes the left factor after taking the
    trace. The primitive assumptions include positivity, invariance,
    nonzero trace, and trace preservation.
\end{proof}

\begin{lemma}[Limiting norm of a primitive-sector excitation]
    \label{thm:primitive_bulk_excitation_norm_limit}
    \lean{MPSTensor.IsPrimitiveMPS.bulkObservable_groundState_norm_sq_tendsto}
    \leanok
    Let $A$ be a normalized primitive tensor with positive-definite invariant
    matrix $\rho$. Let $X$ act on $k$ sites and let $\ell,r\to\infty$
    along any filter. For any family of unit vectors $\psi_N$ eventually
    belonging to the open-boundary MPS space on $N=\ell+k+r$ sites,
    $\|X_{\ell,r}\psi_N\|^2\to\operatorname{Re}\eta_A(X^*X)$.
\end{lemma}
\begin{proof}\leanok
    \uses{thm:primitive_bulk_ground_vector_limits, def:primitive_bulk_observable}
    The interior inclusion preserves products and adjoints, so
    $\langle\psi_N,(X^*X)_{\ell,r}\psi_N\rangle
        =\|X_{\ell,r}\psi_N\|^2$.
    Apply local expectation convergence to $X^*X$ and take real parts.
\end{proof}

\begin{lemma}[Primitive-sector commutator gap reduction]
    \label{thm:primitive_bulk_commutator_gap_reduction}
    \lean{MPSTensor.IsPrimitiveMPS.commutator_gap_of_bulkObservable_energy_limit}
    \leanok
    Let $A$ be a normalized primitive tensor with positive-definite invariant
    matrix $\rho$, and let $X$ be a fixed observable with $\eta_A(X)=0$.
    Consider a sequence $N_n=\ell_n+k+r_n$ with both $\ell_n,r_n\to\infty$.
    Suppose the operators $H_n$ are eventually positive, have kernel equal
    to the open-boundary MPS space, and have a common norm gap
    $\gamma\geq0$ on its orthogonal complement.
    Let $\psi_n$ be unit vectors eventually in that ground space.
    If
    \begin{align}
        \operatorname{Re}\langle X_{\ell_n,r_n}^*
        [H_n,X_{\ell_n,r_n}]\psi_n,\psi_n\rangle & \longrightarrow e,
    \end{align}
    then $\gamma\operatorname{Re}\eta_A(X^*X)\leq e$.
    The energy limit is an explicit hypothesis here; the passage to the
    infinite-volume state in \cite[Section~6]{Nachtergaele1996Spectral}
    remains separate, as discussed in
    \cite{gap:nachtergaele96_infinite_volume_ground_projection}.
\end{lemma}
\begin{proof}\leanok
    \uses{thm:primitive_bulk_excitation_norm_limit, thm:primitive_bulk_ground_vector_limits,
        thm:commutator_gap_projection_limit}
    The excitation norm squared tends to $\operatorname{Re}\eta_A(X^*X)$,
    while the squared norm of its full ground-space projection tends to zero.
    Take limits in the finite-volume commutator gap inequality.
\end{proof}

\begin{lemma}[Interval inclusions and interior observable placement]
    \label{thm:local_interval_bulk_observable}
    \lean{SpinChain.intervalCoordinates_localInclusion_eq_bulkObservable}
    \leanok
    With $I,J$ as above,
    $C_{a-b,b+N+c}(\iota_{I,J}(X))=I_b\otimes C_{a,N}(X)\otimes I_c$.
    Thus the integer-interval inclusion agrees with the interior placement
    used in the local GVBS expectations of
    \cite[Section~3, equations (3.1)--(3.2b)]{Nachtergaele1996Spectral}.
\end{lemma}
\begin{proof}\leanok
    \uses{thm:local_interval_coordinate_inclusion, def:primitive_bulk_observable}
    Write $\sigma=(\sigma_1,\sigma_2,\sigma_3)$ and $\tau=(\tau_1,\tau_2,\tau_3)$
    for the restrictions of configurations to the first $b$ sites, the middle
    $N$ sites, and the last $c$ sites. The interval inclusion formula gives
    \begin{align}
        C_{a-b,b+N+c}(\iota_{I,J}(X))_{\sigma,\tau}
        =\delta_{\sigma_1,\tau_1}\,C_{a,N}(X)_{\sigma_2,\tau_2}\,
        \delta_{\sigma_3,\tau_3},
        \notag
    \end{align}
    which is the entry of $I_b\otimes C_{a,N}(X)\otimes I_c$.
\end{proof}

\begin{lemma}[Norm of a finite local state]
    \label{thm:primitive_insertion_state_norm}
    \lean{MPSTensor.observableInsertionExpectationₗ,
        MPSTensor.observableInsertionExpectationPositive,
        MPSTensor.observableInsertionExpectationPositive_apply,
        MPSTensor.norm_observableInsertionExpectationPositive,
        MPSTensor.norm_observableInsertionExpectation_le,
        MPSTensor.observableInsertionExpectation_conjTranspose}
    \leanok
    Let $\rho\geq0$ satisfy $E_A(\rho)=\rho$ and
    $\operatorname{tr}\rho\neq0$. For each finite interval, the insertion
    expectation $\eta_A$ is a positive continuous linear functional of
    operator norm one. Consequently
    \begin{align}
        |\eta_A(X)| & \leq\|X\|,
                    & \eta_A(X^\dagger) & =\overline{\eta_A(X)}.
    \end{align}
    The adjoint identity requires only $\rho\geq0$; the norm-one assertion
    also uses invariance and nonzero trace. Primitivity is unnecessary.
    This is the finite local state of
    \cite[Section~3, equations (3.1)--(3.2b)]{Nachtergaele1996Spectral}.
\end{lemma}
\begin{proof}\leanok
    \uses{thm:primitive_insertion_local_state, thm:primitive_insertion_positive}
    Positivity and complex linearity give a positive functional on the
    finite matrix algebra. Its norm is its value at the identity, which is
    one. A positive complex-linear functional preserves adjoints.
\end{proof}

\begin{definition}[Finite open interaction sum in matrix coordinates]
    \label{def:finite_open_interaction_matrix}
    \lean{MPSTensor.openInteractionMatrix,
        MPSTensor.openInteractionHamiltonianES_eq_toEuclideanLin_openInteractionMatrix,
        MPSTensor.toEuclideanCLM_openInteractionMatrix_parentInteraction}
    \leanok
    For an arbitrary range-$R$ matrix $h$, let
    $H_N(h)=\sum_{a=0}^{N-R}\iota_a(h)$, where $\iota_a$ places $h$
    on the sites $a,\ldots,a+R-1$ of an $N$-site chain, with the sum
    empty when $N<R$. For $R\geq1$ and $N\geq R$, its Euclidean matrix
    action is the open interaction Hamiltonian. If $h$ is the complement
    projection of $G_R(A)$, it is the canonical open parent Hamiltonian.
    This is the nonwrapping interaction sum in
    \cite[equation (3.12)]{Nachtergaele1996Spectral}.
\end{definition}
\begin{proof}\leanok
    \uses{def:open_parent_hamiltonian_es}
    The admissible starts are exactly $0,\ldots,N-R$. Each matrix inclusion
    acts as the corresponding Euclidean local interaction, so their sums
    agree. Substitute the local complement projection for the canonical case.
\end{proof}

\begin{lemma}[Commutation on disjoint intervals]
    \label{thm:disjoint_interval_observables_commute}
    \lean{MPSTensor.chainWindowOperator_commute_of_disjoint}
    \leanok
    Observables supported on two disjoint nonwrapping intervals of a finite
    chain commute. No self-adjointness assumption is required.
\end{lemma}
\begin{proof}\leanok
    Operators on the two tensor factors commute; each acts as the identity
    on the support of the other.
\end{proof}

\begin{definition}[The fixed local commutator observable]
    \label{def:fixed_local_commutator_observable}
    \lean{MPSTensor.localCommutatorObservable}
    \leanok
    Let $R,k\geq1$, put $t=R-1$ and $w=2t+k$, and let $X$ be supported
    on $k$ sites. On $w$ sites set $B=I_t\otimes X\otimes I_t$ and
    \begin{align}
        Y_h(X) & =B^\dagger\bigl(H_w(h)B-BH_w(h)\bigr).
    \end{align}
    Neither positivity nor self-adjointness of $h$ is needed for this
    definition. It is the fixed local expression used in the limiting
    argument of \cite[Section~6]{Nachtergaele1996Spectral}.
\end{definition}

\begin{lemma}[Finite-range commutator localization]
    \label{thm:finite_open_commutator_locality}
    \lean{MPSTensor.openInteractionMatrix_commutator_chainWindowOperator,
        MPSTensor.openInteractionMatrix_commutator_bulkObservable,
        MPSTensor.chainWindowOperator_adjoint_mul_openInteractionMatrix_commutator,
        MPSTensor.bulkObservable_adjoint_mul_openInteractionMatrix_commutator}
    \leanok
    Let $R,k\geq1$ and $X_N=I_\ell\otimes X\otimes I_r$, where
    $\ell,r\geq t=R-1$ and $N=\ell+k+r$. Then
    \begin{align}
        [H_N(h),X_N]
         & =I_{\ell-t}\otimes[H_{2t+k}(h),B]\otimes I_{r-t}, \\
        X_N^\dagger[H_N(h),X_N]
         & =I_{\ell-t}\otimes Y_h(X)\otimes I_{r-t}.
    \end{align}
    Equivalently, both identities hold for an arbitrary placement of the
    enlarged interval inside a longer chain. These are locality identities;
    no spectral-gap assumption is made.
\end{lemma}
\begin{proof}\leanok
    \uses{def:finite_open_interaction_matrix,
        def:fixed_local_commutator_observable,
        thm:disjoint_interval_observables_commute, thm:bulk_observable_algebra}
    Terms disjoint from the support of $X$ have zero commutator. The remaining
    terms are precisely the translates in the enlarged interval. Reindex
    their sum and use multiplicativity of the interval inclusion. Multiplying
    by $X_N^\dagger$ gives the second identity.
\end{proof}

\begin{lemma}[Primitive-sector local commutator energy limit]
    \label{thm:primitive_local_commutator_energy_limit}
    \lean{MPSTensor.IsPrimitiveMPS.localCommutatorObservable_expectation_tendsto}
    \leanok
    Let $A$ be a normalized primitive tensor with positive-definite invariant
    matrix $\rho$, let $h$ act on $R>0$ sites, and let $X$ act on $k>0$
    sites. Put $t=R-1$, $w=2t+k$, $B=I_t\otimes X\otimes I_t$, and
    $Y=B^*[H_w(h),B]$. On chains of length $N_n=2n+w$, place $X$ at
    start $n+t$. For any eventually unit vectors $\psi_n$ in the
    open-boundary MPS space,
    $\operatorname{Re}\langle X_n^*[H_{N_n}(h),X_n]\psi_n,\psi_n\rangle
        \to\operatorname{Re}\eta_A(Y)$.
\end{lemma}
\begin{proof}\leanok
    \uses{thm:primitive_bulk_ground_vector_limits, thm:finite_open_commutator_locality}
    The commutator energy operator is the inclusion of $Y$ starting at $n$.
    Apply the primitive local expectation limit to $Y$.
    Conjugate symmetry of the inner product preserves its real part.
\end{proof}

\begin{lemma}[Eventually normalized primitive boundary vectors]
    \label{thm:primitive_unit_boundary_approximants}
    \lean{MPSTensor.IsPrimitiveMPS.exists_eventually_unit_groundSpaceES}
    \leanok
    Let $A$ and $\rho$ have the preceding properties. Along any filter with
    $N_i\to\infty$, there exist vectors $\psi_i$ in the open-boundary MPS
    space at every length $N_i$, whose norms are eventually one.
\end{lemma}
\begin{proof}\leanok
    \uses{thm:primitive_ground_space_map_eventual_inverse_gram_bound,
        def:ground_space_map_es}
    Choose a nonzero matrix-unit boundary vector $u$ and set
    $\psi_i=\|\Gamma_{N_i}u\|^{-1}\Gamma_{N_i}u$, with zero inverse
    interpreted as zero. The boundary map is eventually injective, so these
    vectors are eventually normalized; their ground-space membership holds
    at every length.
\end{proof}

\begin{lemma}[Primitive-sector local commutator inequality]
    \label{thm:primitive_local_commutator_gap}
    \lean{MPSTensor.IsPrimitiveMPS.localCommutator_gap_of_open_gap}
    \leanok
    Let $A$ be a normalized primitive tensor with positive-definite invariant
    matrix $\rho$, and let $X$ act on $k>0$ sites with $\eta_A(X)=0$.
    Fix $R>0$ and let $h=I-P_{G_R(A)}$. Suppose that, on the expanding
    chains $N_n=2n+2(R-1)+k$, the canonical open Hamiltonians eventually
    have kernel $G_{N_n}(A)$ and a common norm gap $\gamma\geq0$.
    Suppose also that a family of unit vectors eventually belongs to these
    ground spaces. With $B=I_{R-1}\otimes X\otimes I_{R-1}$,
    \begin{align}
        \gamma\operatorname{Re}\eta_A(X^*X)
         & \leq \operatorname{Re}\eta_A\bigl(B^*[H_{2(R-1)+k}(h),B]\bigr).
    \end{align}
    This is a statement about finite-interval primitive-sector expectations.
    Their extension to the quasi-local algebra and the GNS Hamiltonian remains
    separate, as discussed in
    \cite{gap:nachtergaele96_infinite_volume_ground_projection}.
\end{lemma}
\begin{proof}\leanok
    \uses{thm:primitive_local_commutator_energy_limit,
        thm:primitive_bulk_excitation_norm_limit,
        thm:primitive_bulk_ground_vector_limits, thm:commutator_gap_projection_limit,
        thm:finite_open_commutator_locality}
    The excitation norm squared tends to $\operatorname{Re}\eta_A(X^*X)$;
    its full ground-space projection tends to zero. The canonical matrix
    interaction sum is the open parent Hamiltonian. Apply the commutator
    energy limit and take limits in the finite-volume gap inequality.
\end{proof}

\begin{lemma}[A positive local commutator bound at the Wielandt range]
    \label{thm:primitive_local_commutator_wielandt}
    \lean{MPSTensor.IsPrimitiveMPS.exists_pos_localCommutator_gap_of_wielandt_range}
    \leanok
    Let $A$ be a normalized primitive tensor of bond dimension $D>0$ with
    positive-definite invariant matrix $\rho$. At each fixed range
    $R\geq D^4+1$, there exists $\gamma>0$ such that the preceding local
    commutator inequality holds for every observable on every nonempty
    finite interval with $\eta_A(X)=0$.
    The interaction is $h=I-P_{G_R(A)}$.
\end{lemma}
\begin{proof}\leanok
    \uses{thm:normal_from_primitive_posdef, thm:normal_open_gap_every_sufficient_range,
        thm:open_parent_hamiltonian_kernel_ground_space,
        thm:primitive_unit_boundary_approximants, thm:primitive_local_commutator_gap}
    Primitivity and positive definiteness imply normality.
    The dimension bound gives injectivity at $D^4$ and hence at $R-1$.
    The open-chain kernel is therefore the MPS ground space, and the normal
    open-chain gap theorem gives one positive constant independent of
    length. Choose eventually normalized boundary vectors and apply the
    preceding inequality.
\end{proof}

\begin{lemma}[Equality of open parent-interaction kernels]
    \label{thm:open_parent_interaction_kernel_equal}
    \lean{MPSTensor.IsParentInteraction.%
        ker_openInteractionHamiltonianES_eq_ker_openParentHamiltonianES}
    \leanok
    Let $h\geq0$ act on $R>0$ sites with kernel $G_R(A)$, and let
    $q=I-P_{G_R(A)}$. For every chain length $N$,
    $\ker H_N(h)=\ker H_N(q)$.
\end{lemma}
\begin{proof}\leanok
    \uses{thm:open_positive_parent_gap}
    There are $\kappa,C>0$ with
    $\kappa H_N(q)\leq H_N(h)\leq C H_N(q)$ at every length.
    For positive operators, either bound gives the corresponding reverse
    inclusion of kernels. Their combination gives equality.
\end{proof}

\begin{lemma}[Primitive-sector commutator bound for an open interaction]
    \label{thm:primitive_local_open_interaction_commutator_gap}
    \lean{MPSTensor.IsPrimitiveMPS.localCommutator_gap_of_openInteraction_gap}
    \leanok
    Let $A$ be a normalized primitive tensor with positive-definite invariant
    matrix $\rho$, let $h\geq0$ act on $R>0$ sites, and let $X$ act on
    $k>0$ sites with $\eta_A(X)=0$. Suppose the open interaction sums
    on chains $N_n=2n+2(R-1)+k$ eventually have kernel $G_{N_n}(A)$
    and common norm gap $\gamma\geq0$. Suppose that a family of unit
    vectors eventually belongs to these ground spaces.
    With $B=I_{R-1}\otimes X\otimes I_{R-1}$,
    \begin{align}
        \gamma\operatorname{Re}\eta_A(X^*X)
         & \leq\operatorname{Re}\eta_A\bigl(B^*[H_{2(R-1)+k}(h),B]\bigr).
    \end{align}
\end{lemma}
\begin{proof}\leanok
    \uses{thm:primitive_bulk_commutator_gap_reduction,
        thm:primitive_local_commutator_energy_limit, thm:finite_open_commutator_locality}
    The open matrix sum represents the open interaction Hamiltonian.
    Its commutator energy tends to the expectation of the fixed local patch.
    Apply the primitive boundary commutator reduction, which derives the
    excitation norm and full ground-space projection limits from the tensor.
\end{proof}

\begin{lemma}[Positive parent interactions satisfy a primitive-sector local gap]
    \label{thm:primitive_local_positive_parent_commutator_gap}
    \lean{MPSTensor.IsPrimitiveMPS.exists_pos_localCommutator_gap_of_isParentInteraction}
    \leanok
    Let $A$ be a normalized primitive tensor of bond dimension $D>0$ with
    positive-definite invariant matrix $\rho$. Let $h\geq0$ act on
    $R\geq D^4+1$ sites with kernel $G_R(A)$. There exists $\gamma>0$
    such that the preceding local commutator bound holds for every
    observable on every nonempty finite interval with $\eta_A(X)=0$.
    This concerns one primitive sector and its finite-interval expectations;
    the quasi-local, GNS, and pure multiblock conclusions remain separate
    \cite{gap:nachtergaele96_infinite_volume_ground_projection}.
\end{lemma}
\begin{proof}\leanok
    \uses{thm:normal_from_primitive_posdef, thm:normal_open_gap_every_sufficient_range,
        thm:open_positive_parent_gap, thm:open_parent_interaction_kernel_equal,
        thm:open_parent_hamiltonian_kernel_ground_space,
        thm:primitive_unit_boundary_approximants,
        thm:primitive_local_open_interaction_commutator_gap}
    The canonical open Hamiltonian has a uniform positive gap and the MPS
    ground-space kernel. Positive local comparison transfers the gap and
    preserves that kernel. Choose eventually normalized boundary vectors
    and apply the preceding local commutator inequality.
\end{proof}

\begin{lemma}[Compression between two injective ranges]
    \label{thm:parent_cross_range_compression}
    \lean{ContinuousLinearMap.norm_starProjection_comp_starProjection_le_inverseGram,
        ContinuousLinearMap.tendsto_starProjection_comp_starProjection_zero_of_inverseGram_bounds}
    \leanok
    Let $T\colon E\to F$ and $S\colon G\to F$ be injective bounded linear
    maps between real or complex Hilbert spaces, with $E$ and $G$
    finite-dimensional. Write $K_T=(T^\dagger T)^{-1}$,
    $K_S=(S^\dagger S)^{-1}$, and let $P_T,P_S$ be the orthogonal
    projections onto their ranges. For every bounded operator $X$ on $F$,
    \begin{align}
        \lVert P_T X P_S\rVert
         & \leq\lVert K_T\rVert^{1/2}
        \lVert T^\dagger X S\rVert\lVert K_S\rVert^{1/2}.
    \end{align}
    Consequently, along any filter, if these maps are eventually injective,
    their inverse Gram norms are eventually bounded by constants, and
    $\lVert T^\dagger X S\rVert\to0$, then
    $\lVert P_T X P_S\rVert\to0$. The Hilbert spaces and maps may vary
    with the index.
\end{lemma}

\begin{proof}
    \leanok
    Since $P_T=TK_TT^\dagger$ and $P_S=SK_SS^\dagger$,
    factor the compression as
    $(TK_T)(T^\dagger X S)(K_SS^\dagger)$.
    The identities $\lVert TK_T\rVert=\lVert K_T\rVert^{1/2}$ and
    $\lVert K_SS^\dagger\rVert=\lVert K_S\rVert^{1/2}$ give the bound.
    The convergence assertion follows by comparison with
    $\sqrt C\lVert T^\dagger X S\rVert\sqrt D$, where $C,D$ are
    eventual bounds for the two inverse Gram norms.
\end{proof}

\begin{definition}[MPS expectations on finite regions]
    \label{def:mps_finite_region_expectation}
    \lean{MPSTensor.localObservableExpectation}
    \leanok
    For a finite region $\Lambda\subset\mathbb Z$, include an observable
    $X\in\mathcal A_\Lambda$ into a consecutive interval $I(a,N)$
    containing $\Lambda$ and put
    \begin{align}
        \eta_{A,\Lambda}(X)
         & =\eta_A\bigl(C_{a,N}(\iota_{\Lambda,I(a,N)}X)\bigr).
    \end{align}
    To define the functional before proving independence, one may use
    $I(-r,2r+1)$ with $r=\max_{x\in\Lambda}|x|$ and $r=0$ when
    $\Lambda=\varnothing$. This gives a complex-linear functional without
    a normalization assumption on $A$ or $\rho$.
\end{definition}

\begin{lemma}[Independence and compatibility of finite-region expectations]
    \label{thm:mps_finite_region_expectation_compatible}
    \lean{MPSTensor.localObservableExpectation_eq_of_subset_interval,
        MPSTensor.localObservableExpectation_localInclusion}
    \leanok
    Suppose $\sum_i A_i^\dagger A_i=I$ and $E_A(\rho)=\rho$.
    Then the preceding definition is independent of the containing interval.
    For $\Lambda\subseteq\Gamma$ and $X\in\mathcal A_\Lambda$,
    $\eta_{A,\Gamma}(\iota_{\Lambda,\Gamma}X)=\eta_{A,\Lambda}(X)$.
    Neither positivity nor nonzero trace of $\rho$ is required for these
    algebraic compatibility identities.
\end{lemma}
\begin{proof}\leanok
    \uses{def:mps_finite_region_expectation,thm:local_interval_bulk_observable,thm:primitive_insertion_local_state,lem:qca_local_inclusion_trans}
    Identity extension on the left preserves expectation by trace preservation;
    identity extension on the right preserves expectation by $E_A(\rho)=\rho$.
    Include two containing intervals in a common larger interval to compare
    their values. Composition of local inclusions gives finite-region compatibility.
\end{proof}

\begin{definition}[The compatible local MPS state]
    \label{def:mps_compatible_local_state}
    \lean{MPSTensor.localMPSState, MPSTensor.localMPSState_functional_interval}
    \leanok
    Suppose $\sum_i A_i^\dagger A_i=I$, $\rho\geq0$,
    $E_A(\rho)=\rho$, and $\operatorname{tr}\rho\neq0$.
    The functionals $\eta_{A,\Lambda}$ form a compatible family of
    normalized positive contractive local states. In particular,
    $\eta_{A,I(a,N)}(X)=\eta_A(C_{a,N}X)$ for every $N\geq0$.
    This is the compatible local interpretation of
    \cite[Section~3, equations (3.1)--(3.2b)]{Nachtergaele1996Spectral}.
\end{definition}
\begin{proof}\leanok
    \uses{thm:mps_finite_region_expectation_compatible,thm:primitive_insertion_state_norm,def:qca_compatible_local_state,thm:primitive_insertion_local_state}
    The inclusion and coordinate maps preserve identities, products, and
    adjoints. Positivity and normalization therefore follow from their
    finite-interval counterparts. Star-algebra homomorphisms between
    complex C$^*$-algebras are contractive, so
    $|\eta_{A,\Lambda}(X)|\leq\|X\|$.
\end{proof}

\begin{definition}[The completed quasi-local MPS expectation]
    \label{def:mps_quasi_local_expectation}
    \lean{MPSTensor.quasiLocalExpectation}
    \leanok
    Under the preceding hypotheses and $d>0$, let $\omega_A$ be the
    continuous linear extension of the compatible local MPS state to the
    completed quasi-local observable algebra.
\end{definition}

\begin{lemma}[Local evaluation and state properties of the MPS extension]
    \label{thm:mps_quasi_local_expectation_state}
    \lean{MPSTensor.quasiLocalExpectation_quasiLocalObservable,
        MPSTensor.quasiLocalExpectation_interval, MPSTensor.quasiLocalExpectation_isState}
    \leanok
    Under the hypotheses of the preceding definition, the restriction of
    $\omega_A$ to $\mathcal A_\Lambda$ is $\eta_{A,\Lambda}$.
    Thus $\omega_A(X)=\eta_A(C_{a,N}X)$ for an observable on $I(a,N)$,
    including $N=0$. The extension satisfies
    $\|\omega_A\|=1$, $\omega_A(I)=1$, and $\omega_A(X^*X)\geq0$
    for every quasi-local observable $X$. No primitivity or purity is asserted.
\end{lemma}
\begin{proof}\leanok
    \uses{def:mps_quasi_local_expectation, def:mps_compatible_local_state,
        thm:qca_compatible_state_evaluation, thm:qca_compatible_state_extension}
    Apply the completion theorem for compatible local states. Its finite-region
    evaluation formula gives the asserted restriction, and interval evaluation
    gives the original trace insertion formula.
\end{proof}

\begin{lemma}[Translation preserves interval coordinates]
    \label{thm:local_interval_translation_coordinates}
    \lean{SpinChain.translateRegion_finiteChainRegion,
        SpinChain.intervalCoordinates_localTranslation}
    \leanok
    For $a,b\in\mathbb Z$ and $N\geq0$, translation by $b$ sends
    $I(a,N)$ to $I(a+b,N)$. If $T_b$ denotes the induced local
    star-algebra equivalence, then
    $C_{a+b,N}(T_bX)=C_{a,N}(X)$.
    The equality includes the empty interval.
\end{lemma}
\begin{proof}\leanok
    \uses{def:local_interval_coordinates,def:qca_finite_local_translation,def:qca_translate_region,def:qca_local_inclusion}
    The site $a+i$ is sent to $a+b+i$ and retains its numbered coordinate
    $i$. Hence the two interval coordinate maps give the same configurations
    and matrix entries. The local inclusion identifying the translated region
    with $I(a+b,N)$ has empty complement.
\end{proof}

\begin{lemma}[Translation invariance of finite-region MPS expectations]
    \label{thm:mps_finite_region_translation_invariance}
    \lean{MPSTensor.localObservableExpectation_localTranslation}
    \leanok
    Suppose $\sum_i A_i^\dagger A_i=I$ and $E_A(\rho)=\rho$.
    For every finite region $\Lambda\subset\mathbb Z$, displacement
    $b\in\mathbb Z$, and observable $X\in\mathcal A_\Lambda$,
    $\eta_{A,\Lambda+b}(T_bX)=\eta_{A,\Lambda}(X)$.
    Positivity and nonzero trace of $\rho$ are unnecessary for this identity.
\end{lemma}
\begin{proof}\leanok
    \uses{thm:mps_finite_region_expectation_compatible,thm:local_interval_translation_coordinates,lem:qca_finite_local_translation_naturality,def:qca_translate_region}
    Choose an interval containing $\Lambda$. Translation commutes with the
    local inclusion and preserves numbered interval coordinates. Evaluate
    the two local expectations in this interval and its translate.
\end{proof}

\begin{lemma}[Translation invariance of the completed MPS state]
    \label{thm:mps_quasi_local_translation_invariance}
    \lean{MPSTensor.quasiLocalExpectation_quasiLocalTranslation}
    \leanok
    Suppose $d>0$, $\sum_i A_i^\dagger A_i=I$, $\rho\geq0$,
    $E_A(\rho)=\rho$, and $\operatorname{tr}\rho\neq0$.
    For every displacement $b\in\mathbb Z$ and quasi-local observable $X$,
    the completed state satisfies $\omega_A(T_bX)=\omega_A(X)$.
    This proves translation invariance of the local GVBS state in
    \cite[Section~3, equations (3.1)--(3.2b)]{Nachtergaele1996Spectral};
    no purity conclusion is asserted.
\end{lemma}
\begin{proof}\leanok
    \uses{thm:mps_finite_region_translation_invariance,
        thm:mps_quasi_local_expectation_state, def:qca_quasi_local_translation}
    The equality holds on every finite-region algebra. Translation and the
    completed expectation are continuous, so it extends from the dense
    algebraic local algebra to its completion.
\end{proof}

\begin{lemma}[Interval evaluation of the completed MPS state]
    \label{thm:mps_quasi_local_interval_expectation}
    \lean{MPSTensor.quasiLocalExpectation_quasiLocalIntervalObservable}
    \leanok
    Suppose $d>0$, $\sum_i A_i^\dagger A_i=I$, $\rho\geq0$,
    $E_A(\rho)=\rho$, and $\operatorname{tr}\rho\neq0$.
    For every $a\in\mathbb Z$, $k\geq0$, and $X\in M_{d^k}(\mathbb C)$,
    the completed state satisfies $\omega_A(j_{a,k}(X))=\eta_A(X)$.
\end{lemma}
\begin{proof}\leanok
    \uses{thm:mps_quasi_local_expectation_state, thm:quasi_local_interval_observable}
    Apply interval evaluation after mapping the matrix observable to
    the interval algebra. The inverse coordinate maps cancel.
\end{proof}

\begin{lemma}[A primitive-sector quasi-local commutator gap]
    \label{thm:primitive_quasi_local_parent_commutator_gap}
    \lean{MPSTensor.IsPrimitiveMPS.exists_pos_quasiLocalCommutator_gap_of_isParentInteraction}
    \leanok
    Let $d,D>0$, and let $A$ be a normalized primitive tensor with a
    positive-definite invariant matrix $\rho$. Suppose that $h\geq0$ has
    kernel $G_R(A)$ and $R\geq D^4+1$. There is a constant $\gamma>0$
    such that, for every $a\in\mathbb Z$, $k\geq1$, and
    $X\in M_{d^k}(\mathbb C)$ with $\omega_A(j_{a,k}(X))=0$,
    \begin{align}
        \gamma\operatorname{Re}\omega_A\bigl(j_{a,k}(X)^*j_{a,k}(X)\bigr)
        \leq \operatorname{Re}\omega_A\bigl(j_{a-t,w}(Y)\bigr),
        \notag
    \end{align}
    where $t=R-1$, $w=2t+k$, $B=I_t\otimes X\otimes I_t$, and
    $Y=B^*[H_w(h),B]$. The constant is independent of $a$, $k$, and $X$.
    This is the local commutator inequality in the constructed primitive
    quasi-local state, at sufficient range. Identification of all pure
    multiblock GVBS states and construction of the GNS Hamiltonian remain
    separate from this consequence of
    \cite[Theorem~1.2 and Section~6]{Nachtergaele1996Spectral}.
\end{lemma}
\begin{proof}\leanok
    \uses{thm:primitive_local_positive_parent_commutator_gap,
        thm:mps_quasi_local_interval_expectation, thm:quasi_local_interval_observable}
    The interval inclusion preserves products and adjoints. Evaluate both
    sides with the trace insertion formula and apply the finite-interval
    primitive-sector bound.
\end{proof}

\begin{lemma}[Prefix ground projection]
    \label{thm:ground_prefix_projection_identity}
    \lean{MPSTensor.chainWindowOperator_groundSpaceProjection_comp_groundSpaceProjection}
    \leanok
    \uses{def:ground_space_es}
    Let $A$ be an MPS tensor, let $1\leq L\leq N$, and denote its open
    ground projection on $n$ sites by $P_A(n)$. If $Q_A(L,N)$ is $P_A(L)$
    placed on the first $L$ sites and extended by the identity, then
    \begin{align}
        Q_A(L,N)P_A(N)=P_A(N)=P_A(N)Q_A(L,N).
        \notag
    \end{align}
    No injectivity or normalization is assumed.
\end{lemma}
\begin{proof}\leanok
    \uses{thm:open_ground_space_le_hamiltonian_kernel}
    Every open MPS vector satisfies the parent constraint on its prefix.
    Thus the complementary prefix projection annihilates the full ground
    space. Taking adjoints gives the identity in the other order.
\end{proof}

\begin{lemma}[Local cross compression]
    \label{thm:bulk_cross_compression_prefix_bound}
    \lean{MPSTensor.norm_groundSpaceES_bulkObservable_cross_compression_le}
    \leanok
    Let $A,B$ be arbitrary MPS tensors with the same physical dimension.
    Let $X$ act on $k\geq1$ sites, and let $\ell\geq1$ and $r\geq0$.
    Write $N=\ell+k+r$ and let $X_{\ell,r}$ be the placement of $X$ after
    the first $\ell$ sites. Then
    \begin{align}
        \|P_A(N)X_{\ell,r}P_B(N)\|
        \leq \|P_A(\ell)P_B(\ell)\|\,\|X\|.
        \notag
    \end{align}
\end{lemma}
\begin{proof}\leanok
    \uses{thm:ground_prefix_projection_identity,thm:disjoint_interval_observables_commute}
    The observable commutes with both prefix projections. The preceding
    identity therefore gives
    $P_AX_{\ell,r}P_B=P_AQ_AQ_BX_{\ell,r}P_B$.
    The norm of $Q_AQ_B$ is bounded by the norm of the prefix product,
    while the remaining ground projections are contractions.
\end{proof}

\begin{lemma}[Cross-sector decay by prefix locality]
    \label{thm:bulk_cross_compression_decay}
    \lean{MPSTensor.IsPrimitiveMPS.groundSpaceES_starProjection_comp_tendsto_zero_of_inequivalent,
        MPSTensor.IsPrimitiveMPS.bulkObservable_cross_groundSpace_compression_tendsto_zero_of_inequivalent}
    \leanok
    Let $A,B$ be trace-preserving primitive MPS tensors with positive
    definite invariant matrices and positive bond dimensions. Suppose that $A$ and $B$ are inequivalent under a bond
    similarity and a phase, including any identification of equal bond
    dimensions. Then
    \begin{align}
        \|P_A(n)P_B(n)\|\longrightarrow0\qquad(n\longrightarrow\infty).
        \notag
    \end{align}
    For any fixed observable $X$ on $k\geq1$ sites, any filter along which
    $\ell\longrightarrow\infty$, and any choice of $r\geq0$, one also has
    \begin{align}
        \|P_A(\ell+k+r)X_{\ell,r}P_B(\ell+k+r)\|
        \longrightarrow0.
        \notag
    \end{align}
    No growth condition is required of the right interval. This is a
    locality consequence of the overlap estimate in
    \cite[proof of the lemma on distinct local support spaces]{Nachtergaele1996Spectral}.
\end{lemma}
\begin{proof}\leanok
    \uses{thm:inequivalent_block_ground_space_overlap_decay,thm:bulk_cross_compression_prefix_bound}
    The relative overlap estimate bounds the product of the two ground
    projections by an arbitrary positive number at sufficiently large
    length. The second assertion follows by multiplying this bound by
    the fixed norm of the observable and applying the preceding lemma.
\end{proof}

\begin{lemma}[Vanishing error of a finite sum of sector projections]
    \label{thm:sector_projection_sum_error_decay}
    \lean{Submodule.tendsto_norm_sum_starProjection_sub_iSup_zero}
    \leanok
    Let $G_{j,n}$ be a finite family of subspaces of finite-dimensional
    complex Hilbert spaces $E_n$, along an arbitrary filter. Suppose that
    for every $\varepsilon>0$, eventually every pair $i\neq j$ satisfies
    $|\langle x,y\rangle|\leq\varepsilon\|x\|\|y\|$ for
    $x\in G_{i,n}$ and $y\in G_{j,n}$. Write $P_{j,n}$ for their
    orthogonal projections and $P_n$ for the projection onto their joint span.
    Then $\|\sum_jP_{j,n}-P_n\|\to0$.
\end{lemma}
\begin{proof}\leanok
    \uses{thm:ph_sum_block_projectors}
    Let $C_{ij}=1$ for $i\neq j$ and $C_{ii}=0$. The overlap matrix
    $B=\varepsilon C$ satisfies $\|B\|\to0$. The finite-family comparison
    bounds the projection error by $\delta/(1-\delta)$, where
    $\delta=\|B\|/(1-\|B\|)$. This bound tends to zero.
\end{proof}

\begin{lemma}[Combining vanishing sector compressions]
    \label{thm:sector_compression_assembly}
    \lean{Submodule.tendsto_norm_iSup_starProjection_comp_zero_of_sector_compressions}
    \leanok
    Under the preceding asymptotic overlap hypothesis, let $W_n\subseteq E_n$
    have orthogonal projection $Q_n$, and let the operators $X_n$ be
    eventually uniformly bounded. If $\|P_{j,n}X_nQ_n\|\to0$ for every
    sector $j$, then $\|P_nX_nQ_n\|\to0$. The input space $W_n$ may be
    arbitrary and need not be a member of the finite family.
\end{lemma}
\begin{proof}\leanok
    \uses{thm:sector_projection_sum_error_decay}
    Since $\|Q_n\|\leq1$,
    \begin{align}
        \|P_nX_nQ_n\|
         & \leq\|P_n-\textstyle\sum_jP_{j,n}\|\,\|X_n\|
        +\sum_j\|P_{j,n}X_nQ_n\|.
    \end{align}
    Both terms tend to zero, by the projection comparison and finiteness
    of the sector family.
\end{proof}

\begin{lemma}[Centered observables and the full multiblock ground projection]
    \label{thm:multiblock_centered_full_ground_compression}
    \lean{MPSTensor.bulkObservable_iSup_groundSpace_compression_tendsto_zero}
    \leanok
    Let $A_j$ be finitely many normalized primitive tensors with positive
    bond dimensions $D_j$, positive-definite invariant matrices $\rho_j$,
    and pairwise gauge-and-phase inequivalence whenever their bond dimensions
    agree. Fix a sector $\alpha$ and an observable $X$ on $k>0$ sites
    with $\eta_{A_\alpha}(X)=0$. Along an arbitrary filter let
    $\ell_n,r_n\to\infty$ and $N_n=\ell_n+k+r_n$.
    If $P_n$ projects onto $\sum_jG_{N_n}(A_j)$ and $P_{\alpha,n}$
    onto $G_{N_n}(A_\alpha)$, then
    \begin{align}
        \bigl\|P_n(I_{\ell_n}\otimes X\otimes I_{r_n})P_{\alpha,n}\bigr\|
         & \longrightarrow0.
    \end{align}
    This gives the full finite-ground-space projection term in
    \cite[Section~6]{Nachtergaele1996Spectral}. No finite-chain uniqueness
    or purity of a completed state is inferred.
\end{lemma}
\begin{proof}\leanok
    \uses{thm:sector_compression_assembly, thm:primitive_bulk_ground_compression,
        thm:bulk_cross_compression_decay, thm:inequivalent_block_ground_space_overlap_decay,
        thm:bulk_observable_algebra}
    The sector overlaps tend uniformly to zero by inequivalence and
    primitivity. The diagonal compression tends to zero because the
    observable is centered in sector $\alpha$. Every off-diagonal
    compression tends to zero by prefix locality. Interior placement is
    contractive, so the finite-sector combination theorem applies.
\end{proof}

\begin{lemma}[Exact local support of the MPS state]
    \label{thm:mps_state_local_support}
    \lean{MPSTensor.observableInsertionExpectation_eq_zero_of_groundSpaceES_le_ker,
        MPSTensor.quasiLocalExpectation_interval_eq_zero_of_groundSpaceES_le_ker,
        MPSTensor.IsParentInteraction.quasiLocalExpectation_interval_eq_zero}
    \leanok
    Let $A$ be an MPS tensor and let $X$ act on $k$ sites. If
    $G_k(A)\subseteq\ker X$, then $\eta_A(X)=0$ for every virtual
    matrix $\rho$, without normalization or positivity assumptions.
    If $d>0$, $\sum_i A_i^\dagger A_i=I$, $\rho\geq0$,
    $E_A(\rho)=\rho$, and $\operatorname{tr}\rho\neq0$, this gives
    $\omega_A(j_{a,k}(X))=0$ for every $a\in\mathbb Z$.
    In particular, the constructed state has expectation zero for every
    translate of a positive local parent interaction.
    This is the local support property in
    \cite[Theorem~1.1 and Section~3]{Nachtergaele1996Spectral}.
\end{lemma}
\begin{proof}\leanok
    \uses{thm:primitive_insertion_boundary_identity,thm:quasi_local_interval_observable,thm:mps_quasi_local_interval_expectation}
    The hypothesis gives $X\Gamma_k=0$, and hence $\Gamma_k^*X\Gamma_k=0$.
    The latter operator is the Choi reshuffling of the insertion transfer
    map. Its matrix-unit coefficients therefore imply that this map is
    zero. Taking the trace at $\rho$ proves the first assertion; interval
    evaluation proves the assertions for the completed state.
\end{proof}

\begin{lemma}[Vanishing full projection of a sector excitation]
    \label{thm:multiblock_sector_excitation_projection_limit}
    \lean{MPSTensor.bulkObservable_iSup_groundSpace_projection_norm_sq_tendsto_zero}
    \leanok
    Let $(A_j,\rho_j)_{j\in J}$ be a finite family of pairwise inequivalent
    trace-preserving primitive tensors with positive bond dimensions and
    positive-definite invariant matrices. Write
    $G(N)=\sum_jG_{A_j}(N)$ and let $P(N)$ be its orthogonal projection.
    Fix a sector $\alpha$, an observable $X$ on $k\geq1$ sites with
    $\eta_\alpha(X)=0$, and free lengths $\ell,r\longrightarrow\infty$.
    For any varying boundary vectors $\psi\in G_{A_\alpha}(\ell+k+r)$
    with eventually $\|\psi\|\leq1$,
    \begin{align}
        \|P(\ell+k+r)X_{\ell,r}\psi\|^2\longrightarrow0.
        \notag
    \end{align}
    The assertion holds along an arbitrary filter.
\end{lemma}
\begin{proof}\leanok
    \uses{thm:multiblock_centered_full_ground_compression}
    The sector projection fixes $\psi$. Thus the vector norm is bounded
    by the norm of the compression from that sector into the full ground
    space, multiplied by $\|\psi\|$. Apply the compression limit and square.
\end{proof}

\begin{lemma}[Multiblock commutator reduction]
    \label{thm:multiblock_commutator_gap_reduction}
    \lean{MPSTensor.multiblock_commutator_gap_of_bulkObservable_energy_limit}
    \leanok
    Under the preceding tensor hypotheses, suppose that $H_n$ are
    eventually positive operators on chains of length $N_n=\ell_n+k+r_n$,
    with $\ell_n,r_n\longrightarrow\infty$, and eventually
    \begin{align}
        \ker H_n=G(N_n),\qquad
        \gamma\|v\|\leq\|H_nv\|\quad(v\perp\ker H_n),
        \qquad \gamma\geq0.
        \notag
    \end{align}
    Let $\psi_n\in G_{A_\alpha}(N_n)$ be eventually unit vectors, and
    let $B_n=X_{\ell_n,r_n}$ with $\eta_\alpha(X)=0$. If
    \begin{align}
        \operatorname{Re}\langle B_n^*[H_n,B_n]\psi_n,\psi_n\rangle
        \longrightarrow e,
        \notag
    \end{align}
    then
    \begin{align}
        \gamma\operatorname{Re}\eta_\alpha(X^*X)\leq e.
        \notag
    \end{align}
\end{lemma}
\begin{proof}\leanok
    \uses{thm:multiblock_sector_excitation_projection_limit,thm:primitive_bulk_excitation_norm_limit,thm:commutator_gap_projection_limit}
    The distinguished sector lies in the full ground space. Primitive
    convergence gives $\|B_n\psi_n\|^2\to\operatorname{Re}\eta_\alpha(X^*X)$,
    and the preceding theorem removes its full ground component. Apply the
    positive finite-volume gap inequality and pass to the limit.
\end{proof}

\begin{lemma}[Local multiblock commutator inequality]
    \label{thm:multiblock_local_open_interaction_commutator_gap}
    \lean{MPSTensor.multiblock_localCommutator_gap_of_openInteraction_gap}
    \leanok
    Under the same finite-family tensor hypotheses, let $h$ be a positive
    interaction on $R\geq1$ sites. Put $t=R-1$,
    $N_n=(n+t+k)+(t+n)$, and let $H_n$ be its open interaction sum.
    Suppose that these operators eventually have kernel $G(N_n)$ and
    norm gap at least $\gamma\geq0$. For eventually unit
    $\psi_n\in G_{A_\alpha}(N_n)$ and every $k\geq1$ observable $X$
    with $\eta_\alpha(X)=0$,
    \begin{align}
        \gamma\operatorname{Re}\eta_\alpha(X^*X)
        \leq\operatorname{Re}\eta_\alpha(C_h(X)),
        \qquad C_h(X)=B^*[H_{t+k+t},B],\quad B=X_{t,t}.
        \notag
    \end{align}
\end{lemma}
\begin{proof}\leanok
    \uses{thm:multiblock_commutator_gap_reduction,thm:primitive_local_commutator_energy_limit}
    Finite-range locality places the commutator energy on the fixed
    interval of length $t+k+t$. Primitive expectation convergence determines
    its limit, so the preceding reduction applies without a given
    energy limit.
\end{proof}

\begin{lemma}[A common local commutator bound for all primitive sectors]
    \label{thm:multiblock_local_parent_commutator_gap}
    \lean{MPSTensor.exists_pos_multiblock_localCommutator_gap_of_isParentInteraction,
        MPSTensor.exists_pos_multiblock_localCommutator_gap_of_wordTupleSpanTop}
    \leanok
    Let $A_j$ be trace-preserving primitive tensors of positive bond
    dimensions, with positive definite invariant matrices $\rho_j$,
    and let $T=\bigoplus_j\mu_jA_j$, where every $\mu_j\neq0$.
    Assume that the simultaneous word span is the full product of the bond
    matrix algebras at a length $S\geq1$. Let $R\geq S+1$, and let $h$ be
    a positive interaction on $R$ sites with $\ker h=G_T(R)$.
    There exists $\gamma>0$ such that, for every sector $\alpha$ and every
    observable $X$ on any $k\geq1$ sites with $\eta_\alpha(X)=0$,
    \begin{align}
        \gamma\operatorname{Re}\eta_\alpha(X^*X)
        \leq\operatorname{Re}\eta_\alpha(C_h(X)).
        \notag
    \end{align}
    In particular this holds for the canonical parent projection
    $h=I-P_T(R)$. These are finite-interval sector expectation statements;
    they do not assert the construction of a GNS Hamiltonian or purity of
    its state. The range is sufficient and need not be optimal.
    \cite[Theorem~1.2 and Section~6]{Nachtergaele1996Spectral},
    \cite{gap:nachtergaele96_infinite_volume_ground_projection}.
\end{lemma}
\begin{proof}\leanok
    \uses{thm:multiblock_local_open_interaction_commutator_gap,thm:block_open_gap_all_lengths_simultaneous_injectivity,thm:block-word-span-separation,thm:block_open_kernel_simultaneous_injectivity,thm:open_positive_parent_gap,thm:open_parent_interaction_kernel_equal,thm:primitive_unit_boundary_approximants}
    Simultaneous injectivity implies sector inequivalence and identifies
    the open-chain kernel with the sum of the sector spaces at this range.
    The open-chain theorem gives a uniform canonical gap; positivity
    and the prescribed local kernel transfer it to $h$.
    Choose eventual unit boundary vectors independently in each sector
    and apply the preceding local inequality. The same gap constant works
    for every sector and every observable.
\end{proof}

\begin{lemma}[One quasi-local commutator gap in all constructed sectors]
    \label{thm:multiblock_quasi_local_parent_commutator_gap}
    \lean{MPSTensor.exists_pos_multiblock_quasiLocalCommutator_gap_of_isParentInteraction}
    \leanok
    Let $d>0$, and let $A_j$ be finitely many normalized primitive tensors
    with positive bond dimensions and positive-definite invariant matrices
    $\rho_j$. Suppose that the simultaneous word span is full at a length
    $S>0$. Let $T=\bigoplus_j\mu_jA_j$ with $\mu_j\neq0$, and let $h\geq0$
    have kernel $G_R(T)$ at range $R\geq S+1$.
    There is a constant $\gamma>0$ such that, for every sector $\alpha$,
    $a\in\mathbb Z$, $k\geq1$, and $X\in M_{d^k}(\mathbb C)$
    with $\omega_{A_\alpha}(j_{a,k}(X))=0$,
    \begin{align}
        \gamma\operatorname{Re}\omega_{A_\alpha}
        \bigl(j_{a,k}(X)^*j_{a,k}(X)\bigr)
        \leq\operatorname{Re}\omega_{A_\alpha}(j_{a-t,w}(Y)),
        \notag
    \end{align}
    where $t=R-1$, $w=2t+k$, $B=I_t\otimes X\otimes I_t$, and
    $Y=B^*[H_w(h),B]$. The same constant applies to all sectors,
    positions, and centered observables. The states are the completed
    expectations of the given primitive blocks; identification of the full
    pure GVBS ground-state face remains a separate assertion.
\end{lemma}
\begin{proof}\leanok
    \uses{thm:multiblock_local_parent_commutator_gap,
        thm:mps_quasi_local_interval_expectation, thm:quasi_local_interval_observable}
    The common finite-interval bound is expressed in the constructed states
    by interval evaluation. The interval inclusions preserve products and
    adjoints, so the centered hypothesis and both terms agree exactly.
\end{proof}

\begin{definition}[Rectangular observable insertion]
    \label{def:mixed_observable_insertion}
    \lean{MPSTensor.mixedPhysicalObservableTransfer, MPSTensor.mixedPhysicalObservableTransfer_apply, MPSTensor.mixedPhysicalObservableTransfer_self}
    \leanok
    For tensors $A$ and $B$ with a common physical alphabet and possibly different
    bond dimensions, define the insertion of a $k$-site matrix $X$ by
    \begin{align}
        E_{B,A;X}(Y)=\sum_{\sigma,\tau}X_{\sigma\tau}B_\tau Y A_\sigma^*.
        \notag
    \end{align}
    Here $Y$ is a rectangular matrix from the virtual space of $A$ to that of $B$.
    For $A=B$, this is the ordinary observable insertion.
\end{definition}

\begin{lemma}[Mixed insertion identities]
    \label{thm:mixed_observable_insertion_identities}
    \lean{MPSTensor.inner_single_mixed_observable_groundSpaceMapES_single, MPSTensor.adjoint_groundSpaceMapES_observable_groundSpaceMapES_eq_mixedTransfer, MPSTensor.mixedPhysicalObservableTransfer_one, MPSTensor.mixedPhysicalObservableTransfer_appendObservable, MPSTensor.mixedPhysicalObservableTransfer_bulkObservable, MPSTensor.adjoint_groundSpaceMapES_bulkObservable_groundSpaceMapES_eq_mixedTransfer}
    \leanok
    The coefficient of $\Gamma_{A,k}^*X\Gamma_{B,k}$ between boundary units
    $(b,a)$ and $(e,c)$ is $E_{B,A;X}(|c\rangle\langle a|)_{eb}$.
    Also,
    \begin{align}
        E_{B,A;1_k}=E_{B,A}^{\,k},\qquad
        E_{B,A;X\otimes Z}=E_{B,A;X}E_{B,A;Z},
        \notag
    \end{align}
    and
    \begin{align}
        E_{B,A;1_\ell\otimes X\otimes1_r}
        =E_{B,A}^{\,\ell}E_{B,A;X}E_{B,A}^{\,r}.
        \notag
    \end{align}
\end{lemma}
\begin{proof}\leanok
    \uses{def:mixed_observable_insertion}
    Expand the boundary traces in matrix units. Splitting a physical word at a
    prescribed cut gives the composition identity, and identity insertions give
    the two ordinary transfer powers.
\end{proof}

\begin{lemma}[Decay of mixed observable sandwiches]
    \label{thm:mixed_observable_sandwich_decay}
    \lean{MPSTensor.mixedMapLM_pow_tendsto_zero_of_spectralRadius_lt_one, MPSTensor.mixedPhysicalObservableTransfer_sandwich_tendsto_zero}
    \leanok
    Suppose the rectangular mixed transfer map $E_{B,A}$ has spectral radius less
    than one. Then $E_{B,A}^{\,n}\to0$ in operator norm. For a fixed local matrix
    $X$, and along any filter with $\ell\to\infty$ and $r\to\infty$,
    \begin{align}
        E_{B,A}^{\,\ell}E_{B,A;X}E_{B,A}^{\,r}\longrightarrow0
        \notag
    \end{align}
    in operator norm.
\end{lemma}
\begin{proof}\leanok
    \uses{def:mixed_observable_insertion}
    The mixed transfer powers decay pointwise by the strict spectral-radius
    bound. Finite dimension upgrades this to operator norm convergence;
    continuity of composition then gives the inserted sandwich limit.
\end{proof}

\begin{lemma}[Decay of cross-sector boundary pullbacks]
    \label{thm:mixed_bulk_boundary_pullback_decay}
    \lean{MPSTensor.adjoint_groundSpaceMapES_bulkObservable_comp_tendsto_zero}
    \leanok
    Under the preceding strict mixed spectral-radius hypothesis, let
    $N=\ell+k+r$ and $X_{\ell,r}=1_\ell\otimes X\otimes1_r$.
    As $\ell,r\to\infty$ along any common filter,
    \begin{align}
        \bigl\|\Gamma_{A,N}^*X_{\ell,r}\Gamma_{B,N}\bigr\|\longrightarrow0.
        \notag
    \end{align}
    No injectivity of the boundary maps is required. This is a distinct-sector
    estimate, rather than an assertion of the infinite-volume spectral gap.
\end{lemma}
\begin{proof}\leanok
    \uses{thm:mixed_observable_insertion_identities,thm:mixed_observable_sandwich_decay}
    The pullback is the rectangular coefficient reshuffling of the virtual
    sandwich. This reshuffling is linear between fixed finite-dimensional
    spaces and hence continuous. Apply it to the sandwich limit.
\end{proof}

\begin{lemma}[Constructing sector states from simultaneous injectivity]
    \label{thm:normalized_multiblock_quasi_local_parent_gap}
    \lean{MPSTensor.exists_primitiveFamily_quasiLocalCommutator_gap_of_wordTupleSpanTop}
    \leanok
    Let $d>0$, and let $A_j$ have positive bond dimensions and full
    simultaneous word span at $S>0$. Let $T=\bigoplus_j\mu_jA_j$ with
    $\mu_j\neq0$, and let $h\geq0$ have kernel $G_R(T)$ at $R\geq S+1$.
    There exist normalized primitive representatives $B_j$ with
    positive-definite invariant matrices $\rho_j$ such that
    $G_N(A_j)=G_N(B_j)$ for every $j,N$. Their constructed quasi-local
    states satisfy the preceding common commutator inequality for the
    original interaction $h$, with one positive constant for all sectors,
    interval positions, and centered local observables.
    No normalized representatives or invariant matrices are given.
\end{lemma}
\begin{proof}\leanok
    \uses{thm:block_word_span_primitive_normalization,
        thm:multiblock_quasi_local_parent_commutator_gap}
    Primitive normalization preserves each block support and simultaneous
    injectivity. Hence it also preserves the combined local support and
    the parent-interaction kernel. Apply the common bound to the normalized
    family with the original positive interaction.
\end{proof}

\begin{lemma}[Quasi-local images of interior inclusions]
    \label{thm:quasi_local_bulk_inclusion}
    \lean{SpinChain.quasiLocalIntervalObservable_bulkObservable, SpinChain.quasiLocalIntervalObservable_chainWindowOperator}
    \leanok
    Let $j_{a,k}$ denote inclusion of the interval $[a,a+k)$ into the quasi-local
    algebra. Then, including empty intervals,
    \begin{align}
        j_{a-\ell,\ell+k+r}(I_\ell\otimes X\otimes I_r)=j_{a,k}(X).
        \notag
    \end{align}
    For $k>0$, the nonwrapping inclusion of $X$ at position $b$ in a chain of
    length $N$, where $b+k\leq N$, has quasi-local image $j_{a+b,k}(X)$ when the
    chain begins at $a$.
\end{lemma}
\begin{proof}\leanok
    \uses{thm:local_interval_bulk_observable,thm:quasi_local_interval_observable,lem:qca_quasi_local_observable}
    Interval coordinates identify identity extensions with finite-region
    inclusions. The quasi-local inclusions agree under enlargement of the finite
    region. The nonwrapping statement is obtained by adjoining the remaining
    sites on either side.
\end{proof}

\begin{lemma}[Fixed quasi-local support of commutator energy]
    \label{thm:quasi_local_commutator_locality}
    \lean{SpinChain.quasiLocalIntervalObservable_bulkObservable_commutator, SpinChain.quasiLocalIntervalObservable_commutator_locality}
    \leanok
    Let $R,k>0$, $t=R-1$, $w=k+2t$, and let $h$ be an arbitrary $R$-site matrix.
    Write $H_N$ for its open-chain sum and $X_N=I_\ell\otimes X\otimes I_r$, with
    $\ell,r\geq t$ and $N=\ell+k+r$. Let
    \begin{align}
        Y=(I_t\otimes X\otimes I_t)^*
        [H_w,I_t\otimes X\otimes I_t].
        \notag
    \end{align}
    Then
    \begin{align}
        j_{a-\ell,N}(X_N^*[H_N,X_N])=j_{a-t,w}(Y).
        \notag
    \end{align}
    Equivalently, putting $x=j_{a,k}(X)$ and
    $\widehat H_N=j_{a-\ell,N}(H_N)$ gives
    $x^*[\widehat H_N,x]=j_{a-t,w}(Y)$.
    Neither positivity nor a state assumption is required.
\end{lemma}
\begin{proof}\leanok
    \uses{thm:finite_open_commutator_locality,thm:quasi_local_bulk_inclusion}
    The finite-chain energy is a nonwrapping inclusion of the fixed patch
    observable $Y$. Its quasi-local image is therefore $j_{a-t,w}(Y)$.
    The equivalent form follows because the inclusions preserve products and
    adjoints.
\end{proof}

\begin{lemma}[Finite-volume commutator limit]
    \label{thm:quasi_local_commutator_energy_limit}
    \lean{SpinChain.eventually_quasiLocalIntervalObservable_bulkObservable_commutator_eq, SpinChain.tendsto_quasiLocalIntervalObservable_bulkObservable_commutator, SpinChain.tendsto_apply_quasiLocalIntervalObservable_bulkObservable_commutator, SpinChain.tendsto_apply_quasiLocalIntervalObservable_commutator}
    \leanok
    Under the preceding assumptions, if $\ell,r\to\infty$ along any common
    filter, the quasi-local energy is eventually equal to $j_{a-t,w}(Y)$ and
    hence converges to it in norm. Consequently every functional $\omega$ on the
    quasi-local algebra satisfies
    \begin{align}
        \omega(x^*[\widehat H_N,x])\longrightarrow\omega(j_{a-t,w}(Y)).
        \notag
    \end{align}
    Continuity is unnecessary here because the observable is eventually constant.
\end{lemma}
\begin{proof}\leanok
    \uses{thm:quasi_local_commutator_locality}
    Both growing margins eventually exceed $t$. Apply the fixed-support identity
    on that tail; the convergence statements follow from eventual constancy.
\end{proof}

\begin{definition}[The finite-interval ground projection]
    \label{def:mps_ground_space_projection_matrix}
    \lean{MPSTensor.groundSpaceProjectionMatrix,
        MPSTensor.isStarProjection_groundSpaceProjectionMatrix,
        MPSTensor.isStarProjection_quasiLocalIntervalObservable_groundSpaceProjectionMatrix}
    \leanok
    For an MPS tensor $A$, let $p_A(N)$ be the matrix of the orthogonal
    projection onto $G_A(N)$ in numbered coordinates on $N$ sites.
    It is a self-adjoint idempotent. Its inclusion $j_{a,N}(p_A(N))$
    into the quasi-local algebra on the interval beginning at $a$ is also
    a self-adjoint idempotent, including when $N=0$.
    \uses{def:ground_space_es, thm:quasi_local_interval_observable}
\end{definition}
\begin{proof}\leanok
    Coordinate identification and the interval inclusion preserve products
    and adjoints, hence preserve orthogonal projections.
\end{proof}

\begin{lemma}[Support compression of a positive functional]
    \label{thm:positive_functional_support_compression}
    \lean{PositiveLinearMap.apply_projection_compression_eq,
        MPSTensor.norm_quasiLocalState_interval_sub_scalar_le_projection_compression}
    \leanok
    \uses{def:qca_state_space_purity}
    Let $f$ be a positive complex-linear functional on a star-ordered
    complex algebra in which every element admits a decomposition into
    self-adjoint real and imaginary parts. If $p=p^*=p^2$ and $f(p)=f(1)$,
    then, for every $X$,
    \begin{align}
        f(pXp)=f(X).
        \notag
    \end{align}
    In particular, let $\varphi$ be a normalized positive norm-one state
    on the quasi-local algebra. For an interval projection $p$ with
    $\varphi(j_{a,N}(p))=1$, an interval observable $X$, and $c\in\mathbb C$,
    \begin{align}
        |\varphi(j_{a,N}(X))-c|\leq\|pXp-cp\|.
        \notag
    \end{align}
    No translation invariance is assumed.
\end{lemma}
\begin{proof}\leanok
    \uses{thm:quasi_local_interval_observable}
    Put $q=1-p$. Then $f(q)=f(q^*q)=0$. Cauchy--Schwarz gives
    $f(qX)=f(Xq)=0$, which proves the compression identity.
    Restrict the quasi-local state to the numbered interval matrix algebra,
    apply this identity, and use its norm-one bound on $pXp-cp$.
    The interval inclusion preserves the operator norm.
\end{proof}

\begin{lemma}[Uniqueness of a state with primitive MPS support]
    \label{thm:primitive_supported_quasi_local_state_unique}
    \lean{MPSTensor.IsPrimitiveMPS.quasiLocalState_interval_eq_insertion_of_groundSpace_support,
        MPSTensor.IsPrimitiveMPS.eq_quasiLocalExpectation_of_groundSpace_support}
    \leanok
    \uses{def:mps_ground_space_projection_matrix}
    Let $A$ be a trace-preserving primitive MPS tensor of positive physical
    and bond dimensions, with positive definite invariant matrix $\rho$. Let $\omega_A$ be its constructed quasi-local state.
    Suppose that $\varphi$ is any normalized positive norm-one continuous
    complex-linear state on the quasi-local algebra and that
    \begin{align}
        \forall\,a\in\mathbb Z,\ N\geq0:\ \varphi(j_{a,N}(p_A(N)))=1.
        \notag
    \end{align}
    Then every interval observable has the prescribed insertion expectation:
    \begin{align}
        \varphi(j_{a,k}(X))=\eta_A(X),
        \notag
    \end{align}
    and consequently $\varphi=\omega_A$. The competing state need not be
    translation invariant. This uniqueness statement is a consequence of
    primitive boundary convergence, as used in
    \cite[Sections~3 and~6]{Nachtergaele1996Spectral}; purity is a separate
    consequence of uniqueness and inheritance of support in convex
    decompositions.
\end{lemma}
\begin{proof}\leanok
    \uses{thm:positive_functional_support_compression,thm:primitive_bulk_ground_compression,thm:qca_compatible_state_extension,lem:qca_quasi_local_observable}
    Include a fixed interval observable into intervals expanded by $n$ sites
    on either side. Write $P_n$ for the ground projection of the expanded
    interval and $X_n$ for the included observable. The preceding estimate
    gives
    \begin{align}
        |\varphi(j_{a,k}(X))-\eta_A(X)|
        \leq\|P_nX_nP_n-\eta_A(X)P_n\|.
        \notag
    \end{align}
    Primitive scalar compression makes the right-hand side tend to zero.
    Thus the states agree on every interval observable. Every finite-region
    observable embeds into a containing interval, and their union is dense
    in the quasi-local algebra. Continuity proves equality on the completion.
\end{proof}

\begin{lemma}[Expectation one of MPS support projections]
    \label{thm:mps_support_projection_expectation_one}
    \lean{MPSTensor.quasiLocalExpectation_groundSpaceProjection_eq_one}
    \leanok
    Let $d>0$ and let $A$ be trace preserving. Let $\rho\ge0$ be an invariant
    virtual matrix of nonzero trace. For every $a\in\Z$ and $N\ge0$, the
    constructed state satisfies
    $\omega_A(j_{a,N}(P_{G_N(A)}))=1$.
    Primitivity and faithfulness of $\rho$ are not required.
\end{lemma}
\begin{proof}
    \uses{def:mps_ground_space_projection_matrix,thm:mps_state_local_support,thm:mps_quasi_local_expectation_state}
    \leanok
    The matrix $I-P_{G_N(A)}$ annihilates $G_N(A)$, so its translated
    expectation is zero by the local support theorem. Since $\omega_A(I)=1$,
    linearity gives expectation one on the projection.
\end{proof}

\begin{lemma}[Purity of the primitive quasi-local MPS state]
    \label{thm:primitive_quasi_local_purity}
    \lean{MPSTensor.IsPrimitiveMPS.isPureQuasiLocalState_quasiLocalExpectation}
    \leanok
    Let $d,D>0$, let $A$ be a trace-preserving primitive tensor, and let
    $\rho>0$ be its invariant virtual matrix. The completed MPS state
    $\omega_A$ is pure among all states of the quasi-local algebra.
    The states in a proposed convex decomposition need not be translation
    invariant. This proves the primitive canonical case of the purity
    assertion in \cite[Section~3]{Nachtergaele1996Spectral}; no classification
    of the multiblock ground-state face is asserted here.
\end{lemma}
\begin{proof}
    \uses{thm:qca_purity_of_supported_uniqueness,thm:mps_support_projection_expectation_one,thm:primitive_supported_quasi_local_state_unique,thm:mps_quasi_local_expectation_state,def:mps_ground_space_projection_matrix}
    \leanok
    Apply the supported-state extremality criterion to the family of
    all translated finite-interval MPS projections. They have expectation
    one in $\omega_A$, and every state satisfying the same support constraints
    equals $\omega_A$ by primitive supported-state uniqueness.
\end{proof}

\begin{lemma}[Positivity of a locally supported commutator expectation]
    \label{thm:local_commutator_expectation_positive}
    \lean{MPSTensor.observableInsertionExpectation_adjoint_commutator_nonneg}
    \leanok
    Let $A$ be a tensor, let $\rho\geq0$ be a virtual matrix, and let $H\geq0$
    act on $N$ physical sites with $G_N(A)\subseteq\ker H$. For every $N$-site
    matrix $B$,
    \begin{align}
        \eta_{A,\rho}(B^*[H,B])\geq0.
        \notag
    \end{align}
    In particular, the commutator expectation is real. Neither primitivity nor a
    spectral gap is needed.
\end{lemma}
\begin{proof}\leanok
    \uses{thm:mps_state_local_support,thm:primitive_insertion_local_state}
    The matrix $B^*BH$ annihilates $G_N(A)$, so its expectation vanishes.
    The remaining matrix $B^*HB$ is positive, and the insertion expectation is a
    positive linear functional.
\end{proof}

\begin{lemma}[Reality and positivity of sector commutator patches]
    \label{thm:sector_commutator_patch_positive}
    \lean{MPSTensor.IsParentInteraction.localCommutatorObservable_nonneg, MPSTensor.IsParentInteraction.localCommutatorObservable_im_eq_zero, MPSTensor.IsParentInteraction.quasiLocalExpectation_localCommutatorObservable_nonneg}
    \leanok
    Let $T$ be the weighted sum of tensors $A_\alpha$, with nonzero weights.
    Suppose $h\geq0$ acts on $R>0$ sites and has kernel $G_R(T)$. For a $k$-site
    matrix $X$ with $k>0$, put $t=R-1$, $w=k+2t$, and
    \begin{align}
        Y=(I_t\otimes X\otimes I_t)^*
        [H_w,I_t\otimes X\otimes I_t],
        \notag
    \end{align}
    where $H_w$ is the open-chain sum of $h$. For every block $\alpha$ and every
    virtual matrix $\rho_\alpha\geq0$,
    \begin{align}
        \eta_{A_\alpha,\rho_\alpha}(Y)\geq0,
        \qquad \operatorname{Im}\eta_{A_\alpha,\rho_\alpha}(Y)=0.
        \notag
    \end{align}
    If $A_\alpha$ is trace preserving and $\rho_\alpha$ is an invariant matrix
    of nonzero trace, the corresponding constructed quasi-local state satisfies
    $\omega_\alpha(j_{a,w}(Y))\geq0$ for every interval position $a$.
    No simultaneous injectivity bound is imposed.
\end{lemma}
\begin{proof}\leanok
    \uses{thm:local_commutator_expectation_positive,thm:open_parent_interaction_kernel_equal,thm:mps_quasi_local_expectation_state,thm:block_ground_space_le_joint,def:finite_open_interaction_matrix,lem:open_local_interaction_sum_properties,thm:open_ground_space_le_hamiltonian_kernel}
    The positive parent sum and the canonical open Hamiltonian have the same
    kernel. The canonical open Hamiltonian annihilates $G_w(T)$, which contains
    $G_w(A_\alpha)$. Apply the preceding lemma to $H_w$ and the padded observable.
    The quasi-local statement follows from exact local evaluation of the
    constructed state.
\end{proof}

\begin{lemma}[Infinite-volume commutator gap in the primitive sectors]
    \label{thm:multiblock_quasi_local_commutator_limit_gap}
    \lean{MPSTensor.exists_pos_multiblock_quasiLocalCommutator_limit_gap}
    \leanok
    Let $A=\bigoplus_{\alpha=0}^{b-1}\mu_\alpha A_\alpha$, with nonzero
    coefficients, trace-preserving primitive blocks, and faithful invariant
    density matrices. Suppose that the blocks have simultaneous injectivity
    length $S>0$, and that $h\geq0$ is a parent interaction of range
    $R\geq S+1$. Then there is $\gamma>0$ such that, for every constructed
    sector state $\omega_\alpha$ and every local observable $X$ satisfying
    $\omega_\alpha(X)=0$,
    \begin{align}
        \lim_{\ell,r\to\infty}
        \omega_\alpha\bigl(X^*[H_{[a-\ell,a+k+r)},X]\bigr)
        \geq \gamma\,\omega_\alpha(X^*X).
        \notag
    \end{align}
    Here $X$ is supported on $[a,a+k)$, with $k>0$, and the limit may be
    taken along any filter for which both margins tend to infinity.
    The limit is real and nonnegative. This proves the displayed
    infinite-volume estimate of \cite[Theorem~1.2]{Nachtergaele1996Spectral}
    for these constructed sectors at the stated sufficient range;
    identification of the entire source ground-state face is a separate result.
\end{lemma}
\begin{proof}\leanok
    \uses{thm:multiblock_quasi_local_parent_commutator_gap,thm:quasi_local_commutator_energy_limit,thm:local_commutator_expectation_positive}
    Once both margins contain $R-1$ sites, the quasi-local energy observable
    equals the inclusion of the fixed commutator patch. Thus its expectation
    is eventually constant. Local support and positivity make this expectation
    real and nonnegative. The common sector gap bounds its real part by
    $\gamma\omega_\alpha(X^*X)$, and positivity of the state makes the latter
    real as well.
\end{proof}

\begin{lemma}[Joint compression from sector column limits]
    \label{thm:sector_observable_decomposition_from_columns}
    \lean{Submodule.tendsto_norm_joint_compression_sub_sector_sum_zero}
    \leanok
    Let $J$ be finite, let $G_j(n)$ be subspaces of finite-dimensional complex
    Hilbert spaces, and let $P_j(n)$ and $P(n)$ be the projections onto
    $G_j(n)$ and $\sum_jG_j(n)$, respectively. Fix $c_j\in\C$.
    Suppose $\|\sum_jP_j(n)-P(n)\|\to0$, the operators $X_n$ are eventually
    uniformly bounded, and
    $\|P(n)X_nP_j(n)-c_jP_j(n)\|\to0$ for every $j$.
    Then $\|P(n)X_nP(n)-\sum_jc_jP_j(n)\|\to0$ along the same filter.
\end{lemma}
\begin{proof}
    \leanok
    Write $Q=\sum_jP_j$. Expanding the rightmost joint projection gives
    $PXP-\sum_jc_jP_j=PX(P-Q)+\sum_j(PXP_j-c_jP_j)$.
    Since $\|P\|\le1$,
    \begin{align}
        \label{eq:sector_joint_compression_error}
        \left\|PXP-\sum_jc_jP_j\right\|
         & \le \|Q-P\|\,\|X\|+\sum_j\|PXP_j-c_jP_j\|.
    \end{align}
    The finite sum tends to zero, as does the projection error multiplied
    by the uniform norm bound. Apply~(\ref{eq:sector_joint_compression_error}).
\end{proof}

\begin{lemma}[Sector limits of joint MPS compressions]
    \label{thm:multiblock_bulk_observable_column_limit}
    \lean{MPSTensor.bulkObservable_iSup_groundSpace_column_tendsto_scalar}
    \leanok
    Let $(A_j)_{j\in J}$ be a finite family of trace-preserving primitive
    tensors of positive bond dimensions, with faithful invariant matrices
    $\rho_j$. Suppose distinct tensors are inequivalent up to virtual gauge
    and phase. Fix a sector $\alpha$, an observable $X$ on $k>0$ sites,
    and margins $\ell_n,r_n\to\infty$ along any filter. Put
    $N_n=\ell_n+k+r_n$, let $P_j(n)=P_{G_{N_n}(A_j)}$, let $P(n)$ project
    onto $\sum_jG_{N_n}(A_j)$, and put
    $X_n=I_{\ell_n}\otimes X\otimes I_{r_n}$.
    Then $\|P(n)X_nP_\alpha(n)-\eta_{A_\alpha}(X)P_\alpha(n)\|\to0$.
\end{lemma}
\begin{proof}
    \uses{thm:multiblock_centered_full_ground_compression,thm:primitive_insertion_local_state,thm:bulk_observable_algebra}
    \leanok
    Center $X$ as $Y=X-\eta_{A_\alpha}(X)I$. Linearity and normalization
    give $\eta_{A_\alpha}(Y)=0$. Interior inclusion preserves this
    centering, and $P(n)P_\alpha(n)=P_\alpha(n)$.
    The centered full compression limit therefore gives the claimed limit.
\end{proof}

\begin{lemma}[Asymptotic sector decomposition of the compressed observable]
    \label{thm:multiblock_bulk_observable_sector_decomposition}
    \lean{MPSTensor.bulkObservable_iSup_groundSpace_sector_decomposition_tendsto_zero}
    \leanok
    Under the hypotheses and notation of
    Lemma~\ref{thm:multiblock_bulk_observable_column_limit}, with no
    distinguished sector, one has
    $\|P(n)X_nP(n)-\sum_j\eta_{A_j}(X)P_j(n)\|\to0$.
    The finite-length projections need not be mutually orthogonal.
\end{lemma}
\begin{proof}
    \uses{thm:sector_observable_decomposition_from_columns,thm:multiblock_bulk_observable_column_limit,thm:sector_projection_sum_error_decay,thm:inequivalent_block_ground_space_overlap_decay,thm:bulk_observable_algebra}
    \leanok
    The sector overlap limits and $N_n\to\infty$ imply
    $\|\sum_jP_j(n)-P(n)\|\to0$. Interior inclusion gives
    $\|X_n\|\le\|X\|$. Apply the joint-compression theorem with
    $c_j=\eta_{A_j}(X)$ and the sector column limits.
\end{proof}

\begin{lemma}[Zero energy and a support projection]
    \label{thm:positive_functional_zero_energy_support}
    \lean{PositiveLinearMap.apply_projection_eq_one_of_le_of_apply_eq_zero,
        MPSTensor.quasiLocalState_groundSpaceProjection_eq_one_of_interval_eq_zero}
    \leanok
    \uses{def:mps_ground_space_projection_matrix,def:qca_state_space_purity}
    Let $f$ be a positive normalized complex-linear functional and
    $p=p^*=p^2$. If $f(H)=0$ and $\kappa(1-p)\leq H$ for some
    $\kappa>0$, then $f(p)=1$.
    In particular, if $H\geq0$ is a finite-interval Hamiltonian with
    $\ker H=G_N(A)$, then every normalized positive quasi-local state
    satisfying $\varphi(j_{a,N}(H))=0$ also satisfies
    $\varphi(j_{a,N}(p_A(N)))=1$. No uniform gap is required.
    This is the finite-dimensional support implication used in
    \cite[Section~3]{Nachtergaele1996Spectral}.
\end{lemma}
\begin{proof}\leanok
    \uses{lem:positive-kernel-projection-comparison}
    Positivity gives $f(1-p)\geq0$, whereas monotonicity gives
    $\kappa f(1-p)\leq f(H)=0$. Thus $f(1-p)=0$ and $f(p)=1$.
    A finite positive Hamiltonian dominates a positive multiple of its
    kernel-complement projection. Restrict the quasi-local state to the
    interval matrix algebra and apply the preceding argument.
\end{proof}

\begin{lemma}[Local zero energy implies joint MPS support]
    \label{thm:quasi_local_parent_zero_energy_joint_support}
    \lean{MPSTensor.%
        quasiLocalState_groundSpaceProjection_eq_one_of_isParentInteraction_wordTupleSpanTop,
        MPSTensor.eventually_quasiLocalState_symmetric_groundSpaceProjection_eq_one}
    \leanok
    Let $A=\bigoplus_\alpha\mu_\alpha A_\alpha$ have nonzero weights and
    simultaneous injectivity length $S>0$. Let $h\geq0$ have
    $\ker h=G_R(A)$ at range $R\geq S+1$.
    If a normalized positive quasi-local state $\varphi$ satisfies
    $\varphi(j_{a,R}(h))=0$ for every $a\in\mathbb Z$, then
    \begin{align}
        \forall\,a\in\mathbb Z,\ N\geq R:\ \varphi(j_{a,N}(p_A(N)))=1.
        \notag
    \end{align}
    Consequently the same support constraint holds eventually on the
    symmetric intervals $[-n,n+1)$. The state need not be translation
    invariant, and primitive normalization of the original blocks is
    unnecessary. This is the local support consequence of the zero-energy
    condition in \cite[Theorem~1.1]{Nachtergaele1996Spectral}, at the stated
    sufficient range.
\end{lemma}
\begin{proof}\leanok
    \uses{thm:quasi_local_open_interaction_zero_expectation,thm:block_open_kernel_simultaneous_injectivity,thm:positive_functional_zero_energy_support}
    Every finite open sum has zero expectation, since each translated local
    term does. Positivity is preserved by the sum, and its kernel is exactly
    $G_N(A)$ when $N\geq R$. The preceding support lemma applies.
    Since $2n+1\geq R$ eventually, it also applies to the symmetric intervals.
\end{proof}

\begin{lemma}[Local zero energy and finite open-chain expectations]
    \label{thm:quasi_local_open_interaction_zero_expectation}
    \lean{SpinChain.quasiLocalIntervalObservable_openInteractionMatrix, SpinChain.apply_quasiLocalIntervalObservable_openInteractionMatrix_eq_zero}
    \leanok
    Let $h$ be an $R$-site matrix, where $R>0$, and let $H_N$ be its open-chain
    sum. For every interval position $a$ and every $N\geq0$,
    \begin{align}
        j_{a,N}(H_N)=\sum_{i=0}^{N-R}j_{a+i,R}(h),
        \notag
    \end{align}
    with the sum understood as empty when $N<R$.
    Consequently, any complex-linear functional $\varphi$ satisfying
    $\varphi(j_{a,R}(h))=0$ for every $a\in\mathbb Z$ also satisfies
    \begin{align}
        \varphi(j_{a,N}(H_N))=0
        \notag
    \end{align}
    for every $a$ and $N$. This applies in particular to arbitrary continuous
    states, without requiring translation invariance.
\end{lemma}
\begin{proof}\leanok
    \uses{thm:quasi_local_bulk_inclusion,def:finite_open_interaction_matrix}
    Every nonwrapping interaction window has quasi-local image $j_{a+i,R}(h)$.
    The interval inclusion preserves the finite sum. Apply linearity of
    $\varphi$ and the vanishing hypothesis to its summands.
\end{proof}

\begin{definition}[The positive finite-interval restriction]
    \label{def:positive_interval_state_restriction}
    \lean{MPSTensor.intervalStateFunctional, MPSTensor.intervalStateFunctional_apply}
    \leanok
    Let $\varphi$ be a continuous complex-linear functional on the quasi-local
    algebra that is nonnegative on every element $Y^*Y$. Its restriction to
    the numbered matrix algebra on an interval is the positive functional
    $f_{a,N}(X)=\varphi(j_{a,N}(X))$.
    \uses{thm:quasi_local_interval_observable}
\end{definition}
\begin{proof}\leanok
    Every positive matrix is a square $Y^*Y$. The interval inclusion preserves
    products and adjoints, so the restriction is positive.
\end{proof}

\begin{lemma}[Projection estimates for a finite-interval state]
    \label{thm:interval_state_projection_estimates}
    \lean{MPSTensor.intervalStateFunctional_projection_mem_unitInterval,
        MPSTensor.intervalStateFunctional_projection_eq_ofReal_re,
        MPSTensor.norm_quasiLocalState_interval_le,
        MPSTensor.quasiLocalState_interval_projection_compression_eq,
        MPSTensor.norm_quasiLocalState_interval_sub_le_projection_compression}
    \leanok
    \uses{def:positive_interval_state_restriction,def:qca_state_space_purity}
    Suppose $\varphi$ is a normalized positive norm-one quasi-local state.
    For an interval projection $p=p^*=p^2$, its expectation
    $\varphi(j_{a,N}(p))$ is real and belongs to $[0,1]$.
    Every interval observable satisfies
    $|\varphi(j_{a,N}(X))|\leq\|X\|$.
    If $\varphi(j_{a,N}(p))=1$, then
    \begin{align}
        \varphi(j_{a,N}(pXp))=\varphi(j_{a,N}(X)),\qquad
        |\varphi(j_{a,N}(X))-\varphi(j_{a,N}(Y))|
        \leq\|pXp-Y\|.
        \notag
    \end{align}
    No translation invariance is assumed.
\end{lemma}
\begin{proof}\leanok
    \uses{thm:positive_functional_support_compression,thm:quasi_local_interval_observable}
    Positivity of $p$ and $1-p$ gives the interval bounds and realness.
    The norm-one condition and the isometric interval inclusion give the
    operator-norm estimate. The support-compression identity reduces the
    final difference to the expectation of $pXp-Y$.
\end{proof}

\begin{lemma}[Compactness of finite convex weights]
    \label{thm:finite_convex_weights_compactness}
    \lean{exists_convex_weights_subsequence,
        exists_convex_weights_of_approximate_evaluations}
    \leanok
    Let $J$ be finite, and let $w_j(n)$ be real weights that eventually lie
    in $[0,1]$, with $\sum_jw_j(n)\to1$. There are a strictly increasing
    subsequence $n_s$ and weights $p_j\geq0$, $\sum_jp_j=1$, such that
    $w_j(n_s)\to p_j$ for every $j$.
    For any family of tests $x$, if
    $f(x)-\sum_jw_j(n)g_j(x)\longrightarrow0$ for every test $x$,
    then one choice of normalized weights satisfies
    $f(x)=\sum_jp_jg_j(x)$ for all tests. The test family need not be countable.
\end{lemma}
\begin{proof}\leanok
    Extract a convergent subsequence in the compact cube $[0,1]^J$.
    Continuity of the finite sum shows that the limit has total mass one.
    For each test, compare the assumed limit with the evaluation along
    this same subsequence. Uniqueness of limits gives the asserted identity.
\end{proof}

\begin{lemma}[Pure sector states and the commutator gap from the original blocks]
    \label{thm:original_blocks_pure_sector_commutator_limit_gap}
    \lean{MPSTensor.%
        exists_pure_primitiveFamily_quasiLocalCommutator_limit_gap_of_wordTupleSpanTop}
    \leanok
    Let the original block family have simultaneous injectivity length $S>0$
    and nonzero weights. For every positive parent interaction of the original
    weighted tensor at range $R\geq S+1$, there are trace-preserving primitive
    representatives, with faithful invariant matrices, preserving every
    block ground space at every length. Their completed sector states are
    pure among all quasi-local states and have a common positive
    infinite-volume commutator gap for the original interaction.
    The original blocks need not be normalized, and the invariant matrices
    and purity are conclusions. The two free margins may tend to infinity
    along any filter, as in the preceding theorem.
\end{lemma}
\begin{proof}\leanok
    \uses{thm:block_word_span_primitive_normalization,thm:multiblock_quasi_local_commutator_limit_gap,thm:primitive_quasi_local_purity}
    Normalize the blocks without changing their local ground spaces or their
    simultaneous span. Equality of the block spaces preserves the local
    kernel of the original positive interaction. Apply the common sector
    commutator theorem to the representatives and the primitive purity theorem
    to each completed sector state.
\end{proof}

\begin{lemma}[Ground-state support from the eventual kernel identity]
    \label{thm:quasi_local_zero_energy_support_eventual_kernel}
    \lean{MPSTensor.%
        eventually_quasiLocalState_symmetric_groundSpaceProjection_eq_one_of_eventual_kernel}
    \leanok
    Let $h\geq0$ have positive finite interaction range. Suppose that
    $\ker H_N(h)=G_N(A)$ for every sufficiently large $N$.
    If a normalized positive quasi-local state $\varphi$ has
    $\varphi(j_{a,R}(h))=0$ for every $a\in\mathbb Z$, then, for every sufficiently large $n$,
    \begin{align}
        \varphi\bigl(j_{-n,2n+1}(p_A(2n+1))\bigr)=1.
        \notag
    \end{align}
    This uses the eventual kernel hypothesis in
    \cite[Theorem~1.2]{Nachtergaele1996Spectral} directly. It imposes no
    local parent-kernel equality and no lower bound on the interaction
    range in terms of injectivity.
\end{lemma}
\begin{proof}\leanok
    \uses{thm:quasi_local_open_interaction_zero_expectation,thm:positive_functional_zero_energy_support}
    Every open interaction sum has zero expectation. For sufficiently large
    symmetric intervals it is positive and has the prescribed MPS kernel.
    Its finite-dimensional kernel-complement comparison therefore forces
    expectation one on the MPS support projection.
\end{proof}

\begin{definition}[The zero-energy state set of a local interaction]
    \label{def:quasi_local_parent_ground_state_face}
    \lean{MPSTensor.parentGroundStateFace}
    \leanok
    Let $d>0$ and let $h$ be an observable on $R$ consecutive sites.
    Define $\mathcal F_h$ to be the set of all states $\phi$ of the
    quasi-local algebra such that $\phi(j_{a,R}(h))=0$ for every $a\in\Z$.
    Translation invariance of $\phi$ is not required.
    For a positive parent interaction, this is its zero-energy state set,
    as in \cite[Section~1 and Theorem~1.1]{Nachtergaele1996Spectral}.
\end{definition}

\begin{lemma}[Convexity and the face property]
    \label{thm:quasi_local_parent_ground_state_face_properties}
    \lean{MPSTensor.convex_parentGroundStateFace,
        MPSTensor.isExtreme_parentGroundStateFace}
    \leanok
    The set $\mathcal F_h$ is real convex for every local observable $h$.
    If $h\ge0$, it is a face of the full state space: whenever
    $\omega=(1-t)\phi+t\psi\in\mathcal F_h$, with $0<t<1$ and
    $\phi,\psi$ states, both constituents belong to $\mathcal F_h$.
\end{lemma}
\begin{proof}
    \uses{thm:quasi_local_state_space_convex,thm:qca_state_convex_support,def:quasi_local_parent_ground_state_face}
    \leanok
    The zero-expectation equations are linear, and the state space is
    convex. For the face assertion factor $h=Y^*Y$. Each translate is
    the adjoint square $j_{a,R}(Y)^*j_{a,R}(Y)$. Vanishing of its
    expectation in a nontrivial convex combination forces vanishing in
    each constituent, by positivity.
\end{proof}

\begin{lemma}[Constructed states in the parent zero-energy face]
    \label{thm:mps_quasi_local_parent_ground_state_face_membership}
    \lean{MPSTensor.quasiLocalExpectation_mem_parentGroundStateFace_of_groundSpaceES_le_ker,
        MPSTensor.quasiLocalExpectation_mem_parentGroundStateFace_toTensorFromBlocks}
    \leanok
    Let $A$ be trace preserving and let $\rho\ge0$ be invariant with
    nonzero trace. If $G_R(A)\subseteq\ker h$, then the constructed
    state $\omega_A$ belongs to $\mathcal F_h$.
    Consequently, if $B=\bigoplus_j\mu_jA_j$, with all $\mu_j\ne0$, and
    $h$ is a positive parent interaction for $B$, every sector for which
    such trace-preserving and invariant-matrix data are given determines
    a state in $\mathcal F_h$. Neither primitivity nor faithfulness of
    the invariant matrix is required.
\end{lemma}
\begin{proof}
    \uses{thm:mps_state_local_support,thm:mps_quasi_local_expectation_state,thm:block_sum_local_ground_space,def:quasi_local_parent_ground_state_face,thm:block_ground_space_le_joint}
    \leanok
    Exact local support gives $\omega_A(j_{a,R}(h))=0$ at every position,
    while normalization and positivity make $\omega_A$ a state.
    For the weighted block tensor,
    $G_R(A_j)\subseteq\sum_iG_R(A_i)=G_R(B)=\ker h$.
    Apply the same support assertion to each sector.
\end{proof}

\begin{lemma}[Convex decomposition of a state with joint MPS support]
    \label{thm:multiblock_supported_quasi_local_state_convex_combination}
    \lean{MPSTensor.exists_convex_combination_quasiLocalExpectation_of_eventually_groundSpace_support}
    \leanok
    \uses{def:mps_ground_space_projection_matrix,def:mps_quasi_local_expectation}
    Let $(A_j)_{j=0}^{b-1}$ be finitely many trace-preserving primitive tensors
    of positive physical and bond dimensions, with positive definite invariant matrices $\rho_j$. Suppose distinct sectors are
    inequivalent up to virtual gauge and phase. Let $T$ be their weighted
    block sum with nonzero coefficients, and let $\omega_j$ be the
    constructed quasi-local state of the $j$th sector.
    Suppose $\varphi$ is a normalized positive norm-one continuous
    complex-linear state on the quasi-local algebra and that, for all
    sufficiently large $n$,
    \begin{align}
        \varphi\bigl(j_{-n,2n+1}(p_T(2n+1))\bigr)=1.
        \notag
    \end{align}
    Then there are real weights $w_j\geq0$, $\sum_jw_j=1$, such that
    \begin{align}
        \varphi=\sum_jw_j\omega_j.
        \notag
    \end{align}
    The competing state need not be translation invariant. This is the
    support-to-convex-hull implication in the argument for
    \cite[Theorem~1.1]{Nachtergaele1996Spectral}.
\end{lemma}
\begin{proof}\leanok
    \uses{thm:interval_state_projection_estimates,thm:finite_convex_weights_compactness,thm:multiblock_bulk_observable_sector_decomposition,thm:sector_projection_sum_error_decay,thm:inequivalent_block_ground_space_overlap_decay,thm:quasi_local_interval_observable,thm:mps_quasi_local_interval_expectation}
    On $I_n=[-n,n+1)$, put
    $w_j(n)=\varphi(j_{-n,2n+1}(p_{A_j}(2n+1)))$.
    Positivity gives real weights in $[0,1]$. The sum of the sector
    projections approaches the joint projection in operator norm.
    The support condition and the norm-one bound therefore give
    $\sum_jw_j(n)\to1$.

    Fix an observable $X$ on $I_m$. For $n\geq m$, include it into $I_n$
    with equal margins $n-m$. Support compression and asymptotic sector
    decomposition give
    \begin{align}
        \varphi(j_{-m,2m+1}(X))-
        \sum_jw_j(n)\eta_{A_j}(X)\longrightarrow0.
        \notag
    \end{align}
    Extract one convergent subsequence of the finite weight vectors.
    Its limit is nonnegative with total mass one. The same subsequence
    identifies all centered interval expectations with the corresponding
    convex combination. Every finite region lies in such an interval,
    and local observables are dense. Continuity proves the state identity
    on the completed algebra.
\end{proof}

\begin{lemma}[Sector zero energy from a finite open kernel]
    \label{thm:sector_local_zero_from_open_kernel}
    \lean{MPSTensor.quasiLocalExpectation_interval_eq_zero_of_open_groundSpaceES_le_ker}
    \leanok
    Let $T=\bigoplus_j\mu_jA_j$ with every $\mu_j\neq0$, and let $A_\alpha$
    be a trace-preserving block with positive invariant density
    $\rho_\alpha$ of nonzero trace, and let $\omega_\alpha$ be its constructed
    quasi-local state. Let $h$ act on $R>0$ sites. Suppose that, for some
    $N\geq R$, the joint MPS space $G_N(T)$ is contained in $\ker H_N$, where $H_N$ is
    the open-chain sum of $h$. Then $\omega_\alpha(j_{a,R}(h))=0$ for every
    interval position $a$. Positivity of $h$ is unnecessary.
\end{lemma}
\begin{proof}\leanok
    \uses{thm:quasi_local_open_interaction_zero_expectation,thm:mps_state_local_support,thm:block_ground_space_le_joint,thm:mps_quasi_local_interval_expectation}
    The sector space is contained in the joint MPS space, so
    $\omega_\alpha(j_{a,N}(H_N))=0$. Exact interval evaluation gives the
    same expectation $\eta_\alpha(h)$ for each translate of $h$. Thus
    $(N-R+1)\eta_\alpha(h)=0$. Since $N\geq R$, the coefficient is positive.
\end{proof}

\begin{lemma}[Block spaces inside the joint space]
    \label{thm:block_ground_space_le_joint}
    \lean{MPSTensor.groundSpaceES_block_le_toTensorFromBlocks}
    \leanok
    Let $T=\bigoplus_j\mu_jA_j$ with every $\mu_j\neq0$. For every block
    $\alpha$ and every length $L$, the boundary space $G_L(A_\alpha)$ is
    contained in $G_L(T)$ \cite[Lemma~commutation~(ii)]{Nachtergaele1996Spectral}.
\end{lemma}
\begin{proof}\leanok
    \uses{thm:block_sum_local_ground_space}
    With nonzero weights, the joint boundary space is the sum of the block
    boundary spaces.
\end{proof}

\begin{lemma}[Sector membership from eventual open kernels]
    \label{thm:sector_local_zero_from_eventual_open_kernel}
    \lean{MPSTensor.quasiLocalExpectation_interval_eq_zero_of_eventually_open_kernel}
    \leanok
    Under the same sector-state assumptions, suppose that
    $\ker H_N=\mathcal G_N$ for every sufficiently large $N$, where
    $\mathcal G_N$ is the joint open-boundary MPS space. Then
    $\omega_\alpha(j_{a,R}(h))=0$ for every $a$ and every sector $\alpha$.
    No injectivity condition at the interaction range is required.
    This is the sector-membership consequence of the finite-volume kernel
    hypothesis in \cite[Theorem~1.2]{Nachtergaele1996Spectral}.
\end{lemma}
\begin{proof}\leanok
    \uses{thm:sector_local_zero_from_open_kernel,def:finite_open_interaction_matrix}
    Choose a volume $N\geq R$ at which the kernel equality holds, and apply
    the finite-volume statement.
\end{proof}

\begin{lemma}[Subwindows in active-block coordinates]
    \label{lem:open_interaction_subwindow_coordinates}
    \lean{MPSTensor.cyclicCfg_join_cyclicActiveBlock,
        MPSTensor.extractWindow_join_cyclicActiveBlock,
        MPSTensor.nonwrappingSuffixStartEquiv,
        MPSTensor.periodicLocalInteractionES_apply,
        MPSTensor.periodicLocalInteractionES_conj_cyclicActiveBlock}
    \leanok
    \uses{def:cyclic_active_block_coordinates,
        def:open_local_interaction_sum}
    Suppose $W\leq N$, fix a cyclic start $s$, and let $q<W$ satisfy
    $q+R\leq W$. Replacing the range-$R$ subwindow beginning at $s+q$
    changes only the active configuration. Restriction to this subwindow
    agrees with restriction at $q$ inside the active block. When $R>0$, any local
    operator, extended by the identity, satisfies
    \[
        U_s h_{s+q}^{(N)}U_s^{-1}
        =\bigl(h_q^{(W)}\bigr)^{\oplus[d]^{N-W}}.
    \]
    Its action is obtained by applying $h$ to the restricted window with
    outside coordinates fixed. If $0<R\leq W\leq n\leq N$, translation
    $q\mapsto n-W+q$ also identifies the nonwrapping starts on $W$ sites
    with the range-$R$ starts contained in $[n-W,n)$.
\end{lemma}
\begin{proof}\leanok
    \uses{def:cyclic_active_block_coordinates}
    In active-block coordinates, subwindow replacement changes only the
    first coordinate. The start bijection has inverse subtraction of
    $n-W$. Applying the local operator separately on each complementary
    configuration gives the conjugation identity.
\end{proof}

\begin{definition}[An interaction on an embedded suffix]
    \label{def:open_suffix_interaction_sum}
    \lean{MPSTensor.openSuffixInteractionHamiltonianES}
    \leanok
    \uses{def:open_local_interaction_sum}
    For a range-$R$ operator $h$, let $H_{[n-W,n)}(h)$ be the sum of its
    nonwrapping translates in an $N$-site chain whose starts satisfy
    $n-W\leq i$ and $i+R\leq n$.
\end{definition}

\begin{lemma}[Counting positive terms in complete open windows]
    \label{lem:positive_open_window_counting}
    \lean{MPSTensor.sum_nonwrappingWindowTerms_le,
        MPSTensor.sum_openSuffixInteractionHamiltonianES_le}
    \leanok
    \uses{def:open_suffix_interaction_sum}
    Let $R\leq W$. For any positive operators $T_i$, indexed by the
    nonwrapping range-$R$ starts in an $N$-site chain,
    \[
        \sum_{n=W}^{N}\ \sum_{\substack{n-W\leq i\\i+R\leq n}}T_i
        \leq (W-R+1)\sum_i T_i.
    \]
    In particular, for a positive interaction $h$,
    $\sum_{n=W}^{N}H_{[n-W,n)}(h)
        \leq(W-R+1)H_N^{\mathrm{open}}(h)$.
    This is the counting argument of condition C1 in
    \cite{Nachtergaele1996Spectral}, applied to arbitrary positive terms.
\end{lemma}
\begin{proof}\leanok
    \uses{def:open_suffix_interaction_sum}
    A term beginning at $i$ can occur only for
    $i+R\leq n\leq i+W$, an interval containing $W-R+1$ integers.
    Reverse the order of summation and use positivity of each term.
\end{proof}

\begin{lemma}[Transfer of a finite-window comparison]
    \label{lem:open_interaction_suffix_comparison}
    \lean{MPSTensor.openSuffixInteractionHamiltonianES_conj_cyclicActiveBlock,
        MPSTensor.openSuffixInteractionHamiltonianES_comparison_of_local}
    \leanok
    \uses{lem:open_interaction_subwindow_coordinates,
        def:open_suffix_interaction_sum}
    Suppose $0<R\leq W\leq n\leq N$. The suffix Hamiltonian is unitarily
    equivalent to $H_W^{\mathrm{open}}(h)$ acting independently on the
    complementary configurations. If
    $\kappa p_W(A)\leq H_W^{\mathrm{open}}(h)$, then the same inequality
    holds after embedding both sides on the suffix $[n-W,n)$.
    Here $p_W(A)$ is the canonical complementary ground-space projection;
    no sign condition on $\kappa$ or positivity condition on $h$ is needed
    for this transfer of order.
\end{lemma}
\begin{proof}\leanok
    \uses{lem:open_interaction_subwindow_coordinates}
    Sum the conjugation identity over the corresponding start bijection.
    Extension by the identity and conjugation by an isometry preserve order.
\end{proof}

\begin{lemma}[Open-chain comparison for an arbitrary positive interaction]
    \label{thm:positive_open_interaction_range_comparison}
    \lean{MPSTensor.openInteractionHamiltonianES_comparison_of_local,
        MPSTensor.openInteractionHamiltonianES_gap_of_long_gap_at_length}
    \leanok
    \uses{lem:positive_open_window_counting,
        lem:open_interaction_suffix_comparison,
        lem:open_parent_range_comparison}
    Let $0<R\leq W\leq N$, let $h$ be positive, and suppose
    $\kappa p_W(A)\leq H_W^{\mathrm{open}}(h)$. Then
    \[
        \kappa H_{N,W}^{\mathrm{open}}(A)
        \leq (W-R+1)H_N^{\mathrm{open}}(h).
    \]
    If additionally $\kappa,\gamma>0$, both global kernels equal $G_N(A)$,
    and $\gamma\lVert v\rVert\leq
        \lVert H_{N,W}^{\mathrm{open}}(A)v\rVert$ for $v\perp G_N(A)$, then
    \[
        \frac{\kappa\gamma}{W-R+1}\lVert v\rVert
        \leq\lVert H_N^{\mathrm{open}}(h)v\rVert
        \qquad(v\perp G_N(A)).
    \]
    No local parent-kernel equality is assumed for $h$. This is a
    finite-window consequence of condition C1 and the comparison argument
    in Section~6 of \cite{Nachtergaele1996Spectral}.
\end{lemma}
\begin{proof}\leanok
    \uses{lem:positive_open_window_counting,
        lem:open_interaction_suffix_comparison}
    Sum the embedded local comparisons. Each full range-$W$ parent term
    occurs once, whereas a range-$R$ term occurs at most $W-R+1$ times.
    For the gap, compare the positive operators
    $\kappa H_{N,W}^{\mathrm{open}}(A)$ and
    $(W-R+1)H_N^{\mathrm{open}}(h)$, whose kernels agree, and divide by
    the positive window multiplicity.
\end{proof}

\begin{lemma}[Uniform finite-chain gap from an eventual kernel identity]
    \label{thm:eventual_kernel_open_interaction_gap}
    \lean{MPSTensor.exists_openInteractionHamiltonianES_gap_of_eventual_kernel,
        MPSTensor.exists_openInteractionMatrix_gap_of_eventual_kernel}
    \leanok
    Let $A_j$ be a finite family of normalized primitive MPS tensors with
    faithful fixed density matrices, pairwise inequivalent under gauge and
    phase, and let $B=\bigoplus_j\mu_jA_j$, where every $\mu_j\ne0$.
    Let $h\geq0$ be an interaction of range $R\geq1$.
    Suppose that $\ker H^{\mathrm{open}}_{h,N}=\mathcal G_N(B)$
    for all sufficiently large $N$. Then there exists $\gamma>0$ such that
    \begin{align}
        \gamma\|v\|\leq\|H^{\mathrm{open}}_{h,N}v\|
        \notag
    \end{align}
    for every $N\in\mathbb N$ and every
    $v\perp\ker H^{\mathrm{open}}_{h,N}$.
    No condition on the local kernel of $h$, or lower bound on $R$
    beyond $R\geq1$, is required. This is the finite-chain conclusion of
    \cite[Theorem~1.2]{Nachtergaele1996Spectral}, under its eventual
    finite-interval kernel hypothesis, for the primitive-sector representation;
    the passage from more general periodic presentations is separate
    \cite{gap:nachtergaele96_infinite_volume_ground_projection}.
\end{lemma}
\begin{proof}\leanok
    \uses{thm:primitive_block_sum_open_gap_all_lengths,
        thm:ph_primitive_blocks_open_kernel,
        lem:parent-finite-interval-comparison,
        thm:positive_open_interaction_range_comparison}
    Choose a canonical interaction range $W=2p$ above the original range
    and both kernel thresholds, with canonical gap $\delta>0$ uniform in
    the volume. Finite dimensionality and positivity give
    $\kappa P_{\mathcal G_W(B)^\perp}\leq H^{\mathrm{open}}_{h,W}$
    for some $\kappa>0$. For $N\geq W$, the finite-window comparison gives
    \begin{align}
        \kappa H^{\mathrm{open}}_{N,W}(B)
        \leq (W-R+1)H^{\mathrm{open}}_{h,N},
        \notag
    \end{align}
    The kernels coincide for these volumes, so the transferred gap is
    $\kappa\delta/(W-R+1)$. Taking a minimum with the positive gaps
    on the kernel complements at the finitely many volumes $N<W$
    gives one constant for every volume.
\end{proof}

\begin{lemma}[Classification of the zero-energy state face for primitive sectors]
    \label{thm:primitive_multiblock_quasi_local_ground_state_classification}
    \lean{MPSTensor.mem_parentGroundStateFace_iff_convex_combination_of_eventually_open_kernel}
    \leanok
    \uses{def:quasi_local_parent_ground_state_face}
    Let $T=\bigoplus_{j=0}^{b-1}\mu_jA_j$, with $\mu_j\ne0$, be a finite
    weighted sum of inequivalent trace-preserving primitive tensors with
    faithful invariant density matrices. Write $\omega_j$ for their
    constructed quasi-local states. Let $h\geq0$ have range $R\geq1$,
    and suppose $\ker H_N^{\mathrm{open}}(h)=G_N(T)$ for all sufficiently
    large $N$. Then, for every normalized positive state $\varphi$,
    \begin{align}
        \varphi\in\mathcal F_h
        \quad\Longleftrightarrow\quad
        \exists (w_j)_j:\ w_j\geq0,\quad\sum_jw_j=1,\quad
        \varphi=\sum_j w_j\omega_j.
        \notag
    \end{align}
    No translation invariance of $\varphi$ is assumed. This is the
    ground-state classification associated with
    \cite[Theorems~1.1--1.2]{Nachtergaele1996Spectral} for the stated
    primitive sector family. More general periodic presentations remain
    separate \cite{gap:nachtergaele96_infinite_volume_ground_projection}.
\end{lemma}
\begin{proof}\leanok
    \uses{thm:quasi_local_zero_energy_support_eventual_kernel,thm:multiblock_supported_quasi_local_state_convex_combination,thm:quasi_local_parent_ground_state_face_properties,thm:sector_local_zero_from_eventual_open_kernel}
    Zero local energy and the eventual kernel identity give expectation
    one on the joint support projections of sufficiently large intervals.
    The supported-state decomposition gives the convex coefficients.
    Conversely, each sector state has zero local interaction expectation:
    its expectation of a sufficiently long open sum is zero, and every
    summand has the same expectation. Convexity preserves this property.
\end{proof}

\begin{lemma}[Pure ground states are precisely the primitive sectors]
    \label{thm:primitive_multiblock_pure_ground_state_classification}
    \lean{MPSTensor.isPure_mem_parentGroundStateFace_iff_sector_of_eventually_open_kernel}
    \leanok
    Under the preceding hypotheses, a state $\varphi$ is a pure element of
    $\mathcal F_h$ if and only if $\varphi=\omega_j$ for some $j$.
    Purity is extremality among all states of the quasi-local algebra.
\end{lemma}
\begin{proof}\leanok
    \uses{thm:primitive_multiblock_quasi_local_ground_state_classification,thm:finite_pure_state_decomposition,thm:primitive_quasi_local_purity}
    A pure state in a finite convex combination equals one of the
    constituent states. Each constructed primitive sector state is pure
    and has zero local energy.
\end{proof}

\begin{lemma}[Positive commutator expectation from eventual joint support]
    \label{thm:eventual_kernel_sector_commutator_positive}
    \lean{MPSTensor.quasiLocalExpectation_localCommutatorObservable_nonneg_of_eventual_kernel}
    \leanok
    Let $h\geq0$ have range $R>0$, and suppose that the kernels of its
    sufficiently long open-chain Hamiltonians equal the joint boundary spaces
    of a weighted tensor family with nonzero weights. Let $\omega_\alpha$ be
    a constructed sector state, with a trace-preserving tensor and a positive
    invariant matrix of nonzero trace. For $X$ on $k>0$ sites,
    \begin{align}
        \omega_\alpha\bigl(j_{a-(R-1),\,2(R-1)+k}
        (X^*[H_{\mathrm{patch}},X])\bigr)\geq0.
        \notag
    \end{align}
    Neither primitivity nor faithfulness of the invariant matrix is needed
    for this assertion. The local commutator patch has real expectation.
\end{lemma}
\begin{proof}\leanok
    \uses{thm:local_commutator_expectation_positive,thm:quasi_local_commutator_locality,thm:mps_quasi_local_interval_expectation}
    Choose one sufficiently large volume containing both interaction
    margins. Its Hamiltonian annihilates the sector boundary space and is
    positive. The expectation of $B^*[H,B]$ therefore equals that of
    $B^*HB$ and is nonnegative. The quasi-local commutator identity identifies
    this expectation with the fixed local patch.
\end{proof}

\begin{lemma}[One finite-chain and local commutator gap for all primitive sectors]
    \label{thm:eventual_kernel_sector_local_gap}
    \lean{MPSTensor.exists_pos_multiblock_localCommutator_gap_of_eventual_kernel,
        MPSTensor.exists_pos_multiblock_quasiLocalCommutator_gap_of_eventual_kernel}
    \leanok
    Let $A_1,\ldots,A_b$ be pairwise gauge-phase inequivalent normalized
    primitive tensors with faithful invariant matrices, and let all weights
    be nonzero. Let $h\geq0$ have arbitrary range $R>0$. Suppose that
    $\ker H_N^{\mathrm{open}}(h)$ equals the joint boundary-condition space
    for all sufficiently large $N$. Then there is $\gamma>0$ such that, for every $N$ and every
    $v\perp\ker H_N^{\mathrm{open}}(h)$,
    \begin{align}
        \gamma\lVert v\rVert\leq\lVert H_N^{\mathrm{open}}(h)v\rVert,
        \notag
    \end{align}
    and, for every sector, position, and centered observable $X$ on $k>0$
    sites,
    \begin{align}
        \gamma\omega_\alpha(X^*X)
        \leq\omega_\alpha(X^*[H_{\mathrm{patch}},X]).
        \notag
    \end{align}
    The same inequality holds for the finite-interval insertion expectations.
    No local parent-kernel equality or lower bound on $R$ in terms of an
    injectivity length is assumed. This is the primitive-sector specialization
    of the argument in \cite[Theorem~1.2 and Section~6]{Nachtergaele1996Spectral};
    the more general GVBS presentations are recorded in
    \cite{gap:nachtergaele96_infinite_volume_ground_projection}.
\end{lemma}
\begin{proof}\leanok
    \uses{thm:eventual_kernel_open_interaction_gap,thm:multiblock_local_open_interaction_commutator_gap,thm:mps_quasi_local_interval_expectation}
    Choose the uniform finite-chain constant. Pull the eventual kernel
    equality along growing buffered volumes. Apply the full joint-ground
    projection decay and the finite-volume commutator estimate in each sector.
    Evaluation on consecutive intervals is the insertion expectation.
\end{proof}

\begin{lemma}[A common finite-chain and infinite-volume sector gap]
    \label{thm:eventual_kernel_sector_literal_gap}
    \lean{MPSTensor.exists_pos_multiblock_quasiLocalCommutator_limit_gap_of_eventual_kernel}
    \leanok
    Under the primitive-family hypotheses of
    Lemma~\ref{thm:eventual_kernel_sector_local_gap}, choose the same
    $\gamma>0$ for the finite-chain bound. Let $\ell,r$ tend to infinity
    along any filter. For every constructed sector state and centered local
    $X$ on $k>0$ sites, the finite-volume energies
    \begin{align}
        \omega_\alpha\bigl(X^*[H_{[a-\ell,\,a+k+r)}(h),X]\bigr)
        \notag
    \end{align}
    converge to a real number $e$ satisfying
    $\gamma\omega_\alpha(X^*X)\leq e$. Thus one constant gives both the
    all-volume finite-chain norm gap and the infinite-volume sector bound.
    This states the primitive-sector case explicitly; classification of all
    pure zero-energy states is a separate assertion.
\end{lemma}
\begin{proof}\leanok
    \uses{thm:eventual_kernel_sector_local_gap,thm:eventual_kernel_sector_commutator_positive,thm:quasi_local_commutator_locality}
    Once both margins exceed $R-1$, locality makes the commutator observable
    equal to its fixed patch. Its expectation is real and nonnegative by
    eventual support. Combine this finite-volume limit with the local bound using
    the same finite-chain constant.
\end{proof}

\begin{lemma}[Ground-state classification after primitive normalization]
    \label{thm:normalized_quasi_local_ground_state_classification}
    \lean{MPSTensor.%
        exists_normalized_parentGroundStateFace_classification_of_eventually_open_kernel}
    \leanok
    Let $A_j$ be a finite family with positive bond dimensions whose
    word tuples at some length $S>0$ span the full product matrix algebra.
    This is the simultaneous spanning condition of
    \cite[Section~IV.C]{Cirac2021Matrix}.
    Let $T_A=\bigoplus_j\mu_jA_j$, with every $\mu_j\ne0$.
    Suppose that $h\geq0$ has range $R\geq1$ and that
    $\ker H^{\mathrm{open}}_{h,N}=G_N(T_A)$ for all sufficiently large $N$.
    There exist primitive trace-preserving tensors $B_j$ with faithful
    invariant matrices $\rho_j$ such that their word tuples still span the
    full product algebra at length $S$. Write $T_B=\bigoplus_j\mu_jB_j$. Then
    \begin{align}
        G_N(A_j) & =G_N(B_j), &
        G_N(T_A) & =G_N(T_B).
        \notag
    \end{align}
    These equalities hold for every block $j$ and every length $N$.
    Let $\omega_j$ be the quasi-local state constructed from $B_j,\rho_j$.
    For every continuous complex-linear functional $\varphi$ on the
    quasi-local algebra,
    \begin{align}
        \varphi\in\mathcal F_h
         & \quad\Longleftrightarrow\quad
        \exists (w_j)_j:\ w_j\geq0,\quad\sum_jw_j=1,\quad
        \varphi=\sum_jw_j\omega_j,\notag \\
        \varphi\in\mathcal F_h\ \text{and }\varphi\text{ is pure}
         & \quad\Longleftrightarrow\quad
        \exists j:\ \varphi=\omega_j.\notag
    \end{align}
    The original interaction range need not exceed $S$.
    This combines the normalization in
    \cite[equations~(3.1)--(3.2b)]{Nachtergaele1996Spectral}
    with the ground-state classification of
    \cite[Theorems~1.1--1.2]{Nachtergaele1996Spectral}, for the original
    simultaneously spanning family. The passage from more general periodic
    presentations is separate
    \cite{gap:nachtergaele96_infinite_volume_ground_projection}.
\end{lemma}
\begin{proof}\leanok
    \uses{thm:block_word_span_primitive_normalization,
        thm:block-word-span-separation,thm:ph_block_replacement_parent_equality,
        thm:primitive_multiblock_quasi_local_ground_state_classification,
        thm:primitive_multiblock_pure_ground_state_classification}
    Normalize each block by a nonzero rescaling and a virtual gauge.
    These operations preserve every block ground space and the simultaneous
    word span. The weighted joint ground spaces therefore coincide at every
    length, and the eventual kernel identity transfers to $T_B$.
    Simultaneous spanning excludes gauge-phase equivalence between distinct
    normalized blocks. Apply the convex-combination and pure-sector
    classifications to this normalized family and the unchanged interaction.
\end{proof}

\begin{lemma}[Finite and infinite-volume gaps in every pure ground state]
    \label{thm:eventual_kernel_pure_ground_state_gap}
    \lean{MPSTensor.exists_pos_pure_ground_state_commutator_limit_gap_of_eventual_kernel}
    \leanok
    Let $T=\bigoplus_j\mu_jA_j$ be a finite weighted family of inequivalent
    trace-preserving primitive tensors with faithful invariant density matrices
    and nonzero coefficients. Let $h\geq0$ have range $R\geq1$, and suppose
    $\ker H_N^{\mathrm{open}}(h)=G_N(T)$ for every sufficiently large $N$.
    Let both free margins tend to infinity. Then there is one $\gamma>0$
    such that
    \begin{align}
        \gamma\|v\|\leq\|H_N^{\mathrm{open}}(h)v\|
        \qquad\bigl(N\geq0,\ v\perp\ker H_N^{\mathrm{open}}(h)\bigr),
        \notag
    \end{align}
    and, for every pure $\eta\in\mathcal F_h$ and every local observable
    $X$ with $\eta(X)=0$, the finite-volume limit exists and satisfies
    \begin{align}
        \lim_{M\to-\infty,\,N\to+\infty}
        \eta\bigl(X^*[H_{[M,N]}(h),X]\bigr)
        \geq\gamma\eta(X^*X).
        \notag
    \end{align}
    The margins may diverge along an arbitrary filter. This is the pair
    of conclusions in \cite[Theorem~1.2]{Nachtergaele1996Spectral} for
    the stated primitive family; the original range has no injectivity
    lower bound. General periodic presentations remain separate
    \cite{gap:nachtergaele96_infinite_volume_ground_projection}.
\end{lemma}
\begin{proof}\leanok
    \uses{thm:eventual_kernel_sector_literal_gap,thm:primitive_multiblock_pure_ground_state_classification}
    The sector gap theorem gives one constant for both conclusions.
    The pure-state classification identifies each pure zero-energy state
    with a constructed sector, to which the same finite-volume limit and bound apply.
\end{proof}

\begin{lemma}[The pure ground-state gap from an original spanning family]
    \label{thm:original_family_eventual_kernel_pure_ground_state_gap}
    \lean{MPSTensor.exists_pos_pure_ground_state_commutator_limit_gap_of_wordTupleSpanTop}
    \leanok
    The preceding finite and infinite-volume conclusions also hold when
    the original blocks are only assumed to have full simultaneous word
    span at some positive length $S$. Their normalization and invariant
    density matrices are derived internally. No condition $R\geq S+1$
    is required; positivity and the eventual open-chain kernel identity
    remain the hypotheses on the original interaction.
\end{lemma}
\begin{proof}\leanok
    \uses{thm:normalized_quasi_local_ground_state_classification,thm:eventual_kernel_pure_ground_state_gap}
    Normalize the blocks, preserving every block ground space and every
    weighted joint ground space. The eventual kernel condition transfers
    to this family. Its simultaneous span excludes gauge-and-phase
    equivalence of different blocks. Apply the preceding theorem.
\end{proof}

\begin{lemma}[Existence of a positive interaction for a primitive family]
    \label{thm:primitive_ground_state_interaction_exists}
    \lean{MPSTensor.exists_positive_parent_interaction_of_isPrimitiveMPS}
    \leanok
    \uses{def:open_local_interaction_sum}
    Let $T=\bigoplus_{j=0}^{b-1}\mu_jA_j$, with nonzero weights, be a finite
    family of pairwise gauge-phase inequivalent trace-preserving primitive
    tensors of positive bond dimensions. Suppose each tensor has a faithful
    invariant matrix. There exist $R>0$ and a positive interaction $h$ on
    $R$ sites such that, for every $N\geq R$,
    \begin{align}
        \ker H_N^{\mathrm{open}}(h)=G_N(T).
        \notag
    \end{align}
    Neither $R$ nor a simultaneous injectivity length is given.
    The inequality $N\geq R$ is stronger than the inclusive-interval bound
    $N\geq R+1$ in
    \cite[Theorem~1.1]{Nachtergaele1996Spectral}.
\end{lemma}
\begin{proof}\leanok
    \uses{thm:ph_primitive_blocks_open_kernel,rem:euclidean_local_projector_ingredients}
    Choose the positive open-kernel threshold $R$ and let $h$ be the
    orthogonal projection onto $G_R(T)^\perp$. Its open sum is the canonical
    range-$R$ parent Hamiltonian. The threshold theorem gives the stated
    kernel equality, including the single-window volume $N=R$.
\end{proof}

\begin{lemma}[An interaction realizing the entire primitive-sector ground-state face]
    \label{thm:primitive_ground_state_face_interaction_exists}
    \lean{MPSTensor.exists_positive_parent_interaction_groundStateFace_of_isPrimitiveMPS}
    \leanok
    For the primitive family of the preceding theorem, one can choose a
    single $R>0$ and $h\geq0$ such that the open kernels equal $G_N(T)$
    for every $N\geq R$, and hence for all sufficiently large $N$, and
    \begin{align}
        \mathcal F_h
        =\left\{\sum_{j=0}^{b-1} w_j\omega_j:
        w_j\geq0,\ \sum_jw_j=1\right\}.
        \notag
    \end{align}
    Here $\omega_j$ is the constructed quasi-local state of $A_j$.
    The pure states in $\mathcal F_h$ are precisely the states $\omega_j$.
    Purity is extremality among all normalized positive quasi-local states;
    no translation invariance of a state being classified is assumed.
    This is the stated primitive-family case of the existence assertion in
    \cite[Theorem~1.1 and its finite-family consequence]{Nachtergaele1996Spectral}.
    More general periodic GVBS presentations remain separate
    \cite{gap:nachtergaele96_infinite_volume_ground_projection}.
\end{lemma}
\begin{proof}\leanok
    \uses{thm:primitive_ground_state_interaction_exists,thm:primitive_multiblock_quasi_local_ground_state_classification,thm:primitive_multiblock_pure_ground_state_classification}
    The all-volume kernel bound implies the eventual kernel hypothesis.
    Apply the two classification theorems to the same positive canonical
    interaction. They characterize the entire zero-energy face and its pure
    states without an invariance assumption on the state.
\end{proof}

\begin{lemma}[Ground-space projection observables under blocking]
    \label{thm:blocked_ground_projection_observable}
    \lean{MPSTensor.groundSpaceProjectionMatrix_blockTensor_reindex,
        MPSTensor.quasiLocalBlocking_groundSpaceProjectionMatrix}
    \leanok
    Let $A$ be any MPS tensor, and let $P_A(N)$ denote the matrix of the
    orthogonal projection onto $G_N(A)$. For $L,N\geq0$,
    \begin{align}
        \mathcal R_{N,L}\bigl(P_{A^{[L]}}(N)\bigr)=P_A(NL).
    \end{align}
    If $d>0$ and $L>0$, then for every $a\in\mathbb Z$ and $N\geq0$,
    \begin{align}
        \overline{\mathcal B}_L
        \bigl(j^{(d^L)}_{a,N}(P_{A^{[L]}}(N))\bigr)
        =j^{(d)}_{aL,NL}(P_A(NL)).
    \end{align}
    No normality or translation invariance is assumed. These are coordinate
    consequences of the site regrouping in
    \cite[lines~825--836]{Nachtergaele1996Spectral} and the blocked local
    spaces of \cite[lines~1815--1820 and 1985--1992]{Cirac2021Matrix}.
\end{lemma}
\begin{proof}\leanok
    \uses{thm:blocked_local_ground_space_transport,
        def:mps_ground_space_projection_matrix,
        thm:qca_blocking_quasi_local_interval}
    Evaluate the finite projection conjugacy on a configuration basis vector.
    This gives the matrix reindexing identity. Apply the interval observable
    blocking formula to obtain the completed identity.
\end{proof}

\begin{lemma}[Expectations of aligned ground-space projections]
    \label{thm:blocked_ground_projection_expectation}
    \lean{MPSTensor.quasiLocalBlockingFunctional_groundSpaceProjectionMatrix,
        MPSTensor.quasiLocalBlockingFunctional_symm_groundSpaceProjectionMatrix}
    \leanok
    Let $d,L>0$ and let $T_L\omega=\omega\circ\overline{\mathcal B}_L^{-1}$
    be functional pullback from the blocked chain to the original chain.
    For every continuous complex-linear functional $\omega$ on the blocked
    algebra,
    \begin{align}
        (T_L\omega)\bigl(j^{(d)}_{aL,NL}(P_A(NL))\bigr)
        =\omega\bigl(j^{(d^L)}_{a,N}(P_{A^{[L]}}(N))\bigr).
    \end{align}
    Equivalently, for every continuous complex-linear functional $\varphi$
    on the original algebra,
    \begin{align}
        (T_L^{-1}\varphi)\bigl(j^{(d^L)}_{a,N}(P_{A^{[L]}}(N))\bigr)
        =\varphi\bigl(j^{(d)}_{aL,NL}(P_A(NL))\bigr).
    \end{align}
    Both identities include $N=0$ and arbitrary integer $a$.
\end{lemma}
\begin{proof}\leanok
    \uses{thm:blocked_ground_projection_observable, def:qca_blocking_functional}
    Substitute the projection observable identity into inverse
    precomposition. The inverse evaluation follows by applying the same
    identity to $T_L^{-1}\varphi$.
\end{proof}

\begin{lemma}[Blocked interval support and aligned original support]
    \label{thm:blocked_ground_projection_support}
    \lean{MPSTensor.groundSpace_support_blockTensor_iff_aligned,
        MPSTensor.groundSpace_support_blockTensor_iff_aligned_of_positive_lengths}
    \leanok
    For $d,L>0$, exact support of $T_L^{-1}\varphi$ in every interval
    space of $A^{[L]}$ is equivalent to exact support of $\varphi$ in every
    aligned interval space of $A$:
    \begin{align}
         & \forall a\in\mathbb Z\ \forall N\geq0,\quad
        (T_L^{-1}\varphi)
        \bigl(j^{(d^L)}_{a,N}(P_{A^{[L]}}(N))\bigr)=1
        \\
         & \hspace{2em}\Longleftrightarrow
        \forall a\in\mathbb Z\ \forall N\geq0,\quad
        \varphi\bigl(j^{(d)}_{aL,NL}(P_A(NL))\bigr)=1.
    \end{align}
    The same equivalence holds when both quantifiers over lengths are
    restricted to $N>0$. These identities hold for every continuous
    complex-linear functional; positivity and normalization are not required.
\end{lemma}
\begin{proof}\leanok
    \uses{thm:blocked_ground_projection_expectation}
    Apply the inverse evaluation identity to each pair $(a,N)$, with the
    same restriction on lengths on both sides.
\end{proof}

\begin{lemma}[MPS support descends to contained intervals]
    \label{thm:mps_support_contained_interval}
    \lean{MPSTensor.chainWindowOperator_groundSpaceProjection_mul,
        MPSTensor.quasiLocalState_groundSpaceProjection_eq_one_of_contained_interval}
    \leanok
    Let $A$ be any MPS tensor. Let $P_A(I)$ be the projection onto its
    open-boundary MPS space on a finite nonempty interval $I$.
    If $I\subseteq J$, the projection $P_A(I)$, extended by the identity
    outside $I$, fixes $P_A(J)$. Consequently, every positive normalized
    quasi-local state $\varphi$ with $\varphi(P_A(J))=1$ satisfies
    $\varphi(P_A(I))=1$. Neither injectivity nor translation invariance
    is required. This is the restriction direction of
    \cite[Section~3, equation~(3.12)]{Nachtergaele1996Spectral}.
\end{lemma}
\begin{proof}\leanok
    Every restriction of an open-boundary MPS vector belongs to the MPS
    space on the smaller interval. Thus the smaller interval's extended
    ground projection absorbs the larger ground projection. Positivity
    and Cauchy--Schwarz imply that a state supported in a projection
    evaluates every observable through its compression to that projection.
    Compressing the extended smaller projection therefore gives expectation
    one.
\end{proof}

\begin{lemma}[Aligned blocked intervals determine all MPS support]
    \label{thm:blocked_aligned_mps_support_iff}
    \lean{MPSTensor.quasiLocalState_groundSpaceProjection_eq_one_iff_aligned}
    \leanok
    Let $A$ be any MPS tensor, let $L\geq1$, and let $\varphi$ be a positive
    normalized quasi-local state. Then $\varphi(P_A([a,a+N)))=1$ for every $a\in\mathbb Z$ and $N\geq1$
    if and only if $\varphi(P_A([aL,(a+N)L)))=1$ for every $a\in\mathbb Z$
    and $N\geq1$.
    No primitivity, translation invariance, or spectral gap is assumed.
    This gives the support-space consequence of finite blocking in
    \cite[Section~3]{Nachtergaele1996Spectral}.
\end{lemma}
\begin{proof}\leanok
    \uses{thm:mps_support_contained_interval}
    One implication follows by restricting the collection of intervals.
    For the converse, Euclidean division writes an arbitrary left endpoint
    as $a=cL+b$ with $0\leq b<L$. An aligned interval starting at $cL$
    and containing $b+N$ blocks contains $[a,a+N)$, even when $a<0$.
    Apply the contained-interval support lemma.
\end{proof}

\begin{lemma}[Open MPS spaces from an isometric sector resolution]
    \label{thm:isometric_sector_ground_space_resolution}
    \lean{MPSTensor.groundSpace_eq_iSup_of_isometric_sector_resolution,
        MPSTensor.groundSpaceES_eq_iSup_of_isometric_sector_resolution}
    \leanok
    Let $C$ be an MPS tensor of bond dimension $D$, and let
    $(B_j)_{j\in J}$ be a finite family of tensors with the same physical
    dimension and bond dimensions $D_j$. Suppose that
    $V_j:\C^{D_j}\to\C^D$ satisfy $V_j^\dagger V_j=\Id_{D_j}$,
    $\sum_jV_jV_j^\dagger=\Id_D$, and $C^iV_j=V_jB_j^i$.
    Then $G_N(C)=\bigvee_jG_N(B_j)$ for every $N\ge0$, both in
    computational coordinates and in the physical Hilbert space.
\end{lemma}
\begin{proof}\leanok
    The letterwise intertwining implies $C^wV_j=V_jB_j^w$ for every word,
    including the empty word. Trace cyclicity gives
    \begin{align}
        \Gamma_N^C(V_jYV_j^\dagger) & =\Gamma_N^{B_j}(Y),\notag                      \\
        \Gamma_N^C(X)               & =\sum_j\Gamma_N^{B_j}(V_j^\dagger XV_j).\notag
    \end{align}
    The first identity includes each sector space in the ambient space;
    the second gives the reverse inclusion in their sum.
\end{proof}

\begin{lemma}[Open MPS spaces under prescribed periodic blocking]
    \label{thm:periodic_blocked_ground_space}
    \lean{MPSTensor.IsPeriodic.exists_stepOrbit_groundSpaceDecomposition}
    \leanok
    Suppose that $A$ is periodic of period $m$ and $p>0$.
    Put $r=\gcd(m,p)$. There are nonzero cyclic projections $P_u$ with
    $\sum_uP_u=\Id_D$ and $P_uA^i=A^iP_{u+1}$, positive dimensions
    $D_\alpha$ with $\sum_\alpha D_\alpha=D$, and left-canonical tensors
    $B_\alpha$ indexed by $\alpha\in\{0,\ldots,r-1\}$.
    Their support isometries satisfy
    \begin{align}
        V_\alpha^\dagger V_\alpha & =\Id_{D_\alpha},\notag                      \\
        V_\alpha V_\alpha^\dagger & =
        \sum_{k=0}^{m/r-1}P_{[\alpha+pk]_m},\notag                              \\
        B_\alpha^I                & =V_\alpha^\dagger(A^{[p]})^IV_\alpha,\notag \\
        (A^{[p]})^IV_\alpha       & =V_\alpha B_\alpha^I.\notag
    \end{align}
    Every $B_\alpha$ is periodic of period $m/r$, and for every $N\ge0$,
    \begin{align}
        G_N(A^{[p]}) & =G_N\left(\bigoplus_\alpha B_\alpha\right).\notag
    \end{align}
    The cyclic compression is that of
    \cite[Lemma~6]{DeLasCuevas2017Irreducible}; the local-space equality
    is its supported word contraction.
\end{lemma}
\begin{proof}\leanok
    \uses{thm:periodic_prescribed_blocking_compression,
        thm:isometric_sector_ground_space_resolution,
        thm:block_sum_local_ground_space}
    The step-orbit projections commute with every blocked letter and resolve
    the virtual identity. The displayed compression therefore intertwines
    letters. Apply the isometric sector resolution to obtain
    $G_N(A^{[p]})=\bigvee_\alpha G_N(B_\alpha)$, and identify the latter
    with the open MPS space of the direct-sum tensor.
\end{proof}

\begin{lemma}[Joint support of a normalized sector state]
    \label{thm:sector_state_joint_mps_support}
    \lean{MPSTensor.quasiLocalExpectation_toTensorFromBlocks_groundSpaceProjection_eq_one}
    \leanok
    Let $T=\bigoplus_j\mu_j A_j$, where every $\mu_j\ne0$.
    If $A_j$ is trace preserving and has a nonzero positive invariant matrix,
    its constructed state $\omega_j$ satisfies
    \begin{align}
        \omega_j\bigl(j_{a,N}(P_T(N))\bigr)=1
    \end{align}
    for every $a\in\mathbb Z$ and $N\geq0$. Primitivity is unnecessary.
\end{lemma}
\begin{proof}\leanok
    \uses{thm:block_sum_local_ground_space,
        thm:mps_state_local_support}
    The sector MPS space is contained in the joint space. Thus its vectors
    are annihilated by $1-P_T(N)$, which has zero expectation in the sector
    state. Normalization gives the assertion.
\end{proof}

\begin{lemma}[All states supported in primitive joint MPS spaces]
    \label{thm:primitive_joint_mps_supported_state_classification}
    \lean{MPSTensor.groundSpace_supported_iff_convex_combination_of_isPrimitiveMPS,
        MPSTensor.isPure_groundSpace_supported_iff_sector_of_isPrimitiveMPS}
    \leanok
    Suppose finitely many pairwise inequivalent primitive tensors $A_j$
    have faithful invariant matrices, and let $T=\bigoplus_j\mu_j A_j$
    with every $\mu_j\ne0$. Let $\omega_j$ be their constructed states.
    A state $\varphi$ satisfies
    \begin{align}
        \forall\,a\in\mathbb Z,\ N>0:\ \varphi\bigl(j_{a,N}(P_T(N))\bigr)=1
    \end{align}
    if and only if
    \begin{align}
        \varphi=\sum_jw_j\omega_j,
        \qquad w_j\geq0,\qquad\sum_jw_j=1.
    \end{align}
    Among these supported states, the pure states are precisely the
    $\omega_j$. No translation invariance of $\varphi$ is assumed.
    This is the primitive-family support classification used in
    \cite[Theorem~1.1 and Section~3]{Nachtergaele1996Spectral}; identification
    of a general periodic GVBS presentation remains separate.
\end{lemma}
\begin{proof}\leanok
    \uses{thm:multiblock_supported_quasi_local_state_convex_combination,
        thm:sector_state_joint_mps_support,
        thm:finite_pure_state_decomposition, thm:primitive_quasi_local_purity}
    Full interval support implies support on every centered interval, so
    the centered-interval argument gives the convex coefficients.
    Conversely, every sector state has joint support, and linearity
    preserves this property under convex combinations. A pure state in a
    finite convex combination equals a constituent, and every constituent
    is pure among all states.
\end{proof}

\begin{lemma}[Primitive sectors after blocking by a multiple of the period]
    \label{thm:periodic_primitive_ground_space}
    \lean{MPSTensor.IsPeriodic.exists_isPrimitiveMPS_groundSpaceDecomposition_of_dvd}
    \leanok
    Let $A$ be periodic of period $m$, and let $p>0$ satisfy $m\mid p$.
    Put $r=\gcd(m,p)=m$. There are positive bond dimensions $D_j$ with
    $\sum_{j=0}^{r-1}D_j=D$, tensors $B_j$ on the blocked physical alphabet,
    and positive definite matrices $\rho_j$, such that each pair
    $(B_j,\rho_j)$ has primitive fixed-point data. In particular,
    \begin{align}
        \sum_I(B_j^I)^\dagger B_j^I & =\Id_{D_j},
                                    & \E_{B_j}(\rho_j) & =\rho_j,
    \end{align}
    and for every $N\geq0$,
    \begin{align}
        G_N(A^{[p]})=G_N\left(\bigoplus_{j=0}^{r-1}B_j\right).
    \end{align}
    No inequivalence assumption is imposed on the sector family.
    This combines the periodic blocking of
    \cite[Lemma~6]{DeLasCuevas2017Irreducible} with the primitive normalization
    used in \cite[Section~3, equations (3.1)--(3.2b)]{Nachtergaele1996Spectral}.
    The peripheral period is the corrected $m/\gcd(m,p)$ of
    \cite{gap:dccsp17_blocking_peripheral_roots}.
\end{lemma}
\begin{proof}\leanok
    \uses{thm:periodic_blocked_ground_space,
        lem:periodic_scale_of_irreducible, def:primitive_mps}
    The compressed sectors have period $m/\gcd(m,p)=1$, hence are
    irreducible, trace preserving, and peripheral primitive.
    Finite-dimensional Perron--Frobenius theory gives their primitive
    fixed-point data, and irreducibility makes each invariant matrix positive
    definite. Retain the compressed tensors themselves, so the exact
    direct-sum support identity is unchanged at every length.
\end{proof}

\begin{lemma}[Primitive sectors after blocking by the period]
    \label{thm:periodic_primitive_ground_space_at_period}
    \lean{MPSTensor.IsPeriodic.exists_isPrimitiveMPS_groundSpaceDecomposition}
    \leanok
    A periodic tensor of period $m$ admits $m$ positive-dimensional
    primitive sectors after blocking by $m$. Their bond dimensions sum to
    $D$, they have faithful invariant matrices, and
    \begin{align}
        \forall\,N\geq0:\ G_N(A^{[m]})=G_N\left(\bigoplus_{j=0}^{m-1}B_j\right).
    \end{align}
\end{lemma}
\begin{proof}\leanok
    \uses{thm:periodic_primitive_ground_space}
    Take $p=m$ and use $m>0$ and $\gcd(m,m)=m$.
\end{proof}

\begin{lemma}[Exact open-boundary support under gauge and phase equivalence]
    \label{thm:primitive_representative_support_gauge}
    \lean{MPSTensor.GaugePhaseEquiv.groundSpaceES_eq,
        MPSTensor.MPVBlockPhaseEquiv.groundSpaceES_eq_of_isNormalTensor}
    \leanok
    Gauge-equivalent MPS tensors, including a common nonzero scalar factor
    on every physical matrix, have the same open-boundary MPS space at every
    length. More generally, normalized normal tensors whose periodic vectors
    differ by powers of one nonzero scalar have identical open-boundary
    spaces at every length, even when their bond dimensions are initially
    allowed to differ. The latter assertion uses the dimension equality and
    invertible gauge given by the fundamental theorem
    \cite[Proposition~2.7, Theorem~2.10]{Cirac2016MPDO_arXiv}.
\end{lemma}
\begin{proof}\leanok
    An invertible similarity is absorbed into the boundary matrix, while
    a nonzero length-dependent scalar is absorbed into its normalization.
    For normalized normal tensors, the periodic phase relation first gives
    equal bond dimensions and a gauge relation. Apply the first assertion.
\end{proof}

\begin{lemma}[Primitive representatives preserve the joint MPS spaces]
    \label{thm:primitive_sector_representatives_exact_support}
    \lean{MPSTensor.MPVPhaseClassData.groundSpaceES_toTensorFromBlocks_eq_representatives,
        MPSTensor.exists_primitive_sector_representatives}
    \leanok
    Let $(A_j)_{j=0}^{r-1}$ be a finite family of normalized primitive tensors
    of positive bond dimensions, with faithful fixed-point matrices
    $\rho_j$. Let $\mu_j\neq0$. There is an injective selection of original
    indices $s:\{0,\ldots,g-1\}\to\{0,\ldots,r-1\}$ such that the tensors
    $B_k=A_{s(k)}$ retain the given primitive and faithful fixed points,
    and distinct $B_k$ are not gauge-phase equivalent after any equality
    of their bond dimensions. Every $A_j$ is gauge-phase equivalent to
    one $B_k$, with the required bond-dimension equality given.
    For every $N\geq0$,
    \begin{align}
        \mathcal G_{\bigoplus_j\mu_jA_j}(N)
        =\mathcal G_{\bigoplus_k B_k}(N).
        \notag
    \end{align}
    The finite family may be empty. This is the exact local-support
    consequence of selecting phase-class representatives in
    \cite[Proposition~2.7]{Cirac2016MPDO_arXiv}, used for the
    inequivalent primitive sectors of \cite[Section~3]{Nachtergaele1996Spectral}.
\end{lemma}
\begin{proof}\leanok
    \uses{thm:primitive_representative_support_gauge}
    Select one original index from each periodic phase class. Distinct
    selections are gauge-phase inequivalent. Primitivity and faithfulness
    imply normalized normality, so each original tensor has exactly the
    same local MPS spaces as its selected representative. The local space
    of a weighted direct sum with nonzero weights is the sum of the block
    spaces. Replacing repeated summands by their representatives therefore
    preserves that sum at every length. These fixed points are inherited because
    the representatives are selected original tensors.
\end{proof}

\begin{lemma}[Inequivalent primitive presentation after periodic blocking]
    \label{thm:periodic_ground_state_presentation}
    \lean{MPSTensor.IsPeriodic.exists_inequivalent_primitive_groundStatePresentation_of_dvd}
    \leanok
    Let $A$ be periodic of period $m$, and let $p>0$ satisfy $m\mid p$.
    There is a nonempty finite family $(B_j)_{j=0}^{g-1}$, with positive
    bond dimensions and faithful invariant matrices $\rho_j$, such that
    each $(B_j,\rho_j)$ has normalized primitive fixed-point data, distinct
    sectors are gauge-phase inequivalent after any equality of their bond
    dimensions, and
    \begin{align}
        \forall\,N\geq0:\ G_N(A^{[p]})=G_N\left(\bigoplus_{j=0}^{g-1}B_j\right).
    \end{align}
    The tensors and matrices are selected from the compressed primitive
    family. Inequivalence is derived, rather than assumed.
    This is a finite local-support consequence of the periodic blocking in
    \cite[Lemma~6]{DeLasCuevas2017Irreducible} and the representative selection
    in \cite[Proposition~2.7, Theorem~2.10]{Cirac2016MPDO_arXiv}.
    The compressed peripheral period uses the correction in
    \cite{gap:dccsp17_blocking_peripheral_roots}.
\end{lemma}
\begin{proof}\leanok
    \uses{thm:periodic_primitive_ground_space,
        thm:primitive_sector_representatives_exact_support}
    Obtain primitive compressed sectors with their exact joint local space.
    Select one representative of each gauge-phase class. The given
    primitive and faithful fixed points are inherited, and the joint support
    is preserved at every length. Since the original period is positive,
    the compressed family is nonempty; the gauge cover therefore contains
    at least one selected representative.
\end{proof}

\begin{lemma}[Inequivalent primitive presentation at the period]
    \label{thm:periodic_ground_state_presentation_at_period}
    \lean{MPSTensor.IsPeriodic.exists_inequivalent_primitive_groundStatePresentation}
    \leanok
    Blocking a periodic tensor by its period gives a nonempty finite
    pairwise gauge-phase inequivalent primitive family with faithful
    invariant matrices and
    \begin{align}
        \forall\,N\geq0:\ G_N(A^{[m]})=G_N\left(\bigoplus_jB_j\right).
    \end{align}
\end{lemma}
\begin{proof}\leanok
    \uses{thm:periodic_ground_state_presentation}
    Apply the theorem with $p=m$, using positivity of the period.
\end{proof}

\begin{lemma}[Expanding symmetric support determines every interval]
    \label{thm:eventual_symmetric_mps_support_all_intervals}
    \lean{MPSTensor.%
        quasiLocalState_groundSpaceProjection_eq_one_of_eventual_symmetric_support}
    \leanok
    Let $A$ be any MPS tensor of positive physical dimension, and let
    $\varphi$ be a normalized positive norm-one quasi-local state.
    If $\varphi(j_{-n,2n+1}(p_A(2n+1)))=1$ for every sufficiently large
    $n$, then $\varphi(j_{a,N}(p_A(N)))=1$ for every $a\in\mathbb Z$
    and every $N>0$. This is the finite-support restriction used in
    \cite[Theorem~1.1 and Section~3]{Nachtergaele1996Spectral}.
\end{lemma}
\begin{proof}\leanok
    \uses{thm:mps_support_contained_interval,def:mps_ground_space_projection_matrix}
    Choose $n$ sufficiently large that the symmetric interval
    $[-n,n]$ contains $[a,a+N-1]$ and has full ground-projection
    expectation. Support descends to the contained interval.
\end{proof}

\begin{lemma}[Zero local energy is equivalent to all interval support]
    \label{thm:eventual_kernel_ground_state_support_iff}
    \lean{MPSTensor.mem_parentGroundStateFace_iff_groundSpace_support_of_eventual_kernel}
    \leanok
    Let $A$ be any MPS tensor of positive physical dimension, and let
    $h\geq0$ be a local interaction of range $R>0$. Suppose that
    $\ker H_N(h)=G_N(A)$ for every sufficiently large $N$.
    For every continuous complex-linear functional $\varphi$,
    \begin{align}
        \varphi\in\mathcal F_h
        \quad\Longleftrightarrow\quad
         & \varphi\text{ is a normalized positive norm-one state and}\notag \\
         & \forall a\in\mathbb Z,\ \forall N>0,\quad
        \varphi(j_{a,N}(p_A(N)))=1.\notag
    \end{align}
    These are the eventual finite-volume kernel and local support
    conditions of \cite[Theorems~1.1--1.2 and Section~3]{Nachtergaele1996Spectral}.
    No translation invariance of $\varphi$ or injectivity bound on $R$
    is assumed.
\end{lemma}
\begin{proof}\leanok
    \uses{def:quasi_local_parent_ground_state_face,
        thm:quasi_local_zero_energy_support_eventual_kernel,
        thm:eventual_symmetric_mps_support_all_intervals,
        thm:positive_functional_support_compression,
        thm:quasi_local_open_interaction_zero_expectation}
    Zero local energy first gives full support on sufficiently large
    symmetric intervals, and hence on every nonempty interval.
    Conversely, fix $a$ and choose $N\geq R$ with the exact kernel.
    Set $p=p_A(N)$ and restrict $\varphi$ to this interval.
    Since $H_N(h)p=0$ and $\varphi(j_{a,N}(p))=1$, support compression
    gives $\varphi(j_{a,N}(H_N(h)))=0$. The finite open sum is the sum of
    its translated interaction observables. Their expectations are all
    nonnegative, so each is zero; its first summand gives
    $\varphi(j_{a,R}(h))=0$.
\end{proof}

\begin{lemma}[All interval support is preserved by blocking]
    \label{thm:mps_support_blocking_iff}
    \lean{MPSTensor.groundSpace_support_blockTensor_iff}
    \leanok
    Let $d,L>0$, let $A$ be a tensor, and let $\varphi$ be a normalized
    positive state on the original chain. Write $T_L$ for the pullback of
    states from the blocked chain. Then
    \begin{align}
        \bigl(\forall\,a\in\mathbb Z,\ N>0:\ \varphi(P_{G_A([a,a+N-1])})=1\bigr)
        \quad\Longleftrightarrow\quad
        \bigl(\forall\,a\in\mathbb Z,\ N>0:\
        (T_L^{-1}\varphi)(P_{G_{A^{[L]}}([a,a+N-1])})=1\bigr).
        \notag
    \end{align}
    This statement requires neither primitivity nor translation invariance.
    It makes the regrouping in
    \cite[lines~825--836 and Section~3]{Nachtergaele1996Spectral}
    explicit at the level of supported states.
\end{lemma}
\begin{proof}\leanok
    \uses{thm:blocked_aligned_mps_support_iff,
        thm:blocked_ground_projection_support, def:qca_blocking_functional}
    Support on all intervals is equivalent to support on the aligned
    intervals. The blocked projection identity carries the latter
    conditions to all intervals of the blocked chain.
\end{proof}

\begin{lemma}[An exact blocked representation preserves support]
    \label{thm:blocked_support_representation_iff}
    \lean{MPSTensor.groundSpace_support_blocked_representation_iff}
    \leanok
    Suppose $G_{A^{[L]}}(N)=G_C(N)$ for every $N\geq0$.
    Then a state $\varphi$ is supported on all nonempty original
    $A$-intervals if and only if $T_L^{-1}\varphi$ is supported on all
    nonempty $C$-intervals. No primitivity assumption is made.
\end{lemma}
\begin{proof}\leanok
    \uses{thm:mps_support_blocking_iff}
    Equal subspaces have equal orthogonal projections. Substitute these
    projections in the support equivalence.
\end{proof}

\begin{lemma}[Supported states from an exact blocked primitive representation]
    \label{thm:blocked_primitive_supported_state_classification}
    \lean{MPSTensor.groundSpace_supported_iff_convex_combination_of_blocked_primitive_family,
        MPSTensor.isPure_groundSpace_supported_iff_sector_of_blocked_primitive_family}
    \leanok
    Let $d,L>0$. Suppose $B_1,\ldots,B_b$ are pairwise inequivalent
    primitive tensors of positive bond dimensions on the blocked alphabet,
    with faithful invariant matrices, and $\mu_j\ne0$. Assume the exact identities
    \begin{align}
        \forall\,N\geq0:\ G_{A^{[L]}}(N)=G_{\bigoplus_j\mu_j B_j}(N).
    \end{align}
    Let $\omega_j$ denote the normalized sector states on the blocked
    chain. For any normalized positive state $\varphi$ on the original
    chain, support on all nonempty $A$-intervals is equivalent to
    \begin{align}
        \varphi=\sum_{j=0}^{b-1} w_j T_L\omega_j,
        \qquad w_j\geq0,\qquad \sum_{j=0}^{b-1} w_j=1.
    \end{align}
    The pure supported states are exactly the states $T_L\omega_j$.
    Purity and the support condition are taken among all states.
    This is the conclusion of
    \cite[Theorem~1.1, lines~854--926]{Nachtergaele1996Spectral}
    for the specified exact representation; its existence and identification
    with a general GVBS construction are separate assertions.
\end{lemma}
\begin{proof}\leanok
    \uses{thm:blocked_support_representation_iff,
        thm:primitive_joint_mps_supported_state_classification,
        thm:qca_blocking_state_purity}
    Apply the primitive-family classification to $T_L^{-1}\varphi$.
    State pullback preserves finite real convex combinations and purity,
    so carry the resulting identities back to the original chain.
\end{proof}

\begin{lemma}[Nonwrapping windows commute with physical blocking]
    \label{thm:blocked_interaction_window}
    \lean{MPSTensor.chainWindowOperator_blocking_reindex}
    \leanok
    Let $d,L,K>0$, let $b+K\leq N$, and let $h$ be any operator on
    $K$ blocked sites. Write $W_{N,b}(h)$ for its nonwrapping placement
    starting at $b$, and $\mathcal R_{N,L}$ for decoding both matrix indices.
    Then
    \begin{align}
        \mathcal R_{N,L}\bigl(W_{N,b}(h)\bigr)
        =W_{NL,bL}\bigl(\mathcal R_{K,L}(h)\bigr).
    \end{align}
    This is an operator identity under the site regrouping in
    \cite[lines~825--836 and Section~3]{Nachtergaele1996Spectral} and
    \cite[Appendix, lines~2308 and 2313--2320]{Cirac2017MPU}.
\end{lemma}
\begin{proof}\leanok
    \uses{thm:qca_blocking_quasi_local_interval,
        thm:quasi_local_bulk_inclusion, thm:quasi_local_interval_observable}
    Apply the injective interval inclusion to both matrices. Their
    quasi-local images agree by the window inclusion and interval blocking
    identities: the blocked start $b$ becomes $bL$ and the length becomes $KL$.
\end{proof}

\begin{lemma}[Blocked open sums have aligned original starts]
    \label{thm:blocked_open_interaction_sum}
    \lean{MPSTensor.openInteractionMatrix_blocking_reindex}
    \leanok
    For $d,L,K>0$ and every $N\geq0$, write $H_h(N)$ for the open sum
    of the range-$K$ interaction $h$ on the blocked chain. Then
    \begin{align}
        \mathcal R_{N,L}\bigl(H_h(N)\bigr)
        =\sum_{\substack{b\geq0 \\b+K\leq N}}
        W_{NL,bL}\bigl(\mathcal R_{K,L}(h)\bigr).
    \end{align}
    The sum is zero when $N<K$. This is the open interaction convention
    of \cite[equation~(3.12)]{Nachtergaele1996Spectral} pulled back under blocking.
\end{lemma}
\begin{proof}\leanok
    \uses{thm:blocked_interaction_window, def:finite_open_interaction_matrix}
    Matrix reindexing is linear. Apply the window identity to each term of
    the finite open sum.
\end{proof}

\begin{lemma}[Aligned positive windows are bounded by the full open sum]
    \label{thm:aligned_positive_window_sum}
    \lean{MPSTensor.sum_aligned_chainWindowOperator_le_openInteractionMatrix}
    \leanok
    Let $d\geq0$, $K,L>0$, and let $Y\geq0$ be an operator on $KL$
    original sites. For every $N\geq0$,
    \begin{align}
        \sum_{\substack{b\geq0 \\b+K\leq N}}W_{NL,bL}(Y)
        \leq H_Y(NL).
    \end{align}
    The coefficient is one. This is a derived positive-term counting
    consequence of \cite[equation~(3.12)]{Nachtergaele1996Spectral}.
\end{lemma}
\begin{proof}\leanok
    \uses{def:finite_open_interaction_matrix}
    The map $b\mapsto bL$ is injective because $L>0$, and every aligned
    start satisfies $bL+KL\leq NL$. Thus the left sum is a subsum of the
    full open sum. Each omitted term is positive, since window placement is
    a star-algebra homomorphism.
\end{proof}

\begin{lemma}[Ground states of a periodic MPS tensor]
    \label{thm:periodic_quasi_local_ground_state_classification}
    \lean{MPSTensor.IsPeriodic.exists_parentGroundStateFace_classification_of_eventually_open_kernel}
    \leanok
    Let $A$ be a periodic MPS tensor of period $m$ and positive physical
    dimension. Let $h\geq0$ have interaction range $R>0$, and suppose
    $\ker H_N(h)=G_N(A)$ for every sufficiently large $N$.
    There is a nonempty finite family of pairwise gauge-phase inequivalent
    primitive tensors $B_j$, with positive bond dimensions and faithful
    invariant matrices $\rho_j$, such that $G_N(A^{[m]})=
        G_N(\bigoplus_j B_j)$ at every length $N\geq0$.
    Let $\omega_j$ denote the primitive state of $(B_j,\rho_j)$ pulled back
    to the original lattice by grouping $m$ consecutive sites. For every
    continuous complex-linear functional $\varphi$,
    \begin{align}
        \varphi\in\mathcal F_h
         & \quad\Longleftrightarrow\quad
        \exists (w_j)_j,\quad w_j\geq0,\quad
        \sum_j w_j=1,\quad \varphi=\sum_j w_j\omega_j,
        \label{eq:ph_periodic_ground_face}                             \\
        \varphi\in\mathcal F_h\ \text{and }\varphi\text{ is pure}
         & \quad\Longleftrightarrow\quad \exists j,\ \varphi=\omega_j.
        \label{eq:ph_periodic_pure_face}
    \end{align}
    The sector data and their inequivalence are derived from periodicity.
    The state is not assumed translation invariant, and $R$ is not subject
    to an injectivity bound. This is the periodic-tensor case of the
    ground-state face classification in
    \cite[Theorems~1.1--1.2 and Section~3]{Nachtergaele1996Spectral}; the
    identification of a general periodic GVBS presentation remains separate,
    as described in \cite{gap:nachtergaele96_infinite_volume_ground_projection}.
\end{lemma}
\begin{proof}\leanok
    \uses{thm:periodic_ground_state_presentation_at_period,thm:eventual_kernel_ground_state_support_iff,thm:blocked_primitive_supported_state_classification,thm:mps_quasi_local_expectation_state,thm:qca_blocking_state_purity,thm:sector_separation_bond_transport}
    Blocking by the positive period gives inequivalent primitive sectors
    with faithful invariant matrices and exact joint support at every
    length. The eventual kernel identity makes membership in
    $\mathcal F_h$ equivalent to full expectation on every nonempty
    interval support projection. Carrying this condition to the blocked
    lattice identifies the supported states with the convex hull of the
    primitive sector states. The blocking equivalence preserves purity
    among all states, giving (\ref{eq:ph_periodic_pure_face}).
\end{proof}

\begin{lemma}[Blocking distributes over the joint boundary space]
    \label{thm:blocking_joint_boundary_space}
    \lean{MPSTensor.groundSpaceES_blockTensor_toTensorFromBlocks_eq_iSup}
    \leanok
    Let $A_j$ be a finite family of tensors and let $\mu_j\neq0$.
    For all integers $L,N\geq0$, the open-boundary space of the blocked
    weighted sum satisfies
    \begin{align}
        \mathcal G_{(\bigoplus_j\mu_jA_j)^{(L)}}(N)
        =\sum_j\mathcal G_{A_j^{(L)}}(N).
    \end{align}
    No periodicity, positivity, or nonempty-family assumption is required.
    These are the boundary spaces of
    \cite[Section~3]{Nachtergaele1996Spectral} under the blocking coordinates
    of \cite[Lemma~6]{DeLasCuevas2017Irreducible}.
\end{lemma}
\begin{proof}\leanok
    The configuration isometry identifies each blocked boundary space
    with the original boundary space at length $NL$. Under this isometry,
    the space of the weighted sum is the sum of the original block spaces.
    A linear isomorphism preserves sums of subspaces, giving the asserted
    identity in the blocked coordinates.
\end{proof}

\begin{lemma}[Primitive presentation of a finite collection of families]
    \label{thm:primitive_nested_family_presentation}
    \lean{MPSTensor.exists_primitive_nestedFamilyPresentation}
    \leanok
    Let $(B_{j,a})_a$ be finitely many finite families of primitive tensors
    of positive bond dimensions, with faithful invariant matrices.
    There is a finite pairwise gauge-phase inequivalent primitive faithful
    family $(C_b)_b$ such that, for every $N\geq0$,
    \begin{align}
        \sum_j\sum_a\mathcal G_{B_{j,a}}(N)
        =\mathcal G_{\bigoplus_b C_b}(N).
    \end{align}
    Any of the families may be empty. This is representative selection
    in \cite[Proposition~2.7]{Cirac2016MPDO_arXiv}
    applied to a finite union of primitive families.
\end{lemma}
\begin{proof}\leanok
    \uses{thm:primitive_sector_representatives_exact_support}
    Index all pairs $(j,a)$ by one finite set and apply primitive
    representative selection. A sum over a dependent finite index set
    is the iterated sum over its two indices, so the selected family
    retains the asserted boundary space at every length.
\end{proof}

\begin{lemma}[Primitive presentation at a common period multiple]
    \label{thm:periodic_block_family_primitive_presentation}
    \lean{MPSTensor.exists_primitive_blockFamilyPresentation_of_period_dvd}
    \leanok
    Let $A_j$ be a finite family of normalized periodic tensors with periods
    $m_j$, and let $\mu_j\neq0$. Suppose $L>0$ and $m_j\mid L$ for every $j$.
    There are finitely many primitive tensors $B_k$ of positive bond
    dimensions, with faithful invariant matrices $\rho_k$, such that
    distinct $B_k$ are not gauge-phase equivalent after any equality of
    their bond dimensions. For every $N\geq0$,
    \begin{align}
        \mathcal G_{(\bigoplus_j\mu_jA_j)^{(L)}}(N)
        =\mathcal G_{\bigoplus_k B_k}(N).
    \end{align}
    The original family may be empty, and the representative family may
    consequently be empty. This is the finite-family support consequence
    of \cite[Lemma~6]{DeLasCuevas2017Irreducible}
    and representative selection in
    \cite[Proposition~2.7]{Cirac2016MPDO_arXiv}.
\end{lemma}
\begin{proof}\leanok
    \uses{thm:blocking_joint_boundary_space,
        thm:periodic_primitive_ground_space,
        thm:primitive_nested_family_presentation}
    The cyclic-sector decomposition of $A_j^{(L)}$ has compressed period
    $m_j/\gcd(m_j,L)=1$, so each compressed tensor is primitive and has a
    faithful invariant matrix. Flatten these finitely many sector families.
    The joint boundary-space identity identifies their sum with the
    blocked original space. Selecting one representative from each
    gauge-phase class preserves this sum at every length and inherits
    the primitive and faithful fixed points.
\end{proof}

\begin{lemma}[Existence of a primitive presentation after blocking]
    \label{thm:periodic_block_family_primitive_presentation_exists}
    \lean{MPSTensor.exists_primitive_blockFamilyPresentation_of_isPeriodic}
    \leanok
    Every finite weighted family of normalized periodic tensors with
    nonzero weights admits a positive blocking length and an inequivalent
    primitive faithful family with identical open-boundary spaces at
    every blocked length. No original inequivalence or nonempty-family
    assumption is required. This is the common-period consequence of
    \cite[Lemma~6]{DeLasCuevas2017Irreducible}.
\end{lemma}
\begin{proof}\leanok
    \uses{thm:periodic_block_family_primitive_presentation}
    Take $L=\prod_jm_j$. Every period is positive, and every $m_j$ divides
    this product. For an empty family, the product equals one. Apply the
    common-multiple presentation theorem.
\end{proof}

\begin{definition}[Interaction obtained by grouping consecutive sites]
    \label{def:blocked_interaction_matrix}
    \lean{MPSTensor.blockedInteractionMatrix,
        MPSTensor.reindex_blockedInteractionMatrix}
    \leanok
    For a range-$R$ matrix $h$, let $H_h(n)$ be its finite open interaction
    sum. Group $L$ sites into one site and let $\mathcal R_{K,L}$ decode
    both indices of a matrix on $K$ grouped sites. Define
    \begin{align}
        h_B=\mathcal R_{K,L}^{-1}\bigl(H_h(KL)\bigr).
    \end{align}
    Decoding $h_B$ recovers $H_h(KL)$. The definition allows $d,L,K\geq0$. This is a regrouping construction for
    \cite[Section~3, equations~(3.12)--(3.14)]{Nachtergaele1996Spectral}.
\end{definition}

\begin{lemma}[Positivity of the grouped interaction]
    \label{thm:blocked_interaction_matrix_posSemidef}
    \lean{MPSTensor.blockedInteractionMatrix_posSemidef}
    \leanok
    \uses{def:blocked_interaction_matrix}
    If $R>0$ and $h\geq0$, then $h_B\geq0$, including when $KL<R$.
\end{lemma}
\begin{proof}\leanok
    \uses{def:finite_open_interaction_matrix,thm:open_interaction_matrix_posSemidef}
    Reindexing is an equivalence. Each nonwrapping local term is positive
    because its inclusion is a star-algebra homomorphism; finite sums and
    matrix reindexing preserve positivity.
\end{proof}

\begin{lemma}[Expansion of the grouped open interaction sum]
    \label{thm:blocked_interaction_expansion}
    \lean{MPSTensor.openInteractionMatrix_blockedInteractionMatrix_reindex}
    \leanok
    Let $d,L,R,K>0$, $R\leq KL$, and $N\geq K$. If $h_B$ is the grouped
    interaction, then
    \begin{align}
        \mathcal R_{N,L}\bigl(H_{h_B}(N)\bigr)
        =\sum_{i=0}^{N-K}\sum_{j=0}^{KL-R}W_{NL,iL+j}(h),
    \end{align}
    where $W_{n,b}$ places an operator at the nonwrapping start $b$.
    This is the aligned interaction sum of
    \cite[equation~(3.12)]{Nachtergaele1996Spectral} with each grouped
    interaction expanded into its original terms.
\end{lemma}
\begin{proof}\leanok
    \uses{def:blocked_interaction_matrix, thm:blocked_open_interaction_sum}
    Decoding sends each grouped window to the corresponding original
    window. Inclusion is linear, and nested nonwrapping inclusions add
    their starts, so the internal start $j$ becomes $iL+j$.
\end{proof}

\begin{lemma}[Positive comparison and kernel equality under grouping]
    \label{thm:blocked_interaction_comparison}
    \lean{MPSTensor.openInteractionMatrix_blockedInteractionMatrix_comparison,
        MPSTensor.ker_toEuclideanLin_openInteractionMatrix_blockedInteractionMatrix_reindex,
        MPSTensor.ker_openInteractionHamiltonianES_blockedInteractionMatrix_map}
    \leanok
    Let $d,L,R,K>0$, $h\geq0$, $N\geq K$, and $R+L\leq KL+1$.
    Then
    \begin{align}
        H_h(NL)\leq\mathcal R_{N,L}\bigl(H_{h_B}(N)\bigr)
        \leq(KL-R+1)H_h(NL).
    \end{align}
    The decoded operators have equal kernels. Equivalently, the
    configuration isometry maps $\ker H_{h_B}(N)$ onto $\ker H_h(NL)$.
    These are derived positive counting consequences of
    \cite[Section~3, equations~(3.12)--(3.14)]{Nachtergaele1996Spectral};
    no local MPS-kernel condition is required.
\end{lemma}
\begin{proof}\leanok
    \uses{thm:blocked_interaction_expansion, lem:positive-comparison-kernel-gap}
    For an original start $s$, choose
    $i=\min(\lfloor s/L\rfloor,N-K)$ and $j=s-iL$.
    The covering condition gives $0\leq j\leq KL-R$, and the pair
    $(i,j)$ determines $s$, giving the lower bound by a positive subsum.
    For each fixed $j$, the starts $iL+j$ are distinct and nonwrapping,
    so their sum is bounded by $H_h(NL)$. Sum this bound over the
    $KL-R+1$ internal offsets. Two-sided positive comparison gives kernel
    equality, and the configuration isometry carries that kernel.
\end{proof}

\begin{lemma}[Eventual MPS kernels persist under grouping]
    \label{thm:blocked_interaction_eventual_kernel}
    \lean{MPSTensor.eventually_ker_blockedInteractionMatrix_reindex_eq_groundSpace,
        MPSTensor.eventually_ker_openInteractionHamiltonianES_blockedInteractionMatrix_eq_groundSpaceES}
    \leanok
    Let $d,L,R,K>0$, $R+L\leq KL+1$, $h\geq0$, and let $A$ be any
    MPS tensor. Suppose that $\ker H_h(n)=\mathcal G_n(A)$ for all
    sufficiently large $n$. Then, for all sufficiently large $N$,
    \begin{align}
        \ker H_{h_B}(N)=\mathcal G_N(A^{[L]}).
    \end{align}
    After decoding, this is the equality with the isometric image of
    $\mathcal G_N(A^{[L]})$. No normality or primitivity hypothesis is needed.
    This carries the eventual kernel hypothesis of
    \cite[Theorem~1.2]{Nachtergaele1996Spectral} through its regrouping convention.
\end{lemma}
\begin{proof}\leanok
    \uses{thm:blocked_interaction_comparison, thm:blocked_local_ground_space_transport}
    Since $L>0$, the original volumes $NL$ tend to infinity. Apply the
    original eventual kernel equality and the exact isometric image
    $\mathcal G_N(A^{[L]})\mapsto\mathcal G_{NL}(A)$.
    Injectivity of the configuration isometry gives the equality in
    grouped coordinates.
\end{proof}

\begin{lemma}[Positivity of open interaction sums]
    \label{thm:open_interaction_matrix_posSemidef}
    \lean{MPSTensor.openInteractionMatrix_posSemidef}
    \leanok
    Let $R>0$ and let $h$ be a positive semidefinite interaction on $R$
    sites. For every $n\geq0$, the open sum
    $\sum_{a=0}^{n-R}\iota_{a}(h)$ of its nonwrapping translates on $n$
    sites is positive semidefinite; for $n<R$ the sum is empty and equal
    to zero. This is the positivity of the finite-volume Hamiltonians in
    \cite[equation~(3.12)]{Nachtergaele1996Spectral}.
\end{lemma}
\begin{proof}\leanok
    Each translate is the tensor product of $h$ with an identity, hence
    positive semidefinite, and a finite sum of positive semidefinite
    matrices is positive semidefinite.
\end{proof}

\begin{lemma}[A vanishing square annihilates left multiples]
    \label{thm:positive_functional_vanishing_square}
    \lean{PositiveLinearMap.apply_mul_star_mul_self_eq_zero}
    \leanok
    Let $\mathcal A$ be a star-ordered complex algebra admitting
    self-adjoint decomposition, and let $f:\mathcal A\to\mathbb C$ be
    positive and complex linear. If $f(C^*C)=0$, then $f(YC^*C)=0$ for
    every $Y\in\mathcal A$. This is the Cauchy--Schwarz implication used
    in the zero-energy argument of
    \cite[Section~6]{Nachtergaele1996Spectral}.
\end{lemma}
\begin{proof}\leanok
    Apply Cauchy--Schwarz to $u=(YC^*)^*$ and $v=C$.
    Its right-hand side contains $\sqrt{|f(C^*C)|}=0$.
\end{proof}

\begin{lemma}[Positive comparison of commutator energies]
    \label{thm:positive_functional_commutator_comparison}
    \lean{PositiveLinearMap.apply_mul_eq_zero_of_posSemidef}
    \lean{PositiveLinearMap.apply_conjTranspose_commutator_eq}
    \lean{PositiveLinearMap.apply_conjTranspose_commutator_le}
    \leanok
    Let $f$ be a positive complex-linear functional on a finite matrix
    algebra, and let $H\geq0$ satisfy $f(H)=0$. Every left multiple of
    $H$ has zero expectation, and
    \begin{align}
        f(X^*[H,X])=f(X^*HX).
        \label{eq:ph_zero_energy_commutator}
    \end{align}
    If $C\geq0$, $c\in\mathbb R$, and $C\leq cH$, then
    \begin{align}
        f(X^*[C,X])\leq c f(X^*[H,X]).
        \label{eq:ph_positive_commutator_comparison}
    \end{align}
    These are finite-dimensional positive-comparison consequences of the
    zero-energy argument in \cite[Section~6]{Nachtergaele1996Spectral}.
    Normalization of $f$ is unnecessary.
\end{lemma}
\begin{proof}\leanok
    \uses{thm:positive_functional_vanishing_square}
    Factor $H=S^*S$ and apply the vanishing-square implication to obtain
    $f(YH)=0$. Expanding the commutator proves
    (\ref{eq:ph_zero_energy_commutator}). Positivity gives
    $0\leq f(C)\leq c f(H)=0$, hence the same identity for $C$.
    Conjugating $C\leq cH$ by $X$ gives $X^*CX\leq cX^*HX$;
    monotonicity of $f$ proves (\ref{eq:ph_positive_commutator_comparison}).
\end{proof}

\begin{lemma}[Configuration reindexing preserves kernels and norm gaps]
    \label{thm:configuration_reindex_kernel_gap}
    \lean{Matrix.toEuclideanLin_reindex_piLpCongrLeft,
        Matrix.ker_toEuclideanLin_reindex_map,
        Matrix.norm_gap_toEuclideanLin_reindex}
    \leanok
    Let $e:I\to J$ be a bijection of finite configuration sets, and let $H$
    be any square complex matrix on $I$. The associated Euclidean isometry
    $U_e$ satisfies
    \begin{align}
        H^e U_e=U_e H,\qquad \ker H^e=U_e(\ker H).
    \end{align}
    Every estimate $\gamma\|v\|\leq\|Hv\|$ on $(\ker H)^\perp$
    holds with the same constant for $H^e$ on $(\ker H^e)^\perp$.
    Positivity is unnecessary.
\end{lemma}
\begin{proof}\leanok
    Change the summation index in the matrix action. The intertwining
    identity and invertibility give the kernel identity. Isometries
    preserve norms and orthogonal complements, yielding the gap estimate.
\end{proof}

\begin{lemma}[A blocked finite-chain gap gives an original gap]
    \label{thm:blocked_finite_gap_transfer}
    \lean{MPSTensor.openInteractionMatrix_gap_of_blocked_gap}
    \leanok
    Let $d,L,R,K>0$, $N\geq K$, $R+L\leq KL+1$, and $h\geq0$.
    Let $\widehat h$ be the interaction obtained by regrouping $H_h(KL)$.
    If $H_{\widehat h}(N)$ has norm gap $\gamma>0$ on its kernel
    complement, then $H_h(NL)$ has norm gap
    $\gamma/(KL+1-R)$ on its entire kernel complement.
    This is a finite-volume consequence of
    \cite[Section~3, equations~(3.12)--(3.14)]{Nachtergaele1996Spectral}.
\end{lemma}
\begin{proof}\leanok
    \uses{thm:blocked_interaction_comparison,thm:configuration_reindex_kernel_gap,lem:positive-comparison-kernel-gap}
    Carry the blocked norm estimate to original coordinates.
    The two positive order comparisons give coincident kernels and
    transfer the gap with the stated divisor.
\end{proof}

\begin{lemma}[Uniform gaps at all block-divisible volumes]
    \label{thm:blocked_primitive_aligned_uniform_gap}
    \lean{MPSTensor.exists_aligned_openInteractionMatrix_gap_of_blocked_primitive_family}
    \leanok
    Let $d,L>0$. Suppose $A^{[L]}$ has an exact joint-space presentation
    at every length by inequivalent primitive tensors of positive bond
    dimensions, with faithful invariant matrices and nonzero weights. Let $h\geq0$ have positive range $R$
    and eventual original kernel $G_A(n)$. If $K>0$ and
    $R+L\leq KL+1$, there is $\delta>0$ such that
    \begin{align}
        \delta\|v\|\leq\|H_h(NL)v\|
        \qquad\bigl(N\geq0,\ v\perp\ker H_h(NL)\bigr).
    \end{align}
    Only block-divisible original volumes are asserted. This is the
    corresponding restriction of
    \cite[Theorem~1.2]{Nachtergaele1996Spectral}.
\end{lemma}
\begin{proof}\leanok
    \uses{thm:blocked_interaction_eventual_kernel,thm:blocked_finite_gap_transfer,thm:eventual_kernel_open_interaction_gap,thm:sector_separation_bond_transport,lem:positive-kernel-projection-comparison}
    The coarse interaction is positive and has eventual joint primitive
    kernel. Apply the primitive-family uniform gap theorem and transfer
    its estimate for $N\geq K$. Take a positive minimum with the finitely
    many smaller volumes, using their kernels.
\end{proof}

\begin{lemma}[A periodic tensor has uniform gaps at period-divisible volumes]
    \label{thm:periodic_tensor_aligned_uniform_gap}
    \lean{MPSTensor.IsPeriodic.exists_aligned_openInteractionMatrix_gap_of_eventual_kernel}
    \leanok
    Let $A$ be a periodic tensor of period $m>0$, and let $h\geq0$
    have positive range and eventual kernel $G_A(n)$.
    There exists $\delta>0$ such that, for every $N\geq0$,
    \begin{align}
        \forall\,v\perp\ker H_h(Nm):\ \delta\|v\|\leq\|H_h(Nm)v\|.
    \end{align}
    The primitive family, its invariant matrices, and its inequivalence
    are derived from periodicity. Other residue lengths remain separate.
\end{lemma}
\begin{proof}\leanok
    \uses{thm:periodic_ground_state_presentation_at_period,thm:blocked_primitive_aligned_uniform_gap,thm:sector_separation_bond_transport}
    Obtain the exact primitive presentation after blocking by $m$ and
    choose $K=R+1$. Since $m\geq1$, the window coverage condition holds.
\end{proof}

\begin{lemma}[Ground states of a finite periodic tensor family]
    \label{thm:periodic_block_family_ground_states}
    \lean{MPSTensor.exists_parentGroundStateFace_classification_of_periodic_family}
    \leanok
    Let $d\geq1$, and let $(A_j)_{j=0}^{r-1}$ be a finite family of normalized periodic tensors
    with periods $m_j$, and let $T=\bigoplus_j\mu_jA_j$, where $\mu_j\neq0$.
    Let $h\geq0$ have interaction range $R>0$. Suppose that its open-chain
    Hamiltonians satisfy $\ker H_h(N)=\mathcal G_T(N)$ for all sufficiently
    large $N$. Then there are $L>0$ and a finite pairwise gauge-phase
    inequivalent primitive family $(B_k)_{k=0}^{g-1}$, of positive bond
    dimensions with faithful invariant matrices, such that
    $\mathcal G_{T^{[L]}}(N)=\mathcal G_{\bigoplus_kB_k}(N)$ for every
    $N\geq0$. Let $\widetilde\omega_k$ be the corresponding primitive
    sector states pulled back from the blocked lattice to the original
    lattice. For every continuous complex linear functional $\varphi$,
    \begin{align}
        \varphi\in\mathcal F_h
         & \ \Longleftrightarrow\
        \exists w\in[0,\infty)^g:\quad
        \sum_{k=0}^{g-1} w_k=1,\quad
        \varphi=\sum_{k=0}^{g-1} w_k\widetilde\omega_k, \\
        \varphi\in\mathcal F_h\text{ and }\varphi\text{ is pure}
         & \ \Longleftrightarrow\
        \exists k,\ \varphi=\widetilde\omega_k.
    \end{align}
    Here $\mathcal F_h$ is the zero-energy face of all normalized positive
    states, and purity is extremality among all states. The original
    family may be empty and no inequivalence is imposed on it.
    No translation invariance of $\varphi$, given invariant matrices,
    or injectivity bound on $R$ is assumed.

    This is the explicit periodic-tensor-family case of
    \cite[Theorems~1.1--1.2, Section~3]{Nachtergaele1996Spectral}, using
    the site grouping at lines 825--836 and
    \cite[Lemma~6]{DeLasCuevas2017Irreducible}.
    Identification of an arbitrary GVBS presentation with these tensor
    data and finite support spaces remains separate
    \cite{gap:nachtergaele96_infinite_volume_ground_projection}.
\end{lemma}
\begin{proof}\leanok
    \uses{thm:periodic_block_family_primitive_presentation_exists,thm:eventual_kernel_ground_state_support_iff,thm:blocked_primitive_supported_state_classification,thm:mps_quasi_local_expectation_state,thm:qca_blocking_state_purity,thm:sector_separation_bond_transport}
    Blocking by a common period multiple and selecting primitive
    representatives gives the exact boundary-space equality at every
    length. The eventual kernel identity identifies zero local energy
    with support on every original MPS interval. Carrying support
    through blocking identifies these states with convex combinations
    of the primitive sector states. Conversely, such a combination is
    a normalized positive state with all interval supports, hence has
    zero local energy. Blocking preserves the whole state space and
    purity, so the pure-support classification gives the second
    equivalence.
\end{proof}

\begin{lemma}[Blocked stationary expectations]
    \label{thm:blocked_stationary_expectation}
    \lean{MPSTensor.blockedQuasiLocalExpectation_quasiLocalTranslation,
        MPSTensor.blockedQuasiLocalExpectation_interval,
        MPSTensor.blockedQuasiLocalExpectation_window}
    \leanok
    Let $L>0$, and let $B$ be a tensor on the blocked physical alphabet of
    $L$ sites with
    $\sum_i(B^i)^*B^i=I$, and let $\rho\geq0$ be a fixed point of its
    transfer map with $\operatorname{tr}\rho\neq0$. Carry the
    quasi-local state of $B$ and $\rho$ to the original chain through the
    blocking isomorphism. The pulled-back state is invariant under
    translations by integer multiples of $L$. On an aligned interval of
    $NL$ sites it is the insertion expectation of $B$ after reindexing the
    configurations into blocks, and an observable on $k\geq1$ sites placed
    inside such an interval has the insertion expectation of its
    nonwrapping placement \cite[lines~825--836 and
        Section~3]{Nachtergaele1996Spectral}.
\end{lemma}
\begin{proof}\leanok
    \uses{lem:qca_quasi_local_blocking_translation,thm:qca_blocking_quasi_local_interval,thm:mps_quasi_local_translation_invariance,thm:mps_quasi_local_interval_expectation,thm:quasi_local_bulk_inclusion}
    Translation by $bL$ original sites corresponds to translation by $b$
    blocked sites, under which the blocked state is invariant. The interval
    identities follow from the coordinate form of the blocking isomorphism
    and the definition of the quasi-local state on intervals.
\end{proof}

\begin{lemma}[A common translation period of pure ground states]
    \label{thm:periodic_family_pure_ground_state_period}
    \lean{MPSTensor.exists_common_translation_period_of_pure_parentGroundStates_of_periodic_family}
    \leanok
    Let $d\geq1$, and let $A_j$ be finitely many periodic tensors with nonzero weights $\mu_j$,
    and let $h\geq0$ be an interaction of positive range whose open-chain
    kernels eventually equal the boundary spaces of
    $\bigoplus_j\mu_jA_j$. There is $L>0$ such that every pure zero-energy
    state of $h$ is invariant under translation by every integer multiple
    of $L$. No invariance under one-site translation and no minimality of
    the period is asserted \cite[Theorem~1.1]{Nachtergaele1996Spectral}.
\end{lemma}
\begin{proof}\leanok
    \uses{thm:periodic_block_family_ground_states,
        thm:blocked_stationary_expectation}
    Every pure zero-energy state is one of the pulled-back primitive sector
    states of a common blocking by $L$, and each such state is invariant
    under translations by multiples of $L$.
\end{proof}

\begin{lemma}[Every original local observable has blocked interval coordinates]
    \label{thm:original_local_observable_blocked_interval}
    \lean{SpinChain.exists_blocked_quasiLocalIntervalObservable}
    \leanok
    Let $d,L>0$. For every observable $X$ on an original interval
    $[a,a+k)$ with $k>0$, there exist $b\in\mathbb Z$, $N>0$, and an
    observable $Y$ on $N$ blocked sites such that
    \begin{align}
        \mathcal B_L\bigl(j_{b,N}(Y)\bigr)=j_{a,k}(X).
    \end{align}
    Negative interval starts are included. This is the local observable
    form of the regrouping in
    \cite[lines~825--836 and Section~3]{Nachtergaele1996Spectral}.
\end{lemma}
\begin{proof}\leanok
    \uses{thm:qca_blocking_quasi_local_interval,
        thm:quasi_local_interval_observable}
    Put $c=a\bmod L$, $b=\lfloor a/L\rfloor$, and $N=c+k$.
    Since $L\geq1$, the original interval lies inside $[bL,(b+N)L)$.
    Place $X$ at offset $c$ in this enlarged interval, adjoin identities,
    and invert the configuration reindexing. The interval placement and
    blocking identities give the asserted equality.
\end{proof}

\begin{lemma}[Sector separation in identified bond spaces]
    \label{thm:sector_separation_bond_transport}
    \lean{MPSTensor.BlocksNotGaugePhaseEquiv.forall_ne_transport}
    \leanok
    Let $(A_j)$ be a finite family of pairwise gauge-phase inequivalent
    tensors, of bond dimensions $D_j$. If $i\neq j$ and $D_j=D_i$, then
    the tensor $A_j$, carried by this bond-space identification, is
    not gauge-phase equivalent to $A_i$. No normalization, positive
    dimension, or nonempty-family assumption is required. This is the
    separation condition of
    \cite[Proposition~2.7]{Cirac2016MPDO_arXiv}
    written in the bond-space orientation used by the sector support
    and gap arguments.
\end{lemma}
\begin{proof}\leanok
    Apply the separation condition to the ordered pair $(j,i)$ and the
    given dimension equality. The two descriptions of the bond-space
    identification are identical.
\end{proof}

\begin{definition}[Residual boundary contraction after blocking]
    \label{def:blocked_residual_boundary_map}
    \lean{MPSTensor.blockedResidualBoundaryMap,
        MPSTensor.blockedResidualBoundaryMapES,
        MPSTensor.blockedResidualBoundaryMapES_frobeniusEquivEuclidean_apply}
    \leanok
    Let $A$ have physical dimension $d$ and bond dimension $D$, let $B$
    have blocked physical dimension $d^L$ and bond dimension $E$, and let
    $V\colon\C^E\to\C^D$. For $N,r\geq0$, define a linear boundary map
    on $M_{D\times E}(\C)$ by
    \begin{align}
        \mathcal F_{A,B,V}^{N,r}(Y)(u,\tau)
        =\tr(B^u V^\dagger A^\tau Y),
    \end{align}
    where $u$ is the $N$-letter blocked prefix and $\tau$ is the
    $r$-letter original tail. Its physical realization lies on $NL+r$
    original sites. The Euclidean realization uses Hilbert--Schmidt
    virtual coordinates, with the column index preceding the row index.
    All dimensions and lengths may be zero. These are the residual
    word contractions under the compression in
    \cite[Lemma~\texttt{blocking-arbitrary}]{DeLasCuevas2017Irreducible}
    and the local supports of
    \cite[Section~3, lines~1504--1538]{Nachtergaele1996Spectral}.
\end{definition}
\begin{proof}\leanok
    Linearity follows from matrix multiplication and trace. Frobenius
    vectorization and the physical $\ell^2$ identification preserve
    the resulting vector.
\end{proof}

\begin{lemma}[Residual contraction as an original boundary vector]
    \label{thm:blocked_residual_boundary_identity}
    \lean{MPSTensor.blockedResidualBoundaryMap_eq_groundSpaceMap_mul_conjTranspose}
    \leanok
    If $V^\dagger(A^{[L]})^i=B^iV^\dagger$ for every blocked letter $i$,
    then every $N,r\geq0$ and $Y\in M_{D\times E}(\C)$ satisfy
    \begin{align}
        \mathcal F_{A,B,V}^{N,r}(Y)
        =\Gamma_{NL+r}^A(YV^\dagger).
    \end{align}
    This is the direct boundary-vector identity from compressed words,
    including empty prefixes and tails.
\end{lemma}
\begin{proof}\leanok
    \uses{def:blocked_residual_boundary_map, lem:eval_word_block,
        thm:left_boundary_map_factorization}
    The letter intertwiner extends to every blocked word, so
    $B^uV^\dagger=V^\dagger A^{\widetilde u}$. Word concatenation and
    cyclicity of the rectangular trace give
    $\tr(B^uV^\dagger A^\tau Y)
        =\tr(A^{\widetilde u}A^\tau YV^\dagger)$.
\end{proof}

\begin{lemma}[Residual sector resolution of local support]
    \label{thm:blocked_residual_ground_space_resolution}
    \lean{MPSTensor.groundSpace_eq_iSup_range_blockedResidualBoundaryMap,
        MPSTensor.groundSpaceES_eq_iSup_range_blockedResidualBoundaryMapES}
    \leanok
    Let $B_a,V_a$ be a finite family satisfying
    $\sum_aV_aV_a^\dagger=I$ and
    $V_a^\dagger(A^{[L]})^i=B_a^iV_a^\dagger$ for every $a,i$.
    Then, for every $N,r\geq0$,
    \begin{align}
        \mathcal G_{NL+r}(A)
        =\sum_a\operatorname{ran}\mathcal F_{A,B_a,V_a}^{N,r}.
    \end{align}
    The same equality holds in the physical Euclidean space. No
    orthogonality or injectivity of the residual ranges is assumed.
\end{lemma}
\begin{proof}\leanok
    \uses{thm:blocked_residual_boundary_identity, def:ground_space_parent,
        def:ground_space_es}
    Resolving the right virtual boundary gives
    $\Gamma^A_{NL+r}(X)=\sum_a\mathcal F_{A,B_a,V_a}^{N,r}(XV_a)$.
    Conversely, each residual vector is the original boundary vector
    $\Gamma^A_{NL+r}(YV_a^\dagger)$. The Euclidean identifications
    image ranges and their finite sum.
\end{proof}

\begin{lemma}[Periodic residual support decomposition]
    \label{thm:periodic_residual_ground_space}
    \lean{MPSTensor.IsPeriodic.exists_stepOrbit_residualGroundSpaceDecomposition}
    \leanok
    Let $A$ be a normalized irreducible tensor of period $m$ and let $L>0$.
    There exist sector dimensions $D_a>0$, tensors $B_a$ on $d^L$ spins,
    and isometries $V_a\colon\C^{D_a}\to\C^D$, indexed by
    $a\in\{0,\ldots,\gcd(m,L)-1\}$, such that $\sum_aD_a=D$,
    $\sum_aV_aV_a^\dagger=I$, and each $B_a$ is trace-preserving with
    period $m/\gcd(m,L)$. Both letter intertwinings hold:
    \begin{align}
        (A^{[L]})^iV_a         & =V_aB_a^i,         \\
        V_a^\dagger(A^{[L]})^i & =B_a^iV_a^\dagger.
    \end{align}
    For every $N,r\geq0$, including zero,
    $\mathcal G_{NL+r}(A)=\sum_a\operatorname{ran}
        \mathcal F_{A,B_a,V_a}^{N,r}$.
    This extends the finite support decomposition in
    \cite[Lemma~\texttt{blocking-arbitrary}]{DeLasCuevas2017Irreducible}
    to an ungrouped terminal interval. The period formula includes the
    peripheral-root correction in
    \cite{gap:dccsp17_blocking_peripheral_roots}.
\end{lemma}
\begin{proof}\leanok
    \uses{thm:periodic_blocked_ground_space,
        thm:blocked_residual_ground_space_resolution}
    Step-orbit compression gives the normalized tensors, support isometries,
    resolution, and right intertwining. The orbit projections commute
    with every blocked letter. Using $V_a^\dagger V_a=I$ gives the left
    intertwining. Apply the residual range resolution to this derived data.
\end{proof}

\begin{lemma}[Uniform open-chain gap at a common period multiple]
    \label{thm:periodic_block_family_aligned_gap_of_dvd}
    \lean{MPSTensor.exists_aligned_openInteractionMatrix_gap_of_periodic_family_of_dvd}
    \leanok
    Let $d\geq1$, and let $(A_j)$ be a finite family of normalized
    periodic tensors with
    periods $m_j$, and let $T=\bigoplus_j\mu_jA_j$, with $\mu_j\neq0$.
    Let $L>0$ satisfy $m_j\mid L$ for every $j$. Suppose $h\geq0$ has
    range $R>0$ and $\ker H_h(n)=\mathcal G_T(n)$ for all sufficiently
    large $n$. There is $\gamma>0$ such that, for every $N\geq0$ and
    every $v\perp\ker H_h(NL)$,
    \begin{align}
        \gamma\lVert v\rVert\leq\lVert H_h(NL)v\rVert.
    \end{align}
    This includes short volumes with their kernels. The original
    family may be empty; its sectors need not be inequivalent and no
    faithful invariant matrices are given. This is the explicit
    periodic-family, block-divisible finite-chain consequence of
    \cite[Theorem~1.2, Section~3, equations (3.12)--(3.14)]{Nachtergaele1996Spectral}
    and \cite[Lemma~6]{DeLasCuevas2017Irreducible}.
    The other original residue lengths and general GVBS presentations
    remain separate
    \cite{gap:nachtergaele96_infinite_volume_ground_projection}.
\end{lemma}
\begin{proof}\leanok
    \uses{thm:periodic_block_family_primitive_presentation,
        thm:blocked_primitive_aligned_uniform_gap}
    Derive an exact primitive faithful inequivalent presentation after
    blocking by $L$. Take the coarse interaction range $K=R+1$;
    positivity of $L$ gives $R+L\leq KL+1$. The positive window comparison
    transfers the primitive coarse gap to every sufficiently long
    original volume of length $NL$. Taking a finite minimum with the
    positive gaps on the finitely many remaining kernel complements
    gives one constant for all $N$.
\end{proof}

\begin{lemma}[Existence of an aligned open-chain gap for a periodic family]
    \label{thm:periodic_block_family_aligned_gap_exists}
    \lean{MPSTensor.exists_aligned_openInteractionMatrix_gap_of_periodic_family}
    \leanok
    Assume $d\geq1$. Under the periodic-family, nonzero-weight, positivity, positive-range,
    and eventual original kernel hypotheses of
    Lemma~\ref{thm:periodic_block_family_aligned_gap_of_dvd}, there exist
    $L>0$ and $\gamma>0$ such that the same norm gap holds at every
    original length $NL$, on the orthogonal complement of its
    kernel. No common blocking length is given, and the original
    family may be empty. These are the aligned finite-chain lengths
    associated with the regrouping in
    \cite[lines 825--836, Theorem~1.2]{Nachtergaele1996Spectral}.
\end{lemma}
\begin{proof}\leanok
    \uses{thm:periodic_block_family_aligned_gap_of_dvd,thm:sector_separation_bond_transport}
    The product $L=\prod_jm_j$ is positive and divisible by every period.
    For an empty family it equals one. Apply the common-multiple theorem.
\end{proof}

\begin{lemma}[Independence of the local energy representative]
    \label{thm:local_commutator_energy_representation}
    \lean{MPSTensor.localCommutatorObservable_eq_of_quasiLocalIntervalObservable_eq}
    \leanok
    Let $h$ have positive range $R$, and let $X,Y$ be observables on positive-length
    intervals with the same image in the quasi-local algebra. The images of their
    stabilized local commutator energies $X^*[H,X]$ and $Y^*[H,Y]$ are equal.
\end{lemma}
\begin{proof}\leanok
    \uses{thm:quasi_local_commutator_locality}
    Place both observables in a common interval containing the interaction margins
    on either side. The finite-volume commutator energy is the same for the two
    representations. Locality identifies it with each stabilized representative.
\end{proof}

\begin{lemma}[Nonnegative energy under a zero-energy positive functional]
    \label{thm:quasi_local_zero_energy_commutator_nonnegative}
    \lean{MPSTensor.quasiLocalIntervalObservable_localCommutatorObservable_nonneg}
    \leanok
    Let $\varphi$ be a continuous complex-linear functional that is nonnegative on
    squares, and let $h\geq0$ have range $R>0$. If $\varphi$ vanishes on every
    translate of $h$, then every stabilized local commutator energy has real,
    nonnegative expectation. Normalization and translation invariance are not
    required.
\end{lemma}
\begin{proof}\leanok
    \uses{thm:positive_functional_commutator_comparison}
    On the enlarged finite interval, write the interaction sum as $H\geq0$ and
    restrict the functional to its matrix algebra. Its value on $H$ is zero, so
    $\varphi(X^*[H,X])=\varphi(X^*HX)\geq0$.
\end{proof}

\begin{lemma}[Commutator-energy comparison under site grouping]
    \label{thm:blocked_commutator_energy_comparison}
    \lean{MPSTensor.blocking_localCommutator_energy_le}
    \leanok
    Let $d,L,R,K>0$, let $h\geq0$ have range $R$, and assume $R+L\leq KL+1$.
    On the grouped chain, let $h_B$ be the original open interaction sum on $KL$
    sites expressed in grouped coordinates, and put $M=KL+1-R$.
    Let $\varphi$ be positive on squares and vanish on every original translate
    of $h$, and let $\psi=\varphi\circ\mathcal B_L$ be its functional on the grouped
    quasi-local algebra. For every grouped interval observable $Y$, with original
    image $X=\mathcal B_L(Y)$,
    \begin{align}
        \psi\bigl(Y^*[H_B,Y]\bigr)
        \leq M\varphi\bigl(X^*[H,X]\bigr).
        \notag
    \end{align}
    Both commutators denote their stabilized finite-range representatives.
\end{lemma}
\begin{proof}\leanok
    In a sufficiently large grouped interval, the decoded Hamiltonian $C$ satisfies
    $0\leq C\leq MH$. The original functional has zero expectation on $H$.
    Positive comparison therefore bounds the two finite commutator energies.
    Locality identifies their stabilized representatives.
\end{proof}

\begin{lemma}[Infinite-volume gap transfer under site grouping]
    \label{thm:blocked_quasi_local_gap_transport}
    \lean{MPSTensor.quasiLocalCommutator_gap_of_blocking,
        MPSTensor.quasiLocalCommutator_limit_gap_of_blocking}
    \leanok
    Under the preceding hypotheses, suppose that a real constant $\gamma$ satisfies
    $\gamma\psi(Y^*Y)\leq\psi\bigl(Y^*[H_B,Y]\bigr)$
    for every positive-length grouped interval observable $Y$ with $\psi(Y)=0$. Then every original
    local observable $X$ with $\varphi(X)=0$ satisfies
    \begin{align}
        \frac{\gamma}{M}\varphi(X^*X)
        \leq\varphi\bigl(X^*[H,X]\bigr).
        \notag
    \end{align}
    The original commutator energy is also the limit of the finite-volume
    commutator energy as both boundary distances tend to infinity, along any
    filter. In particular, a positive grouped gap yields a positive original gap.
    This is a conditional transfer statement; it does not assert the existence
    of a gap for an arbitrary interaction. The argument is the site-grouping
    reduction used in \cite[Section~6]{Nachtergaele1996Spectral}.
\end{lemma}
\begin{proof}\leanok
    \uses{thm:blocked_commutator_energy_comparison,thm:local_commutator_energy_representation,thm:quasi_local_zero_energy_commutator_nonnegative}
    Every original interval observable has a grouped interval representation.
    Apply the grouped gap and the energy comparison to this representation.
    The variance is preserved by the star-algebra isomorphism, and the local
    energy is independent of its interval representation. Divide by $M>0$.
    Locality makes the finite-volume energy eventually equal to the fixed
    representative; positivity makes both sides of the inequality real.
\end{proof}

\begin{lemma}[Infinite-volume gap for an explicit periodic tensor presentation]
    \label{thm:periodic_quasi_local_ground_state_gap}
    \lean{MPSTensor.IsPeriodic.exists_pos_quasiLocalCommutator_limit_gap_of_eventual_kernel}
    \leanok
    Let the physical dimension be positive and let $A$ be a periodic MPS tensor.
    Let $h\geq0$ have arbitrary positive range $R$, and suppose that for all
    sufficiently large $N$ the open Hamiltonian has kernel $G_N(A)$.
    There is a constant $\gamma>0$ common to every pure state $\varphi$ in the
    zero-energy face of $h$ such that, for every positive-length local observable
    $X$ with $\varphi(X)=0$,
    \begin{align}
        \gamma\varphi(X^*X)
        \leq\lim_{\ell,r\to\infty}\varphi\bigl(X^*[H_{\ell,r},X]\bigr).
        \notag
    \end{align}
    Here $H_{\ell,r}$ is the original open interaction sum on the interval extending
    $\ell$ sites to the left and $r$ sites to the right of the support of $X$.
    The limit exists as a complex number along every filter for which both margins
    tend to infinity. The same $\gamma$ applies to every such filter and pair
    of margins. Both sides are real, and the inequality uses their usual real order.

    The primitive sectors and their faithful invariant matrices are derived from
    periodicity. No translation-invariance assumption is imposed on $\varphi$.
    This proves the infinite-volume conclusion for an explicit periodic tensor
    presentation. The general GVBS presentation of
    \cite[Theorem~1.2]{Nachtergaele1996Spectral} requires a separate identification
    with these finite support spaces; see
    \cite{gap:nachtergaele96_infinite_volume_ground_projection}.
\end{lemma}
\begin{proof}\leanok
    \uses{thm:periodic_ground_state_presentation_at_period,thm:blocked_primitive_quasi_local_gap,thm:sector_separation_bond_transport}
    Group sites by the period. The periodic decomposition, followed by removal of
    repeated sectors, gives finitely many inequivalent primitive tensors with
    faithful invariant matrices and the same joint support at every length.
    Choose a sufficiently long coarse interaction, equal in grouped coordinates
    to the original interaction sum on its expanded interval. Positive comparison
    preserves its eventual kernel identity.

    The primitive-sector theorem gives one positive coarse gap. The classification
    of pure zero-energy states identifies every original pure state with a
    pulled-back primitive sector state. Gap transfer gives the original constant
    $\gamma_B/(KL+1-R)>0$, uniformly over these sectors and all original local
    observables. Locality identifies the finite-volume energy limit.
\end{proof}

\begin{lemma}[Correlated residual overlap]
    \label{thm:correlated_residual_overlap}
    \lean{MPSTensor.norm_inner_blockedResidualBoundaryMapES_le_of_prefix_bound}
    \leanok
    Let $B$ and $C$ be blocked-prefix tensors, and form their residual boundary
    ranges using arbitrary original tails and rectangular support maps. If
    $\lvert\langle u,v\rangle\rvert\leq\varepsilon\lVert u\rVert\lVert v\rVert$
    for $u\in\mathcal G_N(B)$ and $v\in\mathcal G_N(C)$, with $\varepsilon\geq0$,
    the same bound holds between their residual ranges at length $NL+r$,
    for every $r\geq0$.
    \uses{def:blocked_residual_boundary_map}
\end{lemma}
\begin{proof}\leanok
    Fixing a tail configuration gives a prefix ground vector with boundary
    $V^*A^\tau Y$. Sum the prefix inner products over the tail configurations
    and apply Cauchy--Schwarz. The sums of the squared fiber norms are exactly
    the squared norms of the full vectors.
\end{proof}

\begin{lemma}[Residual-sector overlap decay]
    \label{thm:residual_sector_overlap_decay}
    \lean{MPSTensor.IsPrimitiveMPS.eventually_norm_inner_blockedResidualBoundaryMapES_le_of_inequivalent}
    \leanok
    For two inequivalent normalized primitive blocked-prefix tensors of
    positive bond dimensions with faithful invariant matrices, their correlated residual ranges become
    uniformly orthogonal as $N\to\infty$. The estimate is uniform in the
    original tail length $r$; the original tensors giveing the tails need
    no normalization assumption.
\end{lemma}
\begin{proof}\leanok
    \uses{thm:correlated_residual_overlap,thm:inequivalent_block_ground_space_overlap_decay}
    Apply primitive prefix-space overlap decay and the preceding fiberwise
    estimate, which preserves the overlap coefficient.
\end{proof}

\begin{lemma}[Residual-sector projection sum]
    \label{thm:residual_sector_projection_sum}
    \lean{MPSTensor.tendsto_norm_residual_sector_projection_sum_sub_iSup_zero}
    \leanok
    Let $B_j$ be a finite family of pairwise gauge-phase inequivalent normalized
    primitive blocked-prefix tensors of positive bond dimensions with
    faithful invariant matrices.
    Let $F_j(N,r)$ be their correlated residual ranges. Along any filter for
    which $N\to\infty$, and with arbitrary varying $r$, one has
    \begin{align}
        \left\lVert\sum_jP_{F_j(N,r)}-P_{\sum_jF_j(N,r)}\right\rVert\longrightarrow0.
        \notag
    \end{align}
\end{lemma}
\begin{proof}\leanok
    \uses{thm:residual_sector_overlap_decay,thm:sector_projection_sum_error_decay}
    The finite family of residual ranges is asymptotically orthogonal by the
    uniform overlap estimate. Apply the finite-family projection comparison.
\end{proof}

\begin{lemma}[A pure-state gap from an exact blocked primitive presentation]
    \label{thm:blocked_primitive_quasi_local_gap}
    \lean{MPSTensor.exists_pos_quasiLocalCommutator_limit_gap_of_blocked_primitive_family}
    \leanok
    Let $d,L>0$ and let $A$ be an MPS tensor. Suppose that a finite,
    possibly empty, family of pairwise gauge-phase inequivalent normalized
    primitive tensors $B_j$ of positive bond dimensions has faithful invariant
    matrices, nonzero weights $\mu_j$, and the exact identities
    \begin{align}
        \forall\,N\geq0:\ \mathcal G_N(A^{[L]})
        =\mathcal G_N\Bigl(\bigoplus_j\mu_jB_j\Bigr).
        \notag
    \end{align}
    Let $h\geq0$ have range $R>0$, and suppose
    $\ker H_h(n)=\mathcal G_n(A)$ for all sufficiently large $n$.
    There is one $\gamma>0$ with the following property. For every pure
    zero-energy state $\varphi$, every $a\in\mathbb Z$, and every matrix $X$
    on $k>0$ consecutive original sites, put $X_a=j_{a,k}(X)$ and assume
    $\varphi(X_a)=0$. Along every filter for which the two margins
    $\ell,r$ tend to infinity, there is $e\in\mathbb C$ such that
    \begin{align}
        \varphi\bigl(X_a^*[H_{\ell,r}^{a,k},X_a]\bigr) & \longrightarrow e,
                                                       &
        \gamma\varphi(X_a^*X_a)                        & \leq e,
        \label{eq:blocked_primitive_literal_gap}                                                        \\
        H_{\ell,r}^{a,k}
                                                       & =j_{a-\ell,\ell+k+r}\bigl(H_h(\ell+k+r)\bigr).
        \nonumber
    \end{align}
    The energy and variance are real and nonnegative. The same $\gamma$
    works for all states, observables, filters, and diverging margins.
    No translation invariance of $\varphi$ or short-range injectivity
    bound is assumed. This is the exact-presentation version of the
    blocking reduction in \cite[Theorem~1.2 and Section~6]{Nachtergaele1996Spectral};
    identification with a general GVBS presentation remains separate,
    as recorded in \cite{gap:nachtergaele96_infinite_volume_ground_projection}.
\end{lemma}
\begin{proof}\leanok
    \uses{thm:blocked_interaction_eventual_kernel,
        thm:eventual_kernel_sector_local_gap,
        thm:eventual_kernel_ground_state_support_iff,
        thm:blocked_primitive_supported_state_classification,
        thm:blocked_quasi_local_gap_transport}
    Take $K=R+L$ and group the original Hamiltonian on $KL$ sites.
    Coverage gives $R+L\leq KL+1$, so its eventual kernel is the given
    joint primitive space. Obtain one positive coarse gap for all sectors.
    The eventual original kernel identifies every pure zero-energy state
    with a pulled-back sector. Positive comparison transfers the coarse
    inequality with divisor $KL+1-R>0$. The coarse constant is chosen
    before any state, observable, or exhaustion. Finite-range locality
    gives the finite-volume limit and positivity gives its reality.
\end{proof}

\begin{lemma}[Original-chain commutator gaps for a finite periodic family]
    \label{thm:periodic_block_family_quasi_local_gap}
    \lean{MPSTensor.exists_pos_quasiLocalCommutator_limit_gap_of_periodic_family}
    \leanok
    Let $d>0$, let $(A_j)$ be a finite, possibly empty, family of normalized
    periodic tensors of periods $m_j$, and let $\mu_j\neq0$.
    Put $T=\bigoplus_j\mu_jA_j$. Let $h\geq0$ have arbitrary range $R>0$,
    and suppose that $\ker H_h(n)=\mathcal G_n(T)$ for every sufficiently
    large $n$. There is one $\gamma>0$ such that every pure zero-energy
    state and every centered original local observable satisfy
    (\ref{eq:blocked_primitive_literal_gap}), for every filter and every pair
    of diverging boundary margins. The common blocking length, primitive
    tensors, faithful invariant matrices, and sector inequivalence are
    derived from periodicity. Competing states are not assumed translation
    invariant. For an empty tensor family the zero-energy state face is empty.
    This concerns explicit periodic tensor families; the general GVBS
    identification in \cite[Theorem~1.2]{Nachtergaele1996Spectral}
    remains separate; see
    \cite{gap:nachtergaele96_infinite_volume_ground_projection}.
\end{lemma}
\begin{proof}\leanok
    \uses{thm:periodic_block_family_primitive_presentation_exists,thm:sector_separation_bond_transport,thm:blocked_primitive_quasi_local_gap}
    Block by the product of the positive periods, which equals one for an
    empty family. Decompose the blocked sectors, remove gauge-phase repeats,
    and obtain an exact joint primitive presentation at every length.
    Apply the preceding theorem to the original tensor $T$ and interaction $h$.
\end{proof}

\begin{lemma}[Periodic residual projection comparison]
    \label{thm:periodic_residual_projection_sum}
    \lean{MPSTensor.IsPeriodic.exists_residual_primitive_projection_sum}
    \leanok
    Let $A$ be periodic with period $m$. There are $m$ primitive normalized
    compressed tensors $B_j$, positive virtual dimensions $D_j$, support
    isometries $V_j$, and faithful invariant matrices $\rho_j$. The tensors
    $B_j$ are pairwise gauge-phase inequivalent and satisfy
    \begin{align}
        V_j^*V_j=I,\qquad \sum_jV_jV_j^*=I,\qquad
        V_j^*A^{(m),i}=B_j^iV_j^*.
        \notag
    \end{align}
    Let $F_j(N,r)$ be the correlated rectangular boundary range with
    coefficient $\operatorname{tr}(B_j^uV_j^*A^\tau Y)$. Then
    \begin{align}
        \mathcal G_{Nm+r}(A)=\sum_jF_j(N,r)
        \notag
    \end{align}
    for every $N,r\geq0$. For every function $r:\mathbb N\to\mathbb N$,
    \begin{align}
        \left\lVert\sum_jP_{F_j(N,r(N))}-P_{\mathcal G_{Nm+r(N)}(A)}\right\rVert
        \longrightarrow0 \qquad(N\to\infty).
        \notag
    \end{align}
    All primitive normalization and separation data are derived from
    periodicity; none are assumed separately. This comparison does not assert
    an overlapping-window defect estimate or a gap at all residue lengths.
\end{lemma}
\begin{proof}\leanok
    \uses{thm:cyclic_primitive_sector_resolution,thm:blocked_residual_ground_space_resolution,thm:residual_sector_projection_sum}
    Choose the primitive cyclic sector resolution, including its faithful
    invariant matrices, pairwise separation, and both intertwiners. The support
    projections resolve the virtual identity, which gives the exact residual
    support decomposition. Apply the residual-sector projection-sum limit and
    substitute this decomposition.
\end{proof}

\begin{lemma}[Residual boundary inner products]
    \label{thm:residual_boundary_inner_fibers}
    \lean{MPSTensor.inner_blockedResidualBoundaryMapES_eq_sum_groundSpaceMapES}
    \leanok
    For residual boundary vectors $\Gamma_{N,r}^{A,B,V}(Y)$ and
    $\Gamma_{N,r}^{A,B,V}(Z)$, one has
    \begin{align}
        \langle\Gamma_{N,r}^{A,B,V}(Y),\Gamma_{N,r}^{A,B,V}(Z)\rangle
        =\sum_{\tau\in[d]^r}
        \langle\Gamma_N^B(V^*A^\tau Y),\Gamma_N^B(V^*A^\tau Z)\rangle.
        \notag
    \end{align}
    This includes zero prefix and tail lengths and requires no normalization
    or letter intertwining.
    \uses{def:blocked_residual_boundary_map}
\end{lemma}
\begin{proof}\leanok
    Decode the blocked prefix and split the physical configuration into prefix
    and tail. Fixing the tail gives exactly the displayed prefix contraction;
    sum the resulting scalar inner products over both configurations.
\end{proof}

\begin{lemma}[Supported Gram inverses and range projections]
    \label{thm:supported_gram_projector_cancellation}
    \lean{ContinuousLinearMap.comp_supportedGramInverse_adjoint_eq_starProjection}
    \lean{ContinuousLinearMap.starProjection_comp_sub_of_supportedGramFactorizations}
    \leanok
    Let $T:E\to\mathcal H$ be a continuous linear map from a finite-dimensional
    Hilbert space to a complete Hilbert space. Suppose $Q=Q^*$, $TQ=T$, and
    $(T^*T)H=Q$. Then $P_{\operatorname{ran}T}=THT^*$.

    More generally, let $T:E\to\mathcal H$, $L:F\to\mathcal H$, and
    $C:G\to\mathcal H$ have finite-dimensional domains and a common
    factorization $C=TJ_1=LJ_2$. For $S\in\{T,L,C\}$, suppose
    $Q_S=Q_S^*$, $SQ_S=S$, and $(S^*S)H_S=Q_S$. If also $H_T=H_T^*$, then
    \begin{align}
        P_{\operatorname{ran}T}P_{\operatorname{ran}L}-P_{\operatorname{ran}C}
         & =TH_T\bigl[T^*L-(T^*T)J_1H_CJ_2^*(L^*L)\bigr]H_LL^*.
        \label{eq:ph_supported_gram_residual}
    \end{align}
    This is the supported-boundary algebra for the interval projector
    comparison in \cite[Lemma~6.1(ii)]{Nachtergaele1996Spectral}.
\end{lemma}
\begin{proof}\leanok
    The vector $THT^*x$ lies in $\operatorname{ran}T$. Taking adjoints of
    $TQ=T$ gives $QT^*=T^*$, so
    $T^*(x-THT^*x)=T^*x-QT^*x=0$. Thus $THT^*$ is the orthogonal range
    projection. For the three-map identity, self-adjointness of $H_T$ gives
    $H_T(T^*T)=Q_T$. Consequently
    $TH_T(T^*T)=T$ and $(L^*L)H_LL^*=L^*$.
    Expanding the right-hand side of~(\ref{eq:ph_supported_gram_residual})
    and using $TJ_1=C$ and $J_2^*L^*=C^*$ gives
    $P_{\operatorname{ran}T}P_{\operatorname{ran}L}-CH_CC^*$.
\end{proof}

\begin{definition}[Opposite cyclic boundary degree]
    \label{def:cyclic_boundary_degree}
    \lean{MPSTensor.cyclicBoundarySpace,
        MPSTensor.cyclicBoundarySpace_mul_left,
        MPSTensor.cyclicBoundarySpace_mul_right,
        MPSTensor.conjTranspose_mem_cyclicMatrixSubspace_of_cyclicBoundary}
    \leanok
    For matrices $P_u$ indexed by $\mathbb Z/m\mathbb Z$, define
    \begin{align}
        V_n=\{X:\ P_{u+n}X=XP_u\text{ for every }u\}.
        \notag
    \end{align}
    If $P_uA_i=A_iP_{u+1}$, both $A_iX$ and $XA_i$ belong
    to $V_n$ whenever $X\in V_{n+1}$.
    When the $P_u$ are Hermitian, $X\in V_n$ implies that
    $X^*$ has the word degree $n$, namely $P_uX^*=X^*P_{u+n}$.
    These are the cyclic coordinates of
    \cite[Lemma~5]{DeLasCuevas2017Irreducible}.
\end{definition}
\begin{proof}\leanok
    Substitute the cyclic letter identity in each product. For the right
    product, use the boundary relation at $u-1$. Taking adjoints
    exchanges the two degree conventions.
\end{proof}

\begin{lemma}[Compression of a cyclic boundary]
    \label{thm:cyclic_boundary_compression}
    \lean{MPSTensor.sum_cyclicBoundary_mem,
        MPSTensor.groundSpaceMap_sum_cyclicBoundary,
        MPSTensor.exists_cyclicBoundary_of_mem_groundSpace}
    \leanok
    Let $m>0$, let $P_uP_v=\delta_{uv}P_u$, and assume
    $\sum_uP_u=I$. If every length-$n$ word satisfies
    $P_uA^\sigma=A^\sigma P_{u+n}$, then
    \begin{align}
        C_n(X)=\sum_uP_{u+n}XP_u\in V_n,
        \qquad \Gamma_n(C_n(X))=\Gamma_n(X).
        \notag
    \end{align}
    In particular every vector in $G_n(A)$ has a representative in $V_n$.
    This is the cyclic-boundary preparation for the intersection
    property in \cite[equations (3.10)--(3.11)]{Nachtergaele1996Spectral}.
\end{lemma}
\begin{proof}\leanok
    \uses{def:cyclic_boundary_degree}
    Orthogonality gives both sides of the boundary relation for $C_n(X)$.
    Cycle $P_u$ through the trace, use the word relation and idempotence,
    and sum the resulting terms with $\sum_uP_u=I$.
\end{proof}

\begin{lemma}[Restricted injectivity of the boundary map]
    \label{thm:boundary_injectivity_adjoint_span}
    \lean{MPSTensor.groundSpaceMap_injOn_of_adjoint_mem_wordSpan}
    \leanok
    Let $S\subseteq\mathcal M_D$ be a complex linear subspace.
    If $X^*$ belongs to the span of the length-$n$ words for every
    $X\in S$, then $\Gamma_n$ is injective on $S$.
\end{lemma}
\begin{proof}\leanok
    For $X,Y\in S$ with $\Gamma_n(X)=\Gamma_n(Y)$, the trace pairing
    of $X-Y$ with every length-$n$ word vanishes. Test against
    $(X-Y)^*$; then $\operatorname{tr}((X-Y)^*(X-Y))=0$, so $X=Y$.
\end{proof}

\begin{lemma}[The graded one-step intersection]
    \label{thm:cyclic_boundary_intersection}
    \lean{MPSTensor.mem_groundSpace_of_cyclicBoundary_restrictions,
        MPSTensor.groundSpace_eq_restriction_intersection_of_cyclicBoundary,
        MPSTensor.groundSpace_eq_restriction_intersection_of_cyclic_wordSpan}
    \leanok
    Suppose $P_uA_i=A_iP_{u+1}$ and $\sum_iA_i^*A_i=I$.
    If every vector of $G_{n+1}(A)$ has a boundary representative in
    $V_{n+1}$, and $\Gamma_n$ is injective on $V_n$, then
    \begin{align}
        G_{n+2}(A)
        =(G_{n+1}(A)\otimes\mathbb C^d)
        \cap(\mathbb C^d\otimes G_{n+1}(A)).
        \notag
    \end{align}
    In particular these hypotheses follow when $m>0$, the $P_u$
    are Hermitian orthogonal projections summing to $I$, and
    the length-$n$ word span equals
    $\{Z:\ P_uZ=ZP_{u+n}\text{ for every }u\}$.
    This is an algebraic criterion for
    \cite[equations (3.10)--(3.11) and Lemma existenceinteraction]
    {Nachtergaele1996Spectral}; eventual spanning for a periodic tensor
    is a separate step.
\end{lemma}
\begin{proof}\leanok
    \uses{def:cyclic_boundary_degree,thm:cyclic_boundary_compression,
        thm:boundary_injectivity_adjoint_span,
        thm:ground_space_restriction_overlap,thm:ground_space_grow_back}
    Write the two sets of one-site restrictions with boundaries $Y_i,Z_j$
    in $V_{n+1}$. Their overlap gives
    $\Gamma_n(A_jY_i)=\Gamma_n(Z_jA_i)$.
    Both boundaries belong to $V_n$, hence $A_jY_i=Z_jA_i$.
    The left-canonical identity reconstructs the enlarged boundary as
    $\sum_jA_j^*Z_j$.
    Under the stated spanning condition, cyclic compression gives
    the representatives, and the adjoint-span criterion gives injectivity.
\end{proof}

\begin{definition}[Cyclic matrix degrees]
    \label{def:cyclic_word_degree}
    \lean{MPSTensor.cyclicMatrixSubspace,
        MPSTensor.evalWord_mem_cyclicMatrixSubspace,
        MPSTensor.wordSpan_le_cyclicMatrixSubspace}
    \leanok
    For a family $(P_u)_{u\in\mathbb Z/m\mathbb Z}$, set
    \begin{align}
        W_n=\{X:\ P_uX=XP_{u+n}\text{ for every }u\}.
        \notag
    \end{align}
    If $P_uA_i=A_iP_{u+1}$, every word of length $n$ belongs
    to $W_n$. The family need not consist of projections for this conclusion.
    These degrees describe the cyclic sectors of
    \cite[Lemma~5]{DeLasCuevas2017Irreducible}.
\end{definition}
\begin{proof}\leanok
    Induct on the word length and move each $P_u$ through the first letter.
\end{proof}

\begin{lemma}[Subtraction of an adjoint degree]
    \label{thm:cyclic_adjoint_degree}
    \lean{MPSTensor.mul_conjTranspose_mem_cyclicMatrixSubspace}
    \leanok
    Suppose every $P_u$ is Hermitian. If $X\in W_{q+r}$ and
    $Y\in W_r$, then $XY^*\in W_q$.
\end{lemma}
\begin{proof}\leanok
    \uses{def:cyclic_word_degree}
    Take the adjoint of the relation for $Y$ at $u+q$ and multiply
    it by $X$, using the relation for $X$ at $u$.
\end{proof}

\begin{lemma}[Propagation of cyclic word spanning]
    \label{thm:cyclic_word_span_propagation}
    \lean{MPSTensor.wordSpan_eq_cyclicMatrixSubspace_add,
        MPSTensor.wordSpan_eq_cyclicMatrixSubspace_of_ge}
    \leanok
    Suppose that every $P_u$ is Hermitian,
    $P_uA_i=A_iP_{u+1}$, and $\sum_iA_i^*A_i=I$.
    If the words of length $q$ span $W_q$, then the words of
    length $N$ span $W_N$ for every $N\ge q$.
    This algebraic implication gives the later lengths in the
    intersection argument of
    \cite[equations (3.10)--(3.11)]{Nachtergaele1996Spectral};
    the initial spanning equality is assumed here.
\end{lemma}
\begin{proof}\leanok
    \uses{def:cyclic_word_degree,thm:cyclic_adjoint_degree}
    Write $N=q+r$. The left-canonical identity implies
    $\sum_{|w|=r}A_w^*A_w=I$. For $X\in W_{q+r}$, reconstruct
    \begin{align}
        X=\sum_{|w|=r}(XA_w^*)A_w.
        \notag
    \end{align}
    Every coefficient $XA_w^*$ belongs to $W_q$, hence is a
    combination of words of length $q$. The displayed sum is therefore
    a combination of words of length $q+r$. The reverse inclusion
    follows from the cyclic degree of each word.
\end{proof}

\begin{definition}[Sector matrix embedding]
    \label{def:sector_matrix_embedding}
    \lean{MPSTensor.sectorMatrixEmbedding}
    \leanok
    For rectangular matrices $V_j$ with a common row space, the sector
    embedding sends a family $M=(M_j)_j$ of square matrices to
    $\mathcal E(M)=\sum_jV_jM_jV_j^*$.
\end{definition}

\begin{lemma}[Word spans under an isometric sector resolution]
    \label{thm:embedded_sector_word_span}
    \lean{MPSTensor.sectorMatrixEmbedding_wordTuple,
        MPSTensor.sectorMatrixEmbedding_range_le_wordSpan,
        MPSTensor.sectorMatrixEmbedding_compress_eq,
        MPSTensor.mem_wordSpan_of_sector_wordTupleSpanTop}
    \leanok
    \uses{def:sector_matrix_embedding}
    Let $C,B_j$ be tensors on the same physical alphabet and let
    $V_j$ be rectangular matrices satisfying
    $\sum_jV_jV_j^*=I$ and $C_iV_j=V_jB_{j,i}$.
    Let $\mathcal E$ be the sector embedding of
    Definition~\ref{def:sector_matrix_embedding}.
    Then every simultaneous word tuple maps to the corresponding word of $C$.
    If the length-$q$ simultaneous word span is the full product matrix algebra,
    $\operatorname{ran}\mathcal E$ is contained in the length-$q$ word span of $C$.
    If also $V_j^*V_j=I$, every matrix $X$ commuting with all
    $V_jV_j^*$ satisfies
    \begin{align}
        X=\mathcal E\bigl((V_j^*XV_j)_j\bigr)
        \notag
    \end{align}
    and hence belongs to that word span.
    No positivity of the sector dimensions or tensor normalization is assumed.
    These are the embedded cyclic-sector identities of
    \cite[Lemma~5]{DeLasCuevas2017Irreducible}.
\end{lemma}
\begin{proof}\leanok
    Intertwining propagates along words. Resolve the virtual identity to
    identify the image of every word tuple. Taking linear spans gives
    the range inclusion. For a commuting matrix, each compressed term
    equals $XV_jV_j^*$; summing reconstructs $X$.
\end{proof}

\begin{lemma}[Initial spanning of cyclic corners]
    \label{thm:cyclic_blocked_initial_word_span}
    \lean{MPSTensor.wordSpan_eq_cyclicMatrixSubspace_of_blocked_sector_span}
    \leanok
    Let $m>0$, let $P_uA_i=A_iP_{u+1}$, and suppose that the tensor
    $A^{[m]}$ has an isometric sector resolution $B_j,V_j$, with
    $j\in\{0,\ldots,m-1\}$, satisfying
    \begin{align}
        V_j^*V_j=I,\qquad \sum_jV_jV_j^*=I,\qquad
        A^{[m]}_iV_j=V_jB_{j,i}.
        \notag
    \end{align}
    Suppose each support $V_jV_j^*$ is one of the matrices $P_u$;
    their cyclic labels need not be prescribed.
    If the length-$q$ simultaneous word span of $B_j$ is full, then
    \begin{align}
        \operatorname{span}\{A^\sigma:|\sigma|=qm\}
        =\{X:\ P_uX=XP_{u+qm}\text{ for every }u\}.
        \notag
    \end{align}
    This gives the initial graded spanning step for the intersection
    construction in \cite[equations (3.10)--(3.11) and Section~6]
    {Nachtergaele1996Spectral}.
\end{lemma}
\begin{proof}\leanok
    \uses{thm:embedded_sector_word_span}
    Every original word has its length as cyclic degree. Conversely $qm$
    is zero modulo $m$, so a matrix of this degree commutes with every
    sector support. Apply the embedded-sector lemma, then identify the
    length-$q$ blocked word span with the original length-$qm$ span.
\end{proof}

\begin{lemma}[Faithful primitive fixed point at period one]
    \label{thm:period_one_primitive_witness}
    \lean{MPSTensor.exists_isPrimitiveMPS_of_isPeriodic_one}
    \leanok
    \uses{def:primitive_mps}
    Let $A$ be a normalized irreducible tensor of period one with positive
    bond dimension. There is a positive-definite matrix $\rho$ such that
    $A$ and $\rho$ form a primitive normalized MPS pair.
    This is the Perron--Frobenius normalization used in
    \cite[Equations~(3.1)--(3.2b)]{Nachtergaele1996Spectral}.
\end{lemma}
\begin{proof}\leanok
    Period one gives a primitive trace-preserving transfer map and
    irreducibility of its Kraus family. The primitive Perron--Frobenius
    theorem gives an invariant matrix $\rho$; irreducibility makes it
    positive definite.
\end{proof}

\begin{lemma}[Norm control from supported projector cancellation]
    \label{thm:supported_gram_projector_norm}
    \lean{ContinuousLinearMap.norm_starProjection_comp_sub_le_of_supported_factorization,
        ContinuousLinearMap.tendsto_norm_starProjection_comp_sub_zero_of_supported_factorizations}
    \leanok
    Let $T:E\to F$, $L:G\to F$, and $C:W\to F$ be bounded
    maps, with finite-dimensional virtual spaces and a complete physical
    Hilbert space. Suppose
    \begin{align}
        P_TP_L-P_C=T H_T R H_L L^*.
        \notag
    \end{align}
    Then
    \begin{align}
        \|P_TP_L-P_C\|
        \le \|T\|\|H_T\|\|H_L\|\|L\|\|R\|.
        \notag
    \end{align}
    For a family along any filter, assume this factorization eventually
    holds, the product $\|T\|\|H_T\|\|H_L\|\|L\|$ is eventually
    bounded, and $\|R\|\to0$. Then $\|P_TP_L-P_C\|\to0$.
    The spaces may vary with the index. This implication gives the
    norm comparison in \cite[Lemma commutation (ii)]{Nachtergaele1996Spectral};
    the factorization, inverse bounds, and residual decay are explicit
    hypotheses here.
\end{lemma}
\begin{proof}\leanok
    Apply submultiplicativity of the operator norm and $\|L^*\|=\|L\|$.
    The eventual bound then gives a constant multiple of $\|R\|$;
    squeeze the nonnegative projection norm between zero and this quantity.
\end{proof}

\begin{definition}[Correlated virtual tail and residual limiting metric]
    \label{def:correlated_residual_gram_metric}
    \lean{MPSTensor.correlatedTailMapES,
        MPSTensor.boundaryFamilyEquiv_correlatedTailMapES_apply,
        MPSTensor.correlatedTailMapES_adjoint_apply,
        MPSTensor.inner_correlatedTailMapES,
        MPSTensor.residualBoundaryGramLimit,
        MPSTensor.residualBoundaryGramLimit_frobeniusEquivEuclidean_apply,
        MPSTensor.blockedResidualBoundaryMapES_adjoint_comp_self}
    \leanok
    Let $Y\in M_{D\times E}(\C)$, let $V\colon\C^E\to\C^D$, and let
    $A$ be an original tensor. Define $C_rY=(V^\dagger A^\tau Y)_\tau$,
    where $\tau$ runs over the original words of length $r$. For
    $\tr\rho\ne0$, set $K_\rho Z=Z\rho/\tr\rho$ on the square prefix
    boundary space. The residual limiting metric is
    $K_{\infty,r}=C_r^\dagger(\bigoplus_\tau K_\rho)C_r$. If
    $Q_r=\sum_\tau A^{\tau\dagger}VV^\dagger A^\tau$, then
    \begin{align}
        \langle C_rY,C_rZ\rangle & =\tr(Y^\dagger Q_rZ),                  &
        C_r^\dagger(x_\tau)_\tau & =\sum_\tau A^{\tau\dagger}Vx_\tau,       \\
        K_{\infty,r}Y            & =Q_rY\rho/\tr\rho,                     &
        (\Gamma_{N,r}^{A,B,V})^\dagger\Gamma_{N,r}^{A,B,V}
                                 & =C_r^\dagger(\bigoplus_\tau K_N^B)C_r.
    \end{align}
    Here $K_N^B=(\Gamma_N^B)^\dagger\Gamma_N^B$. The rectangular
    coordinates place the column index before the row index. All lengths
    may be zero. No letter intertwining or normalization is assumed.
    \uses{def:blocked_residual_boundary_map,thm:residual_boundary_inner_fibers}
\end{definition}
\begin{proof}\leanok
    Expand the Hilbert--Schmidt inner products and sum the trace pairings.
    The residual physical inner product splits into the same prefix fibers.
    The primitive limiting metric acts by right multiplication, giving the
    stated pullback formula.
\end{proof}

\begin{lemma}[Cyclic row corner of a correlated tail]
    \label{thm:correlated_residual_cyclic_corner}
    \lean{MPSTensor.sum_word_conjTranspose_cyclic_projection_mul_word,
        MPSTensor.norm_correlatedTailMapES_le_one_of_wordGram_eq_projection,
        MPSTensor.norm_correlatedTailMapES_le_one_of_cyclic_support,
        MPSTensor.residualBoundaryGramLimit_apply_of_cyclic_support,
        MPSTensor.stepOrbitProjection_self}
    \leanok
    Suppose that $\sum_iA_i^\dagger A_i=I$, that the $P_a$ are orthogonal
    projections, and that $P_aA_i=A_iP_{a+1}$ on a positive cyclic index.
    If $VV^\dagger=P_a$ and $\operatorname{tr}\rho\neq0$, then every
    $r\geq0$ satisfies
    \begin{align}
        \sum_{\tau\in[d]^r} A^{\tau\dagger}P_aA^\tau & =P_{a+r}, &
        \|C_r\|                                      & \leq1,    & K_{\infty,r}Y & =P_{a+r}Y\rho/\tr\rho.
    \end{align}
    More generally, $C_r$ is contractive whenever its exact Gram matrix
    $Q_r$ is an orthogonal projection. Blocking by the period itself leaves
    each cyclic projection as its own step-orbit projection.
    These are the cyclic word consequences of
    \cite[Lemma~\texttt{blocking-arbitrary}]{DeLasCuevas2017Irreducible}.
\end{lemma}
\begin{proof}\leanok
    \uses{def:correlated_residual_gram_metric}
    Carry $P_a$ through each word, then use
    $\sum_\tau A^{\tau\dagger}A^\tau=I$. The complement $I-Q_r$ is
    positive, so $\tr(Y^\dagger Q_rY)\leq\|Y\|_{\rm HS}^2$.
    The limiting action follows from the exact pullback metric.
\end{proof}

\begin{lemma}[Residual Gram error uniformly in the tail length]
    \label{thm:correlated_residual_gram_error}
    \lean{MPSTensor.norm_blockedResidualBoundaryMapES_gram_sub_limit_le,
        MPSTensor.IsPrimitiveMPS.residualBoundaryGram_sub_limit_norm_sq_le_geometric}
    \leanok
    Whenever $\|C_r\|\leq1$,
    \begin{align}
        \|(\Gamma_{N,r}^{A,B,V})^\dagger\Gamma_{N,r}^{A,B,V}-K_{\infty,r}\|
        \leq\|K_N^B-K_\rho\|.
    \end{align}
    Under the cyclic hypotheses of
    Lemma~\ref{thm:correlated_residual_cyclic_corner}, if the positive-dimensional
    prefix tensor $B$ is primitive with invariant matrix $\rho$, there are
    $C>0$ and $0<q<1$ such that every $N\geq1$ and every $r\geq0$ satisfy
    \begin{align}
        \|(\Gamma_{N,r}^{A,B,V})^\dagger\Gamma_{N,r}^{A,B,V}-K_{\infty,r}\|^2
        \leq E^3(Cq^N)^2.
    \end{align}
    No independent tail condition is assumed. This is an intermediate Gram
    estimate for the projection comparison in
    \cite[Lemma~commutation(i), lines~2442--2531]{Nachtergaele1996Spectral}.
\end{lemma}
\begin{proof}\leanok
    \uses{def:correlated_residual_gram_metric,thm:correlated_residual_cyclic_corner,thm:boundary_fiberwise_map_norm,thm:primitive_ground_space_gram_geometric_fixed_point}
    Subtract the two correlated pullback forms. The fiberwise difference has
    norm at most $\|K_N^B-K_\rho\|$, and the two outer correlated maps
    have norm at most one. Apply the primitive prefix Gram bound.
\end{proof}

\begin{lemma}[Derived residual Gram bounds for a periodic tensor]
    \label{thm:periodic_correlated_residual_gram}
    \lean{MPSTensor.IsPeriodic.exists_cyclic_residualBoundaryGram_bound}
    \leanok
    A normalized irreducible periodic tensor $A$ of positive period $m$
    admits cyclic projections $P_a$, positive sector dimensions $D_a$ summing
    to $D$, support isometries $V_a$, primitive compressed tensors $B_a$,
    and faithful invariant matrices $\rho_a$. They satisfy
    $V_aV_a^\dagger=P_a$, $P_aA_i=A_iP_{a+1}$,
    $B_a^i=V_a^\dagger(A^{[m]})^iV_a$, and
    $(A^{[m]})^iV_a=V_aB_a^i$. The blocked MPS ground space agrees with
    the joint compressed-sector space at every length, including zero.
    For each sector, the limiting residual Gram metric is
    $Y\mapsto P_{a+r}Y\rho_a/\tr\rho_a$, and the geometric estimate of
    Lemma~\ref{thm:correlated_residual_gram_error} holds uniformly over
    all original tail lengths $r$. The primitive sectors, faithful matrices,
    and row corners are derived from periodicity.
\end{lemma}
\begin{proof}\leanok
    \uses{thm:correlated_residual_gram_error,thm:correlated_residual_cyclic_corner,thm:periodic_blocked_ground_space,thm:period_one_primitive_witness}
    Compress the step-orbit sectors after blocking by $m$. Their period is
    one, so irreducibility and trace-preserving normalization give faithful
    primitive fixed points. Each step orbit consists of its individual cyclic
    projection. Apply the cyclic limiting-action identity and the uniform
    primitive residual Gram estimate.
\end{proof}

\begin{definition}[Residual window coordinates]
    \label{def:residual_window_coordinates}
    \lean{MPSTensor.rectangularBoundaryFamilyFiberₗ,
        MPSTensor.residualWindowRightVirtualMapES,
        MPSTensor.residualWindowRightMapES,
        MPSTensor.residualWindowLeftMapES,
        MPSTensor.residualWindowMapES,
        MPSTensor.residualWindowConfigEquiv}
    \leanok
    Let $A$ have physical dimension $d$ and bond dimension $D$, let $B$ have
    physical dimension $d^L$ and bond dimension $E$, and let $V\in\C^{D\times E}$.
    For nonnegative integers $K,M,r$, use physical configurations $(u,v,\tau)$
    of respective lengths $K,M,r$, with the first two in the blocked alphabet.
    Write $\mathbf x=(X_u)_u$ for a family in $\C^{D\times E}$ and
    $\mathbf y=(Y_\tau)_\tau$ for a family in $\MN{E}$. Define
    \begin{align}
        (J_1Y)_u                  & =YB^u,                           &
        (T\mathbf x)_{u,v,\tau}   & =\tr(B^vV^*A^\tau X_u),\nonumber   \\
        (L_1\mathbf y)_{u,v,\tau} & =\tr(B^uB^vY_\tau),              &
        (CY)_{u,v,\tau}           & =\tr(B^uB^vV^*A^\tau Y).
        \label{eq:ph_residual_window_maps}
    \end{align}
    The family coordinates are ordered as $(u,\text{column},\text{row})$.
    Decoding the blocked letters and concatenating $u,v,\tau$ gives a bijection
    with the original configurations of length $(K+M)L+r$, including zero lengths.
    These are the residual coordinates for the overlap argument in
    \cite[Lemma~6.1(ii)]{Nachtergaele1996Spectral}.
    \uses{def:blocked_residual_boundary_map}
\end{definition}

\begin{lemma}[Residual window factorizations]
    \label{thm:residual_window_factorizations}
    \lean{MPSTensor.residualWindowRightMapES_comp_rightVirtualMapES,
        MPSTensor.residualWindowLeftMapES_correlatedTailFamily,
        MPSTensor.residualWindowMapES_eq_reindex_blockedResidualBoundaryMapES}
    \leanok
    With the coordinates of Definition~\ref{def:residual_window_coordinates},
    put $(J_2Y)_\tau=V^*A^\tau Y$. Then $TJ_1=C=L_1J_2$.
    $C$ is the residual boundary map at prefix length $K+M$ and tail
    length $r$, carried by the physical configuration bijection.
    No injectivity, normalization, or intertwining hypothesis is required.
    \uses{def:residual_window_coordinates,def:blocked_residual_boundary_map}
\end{lemma}
\begin{proof}\leanok
    Trace cyclicity gives
    $\tr(B^vV^*A^\tau YB^u)=\tr(B^uB^vV^*A^\tau Y)$;
    associativity gives the second factorization.
    The physical-coordinate identity follows from the word product under
    concatenation.
\end{proof}

\begin{lemma}[Residual window mixed Gram]
    \label{thm:residual_window_mixed_gram}
    \lean{MPSTensor.inner_residualWindowRightMapES_leftMapES}
    \leanok
    In the same coordinates, let $G_B(M)=\Gamma_B(M)^*\Gamma_B(M)$ be the
    ordinary prefix Gram. With the Frobenius inner product on virtual matrices,
    \begin{align}
        \langle T\mathbf x,L_1\mathbf y\rangle
         & =\sum_u\sum_\tau
        \big\langle V^*A^\tau X_u,\,
        G_B(M)(Y_\tau B^u)\big\rangle.
        \label{eq:ph_residual_mixed_gram}
    \end{align}
    This identity is exact at every nonnegative length; it asserts no limiting
    metric or quantitative decay.
    \uses{def:residual_window_coordinates,def:ground_space_gram}
\end{lemma}
\begin{proof}\leanok
    \uses{def:ground_space_gram}
    Expand the physical inner product over $(u,v,\tau)$ and interchange the
    $v$ and $\tau$ sums. Trace cyclicity rewrites the left-window coefficient
    as $\tr(B^vY_\tau B^u)$. The $v$ sum is then the Gram pairing of
    $\Gamma_B(M)$.
\end{proof}

\begin{lemma}[Primitive cyclic sector resolution]
    \label{thm:cyclic_primitive_sector_resolution}
    \lean{MPSTensor.IsPeriodic.exists_cyclic_primitive_sector_resolution}
    \leanok
    Let $A$ be a normalized irreducible periodic tensor of period $m$.
    There are positive sector dimensions, primitive normalized tensors $B_j$
    on blocks of $m$ sites, faithful invariant matrices $\rho_j$, and isometries
    $V_j$ whose mutually orthogonal support projections sum to the identity.
    Distinct $B_j$ are inequivalent up to similarity and phase. In the inverse
    cyclic labelling, the projections satisfy $P_{j+1}A_i=A_iP_j$.
    Both rectangular intertwiners hold:
    $A^{[m]}_iV_j=V_jB_{j,i}$ and
    $V_j^*A^{[m]}_i=B_{j,i}V_j^*$.
    This retains the cyclic coordinates of
    \cite[Lemma~5]{DeLasCuevas2017Irreducible} used in the grouping argument of
    \cite[Section~6]{Nachtergaele1996Spectral}.
\end{lemma}
\begin{proof}\leanok
    Take the cyclic isometric resolution of the periodic tensor.
    Each compressed blocked sector has period one, hence admits a faithful
    primitive invariant matrix. The cyclic sector separation theorem gives
    pairwise inequivalence. Multiplying the corner identity by the isometry
    and its adjoint gives the two intertwiners.
\end{proof}

\begin{lemma}[Eventual cyclic spanning and original intersection]
    \label{thm:periodic_original_intersection}
    \lean{MPSTensor.IsPeriodic.exists_eventually_cyclic_wordSpan,
        MPSTensor.IsPeriodic.exists_eventually_groundSpace_restriction_intersection}
    \leanok
    Let $A$ be a normalized irreducible periodic tensor of period $m$.
    There is an orthogonal resolution $\{P_u\}_{u\in\Z/m\Z}$ satisfying
    $P_uA_i=A_iP_{u+1}$, and a threshold $n_0$, such that for every $n\ge n_0$
    the length-$n$ words span precisely
    $\{X:P_uX=XP_{u+n}\text{ for all }u\}$.
    Consequently,
    \begin{align}
        \forall\,n\ge n_0:\ (G_{n+1}(A)\otimes\C^d)\cap(\C^d\otimes G_{n+1}(A))=G_{n+2}(A).
        \notag
    \end{align}
    The lengths here count original sites, with no divisibility requirement.
    This is the cyclic-sector form of the intersection argument in
    \cite[equations (3.10)--(3.11)]{Nachtergaele1996Spectral}.
\end{lemma}
\begin{proof}\leanok
    \uses{thm:cyclic_primitive_sector_resolution,thm:ph_primitive_blocks_word_span,thm:cyclic_blocked_initial_word_span,thm:cyclic_word_span_propagation,thm:cyclic_boundary_intersection}
    Simultaneous spanning of the primitive blocked sectors gives spanning of
    one degree-zero cyclic corner algebra. Negating the cyclic labels puts the
    projections in the word-degree convention. The adjoint-word resolution
    propagates this span to all later lengths. Restricted boundary injectivity
    then gives the intersection identity.
\end{proof}

\begin{lemma}[Canonical original-chain interaction for a periodic tensor]
    \label{thm:periodic_original_parent_interaction}
    \lean{MPSTensor.IsPeriodic.exists_ker_openParentHamiltonianES_eq_groundSpaceES,
        MPSTensor.IsPeriodic.exists_positive_parent_interaction}
    \leanok
    For the same tensor $A$, there is $R_0>0$ such that, for every $R\ge R_0$
    and every $N\ge R$, the canonical range-$R$ open parent Hamiltonian has
    kernel exactly $G_N(A)$. In particular, there exists a positive interaction
    $h$ of some range $R>0$ with this kernel identity at every $N\ge R$.
    No faithful invariant matrix, cyclic resolution, or short-range intersection
    is given as a hypothesis.
    This concerns the explicit normalized periodic tensor presentation; it does
    not assert a presentation theorem for arbitrary GVBS states.
    It implements the finite-support interaction construction of
    \cite[Lemma existenceinteraction and Theorem~1.1]{Nachtergaele1996Spectral}
    in this tensor setting.
\end{lemma}
\begin{proof}\leanok
    \uses{thm:periodic_original_intersection}
    Choose the interaction range beyond the intersection threshold.
    Iteration over contiguous intervals identifies the simultaneous local
    constraints with the full boundary-condition space. Positivity of the
    local complementary projection identifies these constraints with the open
    Hamiltonian kernel, including the single-window volume $N=R$.
    Writing that projection in physical matrix coordinates gives $h$.
\end{proof}

\begin{definition}[Matrix of the canonical parent interaction]
    \label{def:canonical_parent_interaction_matrix}
    \lean{MPSTensor.canonicalParentInteractionMatrix,
        MPSTensor.toEuclideanLin_canonicalParentInteractionMatrix,
        MPSTensor.canonicalParentInteractionMatrix_posSemidef}
    \leanok
    Let \(h_R\) be the configuration-basis matrix of the orthogonal
    projection onto \(G_R(A)^\perp\). This matrix is positive semidefinite
    and represents the same projection on the Euclidean realization of
    the physical configuration space. It is the canonical choice in
    \cite[Lemma existenceinteraction]{Nachtergaele1996Spectral}.
\end{definition}
\begin{proof}\leanok
    Identify matrices with bounded operators in the configuration basis.
    Under this identification, matrix positive semidefiniteness is equivalent
    to positivity of the represented operator. Orthogonal projections are
    positive.
\end{proof}

\begin{lemma}[An original periodic parent interaction with an infinite-volume gap]
    \label{thm:periodic_positive_parent_infinite_gap}
    \lean{MPSTensor.IsPeriodic.exists_positive_parent_interaction_quasiLocalCommutator_gap}
    \leanok
    Let $A$ be a normalized periodic tensor. There exist $R>0$, a
    positive semidefinite interaction $h$ on $R$ sites, and $\gamma>0$
    such that $\ker H_N(h)=G_N(A)$ for every $N\ge R$.
    If $\varphi$ is any pure zero-energy state for this interaction,
    and $X$ is an observable on an interval of positive length with
    $\varphi(X)=0$, then
    \begin{align}
        \varphi(X^*[H_{\ell,r}(h),X])\longrightarrow e,
        \qquad \gamma\varphi(X^*X)\le e,
        \notag
    \end{align}
    along every filter for which both interval margins $\ell,r$
    tend to infinity. The same constant works for all such states,
    observables, and exhaustions. The physical dimension, interaction,
    eventual kernel identity, and primitive invariant matrices are derived.
    This proves the existence and infinite-volume gap conclusions of
    \cite[Theorems 1.1--1.2]{Nachtergaele1996Spectral} for an explicit
    normalized periodic tensor presentation. Identification with a general
    GVBS presentation, and the all-residue finite-volume gap, remain separate.
\end{lemma}
\begin{proof}\leanok
    \uses{thm:periodic_original_parent_interaction,
        thm:periodic_quasi_local_ground_state_gap}
    Choose the positive canonical interaction from the original-chain
    intersection. Its kernel identity for every $N\ge R$ gives the
    eventual kernel hypothesis of the periodic infinite-volume gap theorem.
    Apply that theorem to obtain one positive constant before all states
    and interval exhaustions are chosen.
\end{proof}

\begin{definition}[Finite Euclidean direct-sum operators]
    \label{def:euclidean_fiberwise_operator}
    \lean{EuclideanSpace.fiberₗ,EuclideanSpace.fiberwiseMap,
        EuclideanSpace.fiberₗ_apply,EuclideanSpace.fiberwiseMap_apply,
        EuclideanSpace.fiberₗ_fiberwiseMap}
    \leanok
    Let $S,I,J$ be finite sets and let the scalar field be $\R$ or $\C$.
    For a vector $x\in\ell^2(S\times I)$, its $s$th fiber is
    $x_s=(x_{s,i})_{i\in I}$. For a linear operator
    $G:\ell^2(I)\to\ell^2(J)$, define $G^{\oplus S}$ by
    $(G^{\oplus S}x)_s=Gx_s$.
    The fiber projection itself is defined without a finiteness assumption.
\end{definition}

\begin{lemma}[Finite direct-sum operator identities]
    \label{thm:euclidean_fiberwise_operator}
    \lean{EuclideanSpace.norm_fiberwiseMap_le,
        EuclideanSpace.inner_fiberwiseMap,
        EuclideanSpace.fiberwiseMap_adjoint,
        EuclideanSpace.fiberwiseMap_comp,
        EuclideanSpace.fiberwiseMap_id,
        EuclideanSpace.fiberwiseMap_sub}
    \leanok
    Finite direct sums preserve adjoints, composition, identity operators, and
    operator differences. They satisfy
    $\lVert G^{\oplus S}\rVert\leq\lVert G\rVert$ and
    $\langle y,G^{\oplus S}x\rangle=\sum_s\langle y_s,Gx_s\rangle$.
    These statements include empty coordinate sets.
    \uses{def:euclidean_fiberwise_operator}
\end{lemma}
\begin{proof}\leanok
    The squared norm is the sum of the squared fiber norms. Applying
    $\lVert Gx_s\rVert\leq\lVert G\rVert\lVert x_s\rVert$ to every fiber gives
    the norm bound. The inner-product identity follows by splitting the sum
    over $S\times I$. Apply the adjoint identity in each fiber to obtain the
    adjoint formula. Composition, identity, and subtraction hold coordinatewise.
\end{proof}

\begin{lemma}[Residual window Grams and boundary norms]
    \label{thm:residual_window_grams}
    \lean{MPSTensor.residualWindowRightMapES_adjoint_comp_self,
        MPSTensor.residualWindowLeftMapES_adjoint_comp_self,
        MPSTensor.norm_residualWindowRightMapES_le,
        MPSTensor.norm_residualWindowLeftMapES_le}
    \leanok
    Use the residual windows $T,L_1$ at full length $(K+M)L+r$ from
    Definition~\ref{def:residual_window_coordinates}. Write $\Gamma_{M,r}$
    for the correlated residual boundary map and $G_B(n)$ for the ordinary
    blocked-prefix Gram. Then
    \begin{align}
        T^*T     & =(\Gamma_{M,r}^*\Gamma_{M,r})^{\oplus u}, &
        L_1^*L_1 & =G_B(K+M)^{\oplus\tau}.
        \label{eq:ph_residual_window_grams}
    \end{align}
    Consequently $\lVert T\rVert\leq\lVert\Gamma_{M,r}\rVert$ and
    $\lVert L_1\rVert\leq\lVert\Gamma_B(K+M)\rVert$.
    These identities require no normalization or injectivity and include
    zero lengths and empty configuration sets.
    \uses{def:residual_window_coordinates,def:euclidean_fiberwise_operator,
        def:ground_space_gram}
\end{lemma}
\begin{proof}\leanok
    \uses{thm:residual_boundary_inner_fibers,thm:euclidean_fiberwise_operator}
    Expand the physical inner products over $(u,v,\tau)$.
    For the right Gram, first sum over the residual physical variables
    $(v,\tau)$, retaining $u$ as a spectator.
    For the left Gram, concatenate $u,v$ and retain $\tau$ as a spectator.
    Apply the finite direct-sum norm bound to each Gram and use
    $\lVert F^*F\rVert=\lVert F\rVert^2$.
\end{proof}

\begin{lemma}[Correlated residual window factorization]
    \label{thm:residual_window_correlated_factorization}
    \lean{MPSTensor.residualWindowLeftMapES_comp_correlatedTailMapES}
    \leanok
    The full residual contraction satisfies $L_1C_r=C$, where
    $(C_rY)_\tau=V^*A^\tau Y$ is the correlated tail map.
    This is an operator identity, without injectivity or normalization.
    \uses{def:residual_window_coordinates,def:correlated_residual_gram_metric}
\end{lemma}
\begin{proof}\leanok
    \uses{thm:residual_window_factorizations}
    Identify each correlated tail fiber and apply the left window factorization.
\end{proof}

\begin{lemma}[Residual window Gram errors]
    \label{thm:residual_window_gram_errors}
    \lean{MPSTensor.norm_residualWindowRightMapES_gram_sub_limit_le,
        MPSTensor.norm_residualWindowLeftMapES_gram_sub_limit_le}
    \leanok
    \uses{def:correlated_residual_gram_metric}
    Let $\rho\in\MN{E}$ have nonzero trace, put
    $G_\infty=\operatorname{gramReshuffle}(\mathcal P_\rho)$, and let
    $K_r=C_r^*G_\infty^{\oplus\tau}C_r$.
    If $\lVert C_r\rVert\leq1$, then
    \begin{align}
        \lVert T^*T-K_r^{\oplus u}\rVert
         & \leq\lVert G_B(M)-G_\infty\rVert,\nonumber \\
        \lVert L_1^*L_1-G_\infty^{\oplus\tau}\rVert
         & \leq\lVert G_B(K+M)-G_\infty\rVert.
        \label{eq:ph_residual_window_errors}
    \end{align}
    The second bound does not require the contractivity assumption. Neither
    bound requires primitivity or a fixed-point equation. Both coefficients
    are one and are independent of the spectator lengths.
\end{lemma}
\begin{proof}\leanok
    \uses{thm:euclidean_fiberwise_operator,thm:correlated_residual_gram_error,thm:residual_window_grams}
    Finite direct sums preserve differences and do not increase operator norm.
    For the right window also apply the coefficient-one correlated Gram bound.
\end{proof}

\begin{definition}[Fixed rectangular metric and row projection]
    \label{def:residual_fixed_metric_row_projection}
    \lean{MPSTensor.rectangularRightGramMetric,
        MPSTensor.rectangularRightGramMetric_apply,
        MPSTensor.rectangularRightGramMetric_isUnit,
        MPSTensor.rectangularRightGramMetric_inverse_apply,
        MPSTensor.rectangularRightGramMetric_commute_rowProjection,
        MPSTensor.rectangularRightGramMetric_isSelfAdjoint,
        MPSTensor.rectangularRowProjection,
        MPSTensor.rectangularRowProjection_apply,
        MPSTensor.rectangularRowProjection_isStarProjection}
    \leanok
    On $M_{D\times E}(\C)$ with its Hilbert--Schmidt inner product, define
    $R_\rho Y=Y\rho/\tr\rho$ and $\mathcal QY=QY$.
    If $E>0$ and $\rho>0$, then $R_\rho$ is self-adjoint and invertible,
    with $R_\rho^{-1}Y=(\tr\rho)Y\rho^{-1}$.
    If $Q$ is an orthogonal projection, then so is $\mathcal Q$.
    Right multiplication commutes with every row operator $\mathcal Q$.
\end{definition}
\begin{proof}\leanok
    \uses{def:correlated_residual_gram_metric}
    In column-first rectangular coordinates the operators are
    $(\tr\rho)^{-1}(\rho^{\mathsf T}\otimes I_D)$ and $I_E\otimes Q$.
    Positivity of $\rho$ gives the first operator's invertibility and
    self-adjointness. The remaining assertions follow from multiplication
    of rectangular matrices.
\end{proof}

\begin{lemma}[Exact residual row support]
    \label{thm:residual_exact_row_support}
    \lean{MPSTensor.correlatedTailMapES_adjoint_comp_apply,
        MPSTensor.correlatedTailMapES_adjoint_comp_eq_rowProjection,
        MPSTensor.correlatedTailMapES_comp_rowProjection,
        MPSTensor.blockedResidualBoundaryMapES_comp_rowProjection,
        MPSTensor.residualBoundaryGramLimit_eq_rightMetric_comp_rowProjection}
    \leanok
    Let $C_rY=(V^\dagger A^\tau Y)_\tau$ and suppose that
    $Q_r=\sum_\tau A^{\tau\dagger}VV^\dagger A^\tau$ is an orthogonal
    projection. Write $\mathcal Q_rY=Q_rY$. For every prefix length $N$,
    \begin{align}
        C_r^\dagger C_r          & =\mathcal Q_r,       &
        C_r\mathcal Q_r          & =C_r,                &
        \Gamma_{N,r}\mathcal Q_r & =\Gamma_{N,r},       &
        K_{\infty,r}             & =R_\rho\mathcal Q_r.
    \end{align}
    The last identity requires $\tr\rho\ne0$. No positivity or
    injectivity of the prefix tensor is required for these identities.
\end{lemma}
\begin{proof}\leanok
    \uses{def:residual_fixed_metric_row_projection,def:correlated_residual_gram_metric}
    The exact tail Gram form is $\tr(Y^\dagger Q_rZ)$.
    A boundary map whose Gram operator vanishes on a vector also vanishes
    on that vector. Applying this fact to $I-\mathcal Q_r$ proves the two
    support identities. The residual Gram factorization and the limiting
    right-multiplication formula give the remaining identities.
\end{proof}

\begin{definition}[Extended Gram operator and supported inverse]
    \label{def:residual_extended_gram_supported_inverse}
    \lean{MPSTensor.residualBoundaryGramExtension,
        MPSTensor.residualBoundaryGramInverse,
        MPSTensor.residualBoundaryGramExtension_sub_rightMetric}
    \leanok
    With $G_{N,r}=\Gamma_{N,r}^\dagger\Gamma_{N,r}$, define
    \begin{align}
        F_{N,r} & =G_{N,r}+R_\rho(I-\mathcal Q_r), &
        H_{N,r} & =\mathcal Q_rF_{N,r}^{-1}.
    \end{align}
    Whenever $Q_r$ is the exact tail Gram matrix,
    $F_{N,r}-R_\rho=G_{N,r}-K_{\infty,r}$.
    The inverse equations for $H_{N,r}$ hold under the hypotheses of the
    following theorem.
\end{definition}
\begin{proof}\leanok
    \uses{thm:residual_exact_row_support}
    Substitute $K_{\infty,r}=R_\rho\mathcal Q_r$ and expand the complementary
    term.
\end{proof}

\begin{lemma}[Quantitative supported residual inverse]
    \label{thm:residual_supported_inverse_small_error}
    \lean{MPSTensor.residualBoundaryGramInverse_properties_of_small_error,
        MPSTensor.norm_blockedResidualBoundaryMapES_sq_le_rightMetric_add_gram_error}
    \leanok
    Assume $D,E>0$, $\rho>0$, and that $Q_r$ is an orthogonal projection
    equal to the exact tail Gram matrix. Let
    $\delta_N=\|K_N^B-K_\rho\|$ and $u=\|R_\rho^{-1}\|$.
    If $u\delta_N\leq a<1$, then $F_{N,r}$ is invertible, $H_{N,r}$ is
    self-adjoint, and
    \begin{align}
        \Gamma_{N,r}\mathcal Q_r & =\Gamma_{N,r},   &
        G_{N,r}H_{N,r}           & =\mathcal Q_r,   &
        \|H_{N,r}\|              & \leq(1-a)^{-1}u.
    \end{align}
    Without the small-error condition,
    $\|\Gamma_{N,r}\|^2\leq\|R_\rho\|+\delta_N$.
    These are intermediate estimates for
    \cite[Lemma~commutation(ii)]{Nachtergaele1996Spectral},
    rather than an assertion of the three-interval projection estimate.
\end{lemma}
\begin{proof}\leanok
    \uses{def:residual_extended_gram_supported_inverse,thm:correlated_residual_gram_error}
    The extended error is at most $\delta_N$, so the Neumann estimate
    at the fixed invertible operator $R_\rho$ bounds $F_{N,r}^{-1}$.
    The exact support identity gives
    $F_{N,r}\mathcal Q_r=\mathcal Q_rF_{N,r}=G_{N,r}$.
    Thus its inverse commutes with $\mathcal Q_r$, proving self-adjointness
    and the supported inverse equation. Finally,
    $\|G_{N,r}\|=\|\Gamma_{N,r}\|^2$ and
    $\|R_\rho\mathcal Q_r\|\leq\|R_\rho\|$ give the boundary-map bound.
\end{proof}

\begin{lemma}[Uniform supported inverse convergence]
    \label{thm:residual_supported_inverse_uniform_tail}
    \lean{MPSTensor.IsPrimitiveMPS.eventually_residualBoundaryGramInverse_properties,
        MPSTensor.norm_residualBoundaryGramInverse_sub_limit_le,
        MPSTensor.IsPrimitiveMPS.residualBoundaryGramInverse_sub_limit_tendsto_zero}
    \leanok
    Assume $D>0$, and let the positive-dimensional prefix tensor $B$ be
    primitive with faithful invariant matrix $\rho$. Suppose that every exact tail Gram matrix $Q_r$
    is an orthogonal projection. For all sufficiently large prefix lengths,
    uniformly over every $r\geq0$, the preceding inverse and support equations
    hold and
    \begin{align}
        \|H_{N,r}\|      & \leq2u,           &
        \|\Gamma_{N,r}\| & \leq\|R_\rho\|+2.
    \end{align}
    Whenever $u\delta_N\leq\tfrac12$,
    \begin{align}
        \|H_{N,r}-\mathcal Q_rR_\rho^{-1}\|\leq2u^2\delta_N.
    \end{align}
    Consequently this inverse error tends to zero along every diverging
    prefix length, with arbitrary tail lengths. Under the cyclic hypotheses
    of Lemma~\ref{thm:correlated_residual_cyclic_corner}, the exact projections
    are $Q_r=P_{a+r}$; no restriction on the tail length is needed.
\end{lemma}
\begin{proof}\leanok
    \uses{thm:residual_supported_inverse_small_error,thm:primitive_ground_space_gram_geometric_fixed_point,thm:correlated_residual_cyclic_corner}
    Primitive prefix convergence gives $\delta_N\to0$, so eventually
    $u\delta_N\leq\tfrac12$ and $\delta_N\leq1$.
    The resolvent identity bounds the inverse displacement by
    $\|F_{N,r}^{-1}\|\,\|F_{N,r}-R_\rho\|\,u$.
    Multiplication by the contractive row projection preserves this bound.
    The estimates are independent of $r$, so the same convergence holds
    for every choice of tails and every filter along which the prefix
    diverges.
\end{proof}

\begin{lemma}[Covering interval energies and kernels]
    \label{thm:open_cover_energy_kernel}
    \lean{MPSTensor.ker_openParentHamiltonianES_le_ker_openPrefix,
        MPSTensor.localTermES_eq_zero_of_openSuffixParentHamiltonianES_eq_zero,
        MPSTensor.ker_openPrefix_inf_ker_openSuffix_eq_ker_openParentHamiltonianES,
        MPSTensor.openPrefix_add_openSuffix_le_two_openParentHamiltonianES,
        MPSTensor.leftBoundary_inf_reassocTail_range_eq_groundSpaceES_of_open_kernel}
    \leanok
    Let $H_N$ be the canonical range-$R$ open Hamiltonian, and let $H_L,H_T$
    be its restrictions to a prefix of length $n$ and a terminal interval of
    length $W$, respectively. Then $H_L+H_T\le2H_N$.
    If $n\le N$ and $N+R\le n+W+1$, every interaction window lies in one
    of these intervals, and
    $\ker H_L\cap\ker H_T=\ker H_N$.
    If the finite kernels equal the MPS boundary spaces and the overlap has
    length at least $R$, the corresponding extended boundary spaces intersect
    exactly in $G_N$.
\end{lemma}
\begin{proof}\leanok
    \uses{thm:fixed-volume-prefix-ground-projections-nested,thm:ph_prefix_kernel_free_right_sites,thm:ph_terminal_interval_kernel,thm:open_parent_overlapping_interval_gap}
    Each interval contains a subset of the positive local terms of $H_N$.
    This proves the energy upper bound. A vector annihilated by an interval
    sum is annihilated by every local projection in that sum. The covering
    condition therefore gives the kernel identity. The boundary-space identity
    follows by the physical coordinate isometries adjoining free sites.
\end{proof}

\begin{lemma}[Finite residue gaps and quotient reconstruction]
    \label{thm:finite_open_residue_gaps}
    \lean{MPSTensor.exists_pos_openParentHamiltonianES_gap_over_residues,
        MPSTensor.openParentHamiltonianES_gap_all_large_of_residue_gaps}
    \leanok
    For fixed $R,W$ and $m>0$, the finitely many canonical open Hamiltonians
    at lengths $W+r$, $0\le r<m$, have a common positive gap on their
    kernel complements. A gap valid at all lengths
    $Km+Mm+r$ with $K\ge K_0$ and $r<m$ is valid at every
    $N\ge(K_0+M)m$.
\end{lemma}
\begin{proof}\leanok
    \uses{lem:positive-kernel-projection-comparison}
    Take the minimum of the finitely many positive finite-dimensional gaps.
    For the second assertion use $r=N\bmod m$ and $K=\lfloor N/m\rfloor-M$.
\end{proof}

\begin{lemma}[Fixed-overlap residue-cover gap criterion]
    \label{thm:open_residue_cover_gap_criterion}
    \lean{MPSTensor.exists_pos_openParentHamiltonianES_gap_of_aligned_cover_defect,
        MPSTensor.IsPeriodic.exists_pos_openParentHamiltonianES_gap_all_of_residue_cover_defect}
    \leanok
    Suppose $m>0$, $R>0$, $Mm\ge R$, and the canonical open kernels equal
    $G_N$ for every $N\ge R$. Cover $N=Km+Mm+r$ by the aligned prefix
    of length $(K+M)m$ and the terminal interval of length $Mm+r$.
    Assume that, uniformly for $K\ge K_0$ and $0\le r<m$,
    \begin{align}
        \|P_{G_N}-P_{G_{(K+M)m}\otimes\C^{d^r}}
        P_{\C^{d^{Km}}\otimes G_{Mm+r}}\|\le\tfrac12.
        \notag
    \end{align}
    If the aligned prefix gaps have a common bound $\gamma>0$, the full open
    Hamiltonians have a common bound $\min(\gamma,\eta)/16>0$, where $\eta$
    is a derived minimum over the terminal residues. For a normalized periodic
    tensor of period $m$, the aligned gap is derived from its primitive blocked
    presentation. Consequently all $N\ge(K_0+M)m$ have a common positive gap.
    The displayed projection estimate remains an explicit hypothesis; this
    criterion does not establish its periodic-tensor decay.
    It is the finite-Hamiltonian consequence of the three-interval estimate in
    \cite[Section~6, Lemma commutation(ii)]{Nachtergaele1996Spectral}.
\end{lemma}
\begin{proof}\leanok
    \uses{thm:open_cover_energy_kernel,thm:finite_open_residue_gaps,thm:open_interval_gap_transport,thm:open_parent_overlapping_interval_gap,thm:periodic_tensor_aligned_uniform_gap}
    Adjoining free sites preserves each interval gap. Their energies sum to at
    most twice the full energy, so the two-projection comparison gives the stated
    bound. The aligned periodic gap follows by blocking and positive comparison;
    the terminal gaps require only finite dimensionality. Quotient reconstruction
    then removes the residue restriction.
\end{proof}

\begin{lemma}[Uniform virtual prefix bound]
    \label{thm:residual_window_virtual_norm}
    \lean{MPSTensor.sum_evalWord_frobenius_norm_sq_of_leftCanonical,
        MPSTensor.norm_residualWindowRightVirtualMapES_le_sqrt}
    \leanok
    Let $B$ have bond dimension $E$ and satisfy $\sum_i B_i^*B_i=I$.
    For every prefix length $K\geq0$,
    \begin{align}
        \sum_{|u|=K}\|B^u\|_{\mathrm F}^2=E.
        \notag
    \end{align}
    Consequently, the map $J_1$ of
    Definition~\ref{def:residual_window_coordinates} satisfies
    $\|J_1\|\leq\sqrt E$, independently of $K$.
    This is the prefix bound in the overlapping-window argument of
    \cite[Lemma~6.1(ii)]{Nachtergaele1996Spectral}.
    \uses{def:residual_window_coordinates}
\end{lemma}
\begin{proof}\leanok
    Taking the trace of $\sum_{|u|=K}(B^u)^*B^u=I$ gives the first identity.
    For a rectangular matrix $Y$, Frobenius submultiplicativity then gives
    \begin{align}
        \|J_1Y\|^2=\sum_{|u|=K}\|YB^u\|_{\mathrm F}^2
        \leq E\|Y\|_{\mathrm F}^2.
        \notag
    \end{align}
\end{proof}

\begin{lemma}[Cross-sector residual window pairing]
    \label{thm:residual_window_cross_sector_pairing}
    \lean{MPSTensor.inner_residualWindowRightMapES_leftMapES_cross_sector,
        MPSTensor.norm_inner_residualWindowRightMapES_leftMapES_le_of_prefix_bound,
        MPSTensor.norm_residualWindowRightMapES_projection_comp_leftMapES_projection_le}
    \leanok
    Let $T_i$ be the right residual window formed from $A,B_i,V_i$, and
    let $L_j$ be the left window formed from $B_j$, in the coordinates of
    Definition~\ref{def:residual_window_coordinates}.
    Their middle interval has length $M$, the outer blocked prefix has length
    $K$, and the original tail has length $r$. For rectangular boundary fibers
    $x_u$ and square boundary fibers $y_\tau$,
    \begin{align}
        \langle T_ix,L_jy\rangle
        =\sum_{u,\tau}\left\langle
        \Gamma_M^{B_i}(V_i^\dagger A^\tau x_u),
        \Gamma_M^{B_j}(y_\tau B_j^u)\right\rangle.
    \end{align}
    The two sector bond dimensions may differ, and every length may be zero.
    If $\varepsilon\geq0$ and
    $|\langle v,w\rangle|\leq\varepsilon\|v\|\|w\|$ for every
    $v\in\mathcal G_M(B_i)$ and $w\in\mathcal G_M(B_j)$, then
    \begin{align}
        |\langle x',y'\rangle|                                 & \leq\varepsilon\|x'\|\|y'\|,
                                                               &                              & x'\in\operatorname{ran}T_i,\quad y'\in\operatorname{ran}L_j, \\
        \|P_{\operatorname{ran}T_i}P_{\operatorname{ran}L_j}\| & \leq\varepsilon.
    \end{align}
    No normalization, injectivity, intertwining, or boundary inverse is assumed.
\end{lemma}
\begin{proof}\leanok
    \uses{def:residual_window_coordinates}
    Fix $(u,\tau)$ and sum over the middle configuration. Trace cyclicity
    moves the blocked prefix word to the right boundary of the left window.
    Both resulting fibers lie in the stated middle ground spaces.
    The sums of their squared norms are exactly the squared norms of the
    full window vectors. The triangle inequality and finite
    Cauchy--Schwarz therefore preserve the coefficient $\varepsilon$.
    The relative overlap estimate gives the projection-product bound.
\end{proof}

\begin{lemma}[Uniform cross-sector residual window decay]
    \label{thm:residual_window_cross_sector_decay}
    \lean{MPSTensor.IsPrimitiveMPS.eventually_norm_residualWindow_cross_projection_le_of_inequivalent,
        MPSTensor.IsPrimitiveMPS.residualWindow_cross_projection_tendsto_zero_of_inequivalent}
    \leanok
    Suppose $B_i$ and $B_j$ are inequivalent normalized primitive tensors
    with positive bond dimensions and faithful invariant matrices.
    For every $\varepsilon>0$, all sufficiently large middle lengths $M$ satisfy, for every $K,r\geq0$,
    \begin{align}
        \|P_{\operatorname{ran}T_i(K,M,r)}
        P_{\operatorname{ran}L_j(K,M,r)}\|\leq\varepsilon.
        \notag
    \end{align}
    Thus the projection product tends to zero along every diverging
    middle length, while both outer lengths may vary arbitrarily.
    The original tensor $A$ and support matrix $V_i$ require no
    normalization or intertwining hypothesis for this conclusion.
    This is the cross-sector part of the projection comparison in
    \cite[Lemmas~disjoint and commutation(ii)]{Nachtergaele1996Spectral}.
\end{lemma}
\begin{proof}\leanok
    \uses{thm:residual_window_cross_sector_pairing,thm:inequivalent_block_ground_space_overlap_decay}
    Primitive inequivalence makes the middle ground-space angle arbitrarily
    small. The preceding comparison transfers that bound uniformly over
    both outer lengths. Apply it with half the desired tolerance to obtain
    convergence along any filter with diverging middle length.
\end{proof}

\begin{lemma}[Finite sector projection estimate]
    \label{thm:finite_sector_projection_estimate}
    \lean{Submodule.norm_starProjection_comp_sub_le_sector_errors}
    \leanok
    Let $U,V,W$ and $S_i,T_i,C_i$ be subspaces of a finite-dimensional
    real or complex Hilbert space, with $i$ ranging over a finite set $I$.
    Write $P_Z$ for the orthogonal projection onto $Z$, and set
    $e_U=\|P_U-\sum_iP_{S_i}\|$,
    $e_V=\|P_V-\sum_iP_{T_i}\|$, and
    $e_W=\|\sum_iP_{C_i}-P_W\|$. Then
    \begin{align}
        \|P_UP_V-P_W\|
        \leq e_U+|I|e_V+
        \sum_{i,j\in I}\|P_{S_i}P_{T_j}-\delta_{ij}P_{C_i}\|+e_W.
    \end{align}
    No orthogonality of the sector subspaces is required.
    This finite-sum estimate assembles the projection comparisons in
    \cite[Lemmas~commutation(ii) and disjoint]{Nachtergaele1996Spectral}.
\end{lemma}
\begin{proof}\leanok
    Put $P=\sum_iP_{S_i}$, $Q=\sum_iP_{T_i}$, and $Z=\sum_iP_{C_i}$.
    Expand
    $P_UP_V-P_W=(P_U-P)P_V+P(P_V-Q)+(PQ-Z)+(Z-P_W)$.
    Every orthogonal projection has norm at most one, so $\|P\|\leq|I|$.
    Expand $PQ-Z$ as the stated double sum and apply the triangle inequality.
\end{proof}

\begin{lemma}[Vanishing sector defects give a vanishing joint projection defect]
    \label{thm:sector_projection_defect_limit}
    \lean{Submodule.tendsto_norm_starProjection_comp_sub_zero_of_sector_errors}
    \leanok
    Let the subspaces in the preceding theorem vary along an arbitrary
    filter, allowing the finite-dimensional ambient Hilbert spaces to vary.
    Suppose the three projection-sum errors tend to zero, and suppose
    $\|P_{S_i}P_{T_i}-P_{C_i}\|\to0$ for each $i$, while
    $\|P_{S_i}P_{T_j}\|\to0$ for each $i\ne j$.
    Then $\|P_UP_V-P_W\|\to0$.
    The finite index set is fixed and may be empty.
\end{lemma}
\begin{proof}\leanok
    \uses{thm:finite_sector_projection_estimate}
    The finite double sum of sector errors tends to zero.
    Apply the finite sector projection estimate and squeeze the nonnegative
    joint defect between zero and its vanishing upper bound.
\end{proof}

\begin{lemma}[Limiting residual mixed-Gram cancellation]
    \label{thm:residual_window_limiting_cancellation}
    \lean{MPSTensor.rectangularRightGramMetric_inverse_comp_correlatedTailMapES_adjoint,
        MPSTensor.inner_residualWindow_limitingMixedGramIntertwining}
    \leanok
    Let the virtual dimension $E$ be positive and let $\rho\in\MN{E}$ be
    positive definite. Put $G_\infty=\operatorname{gramReshuffle}(\mathcal P_\rho)$
    and $R_\rho Y=Y\rho/\tr\rho$. For the correlated tail map $C_r$,
    \begin{align}
        R_\rho^{-1}C_r^*G_\infty^{\oplus\tau} & =C_r^*.
        \label{eq:ph_residual_metric_cancel}
    \end{align}
    Suppose the exact tail Gram matrix is an orthogonal projection $Q$, so that
    $C_r^*C_r=\mathcal Q$ with $\mathcal QY=QY$. Set
    $K_r=C_r^*G_\infty^{\oplus\tau}C_r$ and $I_r=\mathcal Q R_\rho^{-1}$.
    With the residual window coordinates,
    \begin{align}
        \big\langle \mathbf x,
        K_r^{\oplus u}J_1I_rC_r^*G_\infty^{\oplus\tau}\mathbf y\big\rangle
         & =\sum_u\sum_\tau
        \big\langle V^*A^\tau X_u,\,G_\infty(Y_\tau B^u)\big\rangle.
        \label{eq:ph_residual_mixed_limit}
    \end{align}
    No transfer fixed-point equation, injectivity, or normalization is required
    for these identities.
    \uses{def:residual_window_coordinates,def:correlated_residual_gram_metric,
        def:residual_fixed_metric_row_projection}
\end{lemma}
\begin{proof}\leanok
    \uses{thm:residual_exact_row_support,def:residual_fixed_metric_row_projection,def:correlated_residual_gram_metric,thm:euclidean_fiberwise_operator}
    The adjoint is $C_r^*\mathbf y=\sum_\tau A^{\tau*}VY_\tau$.
    Right multiplication by $\rho/\tr\rho$ commutes with this recombination,
    and $R_\rho^{-1}$ cancels it, giving~(\ref{eq:ph_residual_metric_cancel}).
    The projected tail Gram gives $C_r\mathcal Q=C_r$ and hence
    $\mathcal Q C_r^*=C_r^*$. Thus $QZ=Z$ for
    $Z=\sum_\tau A^{\tau*}VY_\tau$.
    Now $J_1Z=(ZB^u)_u$ and $K_r(ZB^u)=ZB^u\rho/\tr\rho$.
    Expand the resulting Frobenius pairing over $\tau$ to obtain
    (\ref{eq:ph_residual_mixed_limit}).
\end{proof}

\begin{lemma}[Centered residual mixed-Gram error]
    \label{thm:residual_window_centered_gram_error}
    \lean{MPSTensor.inner_residualWindow_centeredMixedGram_eq_gramError}
    \leanok
    Under the hypotheses of Lemma~\ref{thm:residual_window_limiting_cancellation},
    define
    $D_M=T^*L_1-K_r^{\oplus u}J_1I_rC_r^*G_\infty^{\oplus\tau}$.
    Then
    \begin{align}
        \langle\mathbf x,D_M\mathbf y\rangle
         & =\sum_u\sum_\tau
        \big\langle V^*A^\tau X_u,\,
        (G_B(M)-G_\infty)(Y_\tau B^u)\big\rangle.
        \label{eq:ph_residual_centered_gram}
    \end{align}
    This is an exact finite-volume identity. A quantitative bound or a limiting
    projection estimate requires further hypotheses on the prefix Gram.
\end{lemma}
\begin{proof}\leanok
    \uses{thm:residual_window_limiting_cancellation,thm:residual_window_mixed_gram}
    Subtract the limiting mixed-Gram identity from the exact mixed-Gram
    pairing, and distribute subtraction through the two finite sums.
\end{proof}

\begin{definition}[Residual configuration bijection]
    \label{def:residual_boundary_config_equiv}
    \lean{MPSTensor.residualBoundaryConfigEquiv}
    \leanok
    Decode a blocked configuration $v$ of length $M$ into its $ML$ original
    sites and append an original tail $\tau$ of length $r$. This gives a
    bijection $e_{M,r}$ from the pairs $(v,\tau)$ to the original configurations
    of length $ML+r$, including zero lengths.
    \uses{def:residual_window_coordinates}
\end{definition}

\begin{lemma}[Right window as a finite family]
    \label{thm:residual_window_right_fiberwise}
    \lean{MPSTensor.residualWindowRightMapES_eq_fiberwise}
    \leanok
    Let $\Gamma_{B,V}(M,r)$ be the residual boundary map and let $U_{M,r}$
    be physical reindexing through $e_{M,r}$. The right window map is the
    finite direct sum of $U_{M,r}\Gamma_{B,V}(M,r)$ over the blocked prefix
    configurations $u$ of length $K$.
    No normalization or intertwining hypothesis is required.
    \uses{def:residual_boundary_config_equiv,def:blocked_residual_boundary_map,
        def:residual_window_coordinates}
\end{lemma}
\begin{proof}\leanok
    Both maps have coefficient $\tr(B^vV^*A^\tau X_u)$ at $(u,v,\tau)$.
\end{proof}

\begin{lemma}[Right residual sector projection approximation]
    \label{thm:residual_window_right_sector_sum}
    \lean{MPSTensor.tendsto_norm_isometric_residual_sector_projection_sum_sub_iSup_zero,
        MPSTensor.tendsto_norm_residualWindowRight_sector_projection_sum_sub_iSup_zero}
    \leanok
    Let $(B_j)_j$ be a finite family of pairwise gauge-phase inequivalent
    primitive tensors with faithful invariant densities, and let $V_j$ be
    arbitrary rectangular embeddings. Write $S_j(K,M,r)$ for the range of the
    right window map of sector $j$. As $M\to\infty$,
    \begin{align}
        \left\|\sum_jP_{S_j(K,M,r)}-P_{\sum_jS_j(K,M,r)}\right\|\longrightarrow0,
        \notag
    \end{align}
    uniformly in arbitrary varying prefix lengths $K$ and original-tail
    lengths $r$. The assertion holds along any filter on which $M$ tends to
    infinity. Residual-sector projection-sum decay also persists under any
    varying physical configuration isometry before adding the spectator.
    This is the right-window comparison in
    \cite[Lemma~6.1(i)]{Nachtergaele1996Spectral}.
\end{lemma}
\begin{proof}\leanok
    \uses{thm:residual_window_right_fiberwise}
    Physical isometries preserve the relative angles between residual sectors.
    The resulting projection-sum error tends to zero by primitive sector
    separation. The right window acts independently on each prefix fiber,
    and for such fiberwise maps the error between the sum of the sector
    range projections and the joint range projection is at most the error
    of a single fiber. Hence adding the prefix spectator does not increase
    the operator-norm error.
\end{proof}

\begin{lemma}[Uniform upper bounds from all selections]
    \label{thm:uniform_upper_bound_all_selections}
    \lean{Filter.eventually_forall_le_of_tendsto_all_selections}
    \leanok
    Let $X$ be a nonempty set, let $\mathcal F$ be a filter on an index set
    $I$, and let $e:I\times X\to\mathbb R$.
    Suppose $e(i,x_i)\to0$ along $\mathcal F$ for every function
    $i\mapsto x_i$ from $I$ to $X$. Then, for every $\varepsilon>0$,
    eventually $e(i,x)\leq\varepsilon$ simultaneously for every $x\in X$.
    No topology or countability assumption on $X$ is required, and $e$
    need not be nonnegative. This elementary implication applies to the
    uniform projection estimates in
    \cite[Lemma~commutation(ii)]{Nachtergaele1996Spectral}.
\end{lemma}
\begin{proof}\leanok
    For each $i$, select a parameter violating the upper bound if one exists;
    otherwise select any element of $X$.
    Convergence along this selection gives $e(i,x_i)<\varepsilon$ eventually.
    The choice of $x_i$ then implies $e(i,x)\leq\varepsilon$ for every $x$.
\end{proof}

\begin{lemma}[Orthogonal projections on finite Euclidean families]
    \label{thm:euclidean_fiberwise_projection}
    \lean{EuclideanSpace.mem_range_fiberwiseMap_iff,
        EuclideanSpace.mem_orthogonal_range_fiberwiseMap_iff,
        EuclideanSpace.fiberwiseMap_range_starProjection,
        EuclideanSpace.iSup_range_fiberwiseMap_eq,
        EuclideanSpace.fiberwiseMap_iSup_range_starProjection,
        EuclideanSpace.fiberwiseMap_sum,
        EuclideanSpace.norm_sum_fiberwiseMap_range_starProjection_sub_iSup_le}
    \leanok
    \uses{def:euclidean_fiberwise_operator}
    Let $S,I,J$ be finite coordinate sets and let $T\colon\C^I\to\C^J$.
    Write $\mathcal F_S(T)$ for the operator applying $T$ to every fiber
    of $\C^{S\times I}$. Then
    \begin{align}
        x\in\operatorname{ran}\mathcal F_S(T)
         & \Longleftrightarrow \forall s,\ x_s\in\operatorname{ran}T,         \\
        x\in(\operatorname{ran}\mathcal F_S(T))^\perp
         & \Longleftrightarrow \forall s,\ x_s\in(\operatorname{ran}T)^\perp, \\
        P_{\operatorname{ran}\mathcal F_S(T)}
         & =\mathcal F_S(P_{\operatorname{ran}T}).
    \end{align}
    For any sector family $T_j\colon\C^{I_j}\to\C^J$ with varying finite
    domains and common finite target, the joint-range projection satisfies
    \begin{align}
        P_{\sum_j\operatorname{ran}\mathcal F_S(T_j)}
         & =\mathcal F_S(P_{\sum_j\operatorname{ran}T_j}).
    \end{align}
    The sector index in this identity may be infinite, with each sum
    interpreted as the linear span. For a finite sector family,
    $\mathcal F_S$ preserves operator sums and
    \begin{align}
        \left\|\sum_jP_{\operatorname{ran}\mathcal F_S(T_j)}
        -P_{\sum_j\operatorname{ran}\mathcal F_S(T_j)}\right\|
        \leq
        \left\|\sum_jP_{\operatorname{ran}T_j}
        -P_{\sum_j\operatorname{ran}T_j}\right\|.
    \end{align}
    Empty spectator sets and empty sector families are included. These
    finite-dimensional identities give the spectator projection conjugation
    used in \cite[Lemma~commutation(i)]{Nachtergaele1996Spectral}.
\end{lemma}
\begin{proof}\leanok
    \uses{thm:euclidean_fiberwise_operator}
    Choose a preimage independently in each fiber to characterize the range.
    The adjoint acts fiberwise, which gives the orthogonal-range criterion.
    The orthogonal projection is characterized by its range and orthogonal
    complement. For a joint span, orthogonality to every sector and membership
    of every fiber commute as universal quantifiers. Taking orthogonal
    complements again proves the joint-range identity. The final estimate
    follows from linearity of $\mathcal F_S$ and its norm bound.
\end{proof}

\begin{lemma}[Open gap from a residual cover limit]
    \label{thm:open_gap_residue_cover_limit}
    \lean{MPSTensor.IsPeriodic.exists_pos_openParentHamiltonianES_gap_of_residue_cover_limit}
    \leanok
    Let $A$ be a periodic tensor of period $m>0$ on a nonzero physical space.
    For $K,M,r\geq0$, let $P$ denote the ground-space projection on
    $Km+Mm+r$ sites, let $P_{\mathrm L}$ project onto the prefix ground
    space on $Km+Mm$ sites tensored with the free $r$ sites, and let
    $P_{\mathrm R}$ project onto the terminal ground space on $Mm+r$ sites
    tensored with the free $Km$ sites.
    Suppose that, for every pair of sequences $K_M,r_M$,
    $\|P-P_{\mathrm L}P_{\mathrm R}\|\to0$ as $M\to\infty$.
    Then there exist $R>0$, $N_0\geq R$, and $\delta>0$ such that every
    $N\geq N_0$ and every $v\perp\ker H_{A,R,N}$ satisfy
    $\delta\|v\|\leq\|H_{A,R,N}v\|$, where $H_{A,R,N}$ is the
    canonical open parent Hamiltonian.
    At the same range, $\ker H_{A,R,N}=G_A(N)$ for every $N\geq R$.
    The convergence assumption is explicit; this is an auxiliary limit
    criterion for the projection comparison in
    \cite[Section~6, Lemma~commutation(ii)]{Nachtergaele1996Spectral}.
\end{lemma}
\begin{proof}\leanok
    \uses{thm:uniform_upper_bound_all_selections,thm:open_residue_cover_gap_criterion,thm:periodic_original_parent_interaction}
    Periodicity gives a range $R$ for which every canonical open kernel at
    volume $N\geq R$ is the original boundary space.
    Apply the selection uniformization to the parameter $(K,r)$ and tolerance
    $1/2$. Choose one $M\geq R$ for which all these defects are at most
    $1/2$. Since $m\geq1$, the overlap $Mm$ contains every interaction window.
    The residue cover criterion with prefix threshold zero gives a common
    positive gap for every $N\geq Mm$. Take $N_0=Mm$.
\end{proof}

\begin{definition}[Left-window coordinate bijection]
    \label{def:residual_window_left_config_equiv}
    \lean{MPSTensor.residualWindowLeftConfigEquiv}
    \leanok
    For lengths $K,M,r\geq0$, the left-window coordinate bijection sends a
    pair $(\tau,\sigma)$ of an original tail $\tau$ of length $r$ and a
    blocked configuration $\sigma=u\mathbin{\frown}v$ of length $K+M$ to
    $(u,(v,\tau))$.
\end{definition}

\begin{lemma}[The aligned left window as a finite family of prefix boundaries]
    \label{thm:residual_window_left_fiberwise_coordinates}
    \lean{MPSTensor.residualWindowLeftMapES_eq_fiberwise}
    \leanok
    Let $L_B(K,M,r)$ denote the left residual window of
    Definition~\ref{def:residual_window_coordinates}.
    The coordinate bijection of
    Definition~\ref{def:residual_window_left_config_equiv},
    \begin{align}
        (\tau,\sigma)\longmapsto(u,(v,\tau)),\qquad
        \sigma=u\mathbin{\frown}v,
        \notag
    \end{align}
    splits the blocked configuration into lengths $K,M$ and moves the
    original tail to the last coordinate. If $U_{K,M,r}$ is the associated
    unitary permutation, then
    \begin{align}
        L_B(K,M,r)=U_{K,M,r}\mathcal F_{[d]^r}(\Gamma_B(K+M)).
        \notag
    \end{align}
    All lengths may be zero. No normalization or injectivity is assumed.
    \uses{def:residual_window_coordinates,def:euclidean_fiberwise_operator,
        def:residual_window_left_config_equiv}
\end{lemma}
\begin{proof}\leanok
    The coefficient of the left side is $\tr(B^uB^v y_\tau)$.
    Concatenating the blocked words gives the same coefficient for the
    ordinary prefix boundary map applied to the tail fiber $y_\tau$.
\end{proof}

\begin{lemma}[Uniform left-window sector projection-sum approximation]
    \label{thm:residual_window_left_sector_sum}
    \lean{MPSTensor.norm_residualWindowLeft_sector_projection_sum_sub_iSup_le,
        MPSTensor.eventually_norm_residualWindowLeft_sector_projection_sum_sub_iSup_le,
        MPSTensor.tendsto_norm_residualWindowLeft_sector_projection_sum_sub_iSup_zero}
    \leanok
    For every finite sector family $(B_j)_j$ with possibly different bond
    dimensions,
    \begin{align}
         & \left\|\sum_jP_{\operatorname{ran}L_{B_j}(K,M,r)}
        -P_{\sum_j\operatorname{ran}L_{B_j}(K,M,r)}\right\|  \\
         & \qquad\leq
        \left\|\sum_jP_{\mathcal G_{K+M}(B_j)}
        -P_{\sum_j\mathcal G_{K+M}(B_j)}\right\|.
    \end{align}
    This comparison includes empty families and requires no normalization.
    Suppose in addition that the sectors are pairwise inequivalent normalized
    primitive tensors with positive bond dimensions and faithful invariant
    matrices. Then for every $\varepsilon>0$, all sufficiently large $M$ satisfy, for every $K,r\geq0$,
    \begin{align}
        \left\|\sum_jP_{\operatorname{ran}L_{B_j}(K,M,r)}
        -P_{\sum_j\operatorname{ran}L_{B_j}(K,M,r)}\right\|\leq\varepsilon.
        \notag
    \end{align}
    Consequently the left-window projection-sum error tends to zero along
    every diverging middle interval, with arbitrary outer-prefix and
    original-tail lengths. This is the aligned left-sector comparison for
    \cite[Lemma~commutation(i)]{Nachtergaele1996Spectral}.
\end{lemma}
\begin{proof}\leanok
    \uses{thm:residual_window_left_fiberwise_coordinates,thm:euclidean_fiberwise_projection,thm:sector_projection_sum_error_decay,thm:inequivalent_block_ground_space_overlap_decay}
    A unitary permutation conjugates orthogonal projections and preserves
    the operator norm. The finite-family projection identity then bounds
    the spectator projection-sum error by the ordinary prefix error.
    Primitive inequivalence makes that prefix error arbitrarily small at
    all sufficiently large lengths. Since $K+M\geq M$, the threshold is
    uniform in $K$, and the comparison is independent of $r$.
\end{proof}

\begin{definition}[Residual overlap factors]
    \label{def:residual_window_overlap_factors}
    \lean{MPSTensor.residualWindowLeftOverlapFactorES,
        MPSTensor.residualWindowRightOverlapFactorES}
    \leanok
    The two maps to the common family of square virtual boundaries have fibers
    $\mathcal A\mathbf x=(V^*A^\tau X_u)_{u,\tau}$ and
    $\mathcal B\mathbf y=(Y_\tau B^u)_{u,\tau}$.
    They are finite direct sums followed by isometric coordinate rearrangements.
    \uses{def:residual_window_coordinates,def:correlated_residual_gram_metric,
        def:euclidean_fiberwise_operator}
\end{definition}

\begin{lemma}[Residual overlap factor estimates]
    \label{thm:residual_window_overlap_factors}
    \lean{MPSTensor.boundaryFamilyEquiv_residualWindowLeftOverlapFactorES_apply,
        MPSTensor.boundaryFamilyEquiv_residualWindowRightOverlapFactorES_apply,
        MPSTensor.norm_residualWindowLeftOverlapFactorES_le,
        MPSTensor.norm_residualWindowRightOverlapFactorES_le,
        MPSTensor.residualWindow_centeredMixedGram_eq_overlapFactors,
        MPSTensor.norm_residualWindow_centeredMixedGram_le}
    \leanok
    \uses{def:residual_window_overlap_factors}
    The fibers are the matrices in Definition~\ref{def:residual_window_overlap_factors},
    and $\|\mathcal A\|\le\|C_r\|$, $\|\mathcal B\|\le\|J_1\|$.
    Under the hypotheses of Lemma~\ref{thm:residual_window_limiting_cancellation},
    \begin{align}
        D_M=\mathcal A^*(G_B(M)-G_\infty)^{\oplus(u,\tau)}\mathcal B,
        \qquad
        \|D_M\|\le\|C_r\|\,\|J_1\|\,\|G_B(M)-G_\infty\|.
        \notag
    \end{align}
    These are finite-dimensional estimates underlying the mixed-pairing step
    of Nachtergaele's commutation lemma~\cite{Nachtergaele1996Spectral}; no transfer
    normalization or injectivity is assumed here.
\end{lemma}
\begin{proof}\leanok
    \uses{thm:euclidean_fiberwise_operator,thm:residual_window_centered_gram_error}
    The coordinate rearrangements preserve norms, and finite direct sums do not
    increase operator norms. The fiber formulas identify the centered pairing
    with $\mathcal A^*(G_B(M)-G_\infty)^{\oplus(u,\tau)}\mathcal B$.
    The norm bound follows from submultiplicativity.
\end{proof}

\begin{lemma}[Full residual window sector projection sum]
    \label{thm:residual_window_full_sector_sum}
    \lean{MPSTensor.tendsto_norm_residualWindow_sector_projection_sum_sub_iSup_zero}
    \leanok
    Let $B_j$ be a fixed finite family of pairwise inequivalent normalized
    primitive tensors with positive bond dimensions and faithful invariant
    matrices. Let $A$ and the rectangular matrices $V_j$ be arbitrary.
    For the full residual window maps $C_j(K,M,r)$, if $M$ tends to infinity
    along an arbitrary filter, then
    \begin{align}
        \left\|\sum_jP_{\operatorname{ran}C_j(K,M,r)}
        -P_{\sum_j\operatorname{ran}C_j(K,M,r)}\right\|\longrightarrow0.
    \end{align}
    The prefix $K$ and tail $r$ may vary arbitrarily.
    No normalization or intertwining is required for $A$ or $V_j$.
    This is the full-window sector sum in the projection comparison of
    \cite[Lemma~commutation(i)]{Nachtergaele1996Spectral}.
\end{lemma}
\begin{proof}\leanok
    \uses{thm:residual_window_factorizations,thm:residual_window_right_sector_sum}
    The physical configuration isometry identifies $C_j(K,M,r)$ with
    the residual boundary map at prefix length $K+M$ and tail length $r$.
    Since $K+M\geq M$, the total prefix tends to infinity.
    Apply the isometric residual-sector projection-sum limit.
\end{proof}

\begin{lemma}[Joint full residual window range]
    \label{thm:residual_window_full_joint_range}
    \lean{MPSTensor.iSup_range_residualWindowMapES_eq_reindex_groundSpaceES}
    \leanok
    Suppose $\sum_jV_jV_j^*=1$ and
    $V_j^*A^{(L)}_i=(B_j)_iV_j^*$.
    Let $U$ be the physical configuration isometry from the original
    $(K+M)L+r$ sites to the three-window coordinates. Then
    $\sum_j\operatorname{ran}C_j(K,M,r)=U G_A((K+M)L+r)$.
    This includes zero lengths and requires no primitivity or injectivity.
\end{lemma}
\begin{proof}\leanok
    \uses{thm:residual_window_factorizations,thm:blocked_residual_ground_space_resolution}
    Carry the exact residual-sector ground-space decomposition by $U$.
    A linear map preserves sums of subspaces, and the range of a composite
    is the image of the original range.
\end{proof}

\begin{lemma}[Full residual window Gram]
    \label{thm:residual_full_window_gram}
    \lean{MPSTensor.residualWindowMapES_adjoint_comp_self}
    \leanok
    Physical reindexing identifies the full window map $C$ with the residual
    boundary map on $K+M$ blocked sites and an original tail of length $r$.
    Consequently,
    \begin{align}
        C^*C=\Gamma_{K+M,r}^*\Gamma_{K+M,r}.
    \end{align}
    This holds for all lengths, including zero, without injectivity or
    normalization hypotheses.

\end{lemma}
\begin{proof}\leanok
    \uses{thm:residual_window_factorizations}
    The configuration reindexing is an isometry. In $C^*C$, its adjoint times
    itself is the identity, leaving the residual Gram.
\end{proof}

\begin{definition}[Finite residual mixed Gram]
    \label{def:residual_window_mixed_gram_residual}
    \lean{MPSTensor.residualWindowMixedGramResidualES}
    \leanok
    For the right, left, and full window maps $T,L_1,C$ and the two virtual
    factors $J_1,C_r$, set
    \begin{align}
        \mathcal D_{K,M,r}
        =T^*L_1-(T^*T)J_1H_{K+M,r}C_r^*(L_1^*L_1),
    \end{align}
    where $H_{K+M,r}$ is the supported full-boundary Gram inverse.
    \uses{def:residual_window_coordinates,def:residual_extended_gram_supported_inverse}
\end{definition}

\begin{lemma}[Primitive residual mixed-Gram decay]
    \label{thm:residual_window_mixed_gram_decay}
    \lean{MPSTensor.IsPrimitiveMPS.residualWindowMixedGramResidual_tendsto_zero}
    \leanok
    Let $D,E>0$, and let the blocked tensor $B$ be primitive with a faithful
    invariant matrix $\rho$. Suppose the correlated tail Gram is an orthogonal
    projection $Q_r$ for every tail length $r$.
    Along any filter with $M\to\infty$, and with arbitrary $K,r$,
    \begin{align}
        \|\mathcal D_{K,M,r}\|\longrightarrow0.
    \end{align}
    This is the same-sector mixed-Gram estimate used in commutation (ii),
    equation (3.15), of~\cite{Nachtergaele1996Spectral}; the algebraic
    row-support hypotheses are stated explicitly here.
    \uses{def:residual_window_mixed_gram_residual}
\end{lemma}
\begin{proof}\leanok
    \uses{thm:residual_window_overlap_factors,
        thm:residual_supported_inverse_uniform_tail,
        thm:residual_window_gram_errors,
        thm:residual_window_virtual_norm,
        thm:three_operator_perturbation_limit}
    The centered limiting mixed Gram is bounded by the common-prefix Gram error
    and hence tends to zero. Replace successively the right-window Gram, the
    supported full-boundary inverse, and the left-window Gram by their limits.
    Their errors tend to zero, since $K+M\ge M$, and the remaining products have
    uniform bounds. Apply the three-operator perturbation estimate and then
    the triangle inequality.
\end{proof}

\begin{lemma}[Three-operator perturbation estimate]
    \label{thm:three_operator_perturbation}
    \lean{ContinuousLinearMap.norm_comp_three_sub_le}
    \leanok
    For bounded linear operators $A,a:G\to W$, $B,b:F\to G$, and
    $C,c:E\to F$ between normed spaces,
    \begin{align}
        \|ABC-abc\|\le
        \|A-a\|\|B\|\|C\|+
        \|a\|\|B-b\|\|C\|+
        \|a\|\|b\|\|C-c\|.
    \end{align}
    No finite-dimensional or completeness hypothesis is required.
\end{lemma}
\begin{proof}\leanok
    Use $ABC-abc=(A-a)BC+a(B-b)C+ab(C-c)$, the triangle inequality,
    and submultiplicativity of the operator norm.
\end{proof}

\begin{lemma}[Vanishing three-operator perturbations]
    \label{thm:three_operator_perturbation_limit}
    \lean{ContinuousLinearMap.tendsto_norm_comp_three_sub_zero_of_bounds}
    \leanok
    Let the normed spaces and operators in Lemma~\ref{thm:three_operator_perturbation}
    vary along an arbitrary filter. Suppose the three individual differences
    tend to zero in norm, and the products $\|B\|\|C\|$, $\|a\|\|C\|$,
    and $\|a\|\|b\|$ have a common eventual upper bound.
    Then $\|ABC-abc\|\to0$.

\end{lemma}
\begin{proof}\leanok
    \uses{thm:three_operator_perturbation}
    Bound the product difference by the common bound times the sum of the
    three individual norm differences, and apply the squeeze theorem.
\end{proof}

\begin{definition}[Physical cover coordinates]
    \label{def:residual_cover_coordinates}
    \lean{MPSTensor.residualCoverConfigEquiv,MPSTensor.residualCoverConfigEquiv_apply}
    \leanok
    Decode the blocked prefix $u$ of length $K$ and blocked middle $v$
    of length $M$ into configurations of lengths $KL$ and $ML$.
    Concatenation with the original tail $\tau$ gives a configuration on
    $KL+ML+r$ sites. This bijection is the full residual-window configuration
    bijection followed by the identification $(K+M)L=KL+ML$, and it sends
    $(u,v,\tau)$ to the concatenation of the three configurations in that order.
    All zero lengths are included.
    These are the physical covering coordinates of
    \cite[Section~6]{Nachtergaele1996Spectral}.
    \uses{def:residual_window_coordinates}
\end{definition}

\begin{lemma}[Joint right residual window range]
    \label{thm:residual_window_right_joint_range}
    \lean{MPSTensor.iSup_range_residualWindowRightMapES_eq_reindex_reassocTailBoundaryMapES}
    \leanok
    Let $A$, $B_j$, and $V_j$ satisfy
    $\sum_jV_jV_j^*=1$ and $V_j^*A^{(L)}_i=(B_j)_iV_j^*$.
    Let $U$ be the configuration isometry from the original $KL+ML+r$
    sites to the physical cover coordinates. Then
    \begin{align}
        \sum_j\operatorname{ran}T_j(K,M,r)
        =U\operatorname{ran}\Gamma^{\mathrm{tail}}_{KL,ML,r}.
    \end{align}
    Here the terminal boundary map leaves the first $KL$ sites free and
    places the original ground space on the last $ML+r$ sites.
    No primitivity, injectivity, or isometry of the virtual embeddings is required.
    This is the right joint-window identification in the projection comparison of
    \cite[Lemma~commutation(ii)]{Nachtergaele1996Spectral}.
\end{lemma}
\begin{proof}\leanok
    \uses{def:residual_cover_coordinates,thm:residual_window_right_fiberwise,thm:blocked_residual_ground_space_resolution}
    Each right residual window is the direct sum, over prefix configurations,
    of its physically reindexed residual boundary map.
    The exact residual-sector decomposition identifies the joint fiber range
    with the original ground space on $ML+r$ sites.
    Taking finite direct sums preserves the joint range.
    In the physical cover coordinates, the terminal boundary map is the same
    direct sum of ordinary boundary maps, after a bijection of its prefix
    boundary coordinates. That bijection is surjective, so it does not change
    the range.
\end{proof}

\begin{lemma}[Original ground space of the aligned left residual window]
    \label{thm:residual_window_left_original_ground_space}
    \lean{MPSTensor.residualWindowLeftMapES_blockTensor_eq_reindex_leftBoundaryMapES,
        MPSTensor.iSup_range_residualWindowLeftMapES_eq_of_groundSpaceES_eq_iSup,
        MPSTensor.iSup_range_residualWindowLeftMapES_eq_reindex_leftBoundaryMapES}
    \leanok
    Let $A$ be an MPS tensor and let $U_{K,M,r}$ be the unitary from original
    configurations to the three-window coordinates of
    Definition~\ref{def:residual_window_coordinates}. Then
    \begin{align}
        L_{\operatorname{block}_L A}(K,M,r)
        =U_{K,M,r}\Gamma^{\mathrm{left},A}_{(K+M)L,r}.
    \end{align}
    Suppose a family $(B_j)_j$ presents the blocked prefix ground space exactly:
    \begin{align}
        \mathcal G_{K+M}(\operatorname{block}_L A)
        =\bigvee_j\mathcal G_{K+M}(B_j).
    \end{align}
    Then its joint left-window range is the original aligned prefix ground space,
    with the tail unconstrained:
    \begin{align}
        \bigvee_j\operatorname{ran}L_{B_j}(K,M,r)
        =U_{K,M,r}\operatorname{ran}\Gamma^{\mathrm{left},A}_{(K+M)L,r}.
    \end{align}
    In particular, this equality follows for a finite family with an isometric
    virtual resolution satisfying
    \begin{align}
        V_j^\dagger V_j=\Id_{D_j},\qquad
        \sum_jV_jV_j^\dagger=\Id_D,\qquad
        (\operatorname{block}_L A)^iV_j=V_jB_j^i.
    \end{align}
    All lengths may be zero. No primitivity, normalization, or injectivity is
    assumed. This is the exact left-window identification used in the
    three-interval comparison of
    \cite[Lemma~commutation(ii)]{Nachtergaele1996Spectral}.
    \uses{def:residual_window_coordinates}
\end{lemma}
\begin{proof}\leanok
    \uses{thm:residual_window_left_fiberwise_coordinates,
        thm:euclidean_fiberwise_projection,
        thm:isometric_sector_ground_space_resolution}
    Concatenating the two blocked words and decoding their physical letters
    identifies the blocked left map with the original prefix boundary map.
    Each tail configuration contributes an independent copy of the prefix
    space. The joint span of these finite families is therefore obtained by
    applying the joint prefix-space projection independently on each tail
    fiber. The isometric virtual resolution gives the stated prefix
    presentation by trace cyclicity.
\end{proof}

\begin{lemma}[Full joint range in physical cover coordinates]
    \label{thm:residual_window_full_joint_range_cover}
    \lean{MPSTensor.iSup_range_residualWindowMapES_eq_reindex_groundSpaceES_cover}
    \leanok
    Under the virtual resolution and blocked cointertwining hypotheses of
    Lemma~\ref{thm:residual_window_full_joint_range}, let $U$ be the
    configuration isometry from the original $KL+ML+r$ sites to the
    physical cover coordinates. Then
    $\sum_j\operatorname{ran}C_j(K,M,r)=U G_A(KL+ML+r)$.
    The statement includes zero lengths and requires no primitivity or injectivity.
\end{lemma}
\begin{proof}\leanok
    \uses{thm:residual_window_full_joint_range,def:residual_cover_coordinates}
    Apply the exact full-window joint-range identity and carry it through
    $(K+M)L+r=KL+ML+r$. The cover configuration bijection is precisely the
    full-window bijection followed by this arithmetic identification.
\end{proof}

\begin{lemma}[Isometric conjugation of a projection defect]
    \label{thm:isometric_projection_defect_transport}
    \lean{Submodule.norm_map_starProjection_comp_sub_eq}
    \leanok
    Let $U,V,W$ be closed subspaces of a Hilbert space and let $F$ be a
    surjective linear isometry to another Hilbert space. Then
    \begin{align}
        \|P_{FU}P_{FV}-P_{FW}\|=\|P_W-P_VP_U\|.
    \end{align}
\end{lemma}
\begin{proof}\leanok
    Orthogonal projections are conjugated by $F$, and conjugation by a
    surjective isometry preserves the operator norm. Taking the adjoint
    reverses the order of the two projections and preserves the norm.
\end{proof}

\begin{lemma}[The joint left window in the original covering intervals]
    \label{thm:residual_window_left_joint_range_cover}
    \lean{MPSTensor.iSup_range_residualWindowLeftMapES_eq_reindex_cover_leftBoundaryMapES}
    \leanok
    Under the isometric virtual resolution and blocked letter intertwining of
    Lemma~\ref{thm:residual_window_left_original_ground_space}, let
    $U^{\mathrm{cover}}_{K,M,r}$ be the unitary from original configurations
    of length $KL+ML+r$ to the three-window coordinates of
    Definition~\ref{def:residual_cover_coordinates}. Then
    \begin{align}
        \bigvee_j\operatorname{ran}L_{B_j}(K,M,r)
        =U^{\mathrm{cover}}_{K,M,r}\operatorname{ran}
        \Gamma^{\mathrm{left},A}_{KL+ML,r}.
    \end{align}
    All lengths may be zero. This is the same physical coordinate isometry
    used for the right and full windows; no additional tensor hypothesis
    is imposed. The equality gives the aligned prefix space in
    \cite[Lemma~commutation(ii)]{Nachtergaele1996Spectral}.
    \uses{def:residual_cover_coordinates}
\end{lemma}
\begin{proof}\leanok
    \uses{thm:residual_window_left_original_ground_space}
    Carry the previous left-window equality along the arithmetic
    identification $(K+M)L=KL+ML$. By its definition, the cover bijection is
    exactly the previous configuration bijection followed by this identification.
\end{proof}

\begin{lemma}[Supported residual window cancellation]
    \label{thm:residual_window_supported_projector_cancellation}
    \lean{MPSTensor.residualWindow_supported_projector_factorization}
    \leanok
    Let $T,L_1,C$ denote the right, left, and full window maps. Suppose that
    $Q_r$ is the exact correlated-tail Gram projection, that the right and full
    supported Gram inverses satisfy $G_{M,r}H_{M,r}=Q_r$ and
    $G_{K+M,r}H_{K+M,r}=Q_r$, with $H_{M,r}$ self-adjoint, and that the ordinary
    Gram $G_B(K+M)$ is invertible. Then
    \begin{align}
        P_{\operatorname{ran}T}P_{\operatorname{ran}L_1}-P_{\operatorname{ran}C}
        =T\widehat H_{M,r}\mathcal D_{K,M,r}
        \widehat {G_B(K+M)^{-1}}L_1^*.
    \end{align}
    Here the hats denote the direct sums over the corresponding exterior
    configurations. This is the supported-boundary cancellation in
    commutation (ii), equation (3.15), of~\cite{Nachtergaele1996Spectral}.
    \uses{def:residual_window_coordinates,def:residual_window_mixed_gram_residual}
\end{lemma}
\begin{proof}\leanok
    \uses{thm:residual_full_window_gram,
        thm:supported_gram_projector_cancellation,
        thm:residual_window_grams}
    The right Gram and its supported inverse act independently on prefix
    configurations. The left Gram and its inverse act independently on tail
    configurations. Both factorizations of the full boundary map have the exact
    row support $Q_r$. Substitute these identities into the supported
    range-projection cancellation formula.
\end{proof}

\begin{lemma}[Primitive residual window projection decay]
    \label{thm:residual_window_primitive_projection_decay}
    \lean{MPSTensor.IsPrimitiveMPS.residualWindow_projection_comp_sub_tendsto_zero}
    \leanok
    Let $D,E>0$, and let $B$ be primitive with a faithful invariant matrix
    $\rho$. Suppose the exact correlated-tail Gram is an orthogonal projection
    $Q_r$ for every $r$. Along any filter with $M\to\infty$ and arbitrary
    exterior lengths $K,r$,
    \begin{align}
        \|P_{\operatorname{ran}T}P_{\operatorname{ran}L_1}
        -P_{\operatorname{ran}C}\|\longrightarrow0.
    \end{align}
    This is the same-sector part of commutation (ii)
    in~\cite{Nachtergaele1996Spectral}; the tail support is retained explicitly.

\end{lemma}
\begin{proof}\leanok
    \uses{thm:residual_window_supported_projector_cancellation,
        thm:residual_window_mixed_gram_decay,
        thm:residual_supported_inverse_uniform_tail,
        thm:supported_gram_projector_norm}
    The primitive supported-inverse estimates give the inverse identities
    and uniform bounds at lengths $M$ and $K+M$. The ordinary primitive Gram
    inverse is uniformly bounded at length $K+M$. The window maps are uniformly
    bounded by their boundary-map norms. Apply the norm estimate for the
    supported cancellation formula to the vanishing mixed residual.
\end{proof}

\begin{lemma}[Joint residual window projection decay]
    \label{thm:residual_window_joint_projection_decay}
    \lean{MPSTensor.tendsto_norm_residualWindow_joint_projection_comp_sub_zero}
    \leanok
    Let $D>0$, and let $(B_j)_{j\in J}$ be a finite family of positive-dimensional
    primitive tensors with faithful invariant matrices, pairwise inequivalent
    up to similarity and phase. Let $T_j,L_j,C_j$ be the right, left, and full
    window maps arising from an original tensor and rectangular embeddings $V_j$.
    Suppose that each correlated-tail Gram is an orthogonal projection at every
    tail length. Write $\mathcal T=\sum_j\operatorname{ran}T_j$,
    $\mathcal L=\sum_j\operatorname{ran}L_j$ and
    $\mathcal C=\sum_j\operatorname{ran}C_j$. Along any filter with $M\to\infty$
    and arbitrary exterior lengths $K,r$,
    \begin{align}
        \|P_{\mathcal T}P_{\mathcal L}-P_{\mathcal C}\|\longrightarrow0.
    \end{align}

\end{lemma}
\begin{proof}\leanok
    \uses{thm:residual_window_primitive_projection_decay,
        thm:residual_window_cross_sector_decay,
        thm:residual_window_left_sector_sum,
        thm:residual_window_right_sector_sum,
        thm:residual_window_full_sector_sum,
        thm:sector_projection_defect_limit}
    Replace each joint projection by the sum of its sector projections.
    The three replacement errors vanish. The paired sector defects vanish
    by the supported same-sector cancellation, and the different-sector
    projection products vanish by primitive inequivalence. The finite-sector
    projection estimate combines these limits.
\end{proof}

\begin{lemma}[Periodic residual window projection decay]
    \label{thm:periodic_residual_window_joint_projection_decay}
    \lean{MPSTensor.IsPeriodic.exists_residualWindow_joint_projection_comp_sub_tendsto_zero}
    \leanok
    Let $A$ be a normalized periodic tensor of period $m$. There exist
    positive-dimensional blocked sectors $B_j$, indexed by $j\in\mathbb Z/m\mathbb Z$,
    and rectangular isometries $V_j$ such that
    \begin{align}
        \sum_jV_jV_j^*=I,\qquad A^{(m)}_iV_j=V_j(B_j)_i,
        \qquad V_j^*A^{(m)}_i=(B_j)_iV_j^*,
    \end{align}
    and, for every pair of sequences of exterior lengths $K(M),r(M)$,
    \begin{align}
        \|P_{\mathcal T_{K(M),M,r(M)}}P_{\mathcal L_{K(M),M,r(M)}}
        -P_{\mathcal C_{K(M),M,r(M)}}\|\longrightarrow0.
    \end{align}
    The spaces are the joint residual window ranges in the triple configuration
    coordinates. This establishes their projection comparison from periodicity
    alone; identifying them with original physical interval spaces is a separate
    isometric conjugation. The estimate is the residual-coordinate form of
    commutation (ii) in~\cite{Nachtergaele1996Spectral}.
    \uses{def:residual_window_coordinates}
\end{lemma}
\begin{proof}\leanok
    \uses{thm:cyclic_primitive_sector_resolution,
        thm:residual_window_joint_projection_decay,
        thm:correlated_residual_cyclic_corner}
    Choose the faithful primitive, inequivalent cyclic sector resolution.
    Its cyclic shift is $P_{j+1}A_i=A_iP_j$. Reversing the labels permits the
    word shift formula to show that the correlated-tail Gram is precisely
    $P_{j-r}$, an orthogonal projection. Apply joint residual window projection
    decay and retain the two intertwiner identities of the same resolution.
\end{proof}

\begin{lemma}[Physical cover projection decay for a periodic tensor]
    \label{thm:periodic_physical_cover_projection_decay}
    \lean{MPSTensor.IsPeriodic.residual_cover_projection_defect_tendsto_zero}
    \leanok

    Let $A$ be a normalized periodic tensor of period $m$. For arbitrary
    sequences $K_M,r_M\in\mathbb N$, put $N_M=K_Mm+Mm+r_M$.
    Let $P_{\mathrm L,M}$ and $P_{\mathrm R,M}$ be the ground-space
    projections for the prefix of length $(K_M+M)m$ and the suffix of length
    $Mm+r_M$, extended to the full interval. Then
    \begin{align}
        \|P_{\mathcal G_A(N_M)}-P_{\mathrm L,M}P_{\mathrm R,M}\|
        \longrightarrow 0 \qquad (M\longrightarrow\infty).
    \end{align}
    This is the explicit periodic tensor form of the cover estimate used in
    \cite[Lemma commutation (ii) and Section 6]{Nachtergaele1996Spectral}.
    The passage from a general GVBS presentation to this tensor presentation
    remains separate; see \cite{gap:nachtergaele96_infinite_volume_ground_projection}.
\end{lemma}
\begin{proof}\leanok
    \uses{thm:periodic_residual_window_joint_projection_decay,
        thm:isometric_projection_defect_transport,
        thm:residual_window_right_joint_range,
        thm:residual_window_left_joint_range_cover,
        thm:residual_window_full_joint_range_cover}
    The cyclic primitive sector resolution gives the joint three-window
    projection limit. The right, left, and full joint ranges are the images
    of the corresponding original interval spaces under one physical isometry.
    Isometric conjugation preserves the norm; taking adjoints reverses the
    projection order and gives the displayed defect.
\end{proof}

\begin{lemma}[Uniform original-chain open-parent gap for a periodic tensor]
    \label{thm:periodic_open_parent_gap_all_volumes}
    \lean{MPSTensor.IsPeriodic.exists_pos_openParentHamiltonianES_gap}
    \leanok

    Let $A$ be a normalized periodic tensor. There exist $R\ge1$ and
    $\delta>0$ such that the canonical range-$R$ open-parent Hamiltonian
    $H_{R,N}$ satisfies
    \begin{align}
        \forall\,N\ge R:\ \ker H_{R,N}=\mathcal G_A(N),\qquad
        \forall\,N\in\mathbb N,\ v\perp\ker H_{R,N}:\ \delta\|v\|\le\|H_{R,N}v\|.
    \end{align}
    There is no divisibility condition on $N$. The interaction range and
    constant are derived from periodicity alone. This is the normalized
    periodic tensor specialization of
    \cite[Theorem 1.2 and Section 6]{Nachtergaele1996Spectral}; the general
    GVBS presentation's identification with these support spaces remains
    separate, as recorded in
    \cite{gap:nachtergaele96_infinite_volume_ground_projection}.
\end{lemma}
\begin{proof}\leanok
    \uses{thm:periodic_physical_cover_projection_decay,thm:open_gap_residue_cover_limit,lem:positive-kernel-projection-comparison}
    The physical projection limit and the selection argument give one common
    cover length and a positive gap at every sufficiently large volume.
    The original cyclic intersection gives the exact kernels at the same
    interaction range. Each of the remaining finitely many volumes has a
    positive finite-dimensional norm gap on its kernel complement;
    taking a finite minimum completes the uniform bound.
\end{proof}

\begin{lemma}[Positive original interaction with a uniform finite-volume gap]
    \label{thm:periodic_positive_parent_gap_all_volumes}
    \lean{MPSTensor.IsPeriodic.exists_positive_parent_interaction_uniform_gap}
    \leanok

    For every normalized periodic tensor $A$, there exist $R\ge1$, a positive
    local matrix $h$ on $R$ sites, and $\delta>0$ such that its open-chain
    interaction sums $H_N$ satisfy
    \begin{align}
        \forall\,N\ge R:\ \ker H_N=\mathcal G_A(N),\qquad
        \forall\,N\in\mathbb N,\ v\perp\ker H_N:\ \delta\|v\|\le\|H_Nv\|.
    \end{align}
    No local interaction, invariant density, or sector separation is given.
    The scope is the explicit normalized periodic tensor setting of
    \cite[Theorem 1.2, Lemma existenceinteraction, and Section 6]{Nachtergaele1996Spectral}.
\end{lemma}
\begin{proof}\leanok
    \uses{thm:periodic_open_parent_gap_all_volumes,
        def:canonical_parent_interaction_matrix}
    Choose the configuration-basis matrix of the canonical local ground-space
    complement projection. It is positive, and its open interaction sum is
    exactly the canonical parent Hamiltonian at every volume.
\end{proof}

\begin{lemma}[One periodic parent interaction with finite and infinite-volume gaps]
    \label{thm:periodic_parent_finite_quasi_local_gap}
    \lean{MPSTensor.IsPeriodic.exists_positive_parent_interaction_finite_and_quasiLocal_gap}
    \leanok
    Let $A$ be a normalized periodic tensor. There exist a range $R\ge1$, a
    positive local matrix $h$ on $R$ sites, and constants $\delta,\gamma>0$
    with the following properties. The open interaction sum $H_N$ satisfies
    \begin{align}
        \forall\,N\ge R:\ \ker H_N=\mathcal G_A(N),\qquad
        \forall\,N\in\mathbb N,\ v\perp\ker H_N:\ \delta\|v\|\le\|H_Nv\|.
    \end{align}
    For every pure zero-energy state $\varphi$ of this same interaction and
    every observable $X$ supported on $[a,a+k)$ with $k>0$ and $\varphi(X)=0$,
    along every filter for which both margins $\ell_n,r_n\to\infty$, there is a limit
    \begin{align}
        e=\lim_n\varphi\bigl(X^*[H_{[a-\ell_n,a+k+r_n)},X]\bigr),\qquad
        \gamma\varphi(X^*X)\le e.
    \end{align}
    Both constants precede every state, observable and interval exhaustion.
    The interaction, its range and all support data are derived from periodicity.
    This is an existence consequence of Theorems~1.1--1.2 of
    \cite{Nachtergaele1996Spectral} in the explicit normalized periodic tensor
    setting. The general GVBS presentation remains separate; see
    \cite{gap:nachtergaele96_infinite_volume_ground_projection}.
\end{lemma}
\begin{proof}\leanok
    \uses{thm:periodic_positive_parent_gap_all_volumes,thm:periodic_quasi_local_ground_state_gap}
    Choose the positive canonical interaction and its range from the all-volume
    finite-chain gap theorem. Its exact kernel identities imply the eventual
    kernel hypothesis of the periodic infinite-volume theorem. Apply that theorem
    to the same interaction, and retain the finite-chain constant and the common
    pure-state commutator constant.
\end{proof}

\begin{lemma}[Faithful finite MPS support]
    \label{thm:faithful_local_mps_support}
    \lean{MPSTensor.observableInsertionExpectation_conjTranspose_mul_self_eq_zero_iff}
    \leanok
    Let $A$ be an MPS tensor of positive bond dimension, let $\rho>0$,
    and let $X$ act on $k\geq0$ sites. Define $\eta_A$ by the trace
    insertion formula, without assuming transfer invariance of $\rho$.
    Then
    \begin{align}
        \eta_A(X^*X)=0
        \quad\Longleftrightarrow\quad
        G_A(k)\subseteq\ker X.
        \label{eq:faithful_local_mps_support}
    \end{align}
    Tensor normalization, injectivity, and primitivity are unnecessary.
    This is the finite support characterization used for the faithful
    generators in \cite[Section~3 and the proof of the overlap lemma in
        Section~4]{Nachtergaele1996Spectral}; the construction of minimal
    faithful generators from a general GVBS presentation remains separate
    \cite{gap:nachtergaele96_infinite_volume_ground_projection}.
\end{lemma}
\begin{proof}\leanok
    \uses{def:primitive_insertion_expectation}
    Put $C_b=\sum_\sigma X_{b\sigma}A^\sigma$. Expansion and cyclicity of
    the trace give
    \begin{align}
        \operatorname{tr}\rho\,\eta_A(X^*X)
         & =\sum_b\operatorname{tr}(C_b\rho C_b^*)
        =\operatorname{tr}\!\left(\rho\sum_b C_b^*C_b\right).
        \label{eq:faithful_local_square_expansion}
    \end{align}
    The sum is positive. Faithfulness of $\rho$ implies that the final
    trace vanishes precisely when every $C_b$ vanishes. Since
    $(X\Gamma_k(Y))_b=\operatorname{tr}(C_bY)$ for every boundary matrix
    $Y$, nondegeneracy of the trace pairing is equivalent to
    $X\Gamma_k=0$. The range of $\Gamma_k$ is $G_A(k)$.
\end{proof}

\begin{lemma}[Periodic phase relations preserve open-boundary support]
    \label{thm:periodic_phase_relation_support}
    \lean{MPSTensor.MPVBlockPhaseEquiv.dim_eq_and_gaugePhaseEquiv_of_isPeriodic,
        MPSTensor.MPVBlockPhaseEquiv.groundSpaceES_eq_of_isPeriodic}
    \leanok
    Let $A$ and $B$ be normalized periodic tensors. Suppose that there is
    $z\ne0$ such that their periodic vectors satisfy
    $\ket{V^{(N)}(B)}=z^N\ket{V^{(N)}(A)}$ for every $N>0$.
    Then their bond dimensions agree, they are gauge and phase equivalent,
    and $\mathcal G_A(N)=\mathcal G_B(N)$ for every $N\ge0$.
\end{lemma}
\begin{proof}\leanok
    \uses{thm:periodic_overlap_dichotomy, thm:periodic_self_overlap_tendsto,
        thm:primitive_representative_support_gauge}
    Along common multiples of the two periods, the self-overlaps tend to
    the positive periods. Proportionality gives the exact identity
    $|\braket{V^{(N)}(A)}{V^{(N)}(B)}|^2
        =\braket{V^{(N)}(A)}{V^{(N)}(A)}\braket{V^{(N)}(B)}{V^{(N)}(B)}$.
    Thus the mixed overlap cannot tend to zero. The periodic overlap
    dichotomy gives the gauge relation. Gauge transformations and nonzero
    scalar rescaling preserve all open-boundary spaces.
\end{proof}

\begin{lemma}[Periodic representatives with exact joint support]
    \label{thm:periodic_sector_representatives_exact_support}
    \lean{MPSTensor.MPVPhaseClassData.groundSpaceES_toTensorFromBlocks_eq_representatives_of_isPeriodic,
        MPSTensor.exists_periodic_sector_representatives}
    \leanok
    Let $(A_j)_{j=0}^{b-1}$ be a finite family of normalized periodic tensors,
    with arbitrary periods, and let $\mu_j\ne0$.
    There is a subfamily selected by an injective map
    $s:\{0,\ldots,g-1\}\to\{0,\ldots,b-1\}$ such that its members are pairwise
    inequivalent under gauge and phase transformations. Every original
    member is gauge and phase equivalent to a selected member, after the
    corresponding bond-dimension identification. \begin{align}
        \forall\,N\ge0:\ \mathcal G_{\bigoplus_j\mu_jA_j}(N)
        =\mathcal G_{\bigoplus_k A_{s(k)}}(N).
    \end{align}
    The selected family inherits the original periods and positive bond
    dimensions. The statement also allows an empty original family.
    This is an algebraic support identity; it does not assert the joint
    intersection property or automatic existence of a parent interaction.
\end{lemma}
\begin{proof}\leanok
    \uses{thm:periodic_phase_relation_support}
    Select one member from each periodic-vector phase class. The preceding
    lemma identifies the open-boundary spaces of all members of a class.
    Taking their sums proves the displayed identity. The phase-class
    separation gives pairwise gauge and phase inequivalence.
\end{proof}

\begin{lemma}[Trace preservation of an isometric corner]
    \label{thm:trace_preserving_isometric_corner}
    \lean{MPSTensor.IsLeftCanonical.of_isometry_intertwine}
    \leanok
    Let $A$ be trace preserving, let $V^*V=I$, and suppose
    $A_iV=VB_i$ for every letter $i$. Then $B$ is trace preserving.
    This is the normalization identity used in the compressed periodic
    sectors of \cite[Section~2.1]{DeLasCuevas2017Irreducible}.
\end{lemma}
\begin{proof}\leanok
    \uses{def:left_canonical_tensor}
    Taking products of adjoints in $A_iV=VB_i$ and summing gives
    \begin{align}
        \sum_i B_i^*B_i
        =V^*\left(\sum_i A_i^*A_i\right)V
        =V^*V=I.
    \end{align}
\end{proof}

\begin{lemma}[Periodic support presentation of faithful normalized generators]
    \label{thm:faithful_periodic_ground_space_presentation}
    \lean{MPSTensor.exists_periodic_groundSpaceDecomposition_of_leftCanonical_of_posDef_fixedPoint}
    \leanok
    Let $A$ have positive bond dimension, assume
    $\sum_iA_i^*A_i=I$, and let $\rho>0$ satisfy $E_A(\rho)=\rho$.
    There exist a nonempty finite family of positive-dimensional periodic
    tensors $B_j$ with periods $m_j\geq1$ and support isometries $V_j$ such
    that
    \begin{align}
        V_j^*V_j                    & =I,                &
        \sum_jV_jV_j^*              & =I,                &
        \forall\,i\neq j:\ V_i^*V_j & =0,\notag            \\
        B_j^i                       & =V_j^*A_iV_j,      &
        A_iV_j                      & =V_jB_j^i,         &
        V_j^*A_i                    & =B_j^iV_j^*.\notag
    \end{align}
    For every $N\geq0$,
    \begin{align}
        G_A(N)=\bigvee_jG_{B_j}(N)
        =G_{\bigoplus_jB_j}(N).
        \label{eq:faithful_periodic_support_presentation}
    \end{align}
    This is a consequence for given normalized faithful generating data
    in \cite[Section~3]{Nachtergaele1996Spectral}. It does not construct
    such data from a general GVBS boundary limit or impose minimality;
    that passage remains separate
    \cite{gap:nachtergaele96_infinite_volume_ground_projection}.
\end{lemma}
\begin{proof}\leanok
    \uses{thm:faithful_stationary_projector_closure,
        thm:projector_closure_irreducible_corners,
        thm:trace_preserving_isometric_corner,
        thm:irreducible_left_canonical_periodic,
        thm:isometric_sector_ground_space_resolution,
        thm:block_sum_local_ground_space}
    Faithful stationarity makes every invariant projection reducing.
    Decompose into irreducible isometric corners; the identity resolution
    and positive ambient bond dimension imply that the family is nonempty.
    Trace preservation passes to each corner. Each normalized irreducible
    corner therefore has a finite period. The letterwise intertwining
    identifies the MPS boundary spaces at every length, and the unit-weight
    direct sum has their joint space.
\end{proof}

\begin{lemma}[Finite MPS densities and their ranges]
    \label{thm:local_mps_density_support}
    \lean{MPSTensor.exists_local_density_range_eq_groundSpaceES}
    \leanok
    \uses{def:primitive_insertion_expectation}
    Let $A$ have positive bond dimension and satisfy
    $\sum_i(A^i)^*A^i=I$, and let $\rho>0$ be a virtual matrix.
    For every $k\geq0$, there is a density matrix $\sigma_k$ such that
    \begin{align}
        \sigma_k                   & \geq0,                         & \operatorname{tr}\sigma_k & =1, &
        \eta_A(X)                  & =\operatorname{tr}(\sigma_kX), &
        \operatorname{ran}\sigma_k & =G_A(k).
        \label{eq:local_mps_density_support}
    \end{align}
    The trace representation holds for every observable $X$ on $k$ sites.
    No transfer invariance of $\rho$, injectivity, or primitivity is assumed.
    These finite densities identify the supports for given faithful
    generating data as in
    \cite[Section~3 and the overlap proof in Section~4]{Nachtergaele1996Spectral}.
    The construction of faithful minimal data from a general GVBS presentation
    remains separate
    \cite{gap:nachtergaele96_infinite_volume_ground_projection}.
\end{lemma}
\begin{proof}\leanok
    \uses{thm:left_canonical_physical_dimension,
        thm:primitive_insertion_local_state, thm:faithful_local_mps_support}
    The physical alphabet is nonempty by the positive physical dimension of a
    left-canonical tensor on a nonzero bond space. Trace preservation gives
    $\operatorname{tr}E^k(\rho)=\operatorname{tr}\rho$, so
    $\eta_A(I_k)=1$. Positivity and complex linearity therefore yield a
    representing positive density matrix of trace one.
    For positive $\sigma_k$, zero trace pairing with $X^*X$ is equivalent to
    $(X^*X)\sigma_k=0$. Since $\ker(X^*X)=\ker X$, this is equivalent to
    $\operatorname{ran}\sigma_k\subseteq\ker X$. The faithful support criterion
    identifies the same condition with $G_A(k)\subseteq\ker X$.
    Testing the two orthogonal-complement projections gives equality of
    the two subspaces.
\end{proof}

\begin{lemma}[Independence from an overlap matrix]
    \label{thm:subspace_independence_overlap_matrix}
    \lean{Submodule.iSupIndep_of_overlapMatrix}
    \leanok
    Let $(V_i)_{i\in I}$ be finitely many subspaces of a complex Hilbert
    space, and let $B$ be a real matrix with zero diagonal and $\|B\|<1$.
    Suppose that, for $i\ne j$, $x\in V_i$, and $y\in V_j$,
    \begin{align}
        |\langle x,y\rangle|\le B_{ij}\|x\|\|y\|.
    \end{align}
    Then the family $(V_i)$ is independent: a vanishing sum of vectors,
    one from each subspace, has every component zero.
\end{lemma}
\begin{proof}\leanok
    \uses{thm:subspace_overlap_lower_gram}
    The lower Gram estimate gives
    \begin{align}
        (1-\|B\|)\sum_i\|x_i\|^2
        \le\left\|\sum_i x_i\right\|^2.
    \end{align}
    A zero sum therefore makes every summand zero.
\end{proof}

\begin{lemma}[Eventual independence of asymptotically orthogonal subspaces]
    \label{thm:asymptotic_subspace_independence}
    \lean{Submodule.eventually_iSupIndep_of_pairwise_overlap}
    \leanok
    Let $V_i(n)$ be a finite family of subspaces of Hilbert spaces that
    may depend on $n$. Suppose that for every $\varepsilon>0$, eventually
    \begin{align}
        \forall\,i\ne j,\ x\in V_i(n),\ y\in V_j(n):\ |\langle x,y\rangle|\le\varepsilon\|x\|\|y\|.
    \end{align}
    Then the family $(V_i(n))_i$ is eventually independent.
\end{lemma}
\begin{proof}\leanok
    \uses{thm:subspace_independence_overlap_matrix}
    Let $C$ have zero diagonal and all off-diagonal entries equal to one.
    Choose $\varepsilon=(1+\|C\|)^{-1}$. Then
    $\|\varepsilon C\|<1$, and the preceding lemma applies whenever
    the stated pairwise estimate holds.
\end{proof}

\begin{lemma}[Eventual independence of primitive boundary spaces]
    \label{thm:primitive_ground_space_eventual_independence}
    \lean{MPSTensor.eventually_groundSpaceES_iSupIndep_of_primitive}
    \leanok
    \uses{def:primitive_mps, def:ground_space_es}
    A finite family of pairwise gauge and phase inequivalent primitive
    tensors, with faithful invariant matrices, has independent boundary
    spaces $G_{B_i}(N)$ for every sufficiently large $N$.
    This is the consequence of the overlap estimates in
    \cite[Lemmas disjoint and overlapestimate]{Nachtergaele1996Spectral}.
\end{lemma}
\begin{proof}\leanok
    \uses{thm:inequivalent_block_ground_space_overlap_decay,
        thm:asymptotic_subspace_independence}
    Inequivalence gives uniformly vanishing pairwise ground-space
    overlaps. Finiteness makes the bounds simultaneous, so the preceding
    theorem proves independence.
\end{proof}

\begin{lemma}[Independent ambient sums and middle restrictions]
    \label{thm:independent_ambient_sum_intersection}
    \lean{Submodule.iSup_inf_iSup_eq_of_le_independent,
        Submodule.iSupIndep_iInf_comap_of_separating,
        MPSTensor.iSupIndep_middleRestriction_comap,
        MPSTensor.groundSpace_iSupIndep_iff_groundSpaceES_iSupIndep}
    \leanok
    Let $(W_i)_{i\in I}$ be an independent family of subspaces, with $I$
    finite, and let $U_i,V_i\subseteq W_i$. Then
    \begin{align}
        \Bigl(\sum_i U_i\Bigr)\cap\Bigl(\sum_i V_i\Bigr)
        =\sum_i(U_i\cap V_i).
    \end{align}
    More generally, suppose a family of linear restrictions $f_t:E\to F$
    determines every vector of $E$. For any independent family $(G_i)$ in
    $F$, the subspaces $\bigcap_t f_t^{-1}(G_i)$ are independent in $E$.
    In particular, adjoining the first and last physical indices preserves
    independence of middle support spaces. Independence is also unchanged
    by passing from the coefficient functions to their Euclidean realization.
\end{lemma}
\begin{proof}\leanok
    Independence makes the decomposition into the ambient subspaces unique,
    so the decompositions from the two sums agree term by term.
    For the restriction assertion, apply independence after every
    restriction, then use that the restrictions determine the original vector.
    The first and last physical restrictions determine all coefficients.
    The Euclidean realization is a linear isomorphism.
\end{proof}

\begin{lemma}[Joint intersection from independent middle spaces]
    \label{thm:joint_restriction_intersection_middle_independence}
    \lean{MPSTensor.restriction_intersection_iSup_eq_of_middle_independent}
    \leanok
    Let $G_i\subseteq(\mathbb C^d)^{\otimes n}$ and
    $H_i\subseteq(\mathbb C^d)^{\otimes(n+1)}$ be finite families.
    Assume that the $G_i$ are independent and that every restriction of a
    vector in $H_i$, obtained by fixing either the first or last physical
    index, lies in $G_i$. Write
    $L_i=H_i\otimes\mathbb C^d$ and $R_i=\mathbb C^d\otimes H_i$.
    Then
    \begin{align}
        \Bigl(\sum_iL_i\Bigr)\cap\Bigl(\sum_iR_i\Bigr)
        =\sum_i(L_i\cap R_i).
    \end{align}
    This is the geometric step of Nachtergaele,
    arXiv:cond-mat/9410110, the proposition and equivalence lemma of Section 4; no tensor representation
    is assumed in this statement.
\end{lemma}
\begin{proof}\leanok
    \uses{thm:independent_ambient_sum_intersection}
    Both $L_i$ and $R_i$ lie in
    $\mathbb C^d\otimes G_i\otimes\mathbb C^d$.
    These ambient spaces are independent by the preceding lemma.
    Finite sums commute with adjoining an exterior physical index, as is
    seen by decomposing each fixed-index coefficient separately.
    Apply the independent ambient-sum intersection identity.
\end{proof}

\begin{lemma}[One-step joint ground-space intersection criterion]
    \label{thm:joint_ground_space_intersection_middle_independence}
    \lean{MPSTensor.iSup_groundSpace_eq_restriction_intersection_of_middle_independent}
    \leanok
    Let $(A_i)_{i\in I}$ be a finite family of tensors, possibly empty.
    Suppose that their ground spaces $\mathcal G_i(n)$ are independent
    and that each tensor satisfies
    \begin{align}
        (\mathcal G_i(n+1)\otimes\mathbb C^d)
        \cap(\mathbb C^d\otimes\mathcal G_i(n+1))
        =\mathcal G_i(n+2).
    \end{align}
    Then the sums $\mathcal G(q)=\sum_i\mathcal G_i(q)$ satisfy
    \begin{align}
        (\mathcal G(n+1)\otimes\mathbb C^d)
        \cap(\mathbb C^d\otimes\mathcal G(n+1))
        =\mathcal G(n+2).
    \end{align}
    The individual intersection identities and middle independence remain
    explicit hypotheses. This criterion does not by itself assert the
    unrestricted intersection theorem for periodic tensor families.
\end{lemma}
\begin{proof}\leanok
    \uses{thm:joint_restriction_intersection_middle_independence}
    Ground-space vectors restrict into the corresponding ground spaces
    on the shorter interval. Apply the preceding geometric lemma, then
    substitute the assumed individual intersection identities.
\end{proof}

\begin{lemma}[Finite restrictions of a faithful MPS state]
    \label{thm:faithful_quasi_local_state_density_support}
    \lean{MPSTensor.exists_quasiLocalExpectation_density_range_eq_groundSpaceES}
    \leanok
    \uses{def:mps_quasi_local_expectation, def:ground_space_es}
    Let $A$ have positive bond dimension, assume
    $\sum_iA_i^*A_i=I$, and let $\rho>0$ satisfy $E_A(\rho)=\rho$.
    Write $\omega_{A,\rho}$ for the resulting quasi-local state.
    For every $k\geq0$, there exists a density matrix $\sigma_k\geq0$,
    of trace one, such that
    \begin{align}
        \forall\,a\in\mathbb Z,\ X:\ \omega_{A,\rho}(X_{[a,a+k)})
                                   & =\operatorname{tr}(\sigma_kX),\notag \\
        \operatorname{ran}\sigma_k & =G_A(k).\notag
    \end{align}
    Thus the same density describes every translated interval of length
    $k$, and its support is exactly the tensor boundary space.
    The generating data are given, as in
    \cite[Section~3 and the overlap proof in Section~4]{Nachtergaele1996Spectral};
    their construction from a general GVBS presentation remains separate
    \cite{gap:nachtergaele96_infinite_volume_ground_projection}.
\end{lemma}
\begin{proof}\leanok
    \uses{thm:local_mps_density_support,thm:mps_quasi_local_interval_expectation}
    The quasi-local state agrees on each consecutive interval with the
    trace insertion functional. The finite-density support theorem
    therefore gives the same matrix $\sigma_k$ at every position.
\end{proof}

\begin{lemma}[Positive physical dimension of a normalized tensor]
    \label{thm:left_canonical_physical_dimension}
    \lean{MPSTensor.IsLeftCanonical.physDim_ne_zero}
    \leanok
    A tensor of positive bond dimension satisfying
    $\sum_iA_i^*A_i=I$ has positive physical dimension.
    This is the left-canonical normalization of
    \cite[Section~2.1]{DeLasCuevas2017Irreducible}.
\end{lemma}
\begin{proof}\leanok
    \uses{def:left_canonical_tensor}
    An empty physical alphabet would make the sum zero, contradicting the
    nonzero identity on the bond space.
\end{proof}

\begin{lemma}[Infinite-volume gap for faithful normalized generators]
    \label{thm:faithful_generators_quasi_local_gap}
    \lean{MPSTensor.exists_pos_quasiLocalCommutator_limit_gap_of_leftCanonical_of_posDef_fixedPoint}
    \leanok
    Let $A$ have positive bond dimension, assume $\sum_iA_i^*A_i=I$, and
    let $\rho>0$ satisfy $E_A(\rho)=\rho$. Let $h\geq0$ have range $R>0$,
    and suppose that $\ker H_h(N)=G_A(N)$ for every sufficiently large $N$.
    There is one $\gamma>0$ with the following property. For every pure
    zero-energy state $\varphi$, every observable $X$ supported on a
    positive-length interval $[a,a+k)$ with $\varphi(X)=0$, and every filter
    along which $\ell_n,r_n\to\infty$, put
    $\Lambda_n=[a-\ell_n,a+k+r_n)$. Then the finite-volume limit
    exists and satisfies
    \begin{align}
        \lim_n\varphi\!\left(X^*[H_{h,\Lambda_n},X]\right)
        \geq\gamma\varphi(X^*X).
        \label{eq:faithful_generators_literal_gap}
    \end{align}
    Neither irreducibility nor periodicity of $A$, nor translation
    invariance of $\varphi$, is assumed. This is a consequence for given
    faithful normalized generating data in the setting of
    \cite[Theorem~1.2 and Sections~3,~6]{Nachtergaele1996Spectral}.
    Constructing such data from an arbitrary GVBS boundary limit remains
    separate; see
    \cite{gap:nachtergaele96_infinite_volume_ground_projection}.
\end{lemma}
\begin{proof}\leanok
    \uses{thm:left_canonical_physical_dimension,
        thm:faithful_periodic_ground_space_presentation,
        thm:periodic_block_family_quasi_local_gap}
    Normalization implies positive physical dimension. The faithful
    stationary matrix yields a finite periodic family $B_j$ with
    $G_A(N)=G_{\bigoplus_jB_j}(N)$ at every length. Hence the same
    interaction has these joint spaces as its eventual kernels. Applying
    the periodic-family gap theorem with unit weights proves
    (\ref{eq:faithful_generators_literal_gap}), with its constant chosen
    before the state, observable, or exhaustion.
\end{proof}

\begin{lemma}[Zero energy of the generated stationary state]
    \label{thm:stationary_state_zero_local_energy}
    \lean{MPSTensor.quasiLocalExpectation_mem_parentGroundStateFace_of_eventual_open_kernel}
    \leanok
    Let $A$ satisfy $\sum_i(A^i)^*A^i=I$, let $\rho\geq0$ be a fixed point
    of its transfer map with $\operatorname{tr}\rho\neq0$, and let $h\geq0$
    be an interaction of positive range whose open-chain kernels equal the
    MPS boundary spaces $G_N(A)$ for all sufficiently large $N$. Then the
    quasi-local state generated by $A$ and $\rho$ has zero local energy for
    $h$. No primitivity, faithfulness of $\rho$, or purity is assumed
    \cite[Theorems~1.1--1.2]{Nachtergaele1996Spectral}.
\end{lemma}
\begin{proof}\leanok
    \uses{thm:eventual_kernel_ground_state_support_iff,thm:mps_quasi_local_expectation_state,thm:mps_support_projection_expectation_one}
    The restriction of the generated state to a long interval is supported
    in $G_N(A)$, which is the kernel of the open Hamiltonian there; hence
    every local interaction term has zero expectation.
\end{proof}

\begin{lemma}[Zero-energy states for faithful stationary generators]
    \label{thm:faithful_generators_ground_state_face}
    \lean{MPSTensor.exists_parentGroundStateFace_classification_of_leftCanonical_of_posDef_fixedPoint}
    \leanok
    Let $A$ have positive bond dimension and satisfy $\sum_i(A^i)^*A^i=I$,
    with a positive definite fixed
    point $\rho$ of its transfer map, and let $h\geq0$ be an interaction of
    positive range whose open-chain kernels eventually equal $G_N(A)$.
    The state generated by $A$ and $\rho$ has zero local energy, and there are a blocking length $L>0$ and a nonempty finite family of
    pairwise inequivalent normalized primitive tensors $B_j$ with positive
    definite fixed points such that the blocked tensor of $A$ and the
    unit-weight direct sum of the $B_j$ have the same boundary spaces at
    every length. The zero-energy states of $h$ are exactly the convex
    combinations of the states generated by the $B_j$, pulled back
    through the blocking, and the pure zero-energy states are exactly these
    pulled-back states \cite[Theorems~1.1--1.2]{Nachtergaele1996Spectral}.
\end{lemma}
\begin{proof}\leanok
    \uses{thm:stationary_state_zero_local_energy,thm:faithful_periodic_ground_space_presentation,thm:periodic_block_family_ground_states}
    The faithful fixed point gives a periodic presentation of the boundary
    spaces; blocking by a common multiple of the periods yields primitive
    sectors with faithful fixed points. The classification of the
    zero-energy face for primitive sectors, carried through the
    blocking isomorphism of quasi-local algebras, gives both assertions.
\end{proof}

\begin{lemma}[Bounded preimages on a finite sum of subspaces]
    \label{thm:sector_bounded_preimages}
    \lean{Submodule.exists_norm_le_preimage_of_mem_iSup_of_lower_gram}
    \leanok
    Let $T:E\to F$ be a complex linear map from a normed space to an
    inner-product space, and let $(W_i)_{i\in I}$ be a finite family of
    subspaces of $F$. Suppose $C_i\geq0$ and
    \begin{align}
        \forall\,x_i\in W_i:\ \frac12\sum_i\|x_i\|^2\leq\Bigl\|\sum_i x_i\Bigr\|^2.
    \end{align}
    Suppose also that each $x_i\in W_i$ admits a preimage $u_i$ with
    $Tu_i=x_i$ and $\|u_i\|\leq C_i\|x_i\|$.
    Every $x\in\sum_iW_i$ then admits a preimage $u$ satisfying
    \begin{align}
        Tu=x,\qquad \|u\|\leq2\Bigl(\sum_iC_i\Bigr)\|x\|.
    \end{align}
    This is a consequence of the given lower Gram bound, of the form
    used in the overlap estimate in
    \cite[Section~4]{Nachtergaele1996Spectral}.
\end{lemma}
\begin{proof}\leanok
    Choose a decomposition $x=\sum_i x_i$ with $x_i\in W_i$.
    The lower Gram bound gives $\|x_i\|\leq2\|x\|$ for every $i$.
    Choose the stated preimages and put $u=\sum_i u_i$.
    Linearity and the triangle inequality give the conclusion.
\end{proof}

\begin{lemma}[Canonical open kernels from eventual intersections]
    \label{thm:canonical_open_kernel_from_eventual_intersection}
    \lean{MPSTensor.ker_openParentHamiltonianES_eq_groundSpaceES_of_restriction_intersection,
        MPSTensor.exists_ker_openParentHamiltonianES_eq_groundSpaceES_of_eventually_restriction_intersection}
    \leanok
    Let $d\ge1$ and let $A$ be any tensor. Suppose that, for every $q\ge R$,
    \begin{align}
        (\mathcal G_A(q)\otimes\mathbb C^d)
        \cap(\mathbb C^d\otimes\mathcal G_A(q))=\mathcal G_A(q+1).
    \end{align}
    If $R>0$, the canonical open Hamiltonian of range $R$ has kernel
    $\mathcal G_A(N)$ for every $N\ge R$.
    Consequently, if the one-step identity holds at every sufficiently
    large length, there is $R_0>0$ such that the same kernel identity holds
    for every $N\ge R\ge R_0$.
    The intersection identity is an explicit hypothesis of this criterion.
\end{lemma}
\begin{proof}\leanok
    \uses{lem:mem_ground_space_es,lem:open_chain_iteration_subspace_sequence,thm:open_parent_hamiltonian_window_ground_space,thm:open_ground_space_le_hamiltonian_kernel,rem:cyclic_window_operators}
    A vector in the Hamiltonian kernel satisfies every contiguous local
    ground-space constraint. Iterating the one-step intersections enlarges
    these local spaces to the whole interval. Conversely, every
    open-boundary vector satisfies every local constraint.
\end{proof}

\begin{lemma}[Eventual joint intersection for primitive sectors]
    \label{thm:primitive_joint_intersection_from_overlap}
    \lean{MPSTensor.eventually_iSup_groundSpace_restriction_intersection_of_independent_middle,
        MPSTensor.eventually_iSup_groundSpace_restriction_intersection_of_isPrimitiveMPS}
    \leanok
    For a finite tensor family, eventual independence of the middle ground
    spaces and the individual eventual one-step intersection identities
    imply the eventual joint one-step intersection identity.
    In particular, this applies to any finite family of pairwise gauge and
    phase inequivalent normalized primitive tensors with positive bond
    dimensions and faithful invariant matrices. Thus, writing
    $\mathcal G(q)=\sum_j\mathcal G_j(q)$, for all sufficiently large $n$,
    \begin{align}
        (\mathcal G(n+1)\otimes\mathbb C^d)
        \cap(\mathbb C^d\otimes\mathcal G(n+1))
        =\mathcal G(n+2).
    \end{align}
    The family may be empty. The primitive assertion concerns the given
    generating tensors, rather than an arbitrary GVBS presentation.
\end{lemma}
\begin{proof}\leanok
    \uses{thm:joint_ground_space_intersection_middle_independence,
        thm:primitive_ground_space_eventual_independence,
        thm:periodic_original_intersection}
    Finiteness gives one eventual set on which all individual identities
    hold. Apply the middle-independence criterion on its intersection with
    the eventual independence set. For primitive tensors, overlap decay
    gives that independence, and period-one periodicity gives the
    individual intersections.
\end{proof}

\begin{lemma}[Canonical primitive-family interaction from overlap]
    \label{thm:primitive_parent_interaction_from_overlap}
    \lean{MPSTensor.exists_ker_openParentHamiltonianES_toTensorFromBlocks_eq_groundSpaceES_of_primitive_overlap,
        MPSTensor.exists_positive_parent_interaction_of_primitive_overlap}
    \leanok
    Let $d\ge1$, let $(A_j)$ be a finite family of pairwise gauge and phase
    inequivalent normalized primitive tensors with positive bond dimensions
    and faithful invariant matrices, and let $\mu_j\ne0$.
    For $T=\bigoplus_j\mu_jA_j$, there exists $R_0>0$ such that
    \begin{align}
        \forall\,N\ge R\ge R_0:\ \ker H_{R,N}^{\mathrm{open}}(T)=\mathcal G_T(N).
    \end{align}
    In particular, the canonical local projection of range $R_0$ is a
    positive parent interaction with this exact kernel at every $N\ge R_0$.
    No simultaneous word span is given in this assertion or used in the
    family intersection argument.
\end{lemma}
\begin{proof}\leanok
    \uses{thm:primitive_joint_intersection_from_overlap,
        thm:canonical_open_kernel_from_eventual_intersection}
    Nonzero weights preserve the joint open-boundary space. Substitute
    the eventual joint intersection into the canonical kernel criterion.
    The orthogonal projection onto the local ground-space complement is
    positive, so choosing that projection gives the interaction.
\end{proof}

\begin{lemma}[Observable insertions under a virtual intertwiner]
    \label{thm:observable_insertion_isometric_intertwiner}
    \lean{MPSTensor.physicalObservableTransfer_intertwine}
    \lean{MPSTensor.observableInsertionExpectation_eq_of_isometric_intertwine}
    \leanok
    Suppose $A_iV=VB_i$ for every letter $i$. For every virtual matrix
    $\sigma$ and every observable $X$ on $k\geq0$ sites,
    \begin{align}
        E_{A,X}(V\sigma V^*)=V E_{B,X}(\sigma)V^*.
        \label{eq:observable_insertion_intertwiner}
    \end{align}
    If additionally $V^*V=I$, $\rho=V\sigma V^*$, and
    $\operatorname{tr}\sigma\neq0$, then $\operatorname{tr}\rho=\operatorname{tr}\sigma$ and
    $\eta_{A,\rho}(X)=\eta_{B,\sigma}(X)$.
    No normalization, stationarity, or positivity is needed for these
    identities. This concerns the finite insertion formula in
    \cite[Section~3]{Nachtergaele1996Spectral}.
\end{lemma}
\begin{proof}\leanok
    \uses{def:physical_observable_transfer,def:primitive_insertion_expectation}
    Letter intertwining gives $A^\tau V=VB^\tau$ for every word $\tau$.
    Expanding the insertion as a sum over words proves
    (\ref{eq:observable_insertion_intertwiner}). When $V^*V=I$, cyclicity
    gives $\operatorname{Tr}(VZV^*)=\operatorname{Tr}Z$; apply this to both
    the inserted matrix and the normalization denominator.
\end{proof}

\begin{lemma}[Faithful compression of stationary generating data]
    \label{thm:stationary_support_faithful_generators}
    \lean{MPSTensor.exists_faithful_stationary_supportCompression,
        MPSTensor.exists_faithful_stationary_supportCompression_of_leftCanonical}
    \leanok
    Let $A$ satisfy $\sum_iA_i^*A_i=I$, and let $\rho\geq0$ have nonzero
    trace with $E_A(\rho)=\rho$. There exist a positive dimension $E$, an
    isometry $V$, a tensor $B$, and a faithful matrix $\sigma>0$ such that
    \begin{align}
        V^*V                    & =I,                           & VV^*   & =\operatorname{supp}\rho, &
        B_i                     & =V^*A_iV,                     & \sigma & =V^*\rho V,\notag           \\
        \sum_iB_i^*B_i          & =I,                           & A_iV   & =VB_i,                    &
        E_B(\sigma)             & =\sigma,                      &
        \operatorname{Tr}\sigma & =\operatorname{Tr}\rho.\notag
    \end{align}
    $\eta_{A,\rho}(X)=\eta_{B,\sigma}(X)$ for every observable
    on every finite interval, including length zero.
    All conclusions except $\sum_iB_i^*B_i=I$ hold without the
    normalization of $A$.
    This is a restriction of given stationary generating data to their
    support, as in \cite{Wolf2012Quantum}, restriction to full-rank fixed
    points, and the normalized generating representation of
    \cite[Section~3]{Nachtergaele1996Spectral}. Constructing generating
    data from an arbitrary GVBS boundary limit remains separate
    \cite{gap:nachtergaele96_infinite_volume_ground_projection}.
\end{lemma}
\begin{proof}\leanok
    \uses{thm:fp_inv_proj,thm:trace_preserving_isometric_corner,
        thm:observable_insertion_isometric_intertwiner}
    The stationary support projection $P$ obeys $(I-P)A_iP=0$.
    Choose $V$ with $V^*V=I$ and $VV^*=P$. Then $PV=V$ gives
    $A_iV=V(V^*A_iV)$. Compression preserves trace preservation and the
    stationary equation. Since $P\rho P=\rho$, the compressed matrix is
    positive definite and expands as $\rho=V\sigma V^*$.
    Cyclicity gives $\operatorname{Tr}\sigma=\operatorname{Tr}\rho$,
    whose nonvanishing forces $E>0$. Apply the insertion intertwiner to
    obtain all finite expectation identities.
\end{proof}

\begin{definition}[Corner boundary space]
    \label{def:corner_boundary_space}
    \lean{MPSTensor.cornerGroundSpaceES, MPSTensor.cornerGroundSpaceES_one}
    \leanok
    For an MPS tensor $A$, a projection $P$, and $L\geq0$, let $G_A^P(L)$
    be the set of open-boundary vectors $\Gamma_L^A(Y)$ with $PYP=Y$.
    For $P=I$ it is $G_A(L)$.
\end{definition}

\begin{lemma}[Boundary space of an isometric compression]
    \label{thm:isometric_compression_corner_boundary}
    \lean{MPSTensor.groundSpace_eq_map_cornerSubmodule_of_isometric_bond_intertwiner,
        MPSTensor.groundSpaceES_eq_cornerGroundSpaceES_of_isometric_bond_intertwiner}
    \leanok
    \uses{def:corner_boundary_space}
    Let $V^*V=I$ and $A_iV=VB_i$ for every $i$. Then for every $L\geq0$,
    \begin{align}
        G_B(L) & =G_A^{VV^*}(L).
        \label{eq:isometric_compression_corner_boundary}
    \end{align}
\end{lemma}
\begin{proof}\leanok
    \uses{thm:spt_isometric_bond_restriction}
    The intertwining gives $\Gamma_L^A(VXV^*)=\Gamma_L^B(X)$, and
    $VXV^*$ lies in the corner of $VV^*$. Conversely, if $Y$ lies in that
    corner, then $Y=V(V^*YV)V^*$.
\end{proof}

\begin{lemma}[Local supports of stationary generating data]
    \label{thm:stationary_local_support}
    \lean{MPSTensor.observableInsertionExpectation_conjTranspose_mul_self_eq_zero_iff_of_stationary,
        MPSTensor.exists_local_density_range_eq_cornerGroundSpaceES_of_stationary}
    \leanok
    \uses{def:primitive_insertion_expectation, def:corner_boundary_space}
    Let $A$ be an MPS tensor, let $\rho\geq0$ have nonzero trace and satisfy
    $\sum_iA_i\rho A_i^*=\rho$, and let $P$ be its support projection.
    The matrix $\rho$ need not be faithful. For every $k\geq0$ and every
    $k$-site matrix $X$,
    \begin{align}
        \eta_{A,\rho}(X^*X)=0
        \quad\Longleftrightarrow\quad
        G_A^P(k)\subseteq\ker X.
        \label{eq:stationary_local_support}
    \end{align}
    If in addition $\sum_iA_i^*A_i=I$, there is a density matrix $\sigma_k$
    with $\eta_{A,\rho}(X)=\operatorname{tr}(\sigma_kX)$ for every $X$ and
    $\operatorname{ran}\sigma_k=G_A^P(k)$.
    This describes the local support spaces of normalized stationary
    generating data, as in \cite[Section~3 and the proof of the overlap
        lemma in Section~4]{Nachtergaele1996Spectral}, without the faithfulness
    that the source obtains from minimality. Passing from an arbitrary
    GVBS boundary limit to normalized stationary data remains separate
    \cite{gap:nachtergaele96_infinite_volume_ground_projection}.
\end{lemma}
\begin{proof}\leanok
    \uses{thm:stationary_support_faithful_generators,
        thm:isometric_compression_corner_boundary,
        thm:faithful_local_mps_support, thm:local_mps_density_support}
    Compress to the support of $\rho$. The compressed data are faithful
    and have the same finite expectations, so the faithful support
    theorem and the finite-density theorem apply to them. The boundary
    space of the compressed tensor is $G_A^P(k)$ by
    (\ref{eq:isometric_compression_corner_boundary}).
\end{proof}

\begin{definition}[Embedding a residual boundary]
    \label{def:residual_boundary_embedding}
    \lean{MPSTensor.residualBoundaryEmbeddingES,
        MPSTensor.residualBoundaryEmbeddingES_apply,
        MPSTensor.blockedResidualBoundaryMapES_eq_groundSpaceMapES_comp}
    \leanok
    For a rectangular map $V:\mathbb C^E\to\mathbb C^D$, let
    $J_V(Y)=YV^*$ on the Hilbert space of rectangular matrices, with the
    Hilbert--Schmidt norm. Suppose that
    $V^*(A^{[L]})^i=B^iV^*$ for every blocked letter $i$.
    Then the residual boundary map satisfies
    $\mathcal F^{N,r}_{A,B,V}=\Gamma^A_{NL+r}J_V$ for every $N,r\geq0$.
    In particular, the boundary trace uses $YV^*$ in the original site order.
\end{definition}
\begin{proof}\leanok
    \uses{thm:blocked_residual_boundary_identity}
    The rectangular intertwiner moves $V^*$ through each blocked letter.
    Thus the residual coefficient is
    $\operatorname{tr}(A^{\widetilde u}A^\tau YV^*)$, which is the original
    boundary coefficient at $YV^*$.
\end{proof}

\begin{lemma}[Uniform boundary lifts for a periodic tensor]
    \label{thm:periodic_boundary_uniform_lifts}
    \lean{MPSTensor.IsPeriodic.exists_eventually_groundSpaceMapES_preimage_norm_le}
    \leanok
    Let $A$ be a normalized irreducible periodic tensor. There are $C>0$
    and $n_0$ such that every $n\geq n_0$ and $x\in\mathcal G_n(A)$ have
    a virtual boundary $u$ with $\Gamma^A_n(u)=x$ and
    $\|u\|\leq C\|x\|$. The bound includes every residue length.
    This is the boundary normalization step in the mixed-transfer argument
    of \cite[Lemma disjoint]{Nachtergaele1996Spectral}, applied to
    periodic generating tensors using
    \cite[Lemma~6]{DeLasCuevas2017Irreducible}.
\end{lemma}
\begin{proof}\leanok
    \uses{thm:cyclic_primitive_sector_resolution,
        thm:blocked_residual_ground_space_resolution,
        thm:residual_supported_inverse_uniform_tail,
        thm:residual_sector_overlap_decay,
        thm:subspace_overlap_lower_gram,thm:sector_bounded_preimages,
        def:residual_boundary_embedding,
        thm:inverse_cyclic_word_gram,thm:supported_gram_projector_cancellation}
    Write $n=Nm+r$ using the tensor's period $m$. Its primitive cyclic
    sectors give residual maps $S_j$ and supported inverses $H_j$ with
    $S_jH_jS_j^*=P_{\operatorname{ran}S_j}$ and uniformly bounded
    $\|H_j\|\|S_j\|$, independently of $r$.
    The original lift of a sector vector $x_j$ is $J_{V_j}H_jS_j^*x_j$.
    Small pairwise residual overlaps give
    $\tfrac12\sum_j\|x_j\|^2\leq\|\sum_jx_j\|^2$.
    The finite sum of these bounded lifts therefore has norm at most
    $2(\sum_jc_j)\|x\|$, where
    $c_j$ bounds $\|J_{V_j}\|\|H_j\|\|S_j\|$.
    Enlarging this constant makes it positive. Since $N=\lfloor n/m\rfloor$
    tends to infinity, the estimate holds at every sufficiently large $n$.
\end{proof}

\begin{lemma}[Vanishing overlap of inequivalent periodic boundary spaces]
    \label{thm:periodic_boundary_overlap_decay}
    \lean{MPSTensor.IsPeriodic.eventually_norm_inner_groundSpaceES_le_of_inequivalent}
    \leanok
    Let $A,B$ be normalized irreducible periodic tensors, possibly of
    different periods and bond dimensions. Suppose that whenever their
    bond dimensions agree, they are inequivalent up to similarity and phase.
    For every $\varepsilon>0$, all sufficiently large $n$ satisfy
    $|\langle x,y\rangle|\leq\varepsilon\|x\|\|y\|$ for every
    $x\in\mathcal G_n(A)$ and $y\in\mathcal G_n(B)$.
    The statement concerns boundary spaces of generating tensors and does
    not assert purity of their associated states.
\end{lemma}
\begin{proof}\leanok
    \uses{thm:periodic_boundary_uniform_lifts,thm:mixed_boundary_gram_decay,
        thm:transfer_operator_gap_irr_tp,thm:transfer_operator_gap_rect_irr_tp}
    Irreducibility, normalization and inequivalence give a strict spectral
    radius bound for the original mixed transfer operator. Consequently
    $\|\Gamma^A_n{}^*\Gamma^B_n\|\to0$.
    Choose boundary lifts $u,v$ with
    $\|u\|\leq C_A\|x\|$ and $\|v\|\leq C_B\|y\|$. Then
    $|\langle x,y\rangle|
        \leq C_AC_B\|\Gamma^A_n{}^*\Gamma^B_n\|\|x\|\|y\|$.
\end{proof}

\begin{lemma}[Eventual independence of periodic boundary spaces]
    \label{thm:periodic_ground_space_eventual_independence}
    \lean{MPSTensor.eventually_groundSpaceES_iSupIndep_of_isPeriodic}
    \leanok
    A finite family of pairwise gauge and phase inequivalent normalized
    irreducible periodic tensors has independent original-chain boundary
    spaces at every sufficiently large length. No common period is prescribed.
\end{lemma}
\begin{proof}\leanok
    \uses{thm:periodic_boundary_overlap_decay,thm:asymptotic_subspace_independence}
    Finitely many pairwise overlap bounds hold simultaneously at all
    sufficiently large lengths. Apply the finite lower Gram criterion
    for independence.
\end{proof}

\begin{lemma}[Periodic-family interaction criterion]
    \label{thm:periodic_family_interaction_independence_criterion}
    \lean{MPSTensor.exists_ker_openParentHamiltonianES_of_periodic_family_independence,
        MPSTensor.exists_positive_parent_interaction_of_periodic_family_independence}
    \leanok
    Let $d\ge1$, let $(A_j)$ be a finite family of normalized periodic
    tensors with arbitrary periods, and let $\mu_j\ne0$.
    Assume explicitly that their original middle ground spaces
    $(\mathcal G_{A_j}(n))_j$ are independent at every sufficiently large
    length $n$. For $T=\bigoplus_j\mu_jA_j$, there exists $R_0>0$ such that
    \begin{align}
        \forall\,N\ge R\ge R_0:\ \ker H_{R,N}^{\mathrm{open}}(T)=\mathcal G_T(N).
    \end{align}
    In particular, the canonical local projection of range $R_0$ gives a
    positive interaction whose open kernel is exactly $\mathcal G_T(N)$
    for every $N\ge R_0$. The family may be empty.
    This is an independence criterion: its middle-space hypothesis remains
    explicit and is not an unrestricted periodic-family existence theorem.
\end{lemma}
\begin{proof}\leanok
    \uses{thm:periodic_original_intersection,thm:primitive_joint_intersection_from_overlap,thm:canonical_open_kernel_from_eventual_intersection,lem:positive_canonical_interaction_exact_open_kernels}
    Each periodic sector has its original one-step intersection identity
    eventually. Finiteness and the assumed middle independence allow the
    geometric joint-intersection criterion to combine these identities.
    Nonzero weights preserve the sum of the sector spaces. Iterating the
    joint intersections gives the exact canonical kernels.
\end{proof}

\begin{lemma}[Constructed periodic-family commutator gap criterion]
    \label{thm:periodic_family_constructed_gap_independence_criterion}
    \lean{MPSTensor.exists_positive_parent_interaction_quasiLocalCommutator_gap_of_periodic_family_independence}
    \leanok
    Under the preceding hypotheses, one can choose a positive interaction $h$ of range $R>0$, with exact open kernels for all $N\ge R$,
    and one constant $\gamma>0$ as follows.
    For every pure normalized zero-energy state $\varphi$ of $h$ and every
    centered local observable $X$ on a positive-length interval beginning
    at $a$, the finite-volume commutator energies converge along
    every interval exhaustion whose left and right margins diverge:
    \begin{align}
        \lim \varphi\bigl(X^*[H_{[a-\ell,a+k+r)},X]\bigr)=e,
        \qquad \gamma\varphi(X^*X)\le e.
    \end{align}
    The constant is chosen before the state, observable, filter, and
    boundary margins. No translation-invariance assumption is imposed on
    the state being classified. The middle-space independence assumption
    is retained in this sufficient criterion.
\end{lemma}
\begin{proof}\leanok
    \uses{thm:periodic_family_interaction_independence_criterion,
        thm:periodic_block_family_quasi_local_gap}
    Apply the existing original-chain periodic-family gap theorem to the
    constructed interaction. Its eventual kernel hypothesis follows from
    the exact identities for all volumes at least the interaction range.
\end{proof}

\begin{lemma}[All-length gap from aligned estimates and residual covers]
    \label{thm:aligned_residue_cover_gap_limit}
    \lean{MPSTensor.exists_pos_openParentHamiltonianES_gap_of_aligned_gap_of_residue_cover_limit}
    \leanok
    Let $A$ be any tensor and $R,L>0$. Assume that its canonical open
    Hamiltonian has kernel $\mathcal G_A(N)$ for every $N\geq R$ and one
    positive norm gap at every sufficiently large length divisible by $L$.
    Assume also that, for every choice of exterior lengths $K(M),r(M)$,
    the projection defect of the cover with left length $(K(M)+M)L$,
    right length $ML+r(M)$, and overlap length $ML$ tends to zero as
    $M\to\infty$. Then there is $\delta>0$ such that
    \begin{align}
        \delta\|v\|\leq\|H_R^{\mathrm{open}}(N)v\|
        \qquad(N\geq0,
        \ v\perp\ker H_R^{\mathrm{open}}(N)).
    \end{align}
    All three estimates are given hypotheses of this criterion.
    This is the covering argument of
    \cite[Section~6]{Nachtergaele1996Spectral}.
\end{lemma}
\begin{proof}\leanok
    \uses{thm:open_residue_cover_gap_criterion,thm:uniform_upper_bound_all_selections,thm:finite_open_residue_gaps,lem:positive-kernel-projection-comparison}
    Convergence along every selection gives one overlap length at which
    every exterior choice has projection defect at most $1/2$.
    Increase that overlap length to lie beyond the aligned-gap threshold
    and the interaction range. The two-interval comparison gives a common
    gap at every sufficiently large original length. A positive finite
    minimum includes the remaining lengths, with their kernels.
\end{proof}

\begin{lemma}[Canonical interactions depend only on local support]
    \label{lem:canonical_interaction_support_transport}
    \lean{MPSTensor.canonicalParentInteractionMatrix_eq_of_groundSpaceES_eq,
        MPSTensor.openParentHamiltonianES_eq_of_groundSpaceES_eq}
    \leanok
    If two tensors have the same $R$-site boundary space, their canonical
    interaction matrices agree. For $R>0$, their canonical open Hamiltonians
    consequently agree at every volume. Their bond dimensions may differ.
\end{lemma}
\begin{proof}\leanok
    \uses{def:canonical_parent_interaction_matrix}
    The canonical matrix represents the orthogonal projection onto the
    local boundary-space complement. Equality of that space therefore gives
    equality of the matrices and of every sum of their local translates,
    as in \cite[Lemma \texttt{existenceinteraction}]{Nachtergaele1996Spectral}.
\end{proof}

\begin{lemma}[Boundary-space inclusion under an isometric intertwiner]
    \label{thm:isometric_cointertwiner_ground_space_inclusion}
    \lean{MPSTensor.groundSpaceES_le_of_isometric_cointertwiner}
    \leanok
    Suppose that $V^*V=I$ and $V^*A_i=B_iV^*$ for every physical letter.
    Then $\mathcal G_N(B)\subseteq\mathcal G_N(A)$ for every $N\geq0$.
    No normalization or injectivity assumption is required.
\end{lemma}
\begin{proof}\leanok
    \uses{def:ground_space_es}
    The letter relation extends to words. Trace cyclicity gives
    $\Gamma^A_N(VYV^*)=\Gamma^B_N(Y)$, proving the range inclusion.
\end{proof}

\begin{lemma}[Separation of corners from distinct periodic tensors]
    \label{thm:distinct_periodic_corner_separation}
    \lean{MPSTensor.IsPeriodic.not_gaugePhaseEquiv_of_primitive_corners}
    \leanok
    Let $A,A'$ be gauge and phase inequivalent normalized irreducible
    periodic tensors, with possibly different periods and bond dimensions.
    Fix $L>0$ and tensors $B,C$ on $L$-site blocks, together with isometries
    $V,W$ satisfying
    $V^*(A^{[L]})_i=B_iV^*$ and
    $W^*(A'^{[L]})_i=C_iW^*$ for every blocked letter.
    Suppose that $B$ is a primitive normalized tensor of positive bond
    dimension with a faithful invariant matrix. Then $B,C$ are inequivalent
    up to similarity and phase whenever their bond dimensions agree.
    The conclusion follows from original-chain separation; separation of
    compressed corners is not assumed.
\end{lemma}
\begin{proof}\leanok
    \uses{thm:isometric_cointertwiner_ground_space_inclusion,
        thm:periodic_boundary_overlap_decay,
        thm:blocked_local_ground_space_transport,
        thm:primitive_representative_support_gauge,
        thm:primitive_ground_space_map_eventual_inverse_gram_bound}
    A hypothetical corner gauge relation gives
    $\mathcal G_N(B)=\mathcal G_N(C)$ at every $N$.
    For large $N$, primitivity and faithfulness make $\Gamma^B_N$ injective,
    so $x=\Gamma^B_N(I)$ is nonzero. The two isometries place $x$ in both
    blocked original spaces. The common physical blocking isometry gives a
    nonzero $y\in\mathcal G_{NL}(A)\cap\mathcal G_{NL}(A')$.
    Original-chain overlap decay at coefficient $1/2$ implies
    $\|y\|^2\leq\tfrac12\|y\|^2$, a contradiction.
\end{proof}

\begin{lemma}[Finite density supports of a stationary MPS state]
    \label{thm:stationary_state_local_density_support}
    \lean{MPSTensor.exists_faithful_quasiLocalExpectation_density_support_of_leftCanonical}
    \leanok
    Let $A$ satisfy $\sum_iA_i^*A_i=I$, and let $\rho\geq0$ have nonzero
    trace with $E_A(\rho)=\rho$. There exist a positive dimension $E$, an
    isometry $V$, a normalized tensor $B$, and a faithful stationary matrix
    $\sigma>0$ such that
    \begin{align}
        V^*V                    & =I,                           & VV^*   & =\operatorname{supp}\rho, &
        B_i                     & =V^*A_iV,                     & \sigma & =V^*\rho V,\notag           \\
        \sum_iB_i^*B_i          & =I,                           & A_iV   & =VB_i,                    &
        E_B(\sigma)             & =\sigma,                      &
        \operatorname{Tr}\sigma & =\operatorname{Tr}\rho.\notag
    \end{align}
    Let $\varphi$ be the quasi-local state defined by the original
    generators $A,\rho$. For every $k\geq0$, there is a matrix
    $\tau_k\geq0$ of trace one such that, for every translated interval
    $[a,a+k)$ and every matrix observable $X$ on that interval,
    \begin{align}
        \varphi(j_{a,k}(X))      & =\operatorname{Tr}(\tau_kX), &
        \operatorname{ran}\tau_k & =G_B(k).
        \label{eq:stationary_state_density_support}
    \end{align}
    The physical dimension is derived from the compressed normalization.
    No faithfulness of $\rho$ or primitivity of the generating tensor is
    required. This concerns given stationary generating data in
    \cite[Section~3]{Nachtergaele1996Spectral} and its finite support
    identification in Section~4; constructing such data from an arbitrary
    GVBS boundary limit remains separate
    \cite{gap:nachtergaele96_infinite_volume_ground_projection}.
\end{lemma}
\begin{proof}\leanok
    \uses{thm:stationary_support_faithful_generators,
        thm:left_canonical_physical_dimension,
        thm:faithful_quasi_local_state_density_support,
        thm:mps_quasi_local_interval_expectation}
    Support compression gives $E>0$, faithful stationary generators
    $B,\sigma$, and $\eta_{A,\rho}(X)=\eta_{B,\sigma}(X)$ at every length.
    Normalization of $B$ implies positive physical dimension. The faithful
    state has a density $\tau_k$ of range $G_B(k)$ on every translated
    interval. The interval evaluation formula for the original quasi-local
    state identifies its restriction with $\eta_{A,\rho}$, proving
    (\ref{eq:stationary_state_density_support}).
\end{proof}

\begin{lemma}[The generated state has zero canonical parent energy]
    \label{thm:canonical_parent_generated_state}
    \lean{MPSTensor.quasiLocalExpectation_mem_parentGroundStateFace_canonicalParentInteraction}
    \leanok
    Let $d\geq1$, let $\sum_iA_i^*A_i=I$, and let $\rho\geq0$ have
    nonzero trace and satisfy $E_A(\rho)=\rho$. For every range $R\geq0$,
    the quasi-local state defined by $A,\rho$ belongs to the zero-energy
    face of the canonical interaction determined by $G_A(R)$.
    No faithfulness or purity is assumed.
\end{lemma}
\begin{proof}\leanok
    \uses{def:canonical_parent_interaction_matrix,thm:mps_quasi_local_parent_ground_state_face_membership,thm:canonical_parent_interaction_is_parent_interaction}
    The canonical interaction annihilates $G_A(R)$. The finite insertion
    expectation of every translate of that interaction is consequently zero,
    so the generated state is a zero-energy state, as in
    \cite[Lemma existenceinteraction]{Nachtergaele1996Spectral}.
\end{proof}

\begin{lemma}[Canonical interaction for a periodic family]
    \label{thm:periodic_family_canonical_interaction}
    \lean{MPSTensor.exists_ker_openParentHamiltonianES_of_periodic_family,
        MPSTensor.exists_positive_parent_interaction_of_periodic_family}
    \leanok
    Let $d\ge1$, let $(A_j)$ be a finite family of normalized periodic
    tensors with arbitrary periods, and let $\mu_j\ne0$.
    For $T=\bigoplus_j\mu_jA_j$, there exists $R_0>0$ such that
    \begin{align}
        \forall\,N\ge R\ge R_0:\ \ker H_{R,N}^{\mathrm{open}}(T)=\mathcal G_T(N).
    \end{align}
    In particular, the canonical local projection of range $R_0$ defines
    a positive interaction with exact open kernels at every $N\ge R_0$.
    No inequivalence or independence of the given sectors is assumed.
    The finite family may be empty. This is the finite periodic-tensor
    presentation of the interaction construction in Section~4 of
    \cite{Nachtergaele1996Spectral}; the passage from general GVBS data to such a
    presentation remains separate.
\end{lemma}
\begin{proof}\leanok
    \uses{thm:periodic_sector_representatives_exact_support,thm:periodic_ground_space_eventual_independence,thm:periodic_family_interaction_independence_criterion,lem:canonical_interaction_support_transport,lem:positive_canonical_interaction_exact_open_kernels}
    Select inequivalent periodic representatives without changing any
    finite joint ground space. Their middle ground spaces are eventually
    independent, so the joint-intersection criterion constructs exact
    canonical kernels. Equality of the local support spaces carries
    the canonical interaction back to the original weighted family.
\end{proof}

\begin{lemma}[Constructed periodic-family commutator gap]
    \label{thm:periodic_family_constructed_commutator_gap}
    \lean{MPSTensor.exists_positive_parent_interaction_quasiLocalCommutator_gap_of_periodic_family}
    \leanok
    Under the preceding hypotheses, there are a positive interaction $h$ of range $R>0$, with exact open kernels for all $N\ge R$,
    and a constant $\gamma>0$ such that every pure normalized zero-energy
    state $\varphi$ and every centered local observable $X$ on a
    positive-length interval beginning at $a$ satisfy
    \begin{align}
        \lim \varphi\bigl(X^*[H_{[a-\ell,a+k+r)},X]\bigr)=e,
        \qquad \gamma\varphi(X^*X)\le e.
    \end{align}
    The limit holds along every interval exhaustion with both margins
    diverging. The constant precedes the state, observable, filter, and
    margins. Neither translation invariance nor given interaction,
    kernel, separation, or independence data are required.
    The statement concerns the finite periodic-tensor presentation;
    the general GVBS presentation step remains separate.
\end{lemma}
\begin{proof}\leanok
    \uses{thm:periodic_family_canonical_interaction,
        thm:periodic_block_family_quasi_local_gap}
    Apply the periodic-family original-chain commutator theorem to the
    constructed interaction. Its eventual kernel identity follows from
    the exact kernels for every volume at least the interaction range.
\end{proof}

\begin{lemma}[Canonical interaction for faithful stationary generators]
    \label{thm:faithful_generators_canonical_interaction}
    \lean{MPSTensor.exists_ker_openParentHamiltonianES_of_leftCanonical_of_posDef_fixedPoint,
        MPSTensor.exists_positive_parent_interaction_of_leftCanonical_of_posDef_fixedPoint}
    \leanok
    Let $A$ have positive bond dimension, satisfy $\sum_iA_i^*A_i=I$,
    and have a faithful stationary matrix $\rho>0$.
    There exists $R_0>0$ such that its canonical open Hamiltonians satisfy
    \begin{align}
        \forall\,N\geq R\geq R_0:\ \ker H^{\mathrm{open}}_{R,N}(A)=G_A(N).
    \end{align}
    In particular, the canonical local projection of range $R_0$ defines
    a positive parent interaction with these exact kernels.
    The tensor need not be irreducible or periodic, and no interaction or
    kernel identity is given. This concerns the normalized generating
    data in \cite[Section~3]{Nachtergaele1996Spectral}; their construction
    from an arbitrary GVBS boundary limit remains separate.
\end{lemma}
\begin{proof}\leanok
    \uses{thm:faithful_periodic_ground_space_presentation,
        thm:periodic_family_canonical_interaction,
        lem:canonical_interaction_support_transport}
    The faithful stationary matrix makes every invariant projection reducing.
    Decompose into normalized irreducible periodic corners, preserving every
    boundary space. Construct the periodic-family canonical interaction.
    Equality of the local supports then transfers the canonical matrices
    and open Hamiltonians back to $A$.
\end{proof}

\begin{lemma}[Constructed gap for faithful stationary generators]
    \label{thm:faithful_generators_constructed_quasi_local_gap}
    \lean{MPSTensor.exists_positive_parent_interaction_quasiLocal_gap_of_leftCanonical_of_posDef_fixedPoint}
    \leanok
    Under the preceding hypotheses, there exist $R>0$, a positive interaction $h$, and $\gamma>0$ with exact open kernels $G_A(N)$ for
    all $N\geq R$. The state defined by $A,\rho$ has zero energy for $h$.
    Every pure normalized zero-energy state $\varphi$ and centered local
    observable $X$ on a positive-length interval satisfy
    \begin{align}
        \lim_{\Lambda\nearrow\mathbb Z}
        \varphi\bigl(X^*[H_{h,\Lambda},X]\bigr)
        \geq\gamma\varphi(X^*X).
    \end{align}
    The finite-volume limit holds along every interval exhaustion with both
    margins diverging. The interaction and constant precede the state,
    observable, filter, and margins. No translation invariance is imposed
    on the competing pure states.
\end{lemma}
\begin{proof}\leanok
    \uses{thm:faithful_generators_canonical_interaction,
        thm:canonical_parent_generated_state,
        thm:faithful_generators_quasi_local_gap}
    Choose the canonical interaction given by the preceding theorem.
    Its local projection annihilates the generated state support, so that
    state has zero energy. Its exact kernels give the eventual kernel
    identity in the faithful-generator commutator theorem.
\end{proof}

\begin{lemma}[Preservation of a primitive stationary fixed point]
    \label{thm:primitive_stationary_witness_transport}
    \lean{MPSTensor.IsPrimitiveMPS.blockTensor,
        MPSTensor.IsPrimitiveMPS.reindexPhysical}
    \leanok
    Let $A$ be a normalized primitive tensor of positive bond dimension with
    stationary matrix $\rho$. If $\rho$ is positive definite, then blocking by any positive length
    retains primitivity with the same stationary matrix $\rho$.
    A bijective relabelling of the physical letters also preserves the
    primitive fixed point with the same $\rho$; this second assertion requires
    only the original primitive fixed-point data.
\end{lemma}
\begin{proof}\leanok
    Positive blocking preserves trace preservation, irreducibility, and
    peripheral primitivity. The same $\rho$ is stationary for every
    positive power of the transfer map. The complementary fixed-point
    spectral gap therefore holds at this given stationary matrix.
    Under a bijective physical relabelling, both the normalization sum
    and the transfer map are unchanged.
\end{proof}

\begin{lemma}[Equality of quasi-local MPS expectations]
    \label{thm:quasi_local_mps_expectation_equality}
    \lean{MPSTensor.quasiLocalExpectation_eq_of_forall_observableInsertionExpectation_eq}
    \leanok
    Let $d>0$, let $A,B$ be trace-preserving tensors with possibly
    different bond dimensions, and let $\rho,\sigma\geq0$ be stationary
    matrices of nonzero trace. If
    $\eta_{A,\rho}(X)=\eta_{B,\sigma}(X)$ for every observable at every
    finite length, then their continuous quasi-local expectation
    functionals satisfy $\omega_{A,\rho}=\omega_{B,\sigma}$.
    Faithfulness and purity are unnecessary. This concerns the generating
    expectations in \cite[Section~3]{Nachtergaele1996Spectral}.
\end{lemma}
\begin{proof}\leanok
    \uses{def:mps_finite_region_expectation,def:mps_compatible_local_state,
        thm:mps_quasi_local_expectation_state,thm:qca_compatible_state_extension}
    Embed any finite region into its canonical containing interval.
    Its expectation is the insertion value of that interval observable,
    so the two finite-region functionals agree. Uniqueness of continuous
    extension from the finite-region algebras gives equality on the
    quasi-local algebra.
\end{proof}

\begin{lemma}[Stationary support compression preserves the quasi-local state]
    \label{thm:stationary_support_quasi_local_state_equality}
    \lean{MPSTensor.exists_faithful_quasiLocalExpectation_supportCompression_of_leftCanonical}
    \leanok
    Let $A$ satisfy $\sum_iA_i^*A_i=I$, and let $\rho\geq0$ have nonzero
    trace with $E_A(\rho)=\rho$. There exist a positive dimension $E$,
    an isometry $V$, a normalized tensor $B$, and a faithful stationary
    matrix $\sigma>0$ with
    \begin{align}
        V^*V                    & =I,                           & VV^*   & =\operatorname{supp}\rho, &
        B_i                     & =V^*A_iV,                     & \sigma & =V^*\rho V,\notag           \\
        \sum_iB_i^*B_i          & =I,                           & A_iV   & =VB_i,                    &
        E_B(\sigma)             & =\sigma,                      &
        \operatorname{Tr}\sigma & =\operatorname{Tr}\rho.\notag
    \end{align}
    The physical dimension is positive, and the continuous
    functionals agree: $\omega_{A,\rho}=\omega_{B,\sigma}$.
    No purity or minimality is assumed. This is a consequence for given
    stationary generating data in \cite[Section~3]{Nachtergaele1996Spectral}.
    Their construction from arbitrary GVBS boundary limits remains separate
    \cite{gap:nachtergaele96_infinite_volume_ground_projection}.
\end{lemma}
\begin{proof}\leanok
    \uses{thm:stationary_support_faithful_generators,
        thm:left_canonical_physical_dimension,
        thm:quasi_local_mps_expectation_equality}
    Compress the generators and stationary matrix to their stationary
    support. The compressed bond dimension is positive, and normalization
    of $B$ implies positive physical dimension. All finite insertion
    expectations agree under compression. The preceding equality theorem
    therefore identifies their continuous extensions.
\end{proof}

\begin{lemma}[Physical residual-cover limit from primitive sectors]
    \label{thm:physical_residual_cover_limit_primitive_resolution}
    \lean{MPSTensor.residual_cover_projection_defect_tendsto_zero_of_primitive_resolution}
    \leanok
    Let $A$ have positive bond dimension and let its blocking by $L$ have
    a finite isometric resolution into primitive tensors $B_j$ of positive
    bond dimensions with
    faithful stationary matrices $\rho_j$ and pairwise inequivalent
    sectors. Suppose the embeddings satisfy
    \begin{align}
        V_j^*V_j=I,\qquad \sum_jV_jV_j^*=I,\qquad
        A^{[L]}_iV_j=V_j(B_j)_i,\qquad
        V_j^*A^{[L]}_i=(B_j)_iV_j^*.
    \end{align}
    Assume that for every $j,r$ the matrix
    \begin{align}
        Q_j(r)=\sum_{|\tau|=r}(A^\tau)^*V_jV_j^*A^\tau
    \end{align}
    is an orthogonal projection. For $N=KL+ML+r$, let $P_N$ project
    onto $\mathcal G_A(N)$, let $P_{\mathrm L}$ project onto
    $\mathcal G_A(KL+ML)\otimes\mathcal H_r$, and let
    $P_{\mathrm R}$ project onto
    $\mathcal H_{KL}\otimes\mathcal G_A(ML+r)$.
    Along every filter for which $M\to\infty$, with $K,r$ arbitrary,
    \begin{align}
        \|P_N-P_{\mathrm L}P_{\mathrm R}\|\longrightarrow0.
    \end{align}
    No periodicity assumption on $A$ is imposed. All sector-resolution
    and tail-Gram hypotheses are given explicitly.
\end{lemma}
\begin{proof}\leanok
    \uses{thm:residual_window_joint_projection_decay,
        thm:residual_window_right_joint_range,
        thm:residual_window_left_joint_range_cover,
        thm:residual_window_full_joint_range_cover,
        thm:isometric_projection_defect_transport}
    The three joint sector spaces equal the corresponding original
    interval spaces under one configuration isometry. Their joint
    projection defect converges to zero by primitive-sector decay and
    supported cancellation. Isometric conjugation preserves its norm;
    taking the adjoint reverses the two interval projections.
\end{proof}

\begin{lemma}[All-length gap criterion from a primitive resolution]
    \label{thm:all_length_gap_primitive_resolution_criterion}
    \lean{MPSTensor.exists_pos_openParentHamiltonianES_gap_of_aligned_gap_of_primitive_resolution}
    \leanok
    Assume the preceding hypotheses and $R,L>0$. Suppose that the
    canonical range-$R$ open Hamiltonian has kernel $\mathcal G_A(N)$
    for every $N\ge R$, and has a common positive norm gap at all
    sufficiently large lengths divisible by $L$. Then there is
    $\delta>0$ such that, for every $N$ and every
    $v\perp\ker H_{R,N}^{\mathrm{open}}(A)$,
    \begin{align}
        \delta\|v\|\le\|H_{R,N}^{\mathrm{open}}(A)v\|.
    \end{align}
    Exact kernels and the aligned gap remain explicit assumptions of
    this criterion. The conclusion uses the kernel also at
    volumes shorter than the interaction range.
\end{lemma}
\begin{proof}\leanok
    \uses{thm:physical_residual_cover_limit_primitive_resolution,
        thm:aligned_residue_cover_gap_limit}
    The first theorem gives the residual-cover limit uniformly along
    all prefix and tail selections. The aligned-cover criterion gives
    an eventual original-length gap. A positive minimum over the
    remaining finite set of volumes gives one bound at every length.
\end{proof}

\begin{lemma}[Coordinate resolution of a finite direct sum]
    \label{thm:block_inclusion_coordinate_resolution}
    \lean{MPSTensor.sum_blockInclusion_mul_conjTranspose,MPSTensor.blockInclusion_conjTranspose_mul_toTensorFromBlocks}
    \leanok
    Let $J_j$ be the coordinate inclusion of the $j$th summand of a finite
    direct sum, and let $T_i=\bigoplus_j\mu_j A_j^i$. Then
    \begin{align}
        \sum_jJ_jJ_j^*=I,\qquad J_j^*T_i=\mu_jA_j^iJ_j^*.
    \end{align}
    The statements permit summands of dimension zero and an empty family.
    They are the coordinate identities underlying the canonical direct-sum
    decomposition in \cite{Cirac2016MPDO_arXiv}, Section~II.
\end{lemma}
\begin{proof}\leanok
    The unflattened coordinate inclusions resolve the identity of the dependent
    direct sum and intertwine every block diagonal matrix with its selected
    block. Carry these identities through the common coordinate equivalence.
\end{proof}

\begin{lemma}[Constructed parent interaction for stationary MPS states]
    \label{thm:stationary_generators_parent_interaction_gap}
    \lean{MPSTensor.exists_positive_parent_interaction_quasiLocal_gap_of_leftCanonical_of_posSemidef_fixedPoint}
    \leanok
    Let $A$ satisfy $\sum_i A_i^*A_i=I$, and let $\rho\geq0$ have
    nonzero trace and satisfy $\sum_i A_i\rho A_i^*=\rho$.
    There exist a tensor $B$ of positive bond dimension and a faithful
    stationary matrix $\sigma>0$ such that
    \begin{align}
        \sum_i B_i^*B_i=I,
        \qquad \sum_i B_i\sigma B_i^*=\sigma,
        \qquad \omega_{A,\rho}=\omega_{B,\sigma}.
    \end{align}
    In particular, the physical dimension is positive. There also exist
    a range $R>0$, a positive interaction $h$, and $\gamma>0$ such that
    \begin{align}
        \forall\,N\geq R:\ \ker H^{\mathrm{open}}_{h,N}=G_B(N),
        \qquad \omega_{A,\rho}\in\mathcal F_h.
    \end{align}
    Here $\mathcal F_h$ denotes the normalized positive states with zero
    expectation of every translate of $h$. For every pure
    $\varphi\in\mathcal F_h$, every local observable $X$ on a
    positive-length interval $[a,a+k)$ with $\varphi(X)=0$, and every
    interval exhaustion whose left and right margins both diverge,
    \begin{align}
        \lim \varphi\bigl(X^*[H_{h,[a-\ell,a+k+r)},X]\bigr)=e,
        \qquad \gamma\varphi(X^*X)\leq e.
    \end{align}
    The interaction and constant are chosen before the pure state,
    observable, filter, and margins. The generated state may be mixed;
    neither its purity nor translation invariance of the competing
    states is required. The finite kernels are $G_B(N)$; equality with
    $G_A(N)$ is not asserted.

    This is a consequence for given stationary generating data in
    \cite[Section~3]{Nachtergaele1996Spectral}. Their construction from
    arbitrary GVBS boundary limits remains separate, as documented in
    \cite{gap:nachtergaele96_infinite_volume_ground_projection}.
\end{lemma}
\begin{proof}\leanok
    \uses{thm:stationary_support_quasi_local_state_equality,
        thm:faithful_generators_constructed_quasi_local_gap}
    Compress the stationary matrix to its support. This gives normalized
    faithful stationary generators for the same quasi-local state.
    Apply the constructed parent-interaction theorem to those generators.
    Equality of the state functionals transfers zero-energy membership
    to the original generated state. The interaction, its exact kernels,
    and its common commutator gap are retained.
\end{proof}

\begin{lemma}[Cyclic word Gram in the inverse convention]
    \label{thm:inverse_cyclic_word_gram}
    \lean{MPSTensor.sum_word_conjTranspose_inverseCyclic_projection_mul_word}
    \leanok
    Suppose $\sum_i A_i^*A_i=\mathbf1$ and
    $P_{j+1}A_i=A_iP_j$, with labels in $\mathbb Z/m\mathbb Z$ for some
    $m>0$.
    Then, for every $r\geq0$,
    \begin{align}
        \sum_{|\tau|=r}A_\tau^*P_jA_\tau=P_{j-r}.
    \end{align}
\end{lemma}
\begin{proof}\leanok
    Replace the label $j$ by $-j$ and apply the forward cyclic word identity.
\end{proof}

\begin{lemma}[Cyclic sectors at a common blocking length]
    \label{thm:powered_cyclic_primitive_resolution}
    \lean{MPSTensor.IsPeriodic.exists_powered_cyclic_primitive_sector_resolution}
    \leanok
    Let $A$ be a normalized irreducible periodic tensor of period $m$, and let
    $L>0$ be divisible by $m$. There are positive-dimensional primitive tensors
    $B_j$ on the $L$-site alphabet, faithful stationary matrices $\rho_j$,
    cyclic projections $P_j$, and isometries $V_j$ such that
    \begin{align}
        V_jV_j^*=P_j,\qquad \sum_jP_j=\mathbf1,\qquad
        P_{j+1}A_i=A_iP_j,
    \end{align}
    \begin{align}
        A^{[L]}_iV_j=V_j(B_j)_i,\qquad
        V_j^*A^{[L]}_i=(B_j)_iV_j^*.
    \end{align}
    The tensors $B_j$ are pairwise inequivalent up to similarity and phase.
    For every tail length $r$, their support projections satisfy
    $\sum_{|\tau|=r}A_\tau^*V_jV_j^*A_\tau=P_{j-r}$.
\end{lemma}
\begin{proof}\leanok
    \uses{thm:cyclic_primitive_sector_resolution,thm:distinct_periodic_corner_separation,thm:primitive_stationary_witness_transport,thm:inverse_cyclic_word_gram}
    Write $L=mq$. Power each minimum-period primitive sector by $q$ and identify
    the iterated physical alphabet with the $L$-site alphabet. This preserves
    the stationary matrices and both intertwiners. Separation of the powered
    sectors follows by applying original-chain separation to the distinct
    period-one sectors. The last identity follows from the cyclic word Gram.
\end{proof}

\begin{lemma}[Two representations of the open interaction sum]
    \label{thm:all_length_open_interaction_matrix}
    \lean{MPSTensor.openInteractionMatrix_eq_zero_of_length_lt,MPSTensor.openInteractionHamiltonianES_eq_zero_of_length_lt,MPSTensor.openInteractionHamiltonianES_eq_toEuclideanLin_openInteractionMatrix_all}
    \leanok
    Let $h$ be a matrix on an $R$-site interval. Its matrix sum over all
    nonwrapping windows and the corresponding linear operator both vanish
    when the chain length satisfies $N<R$. If $R>0$, the matrix and linear
    operator representations of the open interaction sum agree for every
    $N\geq0$. No positivity assumption on $h$ is needed.
    This is the empty-sum convention in \cite{Nachtergaele1996Spectral},
    equation~\textup{(3.12)}.
\end{lemma}
\begin{proof}\leanok
    For $N<R$ there is no nonwrapping window of length $R$, so both sums
    are empty. For $N\geq R$, the coordinate representation intertwines
    each local term, and linearity identifies their sums.
\end{proof}

\begin{lemma}[Primitive resolution of a periodic family]
    \label{thm:periodic_family_literal_primitive_resolution}
    \lean{MPSTensor.exists_commonBlock_primitive_sector_resolution_of_isPeriodic}
    \leanok
    Let $(A_j)$ be a finite family of pairwise inequivalent normalized
    irreducible periodic tensors, with periods $m_j$, and let $L>0$ be divisible
    by every $m_j$. Put $\Phi_i=\bigoplus_j(A_j)_i$. There is a family of
    positive-dimensional primitive tensors $(B_s)$ on the $L$-site alphabet,
    indexed by pairs $s=(j,a)$ with $a\in\mathbb Z/m_j\mathbb Z$, with faithful
    stationary matrices $\rho_s$ and isometries $V_s$ into the original direct sum,
    such that
    \begin{align}
        \sum_sV_sV_s^*=\mathbf1,\qquad
        \Phi^{[L]}_iV_s=V_s(B_s)_i,\qquad
        V_s^*\Phi^{[L]}_i=(B_s)_iV_s^*.
    \end{align}
    The tensors $B_s$ are pairwise inequivalent up to similarity and phase.
    For every $s$ and every tail length $r\geq0$, the matrix
    \begin{align}
        Q_s(r)=\sum_{|\tau|=r}\Phi_\tau^*V_sV_s^*\Phi_\tau
    \end{align}
    is an orthogonal projection. The family may be empty.
\end{lemma}
\begin{proof}\leanok
    \uses{thm:powered_cyclic_primitive_resolution,thm:distinct_periodic_corner_separation,thm:block_inclusion_coordinate_resolution}
    Choose the powered cyclic resolution of each original tensor and
    compose its isometries with the coordinate inclusions of the direct sum.
    The inclusions resolve the identity and preserve both intertwiners.
    Within one original tensor, separation is part of the powered resolution;
    between different original tensors, it follows from original-chain angle
    decay. Finally, if $s=(j,a)$, the word identity gives
    $Q_s(r)=J_jP_{j,a-r}J_j^*$. This is an orthogonal projection because $J_j$
    is an isometry and $P_{j,a-r}$ is a cyclic projection.
\end{proof}

\begin{lemma}[Positive canonical interaction from exact open kernels]
    \label{lem:positive_canonical_interaction_exact_open_kernels}
    \lean{MPSTensor.exists_positive_canonical_parent_interaction_of_exact_open_kernels}
    \leanok
    Let $A$ be a tensor and let $R>0$. Suppose the canonical open
    Hamiltonian has kernel $\mathcal G_A(N)$ for every $N\ge R$.
    Then the range-$R$ canonical orthogonal-complement projection is a
    positive interaction $h$, and the open interaction generated by $h$
    has those same exact kernels. The interaction is the
    configuration-basis matrix of that canonical projection.
\end{lemma}
\begin{proof}\leanok
    \uses{def:canonical_parent_interaction_matrix}
    The canonical orthogonal-complement projection is positive. Its
    configuration-basis matrix represents the same local operator,
    hence gives the same open Hamiltonian and kernel at every volume.
\end{proof}

\begin{lemma}[Transfer of canonical gaps to an eventual-kernel interaction]
    \label{thm:canonical_gap_to_eventual_kernel_interaction}
    \lean{MPSTensor.exists_openInteractionHamiltonianES_gap_of_canonical_gaps_of_eventual_kernel}
    \leanok
    Let $A$ be a tensor. Suppose that, for every sufficiently large range $W$,
    its canonical open Hamiltonians have kernels $\mathcal G_A(N)$ for $N\ge W$
    and a positive uniform norm gap above their kernels at all volumes.
    Let $h\ge0$ have positive range $R$, and suppose
    $\ker H_h(N)=\mathcal G_A(N)$ for all sufficiently large $N$.
    Then there exists $\gamma>0$ such that, for every $N\ge0$ and
    $v\perp\ker H_h(N)$,
    \begin{align}
        \gamma\|v\|\le\|H_h(N)v\|.
    \end{align}
    No local kernel condition or tensor injectivity is imposed.
\end{lemma}
\begin{proof}\leanok
    \uses{lem:positive-kernel-projection-comparison,thm:positive_open_interaction_range_comparison}
    Choose $W\ge R$ large enough for both kernel identities and the canonical
    gap. Positivity and finite dimensionality give
    $H_h(W)\ge\kappa P_{\mathcal G_A(W)^\perp}$ with $\kappa>0$.
    Sum this comparison over all $W$-site windows. Their multiplicity is at most
    $W-R+1$, so the canonical gap $\delta_W$ gives the positive bound
    $\kappa\delta_W/(W-R+1)$ on all $N\ge W$.
    Take its minimum with the finitely many positive gaps at $N<W$, each
    measured above the kernel.
\end{proof}

\begin{definition}[Physical matrix slices of a given map]
    \label{def:supplied_physical_isometry_tensor}
    \lean{MPSTensor.tensorOfPhysicalIsometry}
    \leanok
    Let $V:\mathbb C^D\to\mathbb C^d\otimes\mathbb C^D$ be a linear map,
    with the physical coordinate written first. Define its physical matrix
    slices by
    \begin{align}
        (A_i)_{ab} & =V_{(i,a),b}.
    \end{align}
    This is the coordinate convention for the given representation in
    \cite[equation~(3.6)]{Nachtergaele1996Spectral}.
\end{definition}

\begin{lemma}[Normalization and one-site Heisenberg insertion]
    \label{thm:supplied_isometry_heisenberg_insertion}
    \lean{MPSTensor.sum_conjTranspose_mul_tensorOfPhysicalIsometry,
        MPSTensor.isLeftCanonical_tensorOfPhysicalIsometry,
        MPSTensor.physicalIsometry_heisenbergInsertion}
    \leanok
    For the preceding slices,
    \begin{align}
        \sum_iA_i^*A_i   & =V^*V,
        \\
        V^*(X\otimes B)V & =\sum_{i,j}X_{ij}A_i^*BA_j
        \quad(X\in M_d(\mathbb C),\ B\in M_D(\mathbb C)).
    \end{align}
    Thus an isometry gives trace-preserving normalized tensor matrices.
    The insertion identity holds for any given $V$, without an isometry
    assumption. No dimension, stationarity, or faithfulness assumption is
    needed. This identifies the one-site coordinates of
    \cite[equation~(3.6)]{Nachtergaele1996Spectral}.
\end{lemma}
\begin{proof}\leanok
    \uses{def:supplied_physical_isometry_tensor}
    Expand the matrix products in physical and virtual coordinates.
    The normalization is the Gram matrix of $V$. Permuting the finite sums
    in the insertion gives the displayed two-sided Kraus expression.
\end{proof}

\begin{lemma}[Trace duality of the one-site insertion]
    \label{thm:supplied_isometry_insertion_trace_duality}
    \lean{MPSTensor.physicalObservableTransfer_oneSite_apply,
        MPSTensor.trace_physicalIsometry_heisenbergInsertion}
    \leanok
    Identify configurations on one site with their physical letter. For
    any tensor $A$, its one-site physical observable transfer is
    \begin{align}
        E_X(\rho) & =\sum_{i,j}X_{ij}A_j\rho A_i^*.
    \end{align}
    For the slices of $V$, this is trace-dual to the preceding Heisenberg
    insertion:
    \begin{align}
        \operatorname{Tr}\bigl(\rho V^*(X\otimes B)V\bigr)
         & =\operatorname{Tr}\bigl(B E_X(\rho)\bigr).
    \end{align}
    The matrices $X,B,\rho$ are arbitrary. These identities neither
    construct a given isometry from arbitrary GVBS data nor identify
    multi-site expectations with a particular ordering of source sites.
\end{lemma}
\begin{proof}\leanok
    \uses{thm:supplied_isometry_heisenberg_insertion}
    Reindex the one-site configuration sum by its unique physical letter.
    Cyclicity of the trace turns each term
    $X_{ij}\operatorname{Tr}(\rho A_i^*BA_j)$ into
    $X_{ij}\operatorname{Tr}(BA_j\rho A_i^*)$.
\end{proof}

\begin{lemma}[Uniform open gaps for faithful stationary generators]
    \label{thm:faithful_generators_uniform_open_gap}
    \lean{MPSTensor.exists_openParentHamiltonianES_gap_of_leftCanonical_of_posDef_fixedPoint}
    \leanok
    Let $A$ have positive bond dimension and satisfy
    $\sum_i A_i^*A_i=I$. Suppose that $\rho>0$ and
    $\sum_i A_i\rho A_i^*=\rho$. There is $R_0>0$ such that, for each
    $R\geq R_0$, the canonical open Hamiltonians satisfy
    \begin{align}
        \forall\,N\geq R:\ \ker H_{R,N}=G_A(N),
    \end{align}
    and there is $\delta_R>0$ such that
    \begin{align}
        \forall\,N\geq0,\ v\perp\ker H_{R,N}:\ \delta_R\|v\|\leq\|H_{R,N}v\|.
    \end{align}
    The constant is independent of the original chain length, including its
    residue modulo every blocking length. This is the given generating-data
    case of \cite{Nachtergaele1996Spectral}, Theorems~1.1--1.2.
\end{lemma}
\begin{proof}\leanok
    \uses{thm:faithful_periodic_ground_space_presentation,thm:periodic_family_uniform_original_open_gap,lem:canonical_interaction_support_transport}
    Decompose the faithful stationary tensor into a finite periodic family,
    with equality of the boundary spaces at every length. The uniform family
    gap gives a common range threshold. Equality of the local support spaces
    identifies the canonical interactions and all their open Hamiltonians.
    Carry the exact kernels and the norm gap through these identities.
\end{proof}

\begin{lemma}[One faithful parent interaction with both gaps]
    \label{thm:faithful_generators_constructed_finite_quasi_local_gap}
    \lean{MPSTensor.exists_positive_parent_finite_and_quasiLocal_gap_of_leftCanonical_of_posDef_fixedPoint}
    \leanok
    Under the preceding hypotheses, there are $R>0$, a positive interaction $h$
    of range $R$, and constants $\delta,\gamma>0$ such that
    \begin{align}
        \forall\,N\geq R:\ \ker H_h(N)=G_A(N),\qquad
        \forall\,N\geq0,\ v\perp\ker H_h(N):\ \delta\|v\|\leq\|H_h(N)v\|.
    \end{align}
    The generated stationary state belongs to the zero-energy face of $h$.
    Every pure state $\varphi$ in this face and every observable $X$ on an
    interval of positive length with $\varphi(X)=0$ satisfy
    \begin{align}
        \lim_{\Lambda\nearrow\mathbb Z}
        \varphi\bigl(X^*[H_{h,\Lambda},X]\bigr)
        \geq\gamma\varphi(X^*X).
    \end{align}
    The limit holds along every interval exhaustion for which both margins
    around the support of $X$ diverge. The choices of $R,h,\delta,\gamma$
    precede the state, observable, and exhaustion. No purity of the generated
    stationary state is assumed. This is the given generating-data case of
    \cite{Nachtergaele1996Spectral}, Theorems~1.1--1.2 and Section~6.
\end{lemma}
\begin{proof}\leanok
    \uses{thm:faithful_generators_uniform_open_gap,lem:positive_canonical_interaction_exact_open_kernels,thm:canonical_parent_generated_state,thm:faithful_generators_quasi_local_gap}
    Choose the canonical interaction at the uniform threshold. The first theorem
    gives its exact kernels and finite-volume gap. The generated state has zero
    energy for this canonical interaction. Its eventual exact kernel identity
    also gives the common commutator gap in every pure zero-energy state.
\end{proof}

\begin{lemma}[Uniform original-length open gaps for periodic families]
    \label{thm:periodic_family_uniform_original_open_gap}
    \lean{MPSTensor.exists_openParentHamiltonianES_gap_of_periodic_family}
    \leanok
    Let $d\ge1$, let $(A_j)$ be a finite family of normalized periodic
    tensors with arbitrary periods, and let $\mu_j\ne0$.
    For $T=\bigoplus_j\mu_jA_j$, there exists $R_0>0$ such that every
    $R\ge R_0$ has both of the following properties:
    \begin{align}
        \forall\,N\ge R:\ \ker H_{R,N}^{\mathrm{open}}(T)=\mathcal G_T(N),
    \end{align}
    and there exists $\delta_R>0$ such that, for every $N\ge0$ and
    every $v\perp\ker H_{R,N}^{\mathrm{open}}(T)$,
    \begin{align}
        \delta_R\|v\|\le\|H_{R,N}^{\mathrm{open}}(T)v\|.
    \end{align}
    The finite family may be empty. No inequivalence, invariant matrix,
    primitive presentation, interaction, or kernel identity is given.
    The norm bound uses the kernel also when $N<R$.
    This is the finite periodic-tensor presentation of the gap theorem;
    the passage from general GVBS data to such a presentation remains
    separate.
\end{lemma}
\begin{proof}\leanok
    \uses{thm:periodic_family_canonical_interaction,
        thm:periodic_sector_representatives_exact_support,
        thm:periodic_family_literal_primitive_resolution,
        thm:all_length_gap_primitive_resolution_criterion,
        thm:periodic_block_family_aligned_gap_of_dvd,
        lem:canonical_interaction_support_transport}
    Select inequivalent representatives without changing any finite
    ground space. If the selected family is nonempty, a common period
    multiple gives a primitive isometric resolution and exact
    orthogonal tail Gram projections. The aligned finite gap and the
    residual-cover limit extend the gap to every original length.
    For an empty selected family, the aligned theorem applies with
    blocking length one. Equality of local supports carries the
    canonical gap to the original weighted family.
\end{proof}

\begin{lemma}[A periodic-family interaction with finite and infinite gaps]
    \label{thm:periodic_family_constructed_finite_quasi_local_gap}
    \lean{MPSTensor.exists_positive_parent_interaction_finite_quasiLocal_gap_of_periodic_family}
    \leanok
    Under the preceding hypotheses, one can choose $R>0$ and the
    canonical positive interaction $h$ of range $R$ so that its open
    kernels are $\mathcal G_T(N)$ for every $N\ge R$.
    There exists $\delta>0$ giving a norm gap above the open
    kernel at every original volume, including $N<R$.
    For the same interaction there exists $\gamma>0$ such that every
    pure normalized zero-energy state $\varphi$ and every centered
    local observable $X$ on a positive-length interval satisfy
    \begin{align}
        \lim\varphi\bigl(X^*[H_{[a-\ell,a+k+r)},X]\bigr)=e,
        \qquad \gamma\varphi(X^*X)\le e.
    \end{align}
    The limit holds along every exhaustion with both margins diverging.
    Both constants precede the states, observables, filters, and margins.
    No translation invariance of a classified state is assumed.
    The general GVBS presentation step remains separate.
\end{lemma}
\begin{proof}\leanok
    \uses{thm:periodic_family_uniform_original_open_gap,
        lem:positive_canonical_interaction_exact_open_kernels,
        thm:periodic_block_family_quasi_local_gap}
    Choose any range at the canonical threshold and use its canonical
    local projection matrix. The all-length canonical norm gap transfers
    to the matrix interaction sum, including short volumes. The exact
    kernels give the eventual-kernel hypothesis of the original-chain
    periodic-family commutator theorem.
\end{proof}

\begin{lemma}[Constructed finite and infinite gaps for stationary MPS states]
    \label{thm:stationary_generators_finite_quasi_local_gap}
    \lean{MPSTensor.exists_positive_parent_finite_and_quasiLocal_gap_of_leftCanonical_of_posSemidef_fixedPoint}
    \leanok
    Let $A$ satisfy $\sum_i A_i^*A_i=I$, and let $\rho\geq0$ have
    nonzero trace and satisfy $\sum_i A_i\rho A_i^*=\rho$.
    There exist a tensor $B$ of positive bond dimension and a faithful
    stationary matrix $\sigma>0$ such that
    \begin{align}
        \sum_i B_i^*B_i        & =I,
        \\
        \sum_i B_i\sigma B_i^* & =\sigma,
        \\
        \omega_{A,\rho}        & =\omega_{B,\sigma}.
    \end{align}
    In particular, the physical dimension is positive. There also exist
    a range $R>0$, a positive interaction $h$, and constants
    $\delta,\gamma>0$ such that
    \begin{align}
        \forall\,N\geq R:\ \ker H^{\mathrm{open}}_{h,N} & =G_B(N),
        \\
        \omega_{A,\rho}                                 & \in\mathcal F_h.
    \end{align}
    The same interaction satisfies the finite-volume estimate
    \begin{align}
        \delta\|v\| & \leq\|H^{\mathrm{open}}_{h,N}v\|
        \qquad\bigl(N\geq0,\ v\perp\ker H^{\mathrm{open}}_{h,N}\bigr).
    \end{align}
    Here $\mathcal F_h$ denotes the normalized positive states with zero
    expectation of every translate of $h$. For every pure
    $\varphi\in\mathcal F_h$, every local observable $X$ on a
    positive-length interval $[a,a+k)$ with $\varphi(X)=0$, and every
    interval exhaustion whose left and right margins both diverge,
    \begin{align}
        \lim \varphi\bigl(X^*[H_{h,[a-\ell,a+k+r)},X]\bigr) & =e,
        \\
        \gamma\varphi(X^*X)                                 & \leq e.
    \end{align}
    The interaction and constant are chosen before the pure state,
    observable, filter, and margins. The generated state may be mixed;
    neither its purity nor translation invariance of the competing
    states is required. The finite kernels are $G_B(N)$; equality with
    $G_A(N)$ is not asserted.

    This is a consequence for given stationary generating data in
    \cite[Section~3]{Nachtergaele1996Spectral}. Their construction from
    arbitrary GVBS boundary limits remains separate, as documented in
    \cite{gap:nachtergaele96_infinite_volume_ground_projection}.
\end{lemma}
\begin{proof}\leanok
    \uses{thm:stationary_support_quasi_local_state_equality,
        thm:faithful_generators_constructed_finite_quasi_local_gap}
    Compress the stationary matrix to its support. This gives normalized
    faithful stationary generators for the same quasi-local state.
    Apply the constructed parent-interaction theorem to those generators.
    Equality of the state functionals transfers zero-energy membership
    to the original generated state. The interaction, its exact kernels,
    its uniform finite-volume gap, and its common commutator gap are retained.
\end{proof}

\begin{lemma}[Parent gaps from a given physical isometry]
    \label{thm:supplied_isometry_parent_gap}
    \lean{MPSTensor.exists_stationary_density_tensorOfPhysicalIsometry,
        MPSTensor.exists_positive_parent_finite_and_quasiLocal_gap_of_physical_isometry}
    \leanok
    Let $D>0$ and let $V:\mathbb C^D\to\mathbb C^d\otimes\mathbb C^D$ be an
    isometry, with tensor slices $A_i$. Then $\sum_iA_i^*A_i=I$, and the
    transfer map of $A$ has a positive fixed point $\rho$ of trace one.
    There are a faithful normalized tensor $B$ with a positive
    definite fixed point that generates the same quasi-local state, a range
    $R>0$, and a positive interaction $h$ on $R$ sites whose open-chain
    kernels equal $G_N(B)$ for $N\geq R$, such that the generated state has
    zero energy, the open Hamiltonians of $h$ have one positive gap at every
    length, and one constant $\gamma>0$ bounds the commutator energy below in
    every pure zero-energy state, for every centered observable on an
    interval of positive length and every interval exhaustion whose left
    and right margins both diverge \cite[equation~(3.6), Theorems~1.1--1.2,
        and Section~6]{Nachtergaele1996Spectral}. The isometry is given; its
    construction from a general presentation is not asserted
    \cite{gap:nachtergaele96_infinite_volume_ground_projection}.
\end{lemma}
\begin{proof}\leanok
    \uses{thm:supplied_isometry_heisenberg_insertion,thm:stationary_generators_finite_quasi_local_gap}
    The isometry condition is the trace-preserving normalization, so the
    transfer channel has a fixed density matrix. Compressing to its support
    gives faithful generators of the same state, and the stationary
    faithful-generator theorem gives the interaction and both gaps.
\end{proof}

\begin{lemma}[Finite gaps for a given faithful parent interaction]
    \label{thm:faithful_generators_given_interaction_uniform_gap}
    \lean{MPSTensor.exists_given_interaction_gap_of_leftCanonical_of_posDef_fixedPoint}
    \leanok
    Let $A$ be normalized with positive bond dimension and a faithful stationary
    matrix, as in Lemma~\ref{thm:faithful_generators_uniform_open_gap}.
    Let $h\geq0$ be any given interaction of positive range $R$ such that $\ker H_h(N)=G_A(N)$ for all sufficiently large $N$.
    Then there is $\delta>0$ such that
    \begin{align}
        \forall\,N\geq0,\ v\perp\ker H_h(N):\ \delta\|v\|\leq\|H_h(N)v\|.
    \end{align}
    The local kernel $\ker h$ need not equal $G_A(R)$.
    This is the finite-volume conclusion of \cite{Nachtergaele1996Spectral},
    Theorem~1.2, in the given faithful generating-data setting.
\end{lemma}
\begin{proof}\leanok
    \uses{thm:faithful_generators_uniform_open_gap,thm:canonical_gap_to_eventual_kernel_interaction}
    Choose a sufficiently long window $W$ with exact kernel $G_A(W)$ and
    a uniform canonical gap at range $W$. Positivity gives a local lower bound
    $H_h(W)\geq\kappa P_{G_A(W)^\perp}$ with $\kappa>0$.
    Summing these windows counts each original term at most $W-R+1$ times.
    This transfers the canonical gap to all sufficiently long chains. The
    finitely many remaining volumes are included by a positive finite minimum
    above their kernels.
\end{proof}

\begin{lemma}[Both gaps for the same given faithful parent interaction]
    \label{thm:faithful_generators_given_interaction_finite_quasi_local_gap}
    \lean{MPSTensor.exists_given_interaction_finite_quasiLocal_gap_of_leftCanonical_of_posDef_fixedPoint}
    \leanok
    Under the preceding hypotheses, the same given interaction $h$ has a
    uniform finite-volume gap at every original chain length and one common
    infinite-volume commutator gap in every pure zero-energy state.
    The infinite-volume limit holds along every interval exhaustion with both
    margins around the local observable diverging. Both positive constants
    precede the state, observable, and exhaustion.
    This is the given-interaction conclusion of \cite{Nachtergaele1996Spectral},
    Theorem~1.2, in the given faithful generating-data setting.
\end{lemma}
\begin{proof}\leanok
    \uses{thm:faithful_generators_given_interaction_uniform_gap,thm:faithful_generators_quasi_local_gap}
    Apply the preceding finite-volume theorem and the faithful-generator
    commutator theorem to the same interaction and its eventual kernel identity.
\end{proof}

\begin{definition}[Reversed physical product of one-site observables]
    \label{def:source_reversed_product_observable}
    \lean{MPSTensor.reversedProductObservable}
    \leanok
    For a finite list of one-site observables $(X_1,\ldots,X_n)$, define
    $O_{\mathrm{rev}}=X_n\otimes\cdots\otimes X_1$. The empty product
    is the identity on the zero-site space. Equivalently,
    $O_{\mathrm{rev}}(X_1,X_2,\ldots,X_n)
        =O_{\mathrm{rev}}(X_2,\ldots,X_n)\otimes X_1$.
    This records the physical order required by trace duality of the
    source-written Heisenberg composition in
    \cite[equations~(3.5)--(3.6)]{Nachtergaele1996Spectral}.
\end{definition}

\begin{lemma}[Source-ordered Heisenberg expectation and physical reversal]
    \label{thm:source_ordered_heisenberg_trace_reversal}
    \lean{MPSTensor.sourceOrderedHeisenberg_trace_eq_reversedProduct_transfer}
    \leanok
    Let $V:\mathbb C^D\to\mathbb C^d\otimes\mathbb C^D$ be a given
    map, let $(A_i)_{ab}=V_{(i,a),b}$, and set
    $F_X(B)=V^*(X\otimes B)V$. For every matrix $\rho$ and every finite
    list $(X_1,\ldots,X_n)$,
    \begin{align}
        \operatorname{Tr}\bigl(\rho\,
        (F_{X_1}\circ\cdots\circ F_{X_n})(I)\bigr)
         & =\operatorname{Tr}\bigl(E_{X_n\otimes\cdots\otimes X_1}(\rho)\bigr).
    \end{align}
    Here $E$ is the physical observable transfer with forward matrix words.
    The equality includes $n=0$. It holds in particular for an isometry,
    whose slices satisfy $\sum_iA_i^*A_i=I$; neither normalization,
    positivity, nor stationarity is needed for the trace identity itself.
    This is the explicit order relation for the given representation in
    \cite[equations~(3.5)--(3.6)]{Nachtergaele1996Spectral}, and makes no
    assertion that arbitrary GVBS boundary data have this representation.
\end{lemma}
\begin{proof}\leanok
    \uses{def:source_reversed_product_observable,
        thm:supplied_isometry_insertion_trace_duality,
        lem:ldp_observable_transfer_products}
    Induct on the list, allowing an arbitrary matrix $\rho$. Trace duality
    moves the first Heisenberg insertion to $E_{X_1}(\rho)$.
    The induction hypothesis then gives
    $\operatorname{Tr}(E_{O_{\mathrm{rev}}(X_2,\ldots,X_n)}
        (E_{X_1}(\rho)))$.
    The inserted-transfer product identity identifies this with the
    expectation of $O_{\mathrm{rev}}(X_2,\ldots,X_n)\otimes X_1$.
    At the empty list, both expressions are $\operatorname{Tr}\rho$.
\end{proof}

\begin{lemma}[Given-interaction open gaps for a periodic family]
    \label{thm:periodic_family_given_interaction_uniform_open_gap}
    \lean{MPSTensor.exists_openInteractionMatrix_gap_of_periodic_family_of_eventual_kernel}
    \leanok
    Let $d\ge1$, let $(A_j)$ be a finite family of normalized periodic
    tensors, let $\mu_j\ne0$, and write $T=\bigoplus_j\mu_jA_j$.
    Let $h\ge0$ be a given interaction of range $R>0$ whose open
    Hamiltonians satisfy $\ker H_N(h)=\mathcal G_T(N)$ for every
    sufficiently large $N$.
    Then there exists $\delta>0$ such that every original volume $N$
    and every $v\perp\ker H_N(h)$ satisfy
    \begin{align}
        \delta\|v\| & \le\|H_N(h)v\|.
    \end{align}
    The kernel is used also on short intervals. No local-kernel
    identity at range $R$, lower bound on $R$ beyond positivity, sector
    separation, stationary matrix, or finite gap is given.
    The finite family may be empty. The statement concerns finite
    periodic-tensor presentations; the general GVBS presentation step
    remains separate. This is the given-interaction conclusion of
    \cite[Theorem~1.2]{Nachtergaele1996Spectral} in this setting.
\end{lemma}
\begin{proof}\leanok
    \uses{thm:periodic_family_uniform_original_open_gap,
        thm:canonical_gap_to_eventual_kernel_interaction}
    The canonical interactions have exact kernels and uniform gaps at
    every sufficiently large range. Choose one long interval on which
    both the canonical and given interaction kernels equal the MPS
    space. Finite-dimensional positivity bounds the excitation projection
    below the given Hamiltonian. The window comparison transfers the
    canonical gap, and a positive minimum includes the remaining short
    volumes above their kernels.
\end{proof}

\begin{lemma}[Finite and infinite gaps for a given periodic-family interaction]
    \label{thm:periodic_family_given_interaction_finite_quasi_local_gap}
    \lean{MPSTensor.exists_finite_quasiLocal_gap_of_periodic_family_of_eventual_kernel}
    \leanok
    Under the preceding hypotheses, the same given interaction $h$
    has an all-length finite norm gap with a constant $\delta>0$ and
    a constant $\gamma>0$ such that every pure normalized zero-energy
    state $\varphi$ and every centered local observable $X$ on a
    positive-length interval beginning at $a$ satisfy
    \begin{align}
        \lim\varphi\bigl(X^*[H_{[a-\ell,a+k+r)}(h),X]\bigr) & =e,    \\
        \gamma\varphi(X^*X)                                 & \le e.
    \end{align}
    The limit holds along every exhaustion with both margins diverging.
    Both constants precede the states, observables, filters, and margins.
    No translation invariance of a classified state is assumed. The
    general GVBS presentation step remains separate.
\end{lemma}
\begin{proof}\leanok
    \uses{thm:periodic_family_given_interaction_uniform_open_gap,
        thm:periodic_block_family_quasi_local_gap}
    Combine the given-interaction finite gap with the original-chain
    periodic-family commutator theorem, applied to exactly the same
    given interaction and eventual kernel identity.
\end{proof}
